\documentclass[11pt,a4paper,oneside]{book}

\usepackage[T1]{fontenc}
\usepackage[utf8]{inputenc}
\usepackage{lmodern}
\usepackage{microtype}
\usepackage{amsmath,amssymb,amsthm,amscd,mathtools,mathrsfs}
\usepackage{xspace}
\usepackage{xcolor}
\usepackage{graphicx}
\usepackage{array,booktabs,longtable}
\usepackage{enumitem}
\usepackage{tocloft}
\usepackage{tikz-cd}
\usepackage[all]{xy}
\usepackage[a4paper,margin=27mm,marginparwidth=18mm,marginparsep=4mm]{geometry}
\usepackage[hidelinks,pdfencoding=auto,hypertexnames=false]{hyperref}

% Modern compatibility layer for the translated bodies.  The legacy
% sga2-smf.sty is intentionally not loaded: it is incompatible with the
% current LaTeX kernel.  The local shim retains only the macros actually
% required by the source-aligned English reader.
% BEGIN INCLUDED active_source/sga2_reader_macros.tex
% Modern compatibility layer for the source-aligned English SGA 2 bodies.
% This deliberately replaces, rather than loads, the legacy sga2-smf.sty.

\makeatletter
\@ifundefined{c@pageoriginal}{\newcounter{pageoriginal}}{}
\@ifundefined{c@toto}{\newcounter{toto}}{}
\@ifundefined{c@yequation}{\newcounter{yequation}}{}
\@ifundefined{c@sourcecorollary}{\newcounter{sourcecorollary}[section]}{}
\@ifundefined{c@sourcecondition}{\newcounter{sourcecondition}}{}
\@ifundefined{c@footnoteAG}{\newcounter{footnoteAG}}{}
\@ifundefined{c@footnoteNDE}{\newcounter{footnoteNDE}}{}
\@ifundefined{c@aconjecture}{\newcounter{aconjecture}[section]}{}
\makeatother
\renewcommand{\theaconjecture}{\Alph{aconjecture}}

% The corrected/annotated branch is always selected.
\providecommand{\sisi}[2]{#2}
\providecommand{\pageoriginale}{%
  \refstepcounter{pageoriginal}%
  \marginpar{\raggedright\footnotesize\bfseries\thepageoriginal}%
  \xspace}
\providecommand{\srcpage}[1]{%
  \marginpar{\raggedright\scriptsize\color{gray}Original French p.~#1}}
\providecommand{\srcpdfpage}[1]{%
  \marginpar{\raggedright\scriptsize\color{gray}French source PDF p.~#1}}
\providecommand{\chapterspace}[1]{}

% Original and editorial notes.
\providecommand{\nde}[1]{\footnote{\textit{Editor's note.} #1}}
\providecommand{\de}[1]{\footnote{\textit{Editor's note.} #1}}
\providecommand{\ndemark}{\footnotemark}
\providecommand{\ndetext}[1]{\footnotetext{\textit{Editor's note.} #1}}
\providecommand{\sfootnote}[1]{\footnote{#1}}
\providecommand{\sfootnotemark}{\footnotemark}
\providecommand{\sfootnotetext}[1]{\footnotetext{#1}}
\providecommand{\markedfootnote}[2]{%
  \begingroup\renewcommand{\thefootnote}{#1}\footnote{#2}\endgroup}

% Common prose abbreviations and quotations retained from the source macros.
\providecommand{\og}{``}
\providecommand{\fg}{''\xspace}
\xspaceaddexceptions{\,\markedfootnote}
\providecommand{\cf}{cf.\kern.3em}
\providecommand{\Cf}{Cf.\kern.3em}
\providecommand{\ie}{i.e.\kern.3em}
\providecommand{\resp}{resp.\kern.3em}
\providecommand{\loccit}{\textit{loc. cit.}\xspace}
\providecommand{\EGA}{\textup{EGA}\kern.2em}
\providecommand{\SGA}{\textup{SGA}\kern.3em}
\providecommand{\Exp}{Exp.\kern.3em}
\providecommand{\ptbl}{.\kern.2em}
\providecommand{\numero}{no.\kern.3em}
\providecommand{\no}{no.\kern.3em}
\providecommand{\No}{no.\kern.3em}
\providecommand{\afortiori}{\textit{a fortiori}\xspace}
\providecommand{\quoi}{}
\providecommand{\refname}{Bibliography}

% Calligraphic, bold, and ideal alphabets.
\providecommand{\Aa}{\mathcal A}
\providecommand{\Bb}{\mathcal B}
\providecommand{\Cc}{\mathcal C}
\providecommand{\Dd}{\mathcal D}
\providecommand{\Ee}{\mathcal E}
\providecommand{\Ff}{\mathcal F}
\providecommand{\Gg}{\mathcal G}
\providecommand{\Hh}{\mathcal H}
\providecommand{\Ii}{\mathcal I}
\providecommand{\Jj}{\mathcal J}
\providecommand{\Kk}{\mathcal K}
\providecommand{\Ll}{\mathcal L}
\providecommand{\Mm}{\mathcal M}
\providecommand{\Nn}{\mathcal N}
\providecommand{\Oo}{\mathcal O}
\providecommand{\Pp}{\mathcal P}
\providecommand{\Qq}{\mathcal Q}
\providecommand{\Rr}{\mathcal R}
\providecommand{\Ss}{\mathcal S}
\providecommand{\Tt}{\mathcal T}
\providecommand{\Uu}{\mathcal U}
\providecommand{\Vv}{\mathcal V}
\providecommand{\Ww}{\mathcal W}
\providecommand{\Xx}{\mathcal X}
\providecommand{\Yy}{\mathcal Y}
\providecommand{\Zz}{\mathcal Z}
\providecommand{\CJ}{\underline{C}_J}
\providecommand{\DDJ}{\underline{D}_J}
\providecommand{\bbmu}{\boldsymbol{\mu}}
\providecommand{\CC}{\mathbf C}
\providecommand{\GG}{\mathbf G}
\providecommand{\KK}{\mathbf K}
\providecommand{\LL}{\mathbf L}
\providecommand{\NN}{\mathbf N}
\providecommand{\PP}{\mathbf P}
\providecommand{\QQ}{\mathbf Q}
\providecommand{\ZZ}{\mathbf Z}
\providecommand{\Ab}{\mathbf{Ab}}
\providecommand{\ideal}[1]{\mathfrak{#1}}
\providecommand{\aaa}{\ideal a}
\providecommand{\pp}{\ideal p}
\providecommand{\qq}{\ideal q}
\providecommand{\rr}{\ideal r}
\providecommand{\mm}{\ideal m}
\providecommand{\m}{\ideal m}
\providecommand{\mideal}{\mathfrak m}
\providecommand{\formelF}{\mathfrak F}
\providecommand{\formelG}{\mathfrak G}

% Sheaves, standard objects, and local aliases used by the English units.
\let\oldS\S
\providecommand{\sheaf}[1]{\underline{#1}}
\providecommand{\ffam}{\mathbf f}
\providecommand{\gfam}{\mathbf g}
\providecommand{\hfam}{\mathbf h}
\providecommand{\RR}{\mathrm R}
\providecommand{\uC}{\underline{\mathrm C}}
\providecommand{\Ccat}{\mathcal C}
\providecommand{\Cbar}{\overline{\mathcal C}}
\providecommand{\Ctors}{\mathcal C_{\mideal}}
\providecommand{\CatC}{\mathcal C}
\providecommand{\catSch}{\mathbf{Sch}}
\providecommand{\uGamma}{\mathop{\underline{\Gamma}}\nolimits}
\providecommand{\uHom}{\mathop{\underline{\mathrm{Hom}}}\nolimits}
\providecommand{\uExt}{\mathop{\underline{\mathrm{Ext}}}\nolimits}
\providecommand{\uH}{\mathop{\underline{\mathrm H}}\nolimits}
\providecommand{\UH}{\mathop{\underline{\mathrm H}}\nolimits}
\providecommand{\SheafGamma}{\underline{\Gamma}}
\providecommand{\SheafH}{\mathcal H}
\providecommand{\SheafHom}{\mathcal H\!om}
\providecommand{\SheafExt}{\mathcal E\!xt}
\providecommand{\SheafI}{\mathcal I}
\providecommand{\Sheafphi}{\underline{\varphi}}
\providecommand{\SheafPic}{\mathbf{Pic}}
\providecommand{\cH}{\mathcal H}
\providecommand{\cExt}{\mathcal E\!xt}
\providecommand{\cHom}{\mathcal H\!om}
\providecommand{\sH}{\mathcal H}
\providecommand{\sExt}{\mathcal E\!xt}
\providecommand{\cF}{\mathcal F}
\providecommand{\calF}{\mathcal F}
\providecommand{\Lcal}{\mathcal L}
\DeclareMathOperator{\Proj}{Proj}
\DeclareMathOperator{\SheafProj}{\mathbf{Proj}}
\DeclareMathOperator{\gr}{gr}
\DeclareMathOperator{\E}{E}
\DeclareMathOperator{\degtr}{deg\,tr}
\DeclareMathOperator{\Br}{Br}
\DeclareMathOperator{\cd}{cd}
\DeclareMathOperator{\D}{D}
\DeclareMathOperator{\geom}{g\acute{e}om}
\DeclareMathOperator{\hop}{hop}
\DeclareMathOperator{\Ker}{Ker}
\DeclareMathOperator{\tr}{tr}
\DeclareMathOperator{\profet}{prof\;\acute{e}t}
\DeclareMathOperator{\profgeom}{prof\;g\acute{e}om}
\DeclareMathOperator{\rg}{rg}
\providecommand{\h}{\mathop{\underline{H}}\nolimits}
\providecommand{\cG}{\mathcal G}
\providecommand{\cO}{\mathcal O}
\providecommand{\uF}{\underline{\mathrm F}}
\providecommand{\uT}{\underline{\mathrm T}}
\providecommand{\OX}{\mathcal O_X}
\providecommand{\OXx}{\mathcal O_{X,x}}
\providecommand{\OXY}{\mathcal O_{X,Y}}
\providecommand{\OXZ}[1]{\mathcal O_{X,#1}}
\providecommand{\OoXZ}{\mathcal O_{X,Z}}
\providecommand{\OY}[1]{\mathcal O_{Y_{#1}}}
\providecommand{\omx}{\Omega_{X/k}}
\providecommand{\Et}{\mathbf{Et}}
\providecommand{\Etc}{\mathbf{Et}}
\providecommand{\LLcat}{\mathbf L}
\providecommand{\Ib}{\mathcal I'}
\providecommand{\Jb}{\mathcal J}
\providecommand{\restr}[2]{#1\vert_{#2}}
\providecommand{\rest}[1]{_{|#1}}

% Operators.
\providecommand{\Ann}{\operatorname{Ann}}
\providecommand{\Annop}{\operatorname{Ann}}
\providecommand{\Ass}{\operatorname{Ass}}
\providecommand{\Assop}{\operatorname{Ass}}
\providecommand{\Supp}{\operatorname{Supp}}
\providecommand{\Suppop}{\operatorname{Supp}}
\providecommand{\Spec}{\operatorname{Spec}}
\providecommand{\Specop}{\operatorname{Spec}}
\providecommand{\Frac}{\operatorname{Frac}}
\providecommand{\Homop}{\operatorname{Hom}}
\providecommand{\Hom}{\operatorname{Hom}}
\providecommand{\Ob}{\operatorname{Ob}}
\providecommand{\Ext}{\operatorname{Ext}}
\providecommand{\Extop}{\operatorname{Ext}}
\providecommand{\Pic}{\operatorname{Pic}}
\providecommand{\Cl}{\operatorname{Cl}}
\providecommand{\prof}{\operatorname{prof}}
\providecommand{\codh}{\operatorname{codh}}
\providecommand{\codim}{\operatorname{codim}}
\providecommand{\codimn}{\operatorname{codim}}
\providecommand{\codimop}{\operatorname{codim}}
\providecommand{\depthop}{\operatorname{depth}}
\providecommand{\lengthop}{\operatorname{length}}
\providecommand{\pdop}{\operatorname{pd}}
\providecommand{\Hloc}{\operatorname H}
\providecommand{\HH}{\mathrm H}
\renewcommand{\H}{\mathrm H}
\providecommand{\R}{\mathrm R}
\providecommand{\DD}{\mathrm D}
\providecommand{\DA}{\mathrm D A}
\providecommand{\DhatA}{\mathrm D\widehat A}
\providecommand{\CA}{\mathrm C A}
\providecommand{\ChatA}{\mathrm C\widehat A}
\providecommand{\ellA}{\operatorname{length}_A}
\providecommand{\mA}{\mathfrak m_A}
\providecommand{\mB}{\mathfrak m_B}
\providecommand{\Lef}{\operatorname{Lef}}
\providecommand{\Leff}{\operatorname{Leff}}
\providecommand{\Lefc}{\operatorname{Lef}}
\providecommand{\Leffc}{\operatorname{Leff}}

% Arrows, equation controls, and display conveniences.
\providecommand{\from}{\longleftarrow}
\providecommand{\mto}{\longmapsto}
\providecommand{\To}{\Longrightarrow}
\providecommand{\lto}[1]{\xrightarrow{#1}}
\providecommand{\isomto}{\overset{\sim}{\longrightarrow}}
\providecommand{\isomfrom}{\overset{\sim}{\longleftarrow}}
\providecommand{\SSI}{\unskip~$\Longleftrightarrow$~\ignorespaces}
\providecommand{\ALORS}{\unskip~$\Longrightarrow$~\ignorespaces}
\providecommand{\tbigoplus}{\mathop{\textstyle\bigoplus}\limits}
\providecommand{\tbigotimes}{\mathop{\textstyle\bigotimes}\limits}
\providecommand{\boule}{\raisebox{0.2ex}{$\scriptscriptstyle\bullet$}}
\providecommand{\steco}{%
  \setcounter{equation}{\value{yequation}}\refstepcounter{yequation}}
\providecommand{\skipqed}{\renewcommand{\qedsymbol}{}}
\providecommand{\sourceoddity}[2]{%
  \par\smallskip\noindent
  \fcolorbox{gray}{gray!6}{\parbox{0.94\linewidth}{\footnotesize
  \textbf{Source note [#1].} #2}}\par\smallskip}
\providecommand{\danger}{\marginpar{\Large$\triangle$}}
\let\pageoriginaled\pageoriginale
% SMF's \skpt suppresses punctuation in a statement heading; with the
% standard-class compatibility environments the heading punctuation is fixed.
\providecommand{\skpt}{}
\providecommand{\noo}{\hbox{${}^\circ$}}
\providecommand{\et}{\textup{\'et}}
\providecommand{\Isom}{\textup{Isom}}
\providecommand{\llbracket}{[\![}
\providecommand{\rrbracket}{]\!]}
\definecolor{softgray}{gray}{0.94}
\definecolor{notegray}{gray}{0.25}
\definecolor{noteblue}{RGB}{232,241,250}
\definecolor{notebluetext}{RGB}{20,65,110}

% Human-readable statement environments used by the bounded English units.
\newenvironment{sgastatement}[1]
  {\par\addvspace{0.8\baselineskip}\noindent\textbf{#1.}\ \itshape}
  {\par\addvspace{0.35\baselineskip}}
\newenvironment{sgaremark}[1]
  {\par\addvspace{0.8\baselineskip}\noindent\textbf{#1.}\ }
  {\par\addvspace{0.35\baselineskip}}
\newenvironment{sgatheorem}[1]
  {\par\addvspace{0.8\baselineskip}\noindent\textbf{#1.}\ \itshape}
  {\par\addvspace{0.35\baselineskip}}
\newenvironment{sgalemma}[1]
  {\par\addvspace{0.8\baselineskip}\noindent\textbf{#1.}\ \itshape}
  {\par\addvspace{0.35\baselineskip}}
\newenvironment{sgaexample}[1]
  {\par\addvspace{0.8\baselineskip}\noindent\textbf{#1.}\ \itshape}
  {\par\addvspace{0.35\baselineskip}}
\newenvironment{sgaproof}
  {\par\addvspace{0.4\baselineskip}\noindent\textit{Proof.}\ }
  {\hfill$\square$\par\addvspace{0.4\baselineskip}}
\newenvironment{proofbridge}
  {\par\noindent\textit{Proof (reduction).}\ }
  {\par}
\newenvironment{enumeratea}
  {\begin{enumerate}[label=\alph*)]}
  {\end{enumerate}}
\newenvironment{enumeratei}
  {\begin{enumerate}[label=(\roman*)]}
  {\end{enumerate}}

% Source-like theorem names retained in the translated later exposés.
\theoremstyle{plain}
\newtheorem{theoreme}[subsection]{Theorem}
\newtheorem{theorem}[subsection]{Theorem}
\newtheorem{corollaire}[subsection]{Corollary}
\newtheorem{lemme}[subsection]{Lemma}
\newtheorem{lemma}{Lemma}[section]
\newtheorem{proposition}[subsection]{Proposition}
\newtheorem{propositionblah}[subsection]{Proposition}
\newtheorem{sublemme}[subsubsection]{Lemma}
\newtheorem{subsublemme}[subsubsection]{Lemma}
\newtheorem{conjecture}[subsection]{Conjecture}
\newtheorem{probleme}[subsection]{Problem}
\theoremstyle{definition}
\newtheorem{definition}[subsection]{Definition}
\newtheorem{exemple}[subsection]{Example}
\newtheorem{supdefinition}[subsection]{Definition}
\theoremstyle{remark}
\newtheorem{remarque}[subsection]{Remark}
\newtheorem{remarques}[subsection]{Remarks}
\newtheorem{subremarque}[subsubsection]{Remark}
\newtheorem*{remarquestar}{Remark}
\newtheorem*{theoremIXoneone}{Theorem 1.1}
\newtheorem*{sgaproposition}{Proposition}
\newtheorem*{sgacorollary}{Corollary}

\NewDocumentEnvironment{enonce*}{O{plain} m o}
  {\par\addvspace{0.8\baselineskip}\noindent\textbf{#2%
   \IfValueT{#3}{ (#3)}.}\ %
   \def\sgaenoncestyle{#1}\def\sgaremarkstyle{remark}%
   \ifx\sgaenoncestyle\sgaremarkstyle\normalfont\else\itshape\fi}
  {\par\addvspace{0.35\baselineskip}}
\renewcommand{\thesourcecorollary}{\thesection.\arabic{sourcecorollary}}
\NewDocumentEnvironment{sourcecorollary}{o}
  {\refstepcounter{sourcecorollary}\par\medskip\noindent
   \textbf{Corollary~\thesourcecorollary%
   \IfValueT{#1}{\markedfootnote{(3)}{#1}}.}\ \itshape}
  {\par\medskip}

% LaTeX kernels now define uppercase \Ref.  It must therefore be renewed,
% after hyperref is loaded, rather than introduced with \newcommand.
\renewcommand{\Ref}[1]{\textup{\ref{#1}}}
\renewcommand{\thechapter}{\Roman{chapter}}
\renewcommand{\chaptername}{Expos\'e}

\setlength{\emergencystretch}{2em}
\raggedbottom

% SGA2 reader: bounded LMMathExtension10 extraction support, proposal r1.
% Installed authority: pdfx/glyphtounicode-cmr.tex, whole-glyph base scalars.
% Encoding/TFM binding: lm-mathex.enc, pdftex.map, omxlmex.fd.
% This file is not a wholesale import of the pdfx mappings.
% Size variation selectors are deliberately omitted; size is typographic.
% Scope is only lmex10.tfm. No source formulas or glyph outlines are changed.
% Hats/tildes below are spacing accent identities, not base-order repair.
% No arrow tips, extenders, delimiter assemblies, mapsto pieces, or XY tips.
% Root must input this in the pdfLaTeX preamble and rebuild/recheck extraction.
\ifdefined\pdfglyphtounicode
  \pdfgentounicode=1
  % Whole, sized delimiters observed in the pinned reader.
  \pdfglyphtounicode{tfm:lmex10/parenleftbig}{0028}
  \pdfglyphtounicode{tfm:lmex10/parenrightbig}{0029}
  \pdfglyphtounicode{tfm:lmex10/braceleftbig}{007B}
  \pdfglyphtounicode{tfm:lmex10/bracerightbig}{007D}
  \pdfglyphtounicode{tfm:lmex10/parenleftBig}{0028}
  \pdfglyphtounicode{tfm:lmex10/parenrightBig}{0029}
  \pdfglyphtounicode{tfm:lmex10/parenleftbigg}{0028}
  \pdfglyphtounicode{tfm:lmex10/parenrightbigg}{0029}
  \pdfglyphtounicode{tfm:lmex10/parenleftBigg}{0028}
  \pdfglyphtounicode{tfm:lmex10/parenrightBigg}{0029}
  \pdfglyphtounicode{tfm:lmex10/braceleftBigg}{007B}
  \pdfglyphtounicode{tfm:lmex10/bracketleftBig}{005B}
  \pdfglyphtounicode{tfm:lmex10/bracketrightBig}{005D}
  % Whole n-ary operators; bigoplus is 2A01, not binary oplus 2295.
  \pdfglyphtounicode{tfm:lmex10/circleplustext}{2A01}
  \pdfglyphtounicode{tfm:lmex10/circleplusdisplay}{2A01}
  \pdfglyphtounicode{tfm:lmex10/circlemultiplytext}{2A02}
  \pdfglyphtounicode{tfm:lmex10/summationtext}{2211}
  \pdfglyphtounicode{tfm:lmex10/uniontext}{22C3}
  \pdfglyphtounicode{tfm:lmex10/intersectiontext}{22C2}
  \pdfglyphtounicode{tfm:lmex10/summationdisplay}{2211}
  \pdfglyphtounicode{tfm:lmex10/productdisplay}{220F}
  \pdfglyphtounicode{tfm:lmex10/uniondisplay}{22C3}
  % Spacing accents per pdfx; prevents b/c/e/f/g fallback identities.
  % Does not establish the logical order or scope of the accented expression.
  \pdfglyphtounicode{tfm:lmex10/hatwide}{02C6}
  \pdfglyphtounicode{tfm:lmex10/hatwider}{02C6}
  \pdfglyphtounicode{tfm:lmex10/tildewide}{02DC}
  \pdfglyphtounicode{tfm:lmex10/tildewider}{02DC}
  \pdfglyphtounicode{tfm:lmex10/tildewidest}{02DC}
\else
  \errmessage{SGA2 lmex Unicode support requires pdfTeX primitives}
\fi

% END INCLUDED active_source/sga2_reader_macros.tex


% Roman exposé numbers need wider number boxes than the standard book TOC.
\setlength{\cftchapnumwidth}{3.2em}
\setlength{\cftsecnumwidth}{4.2em}
\setlength{\cftsubsecnumwidth}{5.4em}

\hypersetup{
  pdftitle={SGA 2: Local Cohomology of Coherent Sheaves and Local and Global Lefschetz Theorems},
  pdfauthor={Alexander Grothendieck et al.; English translation project},
  pdfsubject={Source-aligned English reader assembled from the corrected arXiv TeX authority math/0511279v1},
  pdfkeywords={SGA 2, local cohomology, Lefschetz theorems, algebraic geometry, English translation}
}

\begin{document}
\frontmatter

\title{Seminar on Algebraic Geometry at Bois Marie\\
   1962\\[0.4em]
   Local Cohomology of Coherent Sheaves and\\
   Local and Global Lefschetz Theorems\\
   (SGA~2)}

\author{Alexander Grothendieck\\
  {\small (written up by a group of participants)}\\[0.5em]
  With an additional expos\'e by Mich\`ele Raynaud}

\maketitle

\chapter*{Abstract}
\addcontentsline{toc}{chapter}{Abstract}
This volume is a newly updated edition of the book
\emph{Cohomologie locale des faisceaux coh\'erents et th\'eor\`emes de
Lefschetz locaux et globaux (SGA~2)}, Advanced Studies in Pure Mathematics~2,
North-Holland Publishing Company, Amsterdam, 1968, by A.~Grothendieck
\emph{et al.}

The monograph gives necessary and sufficient conditions for the finiteness
of the local cohomology sheaves of coherent sheaves. These results yield
algebraization theorems which, together with purity results also proved in
the text, lead to Lefschetz-type theorems for both the fundamental group and
the Picard group.

\medskip
\noindent\textbf{Keywords.} Algebraization, Lefschetz theorem, local duality,
local cohomology, depth, homotopical depth, fundamental group, Picard group.

\smallskip
\noindent\textbf{2020 Mathematics Subject Classification.}
14-02, 14B10, 14B15, 14B20, 14C22, 14J70, 32S20.

\chapter*{Editor's Preface}
\addcontentsline{toc}{chapter}{Editor's Preface}

The present text is a new updated edition of the book
``Cohomologie locale des faisceaux coh\'erents et th\'eor\`emes de Lefschetz
locaux et globaux (SGA~2)'', Advanced Studies in Pure Mathematics~2,
North-Holland Publishing Company, Amsterdam, 1968, by A.~Grothendieck
\emph{et al.} It is the second part of the SGA project initiated by
B.~Edixhoven, who prepared a new edition of SGA~1. This version is intended
to reproduce the original text as closely as possible, while incorporating
minor typographical corrections and numbered editorial footnotes explaining
the current status of questions raised in the first edition. Further details
have also been supplied for some proofs. To avoid confusion, original notes
retain star-based numbering, whereas new editorial notes use Arabic numbers.
The page numbers of the original edition are shown in the margins.

I thank the mathematicians who carried out most of the initial typesetting
in \LaTeXe: L.~Bayle, N.~Borne, O.~Brinon, J.~Buresi, M.~Chardin,
F.~Ducrot, P.~Graftiaux, F.~Han, P.~Karwasz, L.~Koelblen, D.~Madore,
S.~Morel, D.~Naie, B.~Osserman, J.~Riou, and V.~S\'echerre, as well as
C.~Sabbah for adapting the text to the SMF format. I also thank J.-B.~Bost,
P.~Colmez, O.~Gabber, W.~Fulton, S.~Kleiman, F.~Orgogozo, M.~Raynaud, and
J.-P.~Serre for their comments and advice.

\begin{flushright}
Yves Laszlo, editor
\end{flushright}

\medskip
\noindent\textit{English-reader note.} This reader is translated from the
corrected and annotated TeX distributed as arXiv:math/0511279v1. Mathematical
notation, labels, equations, editorial notes, and source ordering are retained.
External English candidates were comparison witnesses only, not source
authority. The editable English Markdown lineage by Jacob Reinhold
(\url{https://github.com/jcreinhold/sga}, snapshot
\texttt{e7a259f3f8608ad3edf9bf6eead3fd504dd2d23e}) was consulted as one such
comparison witness and is credited under its stated CC BY 4.0 license for
the translation contribution; that license does not cover the underlying
French work.

\cleardoublepage
\tableofcontents

\mainmatter
% BEGIN INCLUDED active_source/parts/00_introduction.tex
\chapter*{Introduction} \label{I.0}
\markboth{INTRODUCTION}{INTRODUCTION}

We\pageoriginale present here, in revised and expanded form, a photo-offset
reissue of the second \sisi{ALGEBRAIC GEOMETRY SEMINAR of the INSTITUT
DES HAUTES \'ETUDES SCIENTIFIQUES}{Algebraic Geometry Seminar of the
Institut des Hautes \'Etudes Scientifiques}, held in 1962 and originally
issued in mimeographed form.

The reader is referred to the Introduction to the first of these seminars
(cited below as \SGA 1) for the aims pursued by these seminars and their
relation to the \sisi{\'EL\'EMENTS DE G\'{E}OM\'ETRIE
ALG\'EBRIQUE}{\'El\'ements de G\'eom\'etrie Alg\'ebrique}.

The text of Expos\'es~\Ref{I} through~\Ref{XI} was written up as the
seminar proceeded, from my oral lectures and handwritten notes, by a group
of participants comprising I\ptbl~\sisi{GIROGUITTI}{Giorgiutti},
J\ptbl~\sisi{GIRAUD}{Giraud}, Mlle M\ptbl~\sisi{JAFFE}{Jaffe} (now
Mme M\ptbl~\sisi{HAKIM}{Hakim}), and A\ptbl~\sisi{LAUDAL}{Laudal}.
These notes were originally regarded as provisional and intended for very
limited circulation, pending their absorption into the EGA---an absorption
that has now become problematic, to say the least, as it has for the other
parts of the SGA. As the warning to the original edition said, the
\og confidential\fg character of the notes was meant to excuse certain
\og stylistic weaknesses\fg, doubtless more apparent in the present
seminar, \SGA 2, than in the others. I have tried, as far as possible, to
remedy them in the present reissue by a fairly thorough revision of the
original text. In particular, I have harmonized the systems used to number
statements in the different Expos\'es, introducing everywhere the same
decimal system, already used in most of the original Expos\'es in the
present \SGA 2 as well as in all the other parts of the SGA. This led me,
in particular, to revise completely the numbering\nde{The original
numbering has been preserved as far as possible, with the adverb
\emph{bis} added when ambiguous duplicates occurred here and there.} of
the statements in Expos\'es~\Ref{III} through~\Ref{VIII}\pageoriginale
(and consequently the references to those Expos\'es)\sfootnote{It goes
without saying that every reference to \SGA 2 appearing in those parts of
the SGA published in the \sisi{SERIES IN PURE MATHEMATICS}{Series in Pure
Mathematics} refers to the present volume, not to the original edition of
\SGA 2!}. I have also tried to remove from the original text the principal
typing and syntactical errors, which were numerous and troublesome.
Moreover, Mme M\ptbl~\sisi{HAKIM}{Hakim} kindly undertook to rewrite
Expos\'e~\Ref{IV} in a less telegraphic style than the original. As in the
other SGA reissues, I have also added a number of footnotes, either to give
additional references or to indicate the status of questions on which
progress has been made since the original text was written. Finally, the
seminar has been enlarged by a new Expos\'e, \sisi{namely}{namely}
Expos\'e~\Ref{XIV}, written by Mme\sisi{.}{}
\sisi{MICH\`ELE RAYNAUD}{Mich\`ele Raynaud} in 1967, which takes up and
completes suggestions contained in the \og Comments on
Expos\'e~\Ref{XIII}\fg (\Ref{XIII}~\Ref{XIII.6}), written in March 1963.
This Expos\'e treats Lefschetz-type theorems from the viewpoint of \'etale
cohomology, using the results on \'etale cohomology presented in \SGA 4
and \SGA 5 (to appear in the same
\sisi{SERIES IN PURE MATHEMATICS}{Series in Pure Mathematics})\nde{In
fact, these seminars were published by Springer (numbers 269, 279, 305,
and 589), but unfortunately are out of print.}; it is therefore less
\og elementary\fg in nature than the other Expos\'es of the present
volume, which use little more than the substance of Chapters~I through~III
of the EGA. Here is an outline of the contents of the present volume.
Expos\'e~\Ref{I} contains the sorites of \og cohomology with support in
$Y$\fg, $H_Y^*(X,F)$, where $Y$ is a closed subset of a space $X$,
cohomology that may be interpreted as the cohomology of $X$ \emph{modulo}
the open subset $X-Y$, and that is the abutment of a very useful
\og \emph{local-to-global spectral sequence}\fg
\Ref{I}~\Ref{I.2.6}, involving the \emph{sheaves} of cohomology
\og with support in $Y$\fg, $\sheaf{H}_Y^*(F)$\nde{The underlined
notation has been retained for sheafified versions of functors, since the
calligraphic analogue of $\Gamma$ is not clear.}. In many questions this
formalism can\pageoriginale play a \og localization\fg role analogous to
that played by tubular neighborhoods of $Y$ in differential geometry.
Expos\'e~\Ref{II} studies the preceding notions for quasi-coherent sheaves
on preschemes, and Expos\'e~\Ref{III} relates them to the classical notion
of \emph{depth} (\Ref{III}~\Ref{III.3.3}).

Expos\'es~\Ref{IV} and~\Ref{V} give forms of \emph{local duality}, which
may be compared with Serre's projective duality theorem
(\Ref{XII}~\Ref{XII.1.1}). We note that both kinds of duality theorem are
substantially generalized in \sisi{HARTSHORNE}{Hartshorne}'s seminar
(cited in the footnote at the end of \Exp \Ref{IV})\refstepcounter{toto}
\de{\label{noteconrad}Hartshorne's book contains sign errors and, more
importantly, does not actually prove the compatibility of the trace with
base change. Conrad completely redid this work and proved this crucial and
highly nontrivial compatibility (\textit{Grothendieck duality and base
change}, Lect. Notes in Math., vol.~1750, Springer-Verlag, Berlin, 2000).
Unfortunately, errors remain (\cf two preprints: Conrad~B.,
\og Clarifications and corrections to ``\textit{Grothendieck duality and
base change}''\fg\ and \og An addendum to Chapter 5 of
``\textit{Grothendieck duality and base change}''\fg). For a more concrete
aspect, with particular attention to the notion of residue, see Lipman's
work, especially Lipman~J., \textit{Dualizing sheaves, differentials and
residues on algebraic varieties}, Ast\'erisque, vol.~117, Soci\'et\'e
math\'ematique de France, 1984. A categorical proof of the duality theorem,
based on Brown's representability theorem, was obtained by Neeman
(Neeman~A., \og The Grothendieck duality theorem via Bousfield's techniques
and Brown representability\fg, \emph{J.~Amer. Math. Soc.} \textbf{9}
(1996), \numero 1, p\ptbl~205--236).}

Expos\'es~\Ref{VI} and~\Ref{VII} give elementary technical notions used in
Expos\'e~\Ref{VIII} to prove the \emph{finiteness theorem}
(\Ref{VIII}~\Ref{VIII.2.3}). This theorem gives necessary and sufficient
conditions, for a coherent sheaf $F$ on a noetherian scheme $X$, for the
local cohomology sheaves $\sheaf{H}_Y^i(F)$ to be coherent for $i\leq n$
(or, equivalently, for the sheaves $\R^if_*(F|X-Y)$ to be coherent for
$i\leq n-1$, where $f\colon X-Y\to X$ is the inclusion). It is one of the
central technical results of the seminar, and in Expos\'e~\Ref{IX} we show
how a theorem of this kind can be used to establish a \og comparison
theorem\fg and an \og existence theorem\fg in formal geometry, following
and generalizing the use made in (\EGA III~\oldS\S 4 and~5) of the
finiteness theorem for a proper morphism.

These last results are applied in~\Ref{X} and~\Ref{XI}, devoted
respectively to Lefschetz-type theorems for the fundamental group and
for\pageoriginale the Picard group. These theorems compare, under certain
conditions, the invariants ($\pi_1$ or $\Pic$) associated respectively
with a scheme $X$ and a subscheme $Y$, which plays the role of a
hyperplane section, and in particular give conditions under which they
are isomorphic. Roughly speaking, the hypotheses allow one to pass from
$Y$ to the formal completion of $X$ along $Y$, and then to apply the
results of~\Ref{IX} in order to pass from there to an open
\emph{neighborhood} $U$ of $Y$ in $X$. To pass from $U$ to $X$, one also
needs information of a \og purity\fg or \og parafactoriality\fg type for
the local rings of $X$ at the points of $Z=X-U$, which is a finite
discrete set in the cases under consideration. This explains the
interaction, in the proofs of Expos\'es~\Ref{X}, \Ref{XI}, and \Ref{XII},
between local and global results, particularly in certain inductions. The
main results obtained in~\Ref{X} and~\Ref{XI} are the theorems of
\emph{local nature} \Ref{X}~\Ref{X.3.4} (the \emph{purity theorem}) and
\Ref{XI}~\Ref{XI.3.14} (the \emph{parafactoriality theorem}). Notice that
these theorems are proved by cohomological techniques of an essentially
global nature. In~\Ref{XII}, using the preceding local results, one
obtains global variants for projective schemes over a field, or more
generally over a more or less arbitrary base scheme; typical statements
include \Ref{XII}~\Ref{XII.3.5} and \Ref{XII}~\Ref{XII.3.7}.

In~\Ref{XIII}, we review some of the many problems and conjectures
suggested by the results and methods of the seminar. Perhaps the most
interesting concern cohomological and homotopical Lefschetz-type theorems
for complex analytic spaces; \cf \Ref{XIII}, pages 26 and following
\de{Essentially all the conjectures stated in~\Ref{XIII} and \Ref{XIV}
have now been proved; see the footnotes in those sections for references
and comments.}. In the context of the \'etale cohomology of
schemes\pageoriginale, the corresponding conjectures are proved
in~\Ref{XIV} by a duality technique that should apply equally in the
complex analytic case (\cf the comments in \Ref{XIII}, p\ptbl~25, and
\Ref{XIV}~\Ref{XIV.6.4}). But the corresponding homotopical statements
for analytic spaces, and more particularly those involving the
fundamental group, seem to require entirely new techniques
(\cf \Ref{XIV}~\Ref{XIV.6.4}).

I am happy to thank all those who, in various capacities, helped bring the
present volume to publication, including the collaborators already named
in this Introduction. In particular, I wish to thank Mlle
\sisi{CHARDON}{Chardon} for the good grace with which she carried out the
thankless task of physically preparing the final manuscript for
photo-offset reproduction.

\bigskip\hfill Bures-sur-Yvette, April~1968
\par\smallskip \hfill A\ptbl~Grothendieck.

% END INCLUDED active_source/parts/00_introduction.tex

% BEGIN INCLUDED active_source/parts/01_expose_I.tex
% SGA 2, Expose I — English mathematical body.
% Authority coverage: official corrected arXiv TeX lines 87--503.
% Mechanically extracted without alteration from the two source-aligned units
% named below. The repeated Expose heading and the Section 1 bibliography copy
% were removed; the complete bibliography from Section 2 is retained.

\makeatletter
\providecommand{\sgareaderlabel}[2]{%
  \begingroup
  \def\@currentlabel{#2}%
  \phantomsection\label{#1}%
  \endgroup}
\makeatother

\chapter[Global and Local Cohomological Invariants]
{Global and Local Cohomological Invariants Relative to a Closed Subspace}
\label{I}

\section*{1. The functors $\Gamma_Z$ and $\uGamma_Z$}
\sgareaderlabel{I.1}{1}
\addcontentsline{toc}{section}{1. The functors Gamma Z and underlined Gamma Z}

\srcpage{6}
Let $X$ be a topological space, and let $\uC_X$ be the category of
sheaves of abelian groups on $X$.  Let $\Phi$ be a family of supports in
the sense of Cartan.  The functor $\Gamma_\Phi$ on $\uC_X$ is defined by
\begin{equation}
  \Gamma_\Phi(F)
  =\{f\in\Gamma(F)\mid \operatorname{supp}(f)\in\Phi\}.
  \tag{1}\label{eq:I.1}
\end{equation}
If $Z$ is a closed subset of $X$, by an abuse of notation we denote by
$\Gamma_Z$ the functor $\Gamma_\Phi$, where $\Phi$ is the family of
closed subsets of $X$ contained in $Z$.\footnote{At this single occurrence
the French printed and TeX sources underline $\Gamma_Z$.  The defining
equation~\eqref{eq:I.2} and the systematic distinction used thereafter
identify the intended functor here as the global, non-underlined $\Gamma_Z$.
The source anomaly and this explicit resolution are recorded in the tranche
comparison ledger.}
Thus
\begin{equation}
  \Gamma_Z(F)
  =\{f\in\Gamma(F)\mid \operatorname{supp}(f)\subset Z\}.
  \tag{2}\label{eq:I.2}
\end{equation}

We wish to extend this definition to the case in which $Z$ is a
\emph{locally closed} subset of $X$, hence closed in a suitable open subset
$V$ of $X$.  In that case we set
\begin{equation}
  \Gamma_Z(F)=\Gamma_Z(\restr{F}{V}).
  \tag{3}\label{eq:I.3}
\end{equation}
We must verify that $\Gamma_Z(F)$ ``does not depend'' on the chosen open
set.  It is enough to show that, if $V'$ is open and
$V\supset V'\supset Z$, then the map
$\rho^V_{V'}:F(V)\to F(V')$ maps
$\Gamma_Z(\restr{F}{V})$ isomorphically onto
$\Gamma_Z(\restr{F}{V'})$.  But
\begin{equation}
  \Gamma_Z(\restr{F}{V})=\ker\rho^V_{V-Z}.
  \tag{4}\label{eq:I.4}
\end{equation}
Consequently, if $f\in\Gamma_Z(\restr{F}{V})$ and
$\rho^V_{V'}(f)=\rho^V_{V-Z}(f)=0$, then $f=0$, since
$(V',V-Z)$ is a covering of $V$.  Likewise, if
$f'\in\Gamma_Z(\restr{F}{V'})$, then $f'\in F(V')$ and
\srcpage{7}$0\in F(V-Z)$ define an element $f\in F(V)$ such that
$\rho^V_{V'}(f)=f'$ and $f\in\Gamma_Z(\restr{F}{V})$.
Thus $\rho^V_{V'}$ induces an isomorphism
\[
  \Gamma_Z(\restr{F}{V})\longrightarrow
  \Gamma_Z(\restr{F}{V'}).
\]

Note that every open subset $W$ of $Z$ is induced by an open subset $U$
of $X$ in which $W$ is closed.  It follows that
$W\mapsto\Gamma_W(F)$ defines a presheaf on $Z$; one verifies that it is
a sheaf, denoted by $i^!(F)$, where $i:Z\to X$ is the canonical
immersion.  We obtain
\begin{equation}
  \Gamma_Z(F)=\Gamma(i^!(F)).
  \tag{5}\label{eq:I.5}
\end{equation}
The sheaf $i^!(F)$ is a subsheaf of $i^*(F)$.  Indeed, the canonical
homomorphism
\[
  \Gamma(\restr{F}{U})=\Gamma(U,F)\longrightarrow
  \Gamma(U\cap Z,i^*(F))
\]
is injective on
$\Gamma_{U\cap Z}(\restr{F}{U})\subset\Gamma(\restr{F}{U})$.
In summary, we have the following result.

\begin{sgastatement}{Proposition 1.1}\sgareaderlabel{I.1.1}{1.1}
There exists a unique subsheaf $i^!(F)$ of $i^*(F)$ such that, for every
open subset $U$ of $X$ for which $U\cap Z$ is closed in $U$, the map
\[
  \Gamma(\restr{F}{U})=\Gamma(U,F)\longrightarrow
  \Gamma(U\cap Z,i^*(F))
\]
induces an isomorphism
\[
  \Gamma_{U\cap Z}(\restr{F}{U})\longrightarrow
  \Gamma(U\cap Z,i^!(F)).
\]
\end{sgastatement}

If $Z$ is open, then simply
\begin{equation}
  i^!(F)=i^*(F)=\restr{F}{Z},
  \qquad \Gamma_Z(F)=\Gamma(Z,F).
  \tag{6}\label{eq:I.6}
\end{equation}

Suppose again that $Z$ is arbitrary.  As $U$ ranges through the open
subsets of $X$, the assignment
\[
  U\longmapsto \Gamma_{U\cap Z}(\restr{F}{U})
  =\Gamma(U\cap Z,i^!(F))
\]
is a sheaf on $X$, which we denote by $\uGamma_Z(F)$.  More precisely,
the preceding formula, which expresses the fact that $i^!$ commutes with
restriction to open subsets, gives an isomorphism
\begin{equation}
  \uGamma_Z(F)=i_*(i^!(F)).
  \tag{7}\label{eq:I.7}
\end{equation}
By definition,\srcpage{8} for every open subset $U$ of $X$,
\begin{equation}
  \Gamma(U,\uGamma_Z(F))
  =\Gamma_{U\cap Z}(\restr{F}{U}).
  \tag{8}\label{eq:I.8}
\end{equation}

Here there is a characteristic difference between the case in which
$Z$ is closed and the case in which $Z$ is open.  In the first case,
formula~\eqref{eq:I.8} shows that $\uGamma_Z(F)$ may be regarded as a
subsheaf of $F$; hence there is a \emph{canonical immersion}
\begin{equation}
  \uGamma_Z(F)\lhook\joinrel\longrightarrow F.
  \tag{$8'$}\label{eq:I.8'}
\end{equation}
When $Z$ is open, on the other hand, equation~\eqref{eq:I.6} shows that
the right-hand side of equation~\eqref{eq:I.8} is
$\Gamma(U\cap Z,F)$ and therefore receives a homomorphism from
$\Gamma(U,F)$.  Thus there is a \emph{canonical homomorphism} in the
opposite direction:\footnote{The original reference was \hyperref[eq:I.8]{(8)}.}
\begin{equation}
  F\longrightarrow\uGamma_Z(F).
  \tag{$8''$}\label{eq:I.8''}
\end{equation}
This is none other than the canonical homomorphism
\[
  F\longrightarrow i_*i^*(F),
\]
in view of the isomorphism
\begin{equation}
  \uGamma_Z(F)\simeq i_*i^*(F)
  \tag{6 bis}\label{eq:I.6bis}
\end{equation}
deduced from equations~\eqref{eq:I.6} and~\eqref{eq:I.7}.

For variable $F$, the objects $\Gamma_Z(F)$, $\uGamma_Z(F)$, and
$i^!(F)$ may of course be regarded as functors in $F$, with values,
respectively, in the category of abelian groups, the category of sheaves
of abelian groups on $X$, and the category of sheaves of abelian groups
on $Z$.  It is sometimes convenient to interpret the functor
\[
  i^!: \uC_X\longrightarrow\uC_Z
\]
as the right adjoint of the familiar functor
\[
  i_!: \uC_Z\longrightarrow\uC_X
\]
defined by the following proposition.

\srcpage{9}
\begin{sgastatement}{Proposition 1.2}\sgareaderlabel{I.1.2}{1.2}
Let $G$ be a sheaf of abelian groups on $Z$.  There exists a unique
subsheaf of $i_*(G)$, denoted $i_!(G)$, such that for every open subset
$U$ of $X$ the (identity) isomorphism
\[
  \Gamma(U\cap Z,G)=\Gamma(U,i_*(G))
\]
induces an isomorphism
\[
  \Gamma_{\Phi_{U\cap Z,U}}(U\cap Z,G)=\Gamma(U,i_!(G)),
\]
where $\Phi_{U\cap Z,U}$ denotes the family of subsets of $U\cap Z$
that are closed in $U$.
\end{sgastatement}

The verification reduces to observing that the left-hand side is a
sheaf as $U$ varies: namely, for a section of $i_*(G)$ on $U$, viewed
as a section of $G$ on $U\cap Z$, the property of having support closed
\emph{in $U$} is \emph{local on $U$}.  The sheaf $i_!(G)$ just defined
is also known as the \emph{sheaf obtained from $G$ by extension by zero}
outside $Z$; see Godement~\cite{Godement}.  In particular, if $Z$ is
closed, then
\begin{equation}
  i_!(G)=i_*(G).
  \tag{9}\label{eq:I.9}
\end{equation}
In general, however, the canonical injection
$i_!(G)\to i_*(G)$ is not an isomorphism, as is already well known when
$Z$ is open.  Evidently $i_!(G)$ depends functorially on $G$ (and is in
fact an exact functor of $G$).  We therefore have:

\begin{sgastatement}{Proposition 1.3}\sgareaderlabel{I.1.3}{1.3}
There is an isomorphism of bifunctors in $G$ and $F$, where $G$ is a
sheaf of abelian groups on $Z$ and $F$ is a sheaf of abelian groups on
$X$:
\begin{equation}
  \operatorname{Hom}(i_!(G),F)
  =\operatorname{Hom}(G,i^!(F)).
  \tag{10}\label{eq:I.10}
\end{equation}
\end{sgastatement}

To define such an isomorphism is equivalent to defining functorial
homomorphisms
\[
  i_!i^!(F)\longrightarrow F,
  \qquad
  G\longrightarrow i^!i_!(G),
\]
satisfying the familiar compatibility conditions (see, for example,
Shih's expos\'e in the Cartan seminar on cohomological operations).

Recalling\srcpage{10} that $i_!$ is exact, and hence takes
monomorphisms to monomorphisms, we conclude:

\begin{sgastatement}{Corollary 1.4}\sgareaderlabel{I.1.4}{1.4}
If $F$ is injective, then $i^!(F)$ is injective; consequently
$\uGamma_Z(F)=i_*i^!(F)$ is also injective.
\end{sgastatement}

Replacing $X$ by a variable open subset $U$ of $X$, \hyperref[I.1.3]{Proposition~1.3}
also gives:

\begin{sgastatement}{Corollary 1.5}\sgareaderlabel{I.1.5}{1.5}
There is an isomorphism, functorial in $F$ and $G$,
\begin{equation}
  \uHom(i_!(G),F)=i_*\bigl(\uHom(G,i^!(F))\bigr).
  \tag{11}\label{eq:I.11}
\end{equation}
\end{sgastatement}

Taking $G$ to be the constant sheaf on $Z$ defined by $\ZZ$, denoted
$\ZZ_Z$, \hyperref[I.1.3]{Proposition~1.3} and \hyperref[I.1.5]{Corollary~1.5} specialize as follows.

\begin{sgastatement}{Corollary 1.6}\sgareaderlabel{I.1.6}{1.6}
There are isomorphisms, functorial in $F$,
\begin{equation}
  \begin{aligned}
    \Gamma_Z(F)&=\operatorname{Hom}(\ZZ_{Z,X},F),\\
    \uGamma_Z(F)&=\uHom(\ZZ_{Z,X},F),
  \end{aligned}
  \tag{12}\label{eq:I.12}
\end{equation}
where $\ZZ_{Z,X}=i_!(\ZZ_Z)$ is the sheaf of abelian groups on $X$
obtained from the constant sheaf $\ZZ_Z$ on $Z$ by extension by zero
outside $Z$.
\end{sgastatement}

\begin{sgaremark}{Remark 1.7}\sgareaderlabel{I.1.7}{1.7}
Suppose that $X$ is a ringed space, equip $Z$ with the sheaf of rings
$\mathcal O_Z=i^{-1}(\mathcal O_X)$, and now let $\uC_X$ and $\uC_Z$
denote the categories of modules on $X$ and $Z$, respectively.  The
preceding considerations extend word for word, with $F$ a module on $X$
and $G$ a module on $Z$, and with \hyperref[I.1.3]{statements 1.3}--\hyperref[I.1.6]{1.6} interpreted
accordingly.
\end{sgaremark}

To conclude these generalities, let us examine what happens when the
locally closed subset $Z$ is changed.  Let $Z'\subset Z$ be another
locally closed subset, and let
\[
  j:Z'\longrightarrow Z,
  \qquad i':Z'\longrightarrow X,
  \qquad i'=ij
\]
be the canonical inclusions.\srcpage{11}  There are functorial
isomorphisms
\begin{equation}
  (ij)^!=j^!i^!,
  \qquad
  (ij)_!=i_!j_!.
  \tag{13}\label{eq:I.13}
\end{equation}
The first isomorphism in~\eqref{eq:I.13} defines a functorial
isomorphism
\begin{equation}
  \Gamma_{Z'}(F)
  =\Gamma(Z',(ij)^!(F))
  \simeq\Gamma(Z',j^!(i^!(F)))
  =\Gamma_{Z'}(i^!(F)).
  \tag{14}\label{eq:I.14}
\end{equation}

Suppose now that $Z'$ is \emph{closed} in $Z$, and let
\[
  Z''=Z-Z'
\]
be its complement in $Z$.  This is open in $Z$ and hence locally closed
in $X$.  Applying the canonical injection~\eqref{eq:I.8'} to
$i^!(F)$ on $Z$ with the subset $Z'$, and using~\eqref{eq:I.14}, gives
an injective functorial canonical homomorphism
\begin{equation}
  \Gamma_{Z'}(F)\longrightarrow\Gamma_Z(F).
  \tag{15}\label{eq:I.15}
\end{equation}
If $Z'$ is replaced by $Z''$ in~\eqref{eq:I.14} and
equation~\eqref{eq:I.8''} is used, one obtains a functorial
canonical homomorphism
\begin{equation}
  \Gamma_Z(F)\longrightarrow\Gamma_{Z''}(F).
  \tag{$15'$}\label{eq:I.15'}
\end{equation}

\begin{sgastatement}{Proposition 1.8}\sgareaderlabel{I.1.8}{1.8}
Under the preceding hypotheses, the sequence of functorial
homomorphisms
\begin{equation}
  0\longrightarrow\Gamma_{Z'}(F)
  \longrightarrow\Gamma_Z(F)
  \longrightarrow\Gamma_{Z''}(F)
  \tag{16}\label{eq:I.16}
\end{equation}
is exact.  If $F$ is flasque, the sequence remains exact after a zero is
added on the right.
\end{sgastatement}

\begin{sgaproof}
Replace $X$ by an open subset $V$ in which $Z$ is closed.  This reduces
the assertion to the case in which $Z$, and hence $Z'$, is closed.
Then $Z''$ is closed in the open subset $X-Z'$, and there is a canonical
injection
\[
  \Gamma_{Z''}(F)\longrightarrow\Gamma(X-Z',F).
\]
The exactness of~\eqref{eq:I.16} then says simply that the sections of
$F$ with support in $Z'$ are precisely those whose restriction to
$X-Z'$ is zero.

\srcpage{12}If $F$ is flasque, every element of $\Gamma_{Z''}(F)$, regarded as a
section of $F$ on $X-Z'$, extends to a section of $F$ on $X$.  The
latter section plainly has support in $Z$, proving that the last
homomorphism in~\eqref{eq:I.16} is then surjective.
\end{sgaproof}

\begin{sgastatement}{Corollary 1.9}\sgareaderlabel{I.1.9}{1.9}
There is a functorial exact sequence
\begin{equation}
  0\longrightarrow\uGamma_{Z'}(F)
  \longrightarrow\uGamma_Z(F)
  \longrightarrow\uGamma_{Z''}(F),
  \tag{16 bis}\label{eq:I.16bis}
\end{equation}
and, if $F$ is flasque, the sequence remains exact after a zero is added
on the right.
\end{sgastatement}

\hyperref[I.1.8]{Proposition~1.8} may be interpreted, by means of \hyperref[I.1.6]{Corollary~1.6}, as a
statement about the functors $\operatorname{Hom}$ and $\uHom$, as
follows.  First note that if $G$ is a sheaf of abelian groups on $Z$,
inducing the sheaves $j^*(G)$ and $k^*(G)$ on $Z'$ and $Z''$,
respectively, where $j:Z'\to Z$ and $k:Z''\to Z$ are the canonical
injections, then there is a canonical exact sequence of sheaves on $X$:
\begin{equation}
  0\longrightarrow k^*(G)_X
  \longrightarrow G_X
  \longrightarrow j^*(G)_X
  \longrightarrow 0.
  \tag{17}\label{eq:I.17}
\end{equation}
Here, to simplify notation, the subscript $X$ denotes the sheaf on $X$
obtained by extension by zero outside the space on which the sheaf in
question is defined.  Sequence~\eqref{eq:I.17} generalizes a familiar
exact sequence for $Z=X$ (see Godement~\cite{Godement}); it follows from
that sequence by writing it on $Z$ and applying the functor $i_!$.
Taking $G=\ZZ_Z$, we obtain in particular:

\begin{sgastatement}{Proposition 1.10}\sgareaderlabel{I.1.10}{1.10}
Under the preceding hypotheses, there is an exact sequence of sheaves of
abelian groups on $X$:
\begin{equation}
  0\longrightarrow\ZZ_{Z'',X}
  \longrightarrow\ZZ_{Z,X}
  \longrightarrow\ZZ_{Z',X}
  \longrightarrow 0.
  \tag{18}\label{eq:I.18}
\end{equation}
The two exact sequences in \hyperref[I.1.8]{Proposition~1.8} and \hyperref[I.1.9]{Corollary~1.9} are exactly
the sequences obtained from~\eqref{eq:I.18} by applying, respectively,
the functors $\operatorname{Hom}(-,F)$ and $\uHom(-,F)$.
\end{sgastatement}

This also gives an immediate proof that sequences~\eqref{eq:I.16} and
\eqref{eq:I.16bis} remain exact after a zero is added on the right,
provided that $F$ is \emph{injective}.
% CONTINUATION: source-aligned Section 2 unit, French lines 280--503.
\section*{2. The functors $\mathrm H_Z^*(X,F)$ and $\uH_Z^*(F)$}
\sgareaderlabel{I.2}{2}
\addcontentsline{toc}{section}{2. The functors H Z star and underlined H Z star}

\srcpage{13}
\begin{sgastatement}{Definition 2.1}\sgareaderlabel{I.2.1}{2.1}
The derived functors, in the variable $F$, of the functors $\Gamma_Z(F)$
and $\uGamma_Z(F)$ are denoted by $\mathrm H_Z^*(X,F)$ and $\uH_Z^*(F)$,
respectively.
\end{sgastatement}

They are cohomological functors with values, respectively, in the category
of abelian groups and in the category of sheaves of abelian groups on $X$.
When $Z$ is closed, $\mathrm H_Z^*(X,F)$ is, by definition, simply
$\mathrm H_\Phi^*(X,F)$, where $\Phi$ denotes the family of closed subsets
of $X$ contained in $Z$. When $Z$ is open, we shall see that
$\mathrm H_Z^*(X,F)$ is simply
$\mathrm H^*(Z,F)=\mathrm H^*(Z,\restr{F}{Z})$, by the following
proposition.

\begin{sgastatement}{\hyperref[I.2.2]{Proposition 2.2} (Excision theorem)}
\sgareaderlabel{I.2.2}{2.2}
Let $V$ be an open subset of $X$ containing $Z$. There is an isomorphism of
cohomological functors in $F$:
\begin{equation}
  \mathrm H_Z^*(X,F)
  \longrightarrow
  \mathrm H_Z^*(V,\restr{F}{V}).
  \tag{19}\label{eq:I.19}
\end{equation}
\end{sgastatement}

Indeed, there is a functorial isomorphism
$\Gamma_Z^X\simeq\Gamma_Z^Vj^!$, where $j\colon V\to X$ is the inclusion
and $j^!$ is therefore the restriction functor (cf. \hyperref[eq:I.14]{equation~(14)}). This
last functor is exact and, by \hyperref[I.1.4]{Corollary~1.4}, takes injectives to injectives;
the isomorphism~\eqref{eq:I.19} follows immediately.

When $Z$ is open, we may take $V=Z$, and obtain:

\begin{sgastatement}{Corollary 2.3}\sgareaderlabel{I.2.3bof}{2.3}
Suppose that $Z$ is open. There is an isomorphism of cohomological functors
\begin{equation}
  \mathrm H_Z^*(X,F)=\mathrm H^*(Z,F).
  \tag{20}\label{eq:I.20}
\end{equation}
\end{sgastatement}

The isomorphisms of \hyperref[I.1.6]{Corollary~1.6} and the definitions give the following
(cf.\ \cite{Tohoku}).

\begin{sgastatement}{Proposition 2.3 bis}\sgareaderlabel{I.2.3bis}{2.3 bis}
There are\srcpage{14} isomorphisms of cohomological functors
\begin{align}
  \mathrm H_Z^*(X,F)
    &\simeq \operatorname{Ext}^*(X;\ZZ_{Z,X},F),
    \tag{21}\label{eq:I.21}\\
  \uH_Z^*(F)
    &\simeq \uExt^*(\ZZ_{Z,X},F).
    \tag{21 bis}\label{eq:I.21bis}
\end{align}
\end{sgastatement}

\noindent\textit{Editorial note.} This proposition is numbered 2.3 in the
original edition; the corrected branch designates it 2.3 bis.

We may therefore apply the results of \cite{Tohoku} on the $\operatorname{Ext}$
of modules. First, we record the following interpretation of the sheaves
$\uH_Z^*(F)$ in terms of the global groups $\mathrm H_Z^*(X,F)$.

\begin{sgastatement}{Corollary 2.4}\sgareaderlabel{I.2.4}{2.4}
$\uH_Z^*(F)$ is canonically isomorphic to the sheaf associated with the
presheaf
\[
  U\longmapsto
  \mathrm H_{Z\cap U}^*(U,\restr{F}{U}).
\]
\end{sgastatement}

In particular, using \hyperref[I.2.3bof]{Corollary~2.3}, we obtain:

\begin{sgastatement}{Corollary 2.5}\sgareaderlabel{I.2.5}{2.5}
Suppose that $Z$ is open. There is an isomorphism of cohomological functors
\begin{equation}
  \uH_Z^*(F)=\mathrm R^*i_*i^*(F),
  \tag{22}\label{eq:I.22}
\end{equation}
where $i\colon Z\to X$ is the inclusion.
\end{sgastatement}

The $\operatorname{Ext}$ spectral sequence gives the following important
spectral sequence.

\begin{sgastatement}{\hyperref[I.2.6]{Theorem 2.6}}
\sgareaderlabel{I.2.6}{2.6}
There is a spectral sequence, functorial in $F$, abutting to
$\mathrm H_Z^*(X,F)$ and having initial term
\begin{equation}
  E_2^{p,q}(F)=\mathrm H^p(X,\uH_Z^q(F)).
  \tag{23}\label{eq:I.23}
\end{equation}
\end{sgastatement}

\begin{sgaremark}{Remarks 2.7}\sgareaderlabel{I.2.7}{2.7}
It follows immediately from \hyperref[I.2.4]{Corollary~2.4} that the sheaves $\uH_Z^q(F)$
vanish on $X-\overline Z$, and also vanish in the interior of $Z$ when
$q\ne0$; hence, for such $q$, $\uH_Z^q(F)$ is supported on the boundary of
$Z$.

Consequently, the right-hand side of~\eqref{eq:I.23} may be interpreted as a
cohomology group on $\overline Z$. We shall use \hyperref[I.2.6]{Theorem~2.6} when $Z$ is
closed\srcpage{15} in $X$; in that case the right-hand side of
\eqref{eq:I.23} may be interpreted as a cohomology group computed on $Z$:
\begin{equation}
  E_2^{p,q}(F)=\mathrm H^p(Z,\uH_Z^q(F)).
  \tag{23 bis}\label{eq:I.23bis}
\end{equation}
The corresponding equation was numbered \hyperref[eq:I.23]{(23)} in the original edition.
\end{sgaremark}

When $Z$ is open, the spectral sequence of \hyperref[I.2.6]{Theorem~2.6} is simply the Leray
spectral sequence for the continuous map $i\colon Z\to X$, in view of the
interpretation in \hyperref[I.2.5]{Corollary~2.5} used to compute its initial term.

We return to the exact sequence~\hyperref[eq:I.18]{(18)}. It gives rise to an exact
$\operatorname{Ext}$ sequence (cf.\ \cite{Tohoku}).

\noindent\textit{Editorial note.} The original edition cited
Proposition~1.10 at this point; the corrected branch cites equation~(18).

\begin{sgastatement}{Theorem 2.8}\sgareaderlabel{I.2.8}{2.8}
Let $Z$ be a locally closed subset of $X$, let $Z'$ be a closed subset of
$Z$, and put $Z''=Z-Z'$. There is a long exact sequence, functorial in $F$:
\begin{equation}
\begin{gathered}
  0\longrightarrow \mathrm H_{Z'}^0(X,F)
  \longrightarrow \mathrm H_Z^0(X,F)
  \longrightarrow \mathrm H_{Z''}^0(X,F)
  \xrightarrow{\partial} \mathrm H_{Z'}^1(X,F)
  \longrightarrow \mathrm H_Z^1(X,F)\longrightarrow\cdots \\
  \cdots\longrightarrow \mathrm H_{Z'}^i(X,F)
  \longrightarrow \mathrm H_Z^i(X,F)
  \longrightarrow \mathrm H_{Z''}^i(X,F)
  \xrightarrow{\partial} \mathrm H_{Z'}^{i+1}(X,F)
  \longrightarrow\cdots .
\end{gathered}
\tag{24}\label{eq:I.24}
\end{equation}
\end{sgastatement}

Recall how this exact sequence is obtained. Let $C(F)$ be an injective
resolution of $F$. Then the exact sequence~\hyperref[eq:I.18]{(18)}\footnote{See the preceding
editorial note.} gives the exact sequence
\begin{equation}
  0\longrightarrow\Gamma_{Z'}(C(F))
  \longrightarrow\Gamma_Z(C(F))
  \longrightarrow\Gamma_{Z''}(C(F))
  \longrightarrow0.
  \tag{25}\label{eq:I.25}
\end{equation}
This is precisely the sequence defined in \hyperref[I.1.8]{Proposition~1.8}. Its cohomology
long exact sequence is~\eqref{eq:I.24}.

\noindent\textit{Source note.} The corrected French TeX prints
$\Gamma(C(F))$ in the middle term of~(25). Equation~(24), the construction
from (18), and Proposition~1.8 require $\Gamma_Z(C(F))$; the support
subscript is restored here and the emendation is not silent.

The most important case for us is the one in which $Z$ is closed (and we may
always reduce to it by replacing $X$ with an open subset $V$ in which $Z$ is
closed). Then $Z'$ is closed, $Z''$ is closed in the open subset $X-Z'$, and
we may write
\begin{equation}
  \mathrm H_{Z''}^i(X,F)
  =\mathrm H_{Z''}^i(X-Z',\restr{F}{X-Z'}).
  \tag{26}\label{eq:I.26}
\end{equation}
This permits us to write the exact sequence~\eqref{eq:I.24} in terms of
cohomology with support\srcpage{16} in a given closed subset. The most common
case is $Z=X$. Writing $Z'=A$ for simplicity, we obtain:

\begin{sgastatement}{Corollary 2.9}\sgareaderlabel{I.2.9}{2.9}
Let $A$ be a closed subset of $X$. There is a long exact sequence,
functorial in $F$:
\begin{equation}
\begin{gathered}
  0\longrightarrow \mathrm H_A^0(X,F)
  \longrightarrow \mathrm H^0(X,F)
  \longrightarrow \mathrm H^0(X-A,F)
  \xrightarrow{\partial} \mathrm H_A^1(X,F)\longrightarrow\cdots \\
  \cdots\longrightarrow \mathrm H_A^i(X,F)
  \longrightarrow \mathrm H^i(X,F)
  \longrightarrow \mathrm H^i(X-A,F)
  \xrightarrow{\partial} \mathrm H_A^{i+1}(X,F)
  \longrightarrow\cdots .
\end{gathered}
\tag{27}\label{eq:I.27}
\end{equation}
\end{sgastatement}

This exact sequence shows that $\mathrm H_A^i(X,F)$ plays the role of a
relative cohomology group of $X$ modulo $X-A$, with coefficients in $F$. It
is in this capacity that it arises naturally in applications. By
``sheafifying''~\eqref{eq:I.24} and~\eqref{eq:I.27}, or by arguing directly,
and taking into account that the sheaf associated with
$U\mapsto\mathrm H^i(U,F)$ is zero for $i>0$, we obtain:

\begin{sgastatement}{Corollary 2.10}\sgareaderlabel{I.2.10}{2.10}
Under the hypotheses of \hyperref[I.2.8]{Theorem~2.8} there is a long exact sequence,
functorial in $F$:
\begin{equation}
  \cdots\longrightarrow \uH_{Z'}^i(F)
  \longrightarrow \uH_Z^i(F)
  \longrightarrow \uH_{Z''}^i(F)
  \xrightarrow{\partial}\uH_{Z'}^{i+1}(F)
  \longrightarrow\cdots .
  \tag{24 bis}\label{eq:I.24bis}
\end{equation}
\end{sgastatement}

\begin{sgastatement}{Corollary 2.11}\sgareaderlabel{I.2.11}{2.11}
Let $A$ be a closed subset of $X$. There is an exact sequence, functorial in
$F$:
\begin{equation}
  0\longrightarrow\uH_A^0(F)
  \longrightarrow F
  \longrightarrow f_*(\restr{F}{X-A})
  \xrightarrow{\partial}\uH_A^1(F)
  \longrightarrow0,
  \tag{28}\label{eq:I.28}
\end{equation}
and, for $i\ge2$, canonical isomorphisms
\begin{equation}
  \uH_A^i(F)
  =\uH_{X-A}^{i-1}(F)
  =\mathrm R^{i-1}f_*(\restr{F}{X-A}),
  \tag{29}\label{eq:I.29}
\end{equation}
where $f\colon X-A\to X$ is the inclusion.
\end{sgastatement}

Thus $\uH_A^0(F)$ and $\uH_A^1(F)$ are defined, respectively, as the kernel
and cokernel of the canonical homomorphism
\[
  F\longrightarrow f_*f^*(F)=f_*(\restr{F}{X-A}),
\]
and $\uH_A^i(F)$ for $i\ge2$ is expressed in terms of the derived functors
of $f_*$.

\noindent\textit{Source note.} The French display~(28) omits the argument
$(F)$ from $\uH_A^1$. The explanatory sentence immediately following the
display supplies it; it is restored here and recorded as an emendation.

\begin{sgastatement}{Corollary 2.12}\sgareaderlabel{I.2.12}{2.12}
Let $F$\srcpage{17} be a sheaf of abelian groups on $X$. If $F$ is flasque,
then for every locally closed subset $Z$ of $X$ and every integer $i\ne0$,
\[
  \mathrm H_Z^i(X,F)=0,
  \qquad
  \uH_Z^i(F)=0.
\]
Conversely, if $\mathrm H_Z^1(X,F)=0$ for every closed subset $Z$ of $X$,
then $F$ is flasque.
\end{sgastatement}

Suppose that $F$ is flasque. It induces a flasque sheaf on every open
subset. To prove that $\mathrm H_Z^i(X,F)=0$ for $i>0$, we may therefore
assume that $Z$ is closed; the assertion then follows from the exact
sequence~\eqref{eq:I.27}. For every locally closed $Z$, ``sheafifying''---that
is, applying \hyperref[I.2.4]{Corollary~2.4}---gives $\uH_Z^i(F)=0$ for $i>0$. Conversely,
suppose that $\mathrm H_Z^1(X,F)=0$ for every closed $Z$. Then
\eqref{eq:I.27} shows that
$\mathrm H^0(X,F)\to\mathrm H^0(X-Z,F)$ is surjective for every such $Z$,
which means precisely that $F$ is flasque.

\noindent\textit{Editorial note.} At both corresponding occurrences in the
French proof, the original edition cited Corollary~2.9; the corrected branch
specifies equation~(27).

Combining \hyperref[I.2.6]{Theorems~2.6} and~\hyperref[I.2.8]{2.8}, we deduce:

\begin{sgastatement}{Proposition 2.13}\sgareaderlabel{I.2.13}{2.13}
Let $F$ be a sheaf of abelian groups on $X$, let $Z$ be a closed subset of
$X$, put $U=X-Z$, and let $N$ be an integer. The following conditions are
equivalent:
\begin{enumerate}[label=\textup{(\roman*)},leftmargin=2.7em]
  \item $\uH_Z^i(F)=0$ for $i\le N$.
  \item For every open subset $V$ of $X$, the canonical homomorphism
  \[
    \mathrm H^i(V,F)\longrightarrow\mathrm H^i(V\cap U,F)
  \]
  is
  \begin{enumerate}[label=\textup{\alph*)},leftmargin=2.5em]
    \item bijective for $i<N$;
    \item injective for $i=N$.
  \end{enumerate}
\end{enumerate}
If $N>0$, it is enough in \textup{(ii)} to require \textup{a)}.
\end{sgastatement}

To prove \textup{(i)} $\Rightarrow$ \textup{(ii)}, the local nature of the
sheaves $\uH_Z^i(F)$ reduces us to proving the following.

\begin{sgastatement}{Corollary 2.14}\sgareaderlabel{I.2.14}{2.14}
If condition \textup{(i)} of \hyperref[I.2.13]{Proposition~2.13} holds, then
\[
  \mathrm H^i(X,F)\longrightarrow\mathrm H^i(U,F)
\]
is bijective for $i<N$ and injective for $i=N$.
\end{sgastatement}

Indeed, by the exact sequence~\eqref{eq:I.27}, this is equivalent to
$\mathrm H_Z^i(X,F)=0$ for $i\le N$, and this\srcpage{18} relation is an
immediate consequence of the spectral sequence~\eqref{eq:I.23bis}.

\noindent\textit{Editorial note.} At both corresponding occurrences in the
French text, the original edition cited Theorem~2.6; the corrected branch
specifies equation~(23 bis).

Conversely, condition \textup{(ii)} of \hyperref[I.2.13]{Proposition~2.13} says that, for every
open subset $V$ of $X$,
\[
  \mathrm H_{Z\cap V}^i(V,\restr{F}{V})=0
  \qquad\text{for }i\le N,
\]
which implies \textup{(i)} by \hyperref[I.2.4]{Corollary~2.4}. Moreover, if $N>0$, condition
\textup{b)} is superfluous. Indeed, if $N=1$, condition \textup{a)} together
with~\eqref{eq:I.28} gives $\uH_Z^i(F)=0$ for $i\le N$. If $N>1$,
condition \textup{a)} for $i=N-1>0$, together with~\eqref{eq:I.29}, gives
$\uH_Z^i(F)=0$ for $i\le N$.

Taking \hyperref[I.2.11]{Corollary~2.11} into account, this also proves condition \textup{(i)}
of \hyperref[I.2.13]{Proposition~2.13}.

\begin{sgaremark}{Remark}
Let $Y\to X$ be a closed immersion, and suppose that locally it has the form
$\{0\}\times Y\subset R^n\times Y$. Suppose that $F$ is a locally constant
sheaf on $X$. Then
\begin{equation}
  \uH_Y^i(F)\simeq
  \begin{cases}
    0,& i\ne n,\\
    F\otimes\uT_{Y,X},& i=n,
  \end{cases}
  \qquad
  \uT_{Y,X}\simeq\uH_Y^n(\ZZ_X),
  \tag{30}\label{eq:I.30}
\end{equation}
where $\uT_{Y,X}$ is the extension to $X$ of a sheaf on $Y$ locally
isomorphic to $\ZZ_Y$; it is called the \emph{normal orientation sheaf} of
$Y$ in $X$.
\end{sgaremark}

\noindent\textit{Source note.} The corrected French prose prints lowercase
$x$ where the ambient space must be $X$, and its sentence defining
$\uT_{Y,X}$ is grammatically incomplete. The reading above preserves the
display and states its evident mathematical syntax explicitly.

Using the spectral sequence~\eqref{eq:I.23bis}, we obtain in this case
\begin{equation}
  \mathrm H_Y^i(X,F)
  \simeq
  \mathrm H^{i-n}(Y,F\otimes\uT_{Y,X}),
  \tag{31}\label{eq:I.31}
\end{equation}
and recover the \emph{Gysin homomorphism}\sgareaderlabel{hgysin}{13}
\begin{equation}
  \mathrm H^j(Y,F\otimes\uT_{Y,X})
  \longrightarrow
  \mathrm H^{j+n}(X,F).
  \tag{32}\label{eq:I.32}
\end{equation}

\begin{thebibliography}{Godement}
\bibitem[Godement]{Godement}
R.~Godement, \emph{Topologie alg\'ebrique et th\'eorie des faisceaux},
Act. Scient. et Ind., vol.~1252, Hermann, Paris.

\bibitem[T\^ohoku]{Tohoku}
A.~Grothendieck, ``Sur quelques points d'alg\`ebre homologique,''
\emph{T\^ohoku Math. J.} \textbf{9} (1957), 119--221.
\end{thebibliography}

% END INCLUDED active_source/parts/01_expose_I.tex

% BEGIN INCLUDED active_source/parts/02_expose_II.tex
% SGA 2, Expose II — English mathematical body.
% Authority coverage: official corrected arXiv TeX lines 505--708.
% Mechanically extracted without alteration from the source-aligned unit.

\makeatletter
\providecommand{\sgareaderlabel}[2]{%
  \begingroup
  \def\@currentlabel{#2}%
  \phantomsection\label{#1}%
  \endgroup}
\makeatother

\chapter[Quasi-coherent Sheaves on Preschemes]
{Application to Quasi-coherent Sheaves on Preschemes}
\label{II}

\begin{sgastatement}{Proposition 1}\sgareaderlabel{II.1}{1}
Let\srcpage{19} $X$ be a prescheme and let $Z$ be a locally closed subset of
the form $Z=U-V$, where $U$ and $V$ are open subsets of $X$ such that
$V\subset U$ and the canonical immersions $U\to X$ and $V\to X$ are
quasi-compact. Then, for every quasi-coherent $\OX$-module $F$, the sheaves
$\uH_Z^i(F)$ are quasi-coherent.
\end{sgastatement}

By Expos\'e~I, \hyperref[eq:I.24]{equation~(24)}, there is an exact sequence of relative
cohomology sheaves
\[
  \longrightarrow \uH_U^{i-1}(F)
  \longrightarrow \uH_V^i(F)
  \longrightarrow \uH_Z^i(F)
  \longrightarrow \uH_U^i(F)
  \longrightarrow \uH_V^{i+1}(F)
  \longrightarrow .
\]
By EGA~III, 1.4.17, for the sheaves $\uH_Z^i(F)$ to be quasi-coherent it
therefore suffices that the sheaves $\uH_U^i(F)$ and $\uH_V^i(F)$ be
quasi-coherent. We may thus suppose that $Z$ is open and that the canonical
immersion $j\colon Z\to X$ is quasi-compact.

Since $Z$ is open, Expos\'e~I, \hyperref[eq:I.22]{equation~(22)}, gives a canonical isomorphism
\[
  \uH_Z^i(F)\simeq \mathrm R^i j_*(F|_Z).
\]
But $j$ is separated (EGA~I, 5.5.1) and quasi-compact; hence, by EGA~III,
1.4.10, the sheaves
$\mathrm R^i j_*(F|_Z)=\uH_Z^i(F)$ are quasi-coherent. This completes the
proof.

\begin{sgastatement}{Corollary 2}\sgareaderlabel{II.2}{2}
Let $Z$ be a closed subset of $X$ such that the canonical immersion
$X-Z\to X$ is quasi-compact. Then the $\OX$-modules $\uH_Z^i(F)$ are
quasi-coherent.
\end{sgastatement}

\begin{sgastatement}{Corollary 3}\sgareaderlabel{II.3}{3}
If $X$ is locally noetherian, then, for every locally closed subset $Z$ of
$X$ and every quasi-coherent $\OX$-module $F$, the sheaves $\uH_Z^i(F)$ are
quasi-coherent.
\end{sgastatement}

This follows immediately from \hyperref[II.2]{Corollary~2} and EGA~I, 6.6.4.

\begin{sgastatement}{Corollary 4}\sgareaderlabel{II.4}{4}
Suppose\srcpage{20} that $X$ is the spectrum of a ring $A$. Let $U$ be a
quasi-compact open subset of $X$, put $Y=X-U$, and let $F$ be a
quasi-coherent $\OX$-module. There is an isomorphism of cohomological
functors in $F$:
\begin{equation}
  \uH_Y^i(F)=\widetilde{\bigl(\mathrm H_Y^i(X,F)\bigr)}.
  \tag{4.1}\label{eq:II.4.1}
\end{equation}
There is moreover an exact sequence, functorial in $F$,
\begin{equation}
  0\longrightarrow \mathrm H_Y^0(X,F)
  \longrightarrow \mathrm H^0(X,F)
  \longrightarrow \mathrm H^0(U,F)
  \longrightarrow \mathrm H_Y^1(X,F)
  \longrightarrow 0,
  \tag{4.2}\label{eq:II.4.2}
\end{equation}
and there are isomorphisms, functorial in $F$,
\begin{equation}
  \mathrm H_Y^i(X,F)\simeq \mathrm H^{i-1}(U,F),
  \qquad i\geq 2.
  \tag{4.3}\label{eq:II.4.3}
\end{equation}
\end{sgastatement}

By \hyperref[II.2]{Corollary~2}, the sheaves $\uH_Y^i(F)$ are quasi-coherent. Since $X$ is
affine, it follows that
$\mathrm H^p(X,\uH_Y^i(F))=0$ for $p>0$. The spectral sequence of
Expos\'e~I, \hyperref[eq:I.23]{equation~(23)}, therefore degenerates, and
\[
  \mathrm H_Y^i(X,F)=\Gamma(\uH_Y^i(F)).
\]
Equation~\eqref{eq:II.4.1} now follows from EGA~I, 1.1.3.7; equations
\eqref{eq:II.4.2} and \eqref{eq:II.4.3} follow from the cohomology exact
sequence of Expos\'e~I, \hyperref[eq:I.27]{equation~(27)}, and from the vanishing
$\mathrm H^i(X,F)=0$ for $i>0$, since $X$ is affine.

\noindent\textit{Editorial note.} The corrected branch replaces the
original citation of Proposition~1 by Corollary~2 in the first sentence of
this proof and regularizes both Expos\'e~I equation references; those
corrected references are used above.

Under the hypotheses of \hyperref[II.4]{Corollary~4},\footnote{The corrected French source
adds: ``For consistency and clarity, only equations have been numbered in
parentheses.''} $U$ is a finite union of affine open subsets $X_f$.
Consequently, one can find an ideal $I$, generated by finitely many elements
$f_\alpha$ and defining $Y$; write $\ffam=(f_\alpha)$. With the notation of
EGA~III, 1,\footnote{Recall that $\mathrm H^\bullet(\ffam,M)$ is the Koszul
cohomology $\mathrm H^\bullet(\operatorname{Hom}(K_\bullet(\ffam),M))$ of
$\ffam$ with values in $M$ (EGA~III, 1.1.2), and that
$\mathrm H^\bullet((\ffam),M)$ is the direct limit (ibid., 1.1.6.5)
$\varinjlim_n\mathrm H^\bullet(\ffam^n,M)$, whose transition morphisms are
induced by the natural morphisms
$K_\bullet(\ffam^{n+1})\to K_\bullet(\ffam^n)$ (ibid., 1.1.6).} we have the
following result.

\begin{sgastatement}{Proposition 5}\sgareaderlabel{II.5}{5}
Suppose that $X$ is the spectrum of a ring $A$. Let
$\ffam=(f_\alpha)$ be a finite family of elements of $A$, let $Y$ be the
closed subset of $X$ that they define, let $M$ be an $A$-module, and let $F$
be the associated sheaf. There are isomorphisms of $\partial$-functors in
$M$:
\begin{equation}
  \mathrm H^i((\ffam),M)\simeq \mathrm H_Y^i(X,F).
  \tag{5.1}\label{eq:II.5.1}
\end{equation}
\end{sgastatement}

We shall also write
$\mathrm H_J^i(M)=\mathrm H_Y^i(X,F)$ when $Y$ is the closed subset of
$X=\Spec A$ defined by an ideal $J$ of $A$.

For $i=0$\srcpage{21} and $i=1$, use the exact sequences
\eqref{eq:II.4.2} and EGA~III, 1.4.3.2; for $i\geq2$, use
\eqref{eq:II.4.3} and EGA~III, 1.4.3.1. This gives isomorphisms functorial in
$M$. One verifies that, up to a sign depending only on $i$, they are
compatible with the boundary operator; hence the isomorphism of
$\partial$-functors in \eqref{eq:II.5.1} exists.

\noindent\textit{Editorial note.} For the case $i\geq2$, the original
edition cited EGA~III, 1.4.3.2; the corrected branch cites 1.4.3.1.

Now let $X$ be a prescheme, let $Y$ be a closed subset of $X$, and let
$f\colon Y\to X$ be the inclusion. Let $I$ be a quasi-coherent ideal defining
$Y$ in $X$, and let $F$ be a sheaf on $X$.

We have seen that there are isomorphisms of $\partial$-functors in $F$
\begin{align}
  \Ext_{\OX}^i(X;f_*f^{-1}(\OX),F)
    &\longrightarrow \mathrm H_Y^i(X,F),
    \tag{$*$}\label{eq:II.5.*}\\
  \uExt_{\OX}^i(f_*f^{-1}(\OX),F)
    &\longrightarrow \uH_Y^i(F).
    \tag{$**$}\label{eq:II.5.**}
\end{align}

\noindent\textit{Source note.} The corrected French source calls these
arrows isomorphisms but displays plain right arrows in both $(*)$ and $(**)$.
The arrows have been preserved exactly; no unmarked replacement by
$\xrightarrow{\sim}$ has been made.

Let $n,m$ be integers such that $m\geq n\geq0$. Denote by $i_{n,m}$ the
canonical map
\[
  \mathcal O_{Y_m}=\OX/I^{m+1}
  \longrightarrow \OX/I^{n+1}=\mathcal O_{Y_n},
\]
and by $j_n$ the map
$f_*f^{-1}(\OX)\to\mathcal O_{Y_n}$. The pairs
$(\mathcal O_{Y_n},i_{n,m})$ form a projective system, and the maps $j_n$ are
compatible with the $i_{n,m}$.

Applying the functor $\Ext_{\OX}^i(X;\mathord\cdot,F)$ yields a morphism
\[
  \varphi'\colon
  \varinjlim_n\Ext_{\OX}^i(X;\mathcal O_{Y_n},F)
  \longrightarrow
  \Ext_{\OX}^i(X;f_*f^{-1}(\OX),F).
\]
It is readily seen to be a morphism of cohomological functors in $F$.
Consequently, the morphism
\[
  \varphi\colon
  \varinjlim_n\Ext_{\OX}^i(X;\mathcal O_{Y_n},F)
  \longrightarrow \mathrm H_Y^i(X,F),
\]
obtained by composing $\varphi'$ with~\eqref{eq:II.5.*}, is itself a
morphism of cohomological functors in $F$.

We define\srcpage{22} similarly
\[
  \underline{\varphi}\colon
  \varinjlim_n\uExt_{\OX}^i(\mathcal O_{Y_n},F)
  \longrightarrow \uH_Y^i(F).
\]

Our goal is the following theorem.

\begin{sgastatement}{Theorem 6}\sgareaderlabel{II.6}{6}
\textup{a)} Let $X$ be a locally noetherian prescheme, let $Y$ be a closed
subset of $X$ defined by a coherent ideal $I$, and let $F$ be a
quasi-coherent $\OX$-module. Then $\boldsymbol{\varphi}$ is an isomorphism.

\textup{b)} If $X$ is noetherian, then $\varphi$ is an isomorphism.
\end{sgastatement}

\noindent\textit{Source note.} The corrected source defines the sheaf-level
morphism immediately above as $\underline{\varphi}$ but denotes it by the
bold symbol $\boldsymbol{\varphi}$ in Theorem~6(a) and Lemma~7. Both graphic
forms are retained rather than silently normalized.

\hyperref[II.6]{Theorem~6} will follow from \hyperref[II.6]{Theorem~6}(a) and the following lemma.

\begin{sgastatement}{Lemma 7}\sgareaderlabel{II.7}{7}
If the underlying topological space of $X$ is noetherian and
$\boldsymbol{\varphi}$ is an isomorphism, then $\varphi$ is also an
isomorphism.
\end{sgastatement}

We first prove \hyperref[II.7]{Lemma~7}. There is a spectral sequence
\begin{equation}
  \mathrm H^p(X,\uH_Y^q(F))
  \Longrightarrow \mathrm H_Y^*(X,F).
  \tag{7.1}\label{eq:II.7.1}
\end{equation}
There is, on the other hand, a direct system of spectral sequences
\begin{equation}
  \mathrm H^p\bigl(X,\uExt_{\OX}^q(\mathcal O_{Y_n},F)\bigr)
  \Longrightarrow \Ext_{\OX}^*(X;\mathcal O_{Y_n},F).
  \tag{7.2$_n$}\label{eq:II.7.2n}
\end{equation}

\noindent\textit{Editorial note.} In both abutments above, the original
edition printed a superscript $m$; the corrected branch replaces it by the
superscript $*$, which is preserved here.

It follows from the definitions of $\varphi$ and $\underline{\varphi}$ that
these morphisms are associated with a homomorphism $\Phi$ of spectral
sequences, from the direct limit of~\eqref{eq:II.7.2n} to
\eqref{eq:II.7.1}. If the underlying space of $X$ is noetherian, then by
Godement, 4.12.1,\footnote{Cf. the first bibliographic reference at the end
of Expos\'e~I.}
\[
  \varinjlim_n
    \mathrm H^p\bigl(X,\uExt_{\OX}^q(\mathcal O_{Y_n},F)\bigr)
  \xrightarrow{\sim}
  \mathrm H^p\!\left(
    X,\varinjlim_n\uExt_{\OX}^q(\mathcal O_{Y_n},F)\right).
\]
Thus $\Phi_2$ may be written as a morphism
\[
  \mathrm H^p\!\left(
    X,\varinjlim_n\uExt_{\OX}^q(\mathcal O_{Y_n},F)\right)
  \longrightarrow \mathrm H^p(X,\uH_Y^q(F)),
\]
which is precisely the morphism deduced from $\underline{\varphi}$.

If $\underline{\varphi}$\srcpage{23} is an isomorphism, then so is
$\Phi_2$, and hence so is $\varphi$ by EGA~$0_{\mathrm{III}}$, 11.1.5.
This proves \hyperref[II.7]{Lemma~7}.

We now prove \hyperref[II.6]{Theorem~6}(a). This is a local question on $X$. By \hyperref[II.4]{Corollary~4}
and EGA~I, 1.3.9 and 1.3.12, we may suppose that $X$ is the spectrum of a
ring $A$. Write $M$ for the $A$-module whose associated sheaf is $F$. It is
therefore enough to prove that, under the hypotheses of \hyperref[II.6]{Theorem~6}(a), the
canonical homomorphism
\begin{equation}
  \varinjlim_n\Ext_A^i(A/I^n,M)
  \longrightarrow \mathrm H_Y^i(X,M)
  \tag{7.3}\label{eq:II.7.3}
\end{equation}
is an isomorphism.

\noindent\textit{Source note.} The corrected French text first uses $M$ in
display~(7.3) without introducing it. The affine correspondence with the
previously given quasi-coherent sheaf $F$ is made explicit above.

Let $f_\alpha$ be finitely many elements of $A$ that generate $I$, and put
$\ffam=(f_\alpha)$. The sequence of ideals $(\ffam^n)$ is decreasing and
cofinal with the sequence of ideals $I^n$, so that~\eqref{eq:II.7.3} is
equivalent to a morphism of $\partial$-functors in $M$:
\begin{equation}
  \varinjlim_n\Ext_A^i(A/(\ffam^n),M)
  \longrightarrow \mathrm H_Y^i(X,M).
  \tag{7.4}\label{eq:II.7.4}
\end{equation}
On the other hand, there are canonical isomorphisms
\begin{equation}
\begin{split}
  \varinjlim_n\operatorname{Hom}_A(A/(\ffam^n),M)
  &\simeq \varinjlim_n
    \bigl(m\in M\mid(\ffam^n)m=0\bigr) \\
  &\simeq \mathrm H^0((\ffam),M).
\end{split}
  \tag{7.5}\label{eq:II.7.5}
\end{equation}
Since $\varinjlim_n\Ext_A^i(A/(\ffam^n),M)$ is a universal
$\partial$-functor in $M$, there is a unique morphism of $\partial$-functors
in $M$
\begin{equation}
  \varinjlim_n\Ext_A^i(A/(\ffam^n),M)
  \longrightarrow \mathrm H^i((\ffam),M)
  \tag{7.6}\label{eq:II.7.6}
\end{equation}
that agrees in degree zero with~\eqref{eq:II.7.5}.

The comparison obtained from~\eqref{eq:II.7.3} and
\eqref{eq:II.5.1} is a morphism of $\partial$-functors in $M$ that agrees
with~\eqref{eq:II.7.6} in degree zero; it therefore agrees with
\eqref{eq:II.7.6} in every degree. \hyperref[II.6]{Theorem~6}(a) is consequently an immediate
consequence of the following lemma.

\noindent\textit{Source note.} The corrected French sentence calls the
comparison just used ``the composite of (7.3) and (5.1).'' As the displayed
arrows in (7.3) and (5.1) have the same target rather than composable source
and target, the intended comparison requires identifying that target through
the isomorphism (5.1), equivalently using its inverse. The mathematical
comparison is stated explicitly above; the source's directional anomaly is
recorded here rather than silently concealed.

\begin{sgastatement}{Lemma 8}\sgareaderlabel{II.8}{8}
Let $A$ be a noetherian ring, let $I$ be an ideal generated by a finite
family $\ffam=(f_\alpha)$ of elements, and let $M$ be an $A$-module. Then the
homomorphisms~\eqref{eq:II.7.6} are isomorphisms.
\end{sgastatement}

\begin{sgastatement}{Lemma 9}\sgareaderlabel{II.9}{9}
Let $A$\srcpage{24} be a ring, let $\ffam=(f_\alpha)$ be a finite family of
elements of $A$, let $I$ be the ideal generated by $\ffam$, and let $i>0$ be
an integer. The following conditions are equivalent:
\begin{enumerate}[label=\textup{\alph*)},leftmargin=2.5em]
  \item The homomorphism~\eqref{eq:II.7.6} is an isomorphism for every $M$.
  \item $\mathrm H^i((\ffam),M)=0$ for every injective $M$.
  \item The projective system
  $(\mathrm H_i(\ffam^n,A))=\mathrm H_{i,n}$ is essentially zero; that is,
  for every $n$ there exists $n'>n$ such that
  $\mathrm H_{i,n'}\to\mathrm H_{i,n}$ is zero.
\end{enumerate}
\end{sgastatement}

\noindent\textit{Source note.} Lemma~9(c), and later Lemma~11, print their
systems with the index $n\in\ZZ$, whereas the proofs explicitly take $n>0$
and use ordinary positive powers. The printed indexing is retained; this
translation does not introduce negative powers.

Condition a) trivially implies b).

Condition b) implies a). Indeed, b) implies that
$M\mapsto\mathrm H^i((\ffam),M)$ is a universal cohomological functor, so
\eqref{eq:II.7.6} is a morphism of universal cohomological functors. It is an
isomorphism in degree zero and hence in every degree.

Condition c) implies b). Indeed, if $M$ is injective, then for every $n$
\[
  \mathrm H^i(\ffam^n,M)
  =\operatorname{Hom}(\mathrm H_i(\ffam^n,A),M)
  =\operatorname{Hom}(\mathrm H_{i,n},M).
\]
Thus c) implies that, for every $i$, the direct system
$(\mathrm H^i(\ffam^n,M))_{n\in\ZZ}$ is essentially zero, whence b).

Condition b) implies c). Indeed, fix $n>0$, and let $j$ be a monomorphism
from $\mathrm H_{i,n}$ into an injective module $M$. For $n'\geq n$, let
$j_{n'}\in\operatorname{Hom}(\mathrm H_{i,n'},M)$ be the composite of $j$
with the transition homomorphism
$t_{n',n}\colon\mathrm H_{i,n'}\to\mathrm H_{i,n}$. The $j_{n'}$ define an
element of $\mathrm H^i((\ffam),M)$, which is zero by hypothesis. There is
therefore an $n_0$ such that $j_{n'}=0$ if $n'>n_0$. Since $j$ is a
monomorphism, $j_{n'}=0$ implies $t_{n',n}=0$, which proves c).

\noindent\textit{Editorial note.} In the last implication, the original
edition shortened the transition-map symbol to $t_{n'}$; the corrected
branch restores $t_{n',n}$, as retained above.

\begin{sgastatement}{Corollary 10}\sgareaderlabel{II.10}{10}
Suppose that the underlying space of $X=\Spec(A)$ is noetherian. The
preceding conditions hold for every finite family of elements of $A$ and
every $i>0$ (or, equivalently, for $i=1$) if and only if, for every injective
$A$-module $M$, the sheaf $F$ associated with $M$ is flasque.
\end{sgastatement}

The condition is necessary. Indeed, let $\ffam=(f_\alpha)$ be a finite family
of elements of $A$, let $Y$ be the closed subset defined by $\ffam$, and put
$U=X-Y$. We then have the exact sequence
\[
  \mathrm H^0(X,F)\longrightarrow \mathrm H^0(U,F)
  \longrightarrow \mathrm H^1((\ffam),M)\longrightarrow0,
\]
and \hyperref[II.9]{Lemma~9}(b) shows that
$\mathrm H^0(X,F)\to\mathrm H^0(U,F)$ is surjective.

The condition is sufficient\srcpage{25} by~\eqref{eq:II.5.1} and the fact
that, for every closed subset $Y$ of $X$ and every flasque sheaf $F$ on $X$,
$\mathrm H_Y^i(X,F)=0$ for $i>0$.

\begin{sgastatement}{Lemma 11}\sgareaderlabel{II.11}{11}
Under the hypotheses of \hyperref[II.9]{Lemma~9}, for every noetherian $A$-module $N$ and
every $i>0$, the projective system
$(\mathrm H_{i,n}(N))_{n\in\ZZ}$, where
$\mathrm H_{i,n}(N)=\mathrm H_i(\ffam^n,N)$, is essentially zero.
\end{sgastatement}

\begin{sgaproof}
We argue by induction on the number $m$ of elements of $\ffam$.

If $m=1$, the family $\ffam$ consists of a single element, say $f$.
Then $\mathrm H_{i,n}(N)=0$ for $i>1$, and $\mathrm H_{1,n}(N)$ is
canonically isomorphic to the annihilator $N(n)$ of $f^n$ in $N$. For
$n'\geq n$, the transition homomorphism $N(n')\to N(n)$ is multiplication
by $f^{n'-n}$. The $N(n)$ form an increasing sequence of submodules of $N$.
Since $N$ is noetherian, there is an $n_0$ such that
$N(n)=N(n_0)$ for $n\geq n_0$. Hence all the $N(n)$ are annihilated by
$f^{n_0}$, and every transition homomorphism $N(n')\to N(n)$ is zero when
$n'\geq n+n_0$. The lemma is thus proved for $m=1$.

Now suppose that $m>1$ and that the lemma has been proved for all integers
$m'<m$. Put
$\gfam=(f_1,\ldots,f_{m-1})$ and $\hfam=f_m$.

For every $n>0$, EGA~III, 1.1.4.1 gives an exact sequence
\[
  0\longrightarrow
  \mathrm H_0(\hfam^n,\mathrm H_i(\gfam^n,N))
  \longrightarrow \mathrm H_i(\ffam^n,N)
  \longrightarrow
  \mathrm H_1(\hfam^n,\mathrm H_{i-1}(\gfam^n,N))
  \longrightarrow0,
\]
and, as $n$ varies, a projective system of exact sequences. The induction
hypothesis shows that, for $i>0$, the modules
$\mathrm H_i(\gfam^n,N)$ form an essentially zero projective system. The same
is therefore true of
$\mathrm H_0(\hfam^n,\mathrm H_i(\gfam^n,N))$, which is identified with a
quotient of $\mathrm H_i(\gfam^n,N)$. For the terms on the right, factor the
transition morphisms from $n'$ to $n$ as
\[
\begin{split}
  \mathrm H_1(\hfam^{n'},\mathrm H_{i-1}(\gfam^{n'},N))
  &\longrightarrow
  \mathrm H_1(\hfam^{n'},\mathrm H_{i-1}(\gfam^n,N)) \\
  &\longrightarrow
  \mathrm H_1(\hfam^n,\mathrm H_{i-1}(\gfam^n,N)).
\end{split}
\]
Since $\mathrm H_{i-1}(\gfam^n,N)$ is a noetherian module, the case $m=1$
shows that, for fixed $n$, there exists $n'>n$ such that the second arrow is
zero. One therefore sees\srcpage{26} that, in this projective system of exact
sequences, the two outer projective systems are essentially zero. The middle
projective system is therefore essentially zero as well.
\end{sgaproof}

We have thus proved \hyperref[II.11]{Lemma~11}, hence \hyperref[II.8]{Lemma~8}, and consequently \hyperref[II.6]{Theorem~6}.

\begin{sgaremark}{Remark}
\hyperref[II.6]{Theorem~6} can also be obtained by proving the condition of \hyperref[II.10]{Corollary~10} with
the aid of the structure theorems for injective modules over a noetherian
ring (Matlis, Gabriel).
\end{sgaremark}

% END INCLUDED active_source/parts/02_expose_II.tex

% BEGIN INCLUDED active_source/parts/03_expose_III.tex
% SOURCE COMPONENT: 01_section_1.tex
\makeatletter
\providecommand{\sgareaderlabel}[2]{%
  \begingroup
  \def\@currentlabel{#2}%
  \phantomsection\label{#1}%
  \endgroup}
\makeatother

\chapter{Cohomological Invariants and Depth}
\label{III}

\section*{1. Review}
\sgareaderlabel{III.1}{1}

We\srcpage{27} shall state a few definitions and results which the reader may
find, for example, in Chapter~I of the course given by J.-P. Serre at the
Collège de France in 1957--58.\footnote{\sgareaderlabel{noteserre}{*}\textit{Editorial note.} The
reissue of Serre's text (J.-P. Serre, \emph{Algèbre locale.
Multiplicités}, course at the Collège de France, 1957--58, written up by
Pierre Gabriel, second edition, Lecture Notes in Mathematics, vol.~11,
Springer-Verlag, 1965) no longer contains proofs of these statements. One may
consult N.~Bourbaki, \emph{Algèbre commutative}, Masson, as Serre himself
suggests.}

\begin{sgastatement}{Definition 1.1}\sgareaderlabel{def:III.1.1}{1.1}
Let $A$ be a ring (commutative with identity, as throughout what follows), and
let $M$ be an $A$-module (unital, as throughout what follows). We call:
\begin{itemize}[leftmargin=2.2em]
  \item the \emph{annihilator of $M$}, denoted $\Ann M$, the set of
  $a\in A$ such that $am=0$ for every $m\in M$;
  \item the \emph{support of $M$}, denoted $\Supp M$, the set of prime
  ideals $\mathfrak p$ of $A$ such that the localization $M_{\mathfrak p}$
  is nonzero;
  \item the “\emph{assassin of $M$},” or the “\emph{set of prime ideals
  associated with $M$},” denoted $\Ass M$, the set of prime ideals
  $\mathfrak p$ for which there exists a nonzero element of $M$ whose
  annihilator is $\mathfrak p$.
\end{itemize}
If $\mathfrak a$ is an ideal of $A$, we denote by $\mathfrak r(\mathfrak a)$
the radical of $\mathfrak a$ in $A$, that is, the set of elements of $A$ some
power of which lies in $\mathfrak a$.
\end{sgastatement}

The following results hold under the assumption that \emph{$A$ is noetherian
and $M$ is finitely generated}.

\begin{sgastatement}{Proposition 1.1}\sgareaderlabel{III.1.1}{1.1}
\begin{enumerate}[label=\textup{(\roman*)},leftmargin=2.5em]
  \item $\Ass M$ is a finite set.
  \item An element of $A$ annihilates a nonzero element of $M$ if and only
  if\srcpage{28} it belongs to one of the prime ideals associated with $M$.
  \item The radical of the annihilator of $M$,
  $\mathfrak r(\Ann M)$, is the intersection of the prime ideals associated
  with $M$ that are minimal with respect to inclusion in $\Ass M$.
\end{enumerate}
\end{sgastatement}

\begin{sgastatement}{Proposition 1.2}\sgareaderlabel{III.1.2}{1.2}
Let $\mathfrak p$ be a prime ideal of $A$. The following assertions are
equivalent:
\begin{enumerate}[label=\textup{(\roman*)},leftmargin=2.8em]
  \item $\mathfrak p\in\Supp M$.
  \item There exists $\mathfrak q\in\Ass M$ such that
  $\mathfrak q\subset\mathfrak p$.
  \item $\mathfrak p\supset\Ann M$.
  \item[\textup{(iii bis)}] $\mathfrak p\supset\mathfrak r(\Ann M)$.
\end{enumerate}
\end{sgastatement}

\begin{sgastatement}{Proposition 1.3}\sgareaderlabel{III.1.3}{1.3}
Let $N$ be a finitely generated $A$-module. Then
\[
  \Ass\operatorname{Hom}_A(N,M)=\Supp N\cap\Ass M.
\]
\end{sgastatement}



% SOURCE COMPONENT: 02_section_2.tex
\section*{2. Depth}
\sgareaderlabel{III.2}{2}

Throughout\srcpage{28} this section, $A$ denotes a commutative ring, $I$ an ideal of
$A$, and $M$ and $N$ two $A$-modules. We denote by $X$ the prime spectrum of
$A$ (its structure sheaf will not be used in this section), and by $Y$ the
closed set defined by $I$, namely
$Y=\Supp(A/I)=\{\mathfrak p\in X\mid \mathfrak p\supset I\}$.

\begin{sgastatement}{Lemma 2.1}\sgareaderlabel{III.2.1}{2.1}
Suppose that $A$ is noetherian, that $M$ and $N$ are finitely generated, and
that $\Supp N=Y$. Then the following assertions are equivalent:
\begin{enumerate}[label=\textup{(\roman*)},leftmargin=2.6em]
\item $\Homop_A(N,M)=0$.
\item $\Supp N\cap\Ass M=\varnothing$.
\item The ideal $I$ is not a zero divisor on $M$, meaning that, for every
$m\in M$, the equality $Im=0$ implies $m=0$.
\item There exists an $M$-regular element in $I$. (An element $a\in A$ is
called $M$-regular if multiplication by $a$ on $M$ is injective.)
\item For every $\mathfrak p\in Y$, the maximal ideal
$\mathfrak pA_{\mathfrak p}$ of the local ring $A_{\mathfrak p}$ is not
associated with $M_{\mathfrak p}$; in symbols,
$\mathfrak pA_{\mathfrak p}\notin\Ass M_{\mathfrak p}$.
\end{enumerate}
\end{sgastatement}

\begin{proof}
The equivalence (i)$\Leftrightarrow$(ii) follows because
$\Ass\Homop_A(N,M)=\varnothing$ is equivalent to (ii) by \hyperref[III.1.3]{Proposition 1.3},
and to (i) by an immediate consequence of \hyperref[III.1.2]{Proposition 1.2}.

The implication\srcpage{29} (iii)$\Rightarrow$(ii) follows by contradiction. If
$\mathfrak p\in\Supp N\cap\Ass M$, then $\mathfrak p\supset I$ and there
is an $m\in M$ whose annihilator is $\mathfrak p$; hence
$Im\subset\mathfrak p m=0$, contrary to (iii).

The implication (iv)$\Rightarrow$(iii) is immediate.

Conditions (ii) and (iv) are equivalent because $\Supp N=Y$, so (ii) says
that $I$ is contained in no prime ideal associated with $M$. Since those
associated ideals are prime and finite in number, this is equivalent to
saying that $I$ is not contained in their union. By \hyperref[III.1.1]{Proposition 1.1}(ii),
that union is precisely the set of elements of $A$ which are not
$M$-regular.

To prove (i)$\Rightarrow$(v), let $\mathfrak p\in Y$. From
$\Homop_A(N,M)=0$ and the formula
\[
  \bigl(\Homop_A(N,M)\bigr)_{\mathfrak p}
  =\Homop_{A_{\mathfrak p}}(N_{\mathfrak p},M_{\mathfrak p})
\]
one obtains $\Homop_{A_{\mathfrak p}}(N_{\mathfrak p},M_{\mathfrak p})=0$.
Thus, by \hyperref[III.1.3]{Proposition 1.3},
\[
  \Supp N_{\mathfrak p}\cap\Ass M_{\mathfrak p}=\varnothing.
\]
But $\mathfrak pA_{\mathfrak p}\in\Supp N_{\mathfrak p}$, so
$\mathfrak pA_{\mathfrak p}\notin\Ass M_{\mathfrak p}$.

Finally, suppose (v), and let $\mathfrak p\in\Ass M$. There is an
$m\in M$ with annihilator $\mathfrak p$. Its canonical image in
$M_{\mathfrak p}$ is nonzero, and its annihilator contains
$\mathfrak pA_{\mathfrak p}$ and hence equals it. Thus
$\mathfrak pA_{\mathfrak p}\in\Ass M_{\mathfrak p}$, so (v) gives
$\mathfrak p\notin Y$. This proves (i).
\end{proof}

We shall now work with these conditions after replacing the
functor $\Homop$ by its derived functors.

\begin{sgastatement}{Theorem 2.2}\sgareaderlabel{III.2.2}{2.2}\srcpage{30}
Let $A$ be a commutative ring, $I$ an ideal of $A$, $M$ an $A$-module, and
$n$ an integer.
\begin{enumerate}[label=\textup{(\alph*)},leftmargin=2.5em]
\item If there is a sequence $f_1,\ldots,f_{n+1}$ of elements of $I$ which
is $M$-regular (that is, $f_1$ is $M$-regular and $f_{i+1}$ is regular on
$M/(f_1,\ldots,f_i)M$ for $i\le n$), then, for every $A$-module $N$
annihilated by a power of $I$,
\[
  \Extop_A^i(N,M)=0\qquad\text{for }i\le n.
\]
\item If, in addition, $A$ is noetherian, $M$ is finitely generated, and
there is a finitely generated $A$-module $N$ with $\Supp N=V(I)$ and
$\Extop_A^i(N,M)=0$ for $i\le n$, then there exists an $M$-regular
sequence $f_1,\ldots,f_{n+1}$ of elements of $I$.
\end{enumerate}
\end{sgastatement}

We first prove (a), by induction. If $n<0$, the assertion is empty. Suppose
$n\ge0$ and that (a) has been proved for every $n'<n$. By hypothesis there
is an $M$-regular element $f_1\in I$. Denote by $f_1^i$ multiplication by
$f_1$ on $\Extop_A^i(N,M)$ and by $f_1^M$ multiplication by $f_1$ on $M$.
The sequence
\begin{equation}\label{eq:III.2.1}
  0\longrightarrow M\xrightarrow{f_1^M}M\longrightarrow M/f_1M
  \longrightarrow0
\end{equation}
is exact, and therefore so is
\[
  \Extop_A^{i-1}(N,M/f_1M)\xrightarrow{\delta}\Extop_A^i(N,M)
  \xrightarrow{f_1^i}\Extop_A^i(N,M).
\]
By hypothesis, $I^rN=0$ for some $r$, so $f_1^0$ is nilpotent. Since
$\Extop^i$ is a universal cohomological functor, the same is true of
$f_1^i$ for every $i$. Moreover, $M/f_1M$ has a regular sequence of length
$n$. The induction hypothesis therefore gives
\[
  \Extop_A^{i-1}(N,M/f_1M)=0\qquad\text{for }i\le n.
\]
\footnote{\textit{Source correction.} The French authority prints
$\mathrm{Ext}_A^{i-1}(N,M)=0$ for $i\le n-1$. The quotient module and the
bound are both defective: applying the induction hypothesis to the length
$n$ sequence on $M/f_1M$ gives the displayed formula, which is exactly what
is needed below.}
Consequently, for $i\le n$, $f_1^i$ is both injective and nilpotent, so
$\Extop_A^i(N,M)=0$.

We\srcpage{31} now prove (b), again by induction. If $n<0$, the assertion is
empty. If $n=0$, it follows from \hyperref[III.2.1]{Lemma 2.1}, (i)$\Rightarrow$(iv). If $n>0$,
the case $n=0$ gives an $M$-regular element $f_1\in I$. From
\eqref{eq:III.2.1} one obtains the exact sequence
\begin{equation}\label{eq:III.2.2}
  \Extop_A^{i-1}(N,M)\longrightarrow
  \Extop_A^{i-1}(N,M/f_1M)\longrightarrow\Extop_A^i(N,M).
\end{equation}
The hypotheses of (b) therefore hold for $M/f_1M$ and $n-1$. By induction
there is a regular sequence of $n$ elements of $I$ on $M/f_1M$, and hence
an $M$-regular sequence of $n+1$ elements of $I$ beginning with $f_1$.

This theorem suggests the following generalization of the classical
definition of the depth of a finitely generated module over a noetherian
ring.

\begin{sgastatement}{Definition 2.3}\sgareaderlabel{III.2.3}{2.3}
Let $A$ be a commutative ring with identity, $M$ an $A$-module, and $I$ an
ideal of $A$. The \emph{$I$-depth of $M$}, denoted $\prof_I M$, is the
supremum in $\mathbb N\cup\{+\infty\}$ of the natural numbers $n$ such
that, for every finitely generated $A$-module $N$ annihilated by a power of
$I$,
\[
  \Extop_A^i(N,M)=0\qquad\text{for every }i<n.
\]
\end{sgastatement}

It follows from the preceding theorem that, if $n$ is the supremum of the
lengths of $M$-regular sequences of elements of $I$, then
$n\le\prof_I M$. More precisely:

\begin{sgastatement}{Proposition 2.4}\sgareaderlabel{III.2.4}{2.4}
Let $A$ be a commutative ring, $I$ an ideal of $A$, $M$ an $A$-module,
and $n\in\mathbb N$. Consider the assertions:
\begin{enumerate}[label=\textup{(\arabic*)},leftmargin=2.7em]
\item $n\le\prof_I M$.
\item For every finitely generated $A$-module $N$ annihilated by a power
of $I$, one has
\[
  \Extop_A^i(N,M)=0\qquad\text{for }i<n.
\]
\item There\srcpage{32} exists a finitely generated $A$-module $N$ such that
$\Supp N=V(I)$, and such that $\Extop_A^i(N,M)=0$ whenever $i<n$.
\item There is an $M$-regular sequence of length $n$ consisting of
elements of $I$.
\end{enumerate}
The following logical implications hold:
\[
\begin{array}{ccccc}
\hyperref[III.2.4]{(1)}&\Longleftrightarrow&\hyperref[III.2.4]{(2)}&\Longleftarrow&\hyperref[III.2.4]{(4)}\\[-0.2em]
&&\Big\Downarrow&&\\[-0.2em]
&&\hyperref[III.2.4]{(3)}&&
\end{array}
\]
If $A$ is noetherian and $M$ is finitely generated, all four conditions are
equivalent.
\end{sgastatement}

\begin{proof}
\hyperref[III.2.4]{Conditions (1)} and \hyperref[III.2.4]{(2)} are equivalent by definition, and
\hyperref[III.2.4]{(2)}$\Rightarrow$\hyperref[III.2.4]{(3)} follows by taking $N=A/I$. The implication
\hyperref[III.2.4]{(4)}$\Rightarrow$\hyperref[III.2.4]{(2)} is \hyperref[III.2.2]{Theorem 2.2}(a). Finally, when $A$ is noetherian and
$M$ is finitely generated, \hyperref[III.2.4]{(3)}$\Rightarrow$\hyperref[III.2.4]{(4)} is \hyperref[III.2.2]{Theorem 2.2}(b).
\end{proof}

From now until the end of the section, we assume that $A$ is
noetherian and $M$ is finitely generated.

\begin{sgastatement}{Corollary 2.5}\sgareaderlabel{III.2.5}{2.5}
If $f\in I$ is $M$-regular, then
\[
  \prof_I M=\prof_I(M/fM)+1.
\]
\end{sgastatement}

Indeed, if $n\le\prof_I(M/fM)$, there is an $(M/fM)$-regular sequence
$f_1,\ldots,f_n$ of elements of $I$. Hence $f,f_1,\ldots,f_n$ is
$M$-regular, so $n+1\le\prof_I M$ and
$\prof_I M\ge\prof_I(M/fM)+1$. Conversely, for every finitely generated
$A$-module $N$ annihilated by a power of $I$, the exact sequence
  \eqref{eq:III.2.2} shows that if $i<\prof_I M$, then
  $\Extop_A^{i-1}(N,M/fM)=0$.\footnote{\textit{Source correction.}
  The French authority prints $i\le\prof_I M$. In the displayed exact
  sequence, vanishing of both adjacent terms is guaranteed for
  $i<\prof_I M$; this is precisely the range needed for the following depth
  inequality.} Thus
$\prof_I M-1\le\prof_I(M/fM)$.

\begin{sgastatement}{Corollary 2.6}\sgareaderlabel{III.2.6}{2.6}
Every finite $M$-regular sequence consisting of elements of $I$ can be
extended to a maximal $M$-regular sequence, whose length is necessarily the
$I$-depth of $M$.
\end{sgastatement}

\begin{sgastatement}{Remark 2.7}\sgareaderlabel{III.2.7}{2.7}\srcpage{33}
One can scarcely refrain from saying that the greater the depth of an
$A$-module, the more beautiful it is. A module whose support does not meet
$V(I)$ is among the most beautiful: indeed, $\prof_I M$ is finite if and
only if $\Supp M\cap V(I)\ne\varnothing$.
\end{sgastatement}

\begin{sgastatement}{Remark 2.8}\sgareaderlabel{III.2.8}{2.8}
If $A$ is a semi-local ring, let $\mathfrak r(A)$ be its radical and
$k=A/\mathfrak r(A)$ its residue ring. The relevant notion of depth is
obtained by taking $I=\mathfrak r(A)$; accordingly, we write simply
$\prof M$ for the $\mathfrak r(A)$-depth of an $A$-module $M$. This is the
notion of ``homological codimension'' (cf. Serre, \emph{op. cit.}), formerly
denoted $\codh_A M$, defined as the infimum of the integers $i$ such that
$\Extop_A^i(k,M)\ne0$. Indeed, $\Supp k=V(\mathfrak r(A))$.
\end{sgastatement}

\begin{sgastatement}{Proposition 2.9}\sgareaderlabel{III.2.9}{2.9}
If $A$ is noetherian and $M$ is finitely generated, then
\[
  \prof_I M=\inf_{\mathfrak p\in V(I)}\prof M_{\mathfrak p}.
\]
\end{sgastatement}

\begin{sgastatement}{Corollary 2.10}\sgareaderlabel{III.2.10}{2.10}
If $A$ is a noetherian semi-local ring and $M$ is a finitely generated
$A$-module, then
\[
  \prof M=\inf_{\mathfrak m}\prof M_{\mathfrak m},
\]
where $\mathfrak m$ ranges over the maximal ideals of $A$.
\end{sgastatement}

The corollary follows immediately from \hyperref[III.2.9]{Proposition 2.9}, since the prime
ideals containing the radical are precisely the maximal ideals.

We now prove \hyperref[III.2.9]{Proposition 2.9}. Let $f\in I$. If $f$ is $M$-regular and
$\mathfrak p\supset I$, then the image $g$ of $f$ in $A_{\mathfrak p}$
belongs to the maximal ideal $\mathfrak pA_{\mathfrak p}$ and is
$M_{\mathfrak p}$-regular, as follows from the exact sequence
\begin{equation}\label{eq:III.2.3}
  0\longrightarrow M_{\mathfrak p}\xrightarrow{g'}M_{\mathfrak p}
  \longrightarrow(M/fM)_{\mathfrak p}\longrightarrow0,
\end{equation}
where $g'$ denotes multiplication by $g$ on $M_{\mathfrak p}$. This also
shows that $(M/fM)_{\mathfrak p}\simeq M_{\mathfrak p}/gM_{\mathfrak p}$.
Applying \hyperref[III.2.5]{Corollary 2.5} to $M$ and to $M_{\mathfrak p}$ gives, by induction,
$\prof_I M\le\nu(M)$,\srcpage{34} where
\[
  \nu(M)=\inf_{\mathfrak p\in V(I)}\prof M_{\mathfrak p}.
\]
More precisely, induction gives $\nu(M)=\nu(M/fM)+1$ whenever $f$ is
$M$-regular. It remains to prove that, if $\nu(M)\ge1$, there is an
$M$-regular element in $I$. Applying \hyperref[III.2.1]{Lemma 2.1} to
$M_{\mathfrak p}$, $A_{\mathfrak p}$, and
$\mathfrak pA_{\mathfrak p}$ for every $\mathfrak p\in V(I)$ shows that
$\mathfrak pA_{\mathfrak p}\notin\Ass M_{\mathfrak p}$. Applying the
same lemma to $A$, $M$, and $I$ gives the conclusion.

\begin{sgastatement}{Proposition 2.11}\sgareaderlabel{III.2.11}{2.11}
Let $u:A\to B$ be a homomorphism of noetherian rings, let $I$ be an ideal
of $A$, and let $M$ be a finitely generated $A$-module. Put
$I_B=I\otimes_A B$ and $M_B=M\otimes_A B$. If $B$ is flat over $A$, then
\[
  \prof_{I_B}M_B\ge\prof_I M;
\]
if $B$ is faithfully flat over $A$, equality holds.
\end{sgastatement}

Indeed, put $N=A/I$. Flatness gives $N\otimes_A B=B/I_B$; set
$N_B=N\otimes_A B$. Again by flatness and the noetherian hypotheses,
\[
  \Extop_B^i(N_B,M_B)=\Extop_A^i(N,M)\otimes_A B.
\]
Thus $\Extop_A^i(N,M)=0$ implies $\Extop_B^i(N_B,M_B)=0$, and the converse
holds when $B$ is faithfully flat over $A$.



% SOURCE COMPONENT: 03_lemma_3_1.tex
\section*{3. Depth and Topological Properties}
\sgareaderlabel{III.3}{3}

\begin{sgastatement}{Lemma 3.1}\sgareaderlabel{III.3.1}{3.1}\srcpage{34}
Let $X$ be a topological space, $Y$ a closed subspace, and $F$ a sheaf of
abelian groups on $X$. Put $U=X-Y$. If $n$ is an integer, the following
conditions are equivalent:
\begin{enumerate}[label=\textup{(\roman*)},leftmargin=2.6em]
\item $\cH_Y^i(X,F)=0$ for $i<n$.\srcpage{35}
\item For every open subset $V$ of $X$, the group homomorphism
\[
  \HH^i(V,F)\longrightarrow\HH^i(V\cap U,F)
\]
is bijective for $i<n-1$ and injective for $i=n-1$.
\item For every open subset $V$ of $X$,
\[
  \HH_{Y\cap V}^i\bigl(V,F|_V\bigr)=0\qquad\text{for }i<n.
\]
\end{enumerate}
\end{sgastatement}

\begin{proof}
First, (ii) and (iii) are equivalent. Indeed, let $V$ be an open subset of
$X$, and put $X'=V$, $Y'=Y\cap V$, $F'=F|_V$, and $U'=X'-Y'$. Since $Y'$ is
closed in $X'$, there is an exact sequence
\[
  \HH_{Y'}^i(X',F')\longrightarrow \HH^i(X',F')
  \xrightarrow{\rho_i}\HH^i(U',F')
  \longrightarrow\HH_{Y'}^{i+1}(X',F').
\]
If the two outer terms vanish, $\rho_i$ is bijective; if the term on the left
vanishes, $\rho_i$ is injective. Thus (iii) implies (ii). Conversely, if
$i<n$, then $\HH_{Y'}^i(X',F')$ vanishes because $\rho_i$ is injective and
$\rho_{i-1}$ is surjective.

Next, (i) implies (iii). The local-to-global spectral sequence gives
\[
  \HH^p\bigl(X,\cH_Y^q(X,F)\bigr)\Longrightarrow
  \HH_Y^{p+q}(X,F).
\]
By hypothesis, $\cH_Y^q(X,F)=0$ for $q<n$; hence
$\HH_Y^{p+q}(X,F)=0$ whenever $p+q<n$.

Finally, (iii) implies (i), since (iii) says that the presheaf
\[
  V\longmapsto\HH_{Y\cap V}^i\bigl(V,F|_V\bigr)
\]
is zero. Its associated sheaf is therefore zero; because $Y$ is closed, that
associated sheaf is $\cH_Y^i(X,F)$.
\end{proof}


% SOURCE COMPONENT: 04_remark_3_2_proposition_3_3.tex
% Source-aligned production fragment: corrected French lines 1022--1060.
% Requires statement/proof environments and operators \prof, \cH, \cExt.

\begin{sgastatement}{Remark 3.2}\sgareaderlabel{III.3.2}{3.2}\srcpage{36}
The equivalence of (i) and (ii) was proved in Proposition \hyperref[I.2.13]{I.2.13}. As was
observed there, if $n\ge2$, the condition that $\rho_{n-1}$ be injective may
be omitted.\footnote{The original edition repeated a proof here, but that
proof was not entirely correct.}
\end{sgastatement}

\begin{sgastatement}{Proposition 3.3}\sgareaderlabel{III.3.3}{3.3}
Let $X$ be a locally noetherian prescheme, $Y$ a closed subprescheme of $X$,
and $F$ a coherent $\mathcal O_X$-module. The conditions of \hyperref[III.3.1]{Lemma 3.1} are
equivalent to each of the following conditions:
\begin{enumerate}[label=\textup{(\roman*)},leftmargin=2.8em]
\setcounter{enumi}{3}
\item For every $x\in Y$, one has $\prof F_x\ge n$.
\item For every coherent $\mathcal O_X$-module $G$ on $X$ whose support is
contained in $Y$,
\[
  \cExt_{\mathcal O_X}^i(G,F)=0\qquad\text{for }i<n.
\]
\item There exists a coherent $\mathcal O_X$-module $G$ whose support equals
$Y$ and such that
\[
  \cExt_{\mathcal O_X}^i(G,F)=0\qquad\text{for }i<n.
\]
\end{enumerate}
\end{sgastatement}

If $X$ is affine, \hyperref[III.2.4]{Proposition 2.4} supplies what is needed to establish the
equivalence of the three conditions in \hyperref[III.3.3]{Proposition 3.3}. These conditions are
local except for (v)$\Rightarrow$(vi); for that implication one may take
$G=\mathcal O_Y$ and invoke \hyperref[III.2.4]{Proposition 2.4} again. It therefore suffices to
prove (i)$\Rightarrow$(vi) and (v)$\Rightarrow$(i).

Let $J$ be the ideal of $Y$; it is a coherent sheaf of ideals. Put
$\mathcal O_{Y_m}=\mathcal O_X/J^{m+1}$. This is a coherent
$\mathcal O_X$-module whose support equals $Y$, and Theorem \hyperref[II.6]{II.6(b)} gives
\[
  \cH_Y^i(X,F)=\varinjlim_m
  \cExt_{\mathcal O_X}^i(\mathcal O_{Y_m},F).
\]
Thus (v) implies (i). Moreover,\srcpage{37} the transition morphisms in the
projective system of the $\mathcal O_{Y_m}$ are epimorphisms.

If the functor $\cExt^i$ is left exact in its first argument---at least when
that argument lies in the category of coherent $\mathcal O_X$-modules with
support contained in $Y$---then the transition morphisms of the inductive
system obtained by applying $\cExt^i$ to the $\mathcal O_{Y_m}$ are
injective. Condition (i) makes the limit zero, and therefore makes
$\cExt_{\mathcal O_X}^i(\mathcal O_{Y_m},F)$ zero for every $m$.

We argue by induction. The assertion is trivial for $n<0$. Suppose that
(i)$\Rightarrow$(vi) has been proved in degrees below $q$. Then
(i)$\Rightarrow$(v), so the long exact $\cExt$ sequence shows that
$\cExt^q$ is left exact in its first argument. Hence
$\cExt_{\mathcal O_X}^q(\mathcal O_{Y_m},F)=0$ for every $m$. Thus
(i)$\Rightarrow$(vi) holds through degree $q$.
\hfill$\square$\par

% SOURCE COMPONENT: 05_example_3_4.tex
\begin{sgastatement}{Example 3.4}\sgareaderlabel{III.3.4}{3.4}\srcpage{37}
Let $A$ be a noetherian local ring, $\mathfrak m$ its maximal ideal, $M$ a
finitely generated $A$-module, and $n$ an integer. Put
$X=\Spec(A)$, $Y=\{\mathfrak m\}$, and $U=X-Y$. Let $F$ be the sheaf associated
with $M$. The following conditions are equivalent:
\begin{enumerate}[label=\textup{(\arabic*)},leftmargin=2.8em]
\item $\prof M\ge n$.
\item The natural homomorphism
\[
  \HH^i(X,F)\longrightarrow\HH^i(U,F)
\]
is injective if $i=n-1$ and bijective if $i<n-1$.
\item $\Extop_A^i(k,M)=0$ if $i<n$, where $k=A/\mathfrak m$.
\item $\HH_Y^i(X,F)=0$ if $i<n$.
\end{enumerate}
\end{sgastatement}

Taking \hyperref[III.3.2]{Remark 3.2} into account, one obtains the following corollary.\srcpage{38}


% SOURCE COMPONENT: 06_corollary_3_5_theorem_3_6.tex
\begin{sgastatement}{Corollary 3.5}\sgareaderlabel{III.3.5}{3.5}\srcpage{38}
Let $X$ be a locally noetherian prescheme, $Y$ a closed subprescheme of $X$,
and $F$ a coherent $\mathcal O_X$-module. The following conditions are equivalent:
\begin{enumerate}[label=\textup{(\arabic*)},leftmargin=2.8em]
\item For every $x\in Y$, $\prof F_x\ge2$.
\item For every open subset $V$ of $X$, the natural homomorphism
\[
  \HH^0(V,F)\longrightarrow\HH^0\bigl(V\cap(X-Y),F\bigr)
\]
is bijective.
\end{enumerate}
\end{sgastatement}

\begin{sgastatement}{Theorem 3.6 (Hartshorne)}\sgareaderlabel{III.3.6}{3.6}
Let $X$ be a locally noetherian prescheme and $Y$ a closed subprescheme of
$X$. Suppose that $\prof\mathcal O_{X,x}\ge2$ for every $x\in Y$. Then the
natural map
\[
  \pi_0(X)\longrightarrow\pi_0(X-Y)
\]
is bijective.
\end{sgastatement}

\begin{proof}
Since $X$ is locally noetherian, it is locally connected. It therefore
suffices to prove that $X$ is connected if and only if $X-Y$ is connected.
A locally ringed space $(X,\mathcal O_X)$ is connected if and only if
$\HH^0(X,\mathcal O_X)$ is not a direct product of two nonzero rings. By
\hyperref[III.3.5]{Corollary 3.5}, applied with $F=\mathcal O_X$, the hypothesis implies that
\[
  \HH^0(X,\mathcal O_X)\longrightarrow
  \HH^0(X-Y,\mathcal O_X)
\]
is an isomorphism, and the conclusion follows.
\end{proof}


% SOURCE COMPONENT: 07_corollary_3_7_reduction.tex
\begin{sgastatement}{Corollary 3.7}\sgareaderlabel{III.3.7}{3.7}
Let $X$ be a locally noetherian prescheme. Let $d$ be an integer such that
$\dim\mathcal O_{X,x}\ge d$ implies $\prof\mathcal O_{X,x}\ge2$. If $X$ is
connected, then it is connected in codimension $d-1$: namely, if $X'$ and
$X''$ are two irreducible components of $X$, there is a sequence of
irreducible components of $X$,
\[
  X'=X_0,X_1,\ldots,X_n=X'',
\]
such that, for every $i$ with $0\le i<n$, the codimension in $X$ of
$X_i\cap X_{i+1}$ is at most $d-1$.
\end{sgastatement}

First note\srcpage{39} that if $X$ is Cohen--Macaulay, then $d=2$ has the
property just invoked. Recall also that the codimension of $Y$ in $X$ is
defined as the infimum of the dimensions of the local rings of $X$ at the
points of $Y$.

\begin{proofbridge}
Let $\mathcal F$ be the collection of closed subsets of $X$ whose codimension
in $X$ is at least $d$. This is an antifilter of closed subsets of $X$. A
closed subset $Y$ belongs to $\mathcal F$ if and only if, for every $y\in Y$,
there is an open neighborhood $V$ of $y$ such that $Y\cap V$ has codimension
at least $d$ in $V$.\footnote{Source correction: the French TeX says “an open
neighborhood $V$ of $X$”; the quantified point $y$ and the local condition
require “of $y$.\!”} Finally, if $X$ is connected and $Y\in\mathcal F$, then
$X-Y$ is connected by Hartshorne's theorem. The corollary therefore follows
from the next lemma, which is purely topological.
\end{proofbridge}


% SOURCE COMPONENT: 08_lemma_3_8.tex
\begin{sgastatement}{Lemma 3.8}\sgareaderlabel{III.3.8}{3.8}
Let $X$ be a connected, locally noetherian topological space, and let
$\mathcal F$ be an antifilter of closed subsets of $X$. Suppose that every
closed subset $Y\subset X$ which belongs locally to $\mathcal F$---that is,
for every $x\in X$ there are an open neighborhood $V$ of $x$ in $X$ and a
$Y'\in\mathcal F$ such that $V\cap Y=V\cap Y'$---belongs to $\mathcal F$.
The following conditions are equivalent:
\begin{enumerate}[label=\textup{(\roman*)},leftmargin=2.7em]
\item For every $Y\in\mathcal F$, the complement $X-Y$ is connected.
\item If $X'$ and $X''$ are distinct irreducible components, there is a
sequence $X_0,X_1,\ldots,X_n$ of irreducible components of $X$ such that
$X'=X_0$, $X''=X_n$, and
\[
  X_i\cap X_{i+1}\notin\mathcal F\qquad(0\le i<n).
\]
\end{enumerate}
\end{sgastatement}

\noindent\textit{Source correction.} The corrected French branch prints
$1\le i<n$ in (ii), omitting the first adjacent pair. Both the chain's
indexing and the proof below require $0\le i<n$; the target displays that
range rather than silently reproducing the defect.

\begin{proof}
Assume (ii), and let $Y\in\mathcal F$. We must prove that $U=X-Y$ is
connected. If $U'$ and $U''$ are two irreducible components of $U$, there are
irreducible components $X'$ and $X''$ of $X$ such that
$X'\cap U=U'$ and $X''\cap U=U''$. Choose a sequence
$X_0,\ldots,X_n$ having the property in (ii), and put
$U_i=X_i\cap U$ for $0\le i\le n$. The $U_i$ are irreducible components of
$U$, and $U_i\cap U_{i+1}$ is nonempty for $0\le i<n$. Otherwise
$X_i\cap X_{i+1}\subset Y$ would belong to $\mathcal F$,\srcpage{40} contrary to the
choice of the $X_i$. Hence $U$ is connected.

Conversely, assume (i). Let
\[
  Y=\bigcup_{X',X''}(X'\cap X''),
\]
where the union is over distinct irreducible components $X',X''$ of $X$ for
which $X'\cap X''\in\mathcal F$. This family of intersections is locally
finite because $X$ is locally noetherian; its members are closed, so $Y$ is
closed. Moreover, $Y$ belongs locally to $\mathcal F$, and hence
$Y\in\mathcal F$. Thus $U=X-Y$ is connected.

Let $X'$ and $X''$ be distinct irreducible components of $X$, and let $U'$
and $U''$ be their traces on $U$; by the construction of $Y$, these traces
are nonempty. They are irreducible components of $U$. Since $U$ is connected
and locally noetherian, there is a sequence $U_0,\ldots,U_n$ of irreducible
components of $U$ such that $U_0=U'$, $U_n=U''$, and, for $0\le i<n$,
\[
  U_i\cap U_{i+1}\ne\varnothing
  \quad\text{and}\quad U_i\cap U_{i+1}\ne U_i.
\]
Let $X_0,\ldots,X_n$ be the irreducible components of $X$ satisfying
$X_i\cap U=U_i$. If $X_i\cap X_{i+1}\in\mathcal F$, then by the construction
of $Y$ one would have either $U_i\cap U_{i+1}=\varnothing$ or
$U_i=U_{i+1}$, contradicting the choice of the $U_i$.
\end{proof}


% SOURCE COMPONENT: 09_corollary_3_9.tex
\begin{sgastatement}{Corollary 3.9}\sgareaderlabel{III.3.9}{3.9}
Let $A$ be a noetherian local ring. Suppose that, for every prime ideal
$\mathfrak p$ of $A$,
\[
  \bigl(\dim A_{\mathfrak p}\ge2\bigr)
  \Longrightarrow
  \bigl(\prof A_{\mathfrak p}\ge2\bigr).
\]
Suppose also that $A$ satisfies the chain condition.\footnote{Cf. EGA
$0_{\mathrm{IV}}$, 14.3.2.} Then every minimal prime ideal $\mathfrak p$ of
$A$ satisfies $\dim(A/\mathfrak p)=\dim A$. Equivalently, all irreducible
components of $\Spec A$ have the same dimension, namely $\dim A$.
\end{sgastatement}

Indeed, put $X=\Spec A$. If $X'$ and $X''$ are two irreducible components of
$X$, \hyperref[III.3.7]{Corollary 3.7} joins them by a chain with the stated properties. It is
therefore enough to prove that two successive components have the same
dimension, which follows from the second hypothesis.


% SOURCE COMPONENT: 10_example_3_10_remark_3_11.tex
\begin{sgastatement}{Example 3.10}\sgareaderlabel{III.3.10}{3.10}\srcpage{41}
Let $X$ be the union of two complementary linear subspaces of respective
dimensions $2$ and $3$ in a vector space of dimension $5$. More precisely,
let $X=\Spec A$, where
\[
  A=B/(\mathfrak p\cap\mathfrak q),\qquad
  B=k[X_1,\ldots,X_5],
\]
$\mathfrak p=(X_1,X_2,X_3)$, and $\mathfrak q=(X_4,X_5)$. Removing the
intersection point $x$ of the two linear subspaces disconnects $X$.
Consequently, the depth of $\mathcal O_{X,x}$ is $1$, since it cannot be at
least $2$ by \hyperref[III.3.6]{Theorem 3.6}. The failure of the preceding corollary's
equidimensionality conclusion gives another reason.

More generally, let $X$ be the union of two linear subspaces of dimensions
$p,q\ge2$ in a vector space of dimension $p+q$. Under no embedding of $X$
into a regular scheme is $X$ even set-theoretically a complete intersection
at the origin. Indeed, after possibly modifying its scheme structure without
changing the underlying topological space near the origin, $X$ would then be
Cohen--Macaulay and hence of depth at least $2$ at the origin, which it is not.
\end{sgastatement}

\begin{sgastatement}{Remark 3.11}\sgareaderlabel{III.3.11}{3.11}
Let $X$ be a locally noetherian prescheme, $Y$ a closed subprescheme of $X$,
and $F$ an $\mathcal O_X$-module. Depth is a purely topological notion,
expressed in terms of the vanishing of $\cH_Y^i(X,F)$ for $i<n$. One also
wishes to study these sheaves for a fixed $i$, or for $i>n$. The following
result addresses this question.
\end{sgastatement}


% SOURCE COMPONENT: 11_lemma_3_12_to_end.tex
\begin{sgastatement}{Lemma 3.12}\sgareaderlabel{III.3.12}{3.12}
Let $m$ be an integer. The equality $\cH_Y^i(X,F)=0$ for every $i>m$ and
every coherent $F$ holds if and only if it holds for $F=\mathcal O_X$.
\end{sgastatement}

By passage to a direct limit, it then holds for every quasi-coherent sheaf.
For example, if $Y$ can be locally described by $m$ equations---or, as one
says, if $Y$ is locally a set-theoretic complete intersection (as happens,
for example, when $X$ and $Y$ are nonsingular)---the computation of
$\cH_Y^i(X,F)$ by the Koszul complex shows that these sheaves vanish for
$i>m$. This fact was used implicitly in \hyperref[III.3.10]{Example 3.10}. The cohomological
condition is nevertheless insufficient, as the following example shows.

\begin{sgastatement}{Example 3.13}\sgareaderlabel{III.3.13}{3.13}\srcpage{42}
Let $X=\Spec A$, where $A$ is a normal noetherian local ring of dimension
$2$, and let $Y$ be a curve in $X$. The complement $X-Y$ can be shown to be
affine. Hence, for $i>1$, the supported-cohomology exact sequence gives
\[
  \HH_Y^i(X,\mathcal O_X)
  \simeq \HH^{i-1}(X-Y,\mathcal O_X)=0.
\]
Consequently the corresponding local-cohomology sheaf
$\cH_Y^i(\mathcal O_X)$ vanishes. Nevertheless, one can construct a curve
which is not described by a single equation.
\end{sgastatement}

\noindent\textit{Source correction.} The corrected French branch and its
rendered PDF write a sheaf-local-cohomology term with support $X-Y$ in the
middle of the displayed isomorphism. The stated reason is instead the
vanishing of $\HH^{i-1}(X-Y,\mathcal O_X)$ because $X-Y$ is affine; the
target displays the supported/global exact-sequence formula that this reason
requires and then states the sheaf consequence separately.

We shall seek\footnote{Cf. Expos\'e VIII.} conditions ensuring that
$\cH_Y^i(X,F)$ is coherent for a fixed $i$. In general it is not: for
example, consider $\HH_{\mathfrak m}^n(A)$ when $A$ is a noetherian local
ring of dimension $n>0$. If $A$ is a discrete valuation ring with fraction
field $K$, then
\[
  \HH_{\mathfrak m}^1(A)\simeq K/A,
\]
which is not a finitely generated $A$-module.

To illuminate the question, it is equivalent to the following problem. Let
$f:U\to X$ be an open immersion and $G$ a coherent sheaf on $U$. Find
criteria ensuring that the higher direct images $\RR^i f_*G$ are coherent
sheaves on $X$ for a fixed $i$. Such conditions are necessary for the use of
formal geometry encountered in Expos\'e IX and the subsequent expos\'es.

% END INCLUDED active_source/parts/03_expose_III.tex

% BEGIN INCLUDED active_source/parts/04_expose_IV.tex
% SOURCE COMPONENT: 01_unit_IV_1_open\part.tex

\makeatletter
\providecommand{\sgareaderlabel}[2]{%
  \begingroup
  \def\@currentlabel{#2}%
  \phantomsection\label{#1}%
  \endgroup}
\makeatother

\chapter{Dualizing Modules and Functors}
\label{IV}
\section*{1. Generalities on Functors of Modules}
\sgareaderlabel{IV.1}{1}

Fix\srcpage{43}\sgareaderlabel{pIV.1}{1} a commutative noetherian ring $A$. Let $\Ccat$ be the
category of finitely generated $A$-modules, $\Ccat'$ the category of all
$A$-modules, and $\Ab$ the category of abelian groups. This section studies
certain properties of additive functors
\[
  T:\Ccat^{\mathrm{op}}\longrightarrow\Ab.
\]

For $M\in\Ob\Ccat$, the group $T(M)$ has a canonical $A$-module structure.
If $f_M:M\to M$ denotes multiplication by $f\in A$, then $f$ acts on $T(M)$
as $T(f_M)$. In other words, $T$ factors as
\[
  \Ccat^{\mathrm{op}}\xrightarrow{T_0}\Ccat'
  \longrightarrow\Ab,
\]
where the second arrow is the forgetful functor. We henceforth regard
$T(M)$ as endowed with this $A$-module structure.

Compose the isomorphism $M\xrightarrow{\sim}\Homop_A(A,M)$ with the map
\[
  \Homop_A(A,M)\longrightarrow\Homop_A\bigl(T(M),T(A)\bigr)
\]
induced by $T$. This gives the following mutually adjoint forms of the same
pairing:
\begin{align*}
  M&\longrightarrow\Homop_A\bigl(T(M),T(A)\bigr),\\
  M\times T(M)&\longrightarrow T(A),\\
  T(M)&\longrightarrow\Homop_A\bigl(M,T(A)\bigr).
\end{align*}
They define a natural morphism of contravariant functors
\[
  \varphi_T:T\longrightarrow\Homop_A\bigl(-,T(A)\bigr).
\]

\begin{sgastatement}{Proposition 1.1}\sgareaderlabel{IV.1.1}{1.1}\srcpage{44}
The following two properties are equivalent:
\begin{enumerate}[label=\textup{\roman*)},leftmargin=2.7em]
\item $\varphi_T$ is an isomorphism of functors.
\item $T$ is left exact.
\end{enumerate}
\end{sgastatement}

The implication (i)$\Rightarrow$(ii) is immediate. For the converse, use the
following fact: if $u:F\to F'$ is a morphism between additive left-exact
functors from $\Ccat^{\mathrm{op}}$ to $\Ab$ and $u(A)$ is an isomorphism,
then $u$ is an isomorphism. Here one uses that $A$ is noetherian, so every
finitely generated $A$-module is finitely presented.

\begin{sgastatement}{Remark 1.2}\sgareaderlabel{IV.1.2}{1.2}
In particular, the representable functors
$T:{\Ccat'}^{\mathrm{op}}\to\Ab$ are precisely those which commute with
arbitrary inverse limits in ${\Ccat'}^{\mathrm{op}}$ (indexed by a
preordered set, not necessarily filtered).
\end{sgastatement}

Let $\operatorname{Fun}(\Ccat^{\mathrm{op}},\Ab)_{\mathrm{lex}}$ denote the
full subcategory of additive functors that are left exact. We have proved an
equivalence of categories
\[
  \Ccat'\xrightarrow{\sim}
  \operatorname{Fun}(\Ccat^{\mathrm{op}},\Ab)_{\mathrm{lex}},
\]
with mutually quasi-inverse functors
\[
  H\longmapsto\Homop_A(-,H),
  \qquad
  T\longmapsto T(A).
\]

% SOURCE COMPONENT: 02_unit_IV_1_3\part.tex

Let $J$ be an ideal of $A$, put $Y=V(J)\subset\Spec A$, and let
$\Ccat_Y$ be the full subcategory of $\Ccat$ consisting of finitely generated
$A$-modules $M$ with $\Supp M\subset Y$. Then
\[
  \Ccat_Y=\bigcup_n\Ccat^{(n)},
\]
where $\Ccat^{(n)}$ is the full subcategory of modules satisfying $J^nM=0$.

\begin{sgastatement}{Proposition 1.3}\sgareaderlabel{IV.1.3}{1.3}\srcpage{45}
With this notation, let
$T:\Ccat_Y^{\mathrm{op}}\to\Ab$ be an additive functor, and put
\[
  H=\varinjlim_n T(A/J^n).
\]
There is a natural morphism
\[
  \varphi_T:T\longrightarrow\Homop_A(-,H),
\]
and the following conditions are equivalent:
\begin{enumerate}[label=\textup{\roman*)},leftmargin=2.7em]
\item $\varphi_T$ is an isomorphism.
\item $T$ is left exact.
\end{enumerate}
\end{sgastatement}

\noindent\textit{Source correction.} The proposition line prints
$T:\Ccat_Y\to\Ab$, while its target, its restrictions, and the concluding
category equivalence are contravariant. The target restores the required
opposite category.

\begin{proof}
First define $\varphi_T$. For $M\in\Ob\Ccat_Y$, choose $n$ with $J^nM=0$.
Then $M$ is an $A/J^n$-module. If $T_n$ is the restriction of $T$ to
${\Ccat^{(n)}}^{\mathrm{op}}$ and $H_n=T(A/J^n)$, the construction of
\hyperref[IV.1.1]{Proposition 1.1} gives
\[
  T_n\longrightarrow\Homop_A(-,H_n).
\]
Composing with $H_n\to H$ gives, for every such $M$,
\[
  T(M)=T_n(M)\longrightarrow\Homop_A(M,H),
\]
and these maps define $\varphi_T$.

It is clear that (i) implies (ii). Conversely, suppose $T$ is left exact and
let $M\in\Ob\Ccat^{(n)}$. \hyperref[IV.1.1]{Proposition 1.1} applied over $A/J^n$ gives
\[
  T_n(M)\xrightarrow{\sim}\Homop_A(M,H_n).
\]
For every $n'\ge n$, \hyperref[IV.1.1]{Proposition 1.1} applied over $A/J^{n'}$ gives a
compatible isomorphism with $H_{n'}$. Hence
\[
  T(M)\simeq\varinjlim_{n'\ge n}\Homop_A(M,H_{n'}).
\]
Because this is a filtered direct limit and $M$ is finitely presented,
\[
  \varinjlim_{n'\ge n}\Homop_A(M,H_{n'})
  \xrightarrow{\sim}
  \Homop_A\!\left(M,\varinjlim_{n'\ge n}H_{n'}\right)
  =\Homop_A(M,H).
\]

Consequently, if $\Ccat'_Y$ denotes the category of all $A$-modules supported
in $Y$, without a finite-generation condition, there is a natural equivalence
\[
  \Ccat'_Y\xrightarrow{\sim}
  \operatorname{Fun}(\Ccat_Y^{\mathrm{op}},\Ab)_{\mathrm{lex}}.
\]
\end{proof}

% SOURCE COMPONENT: 03_unit_IV_1_4\part.tex
\subsection*{Application}

Retain the preceding notation, and let\srcpage{46}
\[
  T^*:\Ccat_Y^{\mathrm{op}}\longrightarrow\Ab
\]
be an exact $\partial$-functor. For every $i\in\mathbb Z$, put
\[
  H_n^i=T^i(A/J^n),\qquad H^i=\varinjlim_n H_n^i.
\]

\begin{sgastatement}{Theorem 1.4}\sgareaderlabel{IV.1.4}{1.4}
Let $n\in\mathbb Z$. Suppose there is an $i_0\in\mathbb Z$ such that
$T^i=0$ for every $i<i_0$. Then the following three conditions are equivalent:
\begin{enumerate}[label=\textup{\roman*)},leftmargin=2.7em]
\item $T^i=0$ for every $i<n$.
\item $H^i=0$ for every $i<n$.
\item There exists an $M_0\in\Ccat_Y$ such that $\Supp M_0=Y$ and
$T^i(M_0)=0$ for every $i<n$.
\end{enumerate}
\end{sgastatement}

\begin{proof}
Clearly (i) implies (ii) and (iii), taking $M_0=A/J$. We prove
(ii)$\Rightarrow$(i) by induction on $n$. It is true for $n=i_0$; suppose it
has been proved through $n$. If $H^i=0$ for every $i<n+1$, the induction
hypothesis gives $T^i=0$ for $i<n$. The vanishing of $T^{n-1}$ makes $T^n$
left exact, so \hyperref[IV.1.3]{Proposition 1.3} gives
\[
  T^n\xrightarrow{\sim}\Homop_A(-,H^n)=0.
\]

It remains to prove (iii)$\Rightarrow$(ii), again by induction. The assertion
is true for $n=i_0$. Suppose it has been proved through $n$, and let
$M_0\in\Ccat_Y$ satisfy $\Supp M_0=Y$ and $T^i(M_0)=0$ for every $i<n+1$.
The induction hypothesis gives $H^i=0$ for $i<n$; it remains to show
$H^n=0$. The preceding implication gives $T^{n-1}=0$, and therefore
\[
  T^n\xrightarrow{\sim}\Homop_A(-,H^n).
\]
Moreover,
\[
  \Ass\Homop_A(M_0,H^n)
  =\Supp M_0\cap\Ass H^n
  =\Ass H^n,
\]
because
\[
  \Ass H^n\subset\Supp H^n\subset Y=\Supp M_0.
\]
Thus
\[
  T^n(M_0)=0
  \quad\Longleftrightarrow\quad
  \Ass H^n=\varnothing
  \quad\Longleftrightarrow\quad
  H^n=0,
\]
which completes the proof.
\end{proof}

% SOURCE COMPONENT: 04_unit_IV_2_1\part.tex

\section*{2. Characterization of Exact Functors}
\sgareaderlabel{IV.2}{2}

The ring $A$ remains commutative and noetherian. Retain the notation of
\hyperref[IV.1.3]{Proposition 1.3}:\srcpage{47}
\[
  Y=V(J),\qquad
  T:\Ccat_Y^{\mathrm{op}}\longrightarrow\Ab,\qquad
  H=\varinjlim_n T(A/J^n).
\]
Assume that $T$ is left exact, so that
\[
  T(M)\xrightarrow{\sim}\Homop_A(M,H).
\]

\begin{sgastatement}{Proposition 2.1}\sgareaderlabel{IV.2.1}{2.1}
The following properties are equivalent:
\begin{enumerate}[label=\textup{\roman*)},leftmargin=2.7em]
\item $T$ is exact.
\item $H$ is injective in $\Ccat'$, the category of all $A$-modules.
\end{enumerate}
\end{sgastatement}

\begin{proof}
Only (i)$\Rightarrow$(ii) needs proof. Suppose the restriction of
$\Homop_A(-,H)$ to $\Ccat_Y^{\mathrm{op}}$ is exact. Since $A$ is
noetherian, it is enough to prove the following extension property: whenever
$N\subset M$ are finitely generated $A$-modules, every homomorphism
$f:N\to H$ extends to a homomorphism $\bar f:M\to H$.

By the definition of $H$ and finite generation of $N$, there is an integer
$n$ such that $J^n f(N)=0$. Give $M$ and $N$ their $J$-adic topologies. By
Krull's theorem, the $J$-adic topology on $N$ is equivalent to the topology
induced from the $J$-adic topology on $M$. Hence there is a
$V=J^kM$ such that
\[
  U=V\cap N\subset J^nN.
\]
Thus $U\subset\ker f$, and $f$ factors as
\[
\begin{tikzcd}[column sep=large]
N \arrow[r,two heads] \arrow[dr,"f"'] & N/U \arrow[d,"u"]\\
& H.
\end{tikzcd}
\]
Both $N/U$ and $M/V$ belong to $\Ccat_Y$. Exactness therefore extends $u$
to a map $\bar u:M/V\to H$:
\[
\begin{tikzcd}[column sep=large]
N/U \arrow[r,hook] \arrow[dr,"u"'] & M/V \arrow[d,"\bar u"]\\
& H.
\end{tikzcd}
\]
The composite $M\to M/V\xrightarrow{\bar u}H$ is the required extension
$\bar f$.\srcpage{48}
\end{proof}

% SOURCE COMPONENT: 05_unit_IV_2_2\part.tex

Retain the notation $Y=V(J)$ and let $\Ccat_Y$ denote the category of
finitely generated $A$-modules whose support is contained in $Y$.

\begin{sgastatement}{Corollary 2.2}\sgareaderlabel{IV.2.2}{2.2}\srcpage{48}
Let $K$ be an injective $A$-module. Then the submodule
\[
  H^0_J(K)=\{x\in K: J^n x=0\text{ for some }n\geq 0\}
\]
of elements annihilated by a suitable power of $J$ is injective.
\end{sgastatement}

\begin{proof}
By \hyperref[IV.2.1]{Proposition 2.1}, it is enough to verify that the restriction to
$\Ccat_Y^{\mathrm{op}}$ of
$\Homop_A(-,H^0_J(K))$ is exact. Let $M\in\Ccat_Y$. There is an integer
$k$ such that $J^kM=0$, and the inclusion
\[
  \Homop_A(M,H^0_J(K))\longrightarrow \Homop_A(M,K)
\]
is an isomorphism: every map $M\to K$ has image annihilated by $J^k$ and
therefore lands in $H^0_J(K)$. The assertion follows because
$\Homop_A(-,K)$ is exact.
\end{proof}

% SOURCE COMPONENT: 06_unit_IV_3_1\part.tex

\section*{3. The Case Where $T$ Is Left Exact and Every $T(M)$ Is Finitely Generated}
\sgareaderlabel{IV.3}{3}

As above, let\srcpage{48}
\[
  T:\Ccat_Y^{\mathrm{op}}\longrightarrow\Ab.
\]
From now on, assume that $T$ is left exact and factors through
$\Ccat_Y$:
\[
\begin{tikzcd}[column sep=large,row sep=small]
\Ccat_Y^{\mathrm{op}} \arrow[rr,"T"] \arrow[dr,swap,"\widetilde T"]
  && \Ab \\
& \Ccat_Y \arrow[ur,swap,"\text{forget}"] &
\end{tikzcd}
\]
Thus $T(T(M))=(T\circ T)(M)$ is defined. Under the representation
$T(M)\simeq\Homop_A(M,H)$, the canonical evaluation map
\[
  M\longrightarrow
  \Homop_A\bigl(\Homop_A(M,H),H\bigr)
\]
defines a natural morphism\srcpage{49}
\[
  \eta_M:M\longrightarrow (T\circ T)(M).
\]

\begin{sgastatement}{Proposition 3.1}\sgareaderlabel{IV.3.1}{3.1}
Assume, as before, that $A$ is noetherian, and assume in addition that
$A/J$ is artinian. The following conditions are equivalent:
\begin{enumerate}[label=\textup{\roman*)},leftmargin=2.7em]
\item $T$ is left exact and, for every $M\in\Ccat_Y$, the module $T(M)$
is finitely generated and $\eta_M:M\to(T\circ T)(M)$ is an isomorphism.

\item $T$ is exact and, for every residue field $k=A/\mathfrak m$ of a
maximal ideal $\mathfrak m\supset J$, one has $T(k)\simeq k$.

\item There is an isomorphism
$T\simeq\Homop_A(-,H)$ with $H$ injective and, for every $k$ as in
\emph{(ii)}, one has $\Homop_A(k,H)\simeq k$.

\item $T$ is exact and $\ellA T(M)=\ellA M$ for every $M\in\Ccat_Y$.
\end{enumerate}
\end{sgastatement}

\begin{proof}
The equivalence of \emph{(ii)} and \emph{(iii)} follows from
\hyperref[IV.2.1]{Proposition 2.1}.

\smallskip\noindent
\emph{(ii)$\Rightarrow$(iv).}
If $M\in\Ccat_Y$, then $M$ is an $A/J^n$-module for some $n$; since
$A/J^n$ is artinian, $M$ has finite length. We argue by induction on
$\ellA M$. When $\ellA M=1$, the module $M$ is one of the residue fields
occurring in \emph{(ii)}, so the assertion holds. If $\ellA M>1$, choose
a nonzero proper submodule $M'\subset M$ and put $M''=M/M'$. Exactness
of $T$ turns
\[
  0\longrightarrow M'\longrightarrow M\longrightarrow M''\longrightarrow0
\]
into
\[
  0\longrightarrow T(M'')\longrightarrow T(M)
   \longrightarrow T(M')\longrightarrow0.
\]
Both end terms have smaller length, and therefore
\[
\begin{aligned}
  \ellA T(M)
    &=\ellA T(M')+\ellA T(M'')\\
    &=\ellA M'+\ellA M''=\ellA M.
\end{aligned}
\]

\smallskip\noindent
\emph{(ii)$\Rightarrow$(i).}
By \emph{(iv)}, $T(M)$ has finite length and hence is finitely generated.
It remains to prove that every $\eta_M$ is an isomorphism. Again use
induction on $\ellA M$. For a residue field, $T(k)$ and $(T\circ T)(k)$
both have length one; evaluation at $1\in k$ shows that $\eta_k$ is
nonzero, hence it is an isomorphism. For an exact sequence as above,
naturality gives the
commutative diagram with exact rows\srcpage{50}
\[
\begin{tikzcd}[column sep=small]
0 \arrow[r] & M' \arrow[r] \arrow[d,"\eta_{M'}"']
  & M \arrow[r] \arrow[d,"\eta_M"']
  & M'' \arrow[r] \arrow[d,"\eta_{M''}"'] & 0\\
0 \arrow[r] & (T\!\circ T)(M') \arrow[r]
  & (T\!\circ T)(M) \arrow[r]
  & (T\!\circ T)(M'') \arrow[r] & 0.
\end{tikzcd}
\]
The outer vertical arrows are isomorphisms by induction, hence so is
$\eta_M$.

\smallskip\noindent
\emph{(i)$\Rightarrow$(ii).}
Let
\[
  0\longrightarrow M'\longrightarrow M\longrightarrow M''\longrightarrow0
\]
be exact in $\Ccat_Y$, and let $Q$ be the cokernel of
$T(M)\to T(M')$. Left exactness gives the exact sequence
\[
  0\longrightarrow T(M'')\longrightarrow T(M)
   \longrightarrow T(M')\longrightarrow Q\longrightarrow0.
\]
Applying $T$ at the left end and using naturality of evaluation yields
\[
\begin{tikzcd}[column sep=small]
0 \arrow[r] & T(Q) \arrow[r] & (T\!\circ T)(M') \arrow[r]
  & (T\!\circ T)(M)\\
&& M' \arrow[u,"\eta_{M'}","\sim"'] \arrow[r,hook]
  & M \arrow[u,"\eta_M"',"\sim"] .
\end{tikzcd}
\]
Thus $T(Q)=0$, and condition \emph{(i)} gives
$Q\simeq(T\circ T)(Q)=0$. Therefore $T(M)\to T(M')$ is surjective and
$T$ is exact.

Now let $k=A/\mathfrak m$ with $\mathfrak m\supset J$. The $A$-module
$T(k)$ is a finite-dimensional $k$-vector space. It is nonzero because
$k\simeq(T\circ T)(k)$. Choose a decomposition $T(k)\simeq k\oplus V$.
Additivity and biduality give
\[
\begin{aligned}
  (T\circ T)(k)
    &\simeq T(k)\oplus T(V)\\
    &\simeq k\oplus V\oplus T(V)\simeq k,
\end{aligned}
\]
so comparison of finite lengths forces $V=0$. Hence $T(k)\simeq k$.

\smallskip\noindent
\emph{(iv)$\Rightarrow$(iii).}
\hyperref[IV.2.1]{Proposition 2.1} supplies an injective module $H$ and an isomorphism
\[
  T\simeq\Homop_A(-,H).
\]
It remains only to check the residue fields.
For $k=A/\mathfrak m$ as above, $T(k)$ has length one, hence is isomorphic
to a residue field $k'=A/\mathfrak m'$. Moreover
\[
  \Supp k'=\Supp\Homop_A(k,H)\subseteq\Supp k=\{\mathfrak m\}.
\]
Thus $\mathfrak m'=\mathfrak m$ and $k'\simeq k$, as required.
\end{proof}

% SOURCE COMPONENT: 07_unit_IV_3_2_4_2\part.tex

\begin{sgastatement}{Remark 3.2}\sgareaderlabel{IV.3.2}{3.2}
One can\srcpage{51} show that condition \emph{(iv)} is equivalent to the condition
\begin{quote}
\emph{(iv)$'$} For every $M\in\operatorname{Ob}\Ccat_Y$, one has
\[
  \lengthop T(M)=\lengthop M.
\]
\end{quote}
\end{sgastatement}

\section*{4. Dualizing Module. Dualizing Functor}
\sgareaderlabel{IV.4}{4}

\begin{sgastatement}{Definition 4.1}\sgareaderlabel{IV.4.1}{4.1}
Let $A$ be a noetherian local ring and let $\mathfrak m$ be its maximal
ideal. A \emph{dualizing functor} for $A$ is a functor
\[
  T:\Ccat_{\mathfrak m}^{\mathrm{op}}\longrightarrow\Ab,
\]
where $\Ccat_{\mathfrak m}$ denotes $\Ccat_Y$ when
$Y=V(\mathfrak m)$, that satisfies the equivalent conditions of
\hyperref[IV.3.1]{Proposition 3.1}. An $A$-module $I\in\operatorname{Ob}\Ccat'_{\mathfrak m}$ is said to be \emph{dualizing for
$A$} if the functor
\[
  M\longmapsto\Homop_A(M,I)
\]
is dualizing.
\end{sgastatement}

\hyperref[IV.4.1]{Definition 4.1} can be generalized without assuming that $A$ is local.

\begin{sgastatement}{Definition 4.2}\sgareaderlabel{IV.4.2}{4.2}
Let $A$ be a noetherian ring, and let $\Cbar$ be the full subcategory of
$\Ccat$ consisting of the $A$-modules of finite length. A
\emph{dualizing functor} is an exact $A$-linear functor
\[
  T:\Cbar^{\mathrm{op}}\longrightarrow\Cbar
\]
such that the morphism of functors
\[
  \mathrm{id}\longrightarrow T\circ T
\]
is an isomorphism.
\end{sgastatement}

We shall prove an existence theorem and also prove that the module $I$
representing such a functor is locally artinian. We shall further show
that, for every maximal ideal $\mathfrak m$ of $A$, the
$\mathfrak m$-primary component of the socle of $I$ has length one.

% SOURCE COMPONENT: 08_unit_IV_4_3_5\part.tex

\begin{sgastatement}{Proposition 4.3}\sgareaderlabel{IV.4.3}{4.3}
Let\srcpage{51} $A$ and $B$ be noetherian local rings with maximal ideals
$\mA$ and $\mB$, and suppose that $B$ is a finite $A$-algebra. If $I$
is a dualizing module for $A$, then $\Homop_A(B,I)$ is a dualizing
module for $B$.
\end{sgastatement}

\begin{proof}
Let\srcpage{52}
\[
  R:\Ccat_{\mB}\longrightarrow\Ccat_{\mA}
\]
be the restriction-of-scalars functor; it is exact. Let $T$ be a
dualizing functor for $A$,
\[
  T:\Ccat_{\mA}^{\mathrm{op}}\longrightarrow\Ab.
\]
It is exact, and for every $M\in\operatorname{Ob}\Ccat_{\mA}$ the
natural morphism
\[
  M\longrightarrow (T\circ T)(M)
\]
is an isomorphism. Consequently, the contravariant composite $T\circ R$
is a dualizing functor for $B$. If $I$ represents $T$, the classical
formula
\[
  \Homop_A(M,I)\simeq
  \Homop_B\bigl(M,\Homop_A(B,I)\bigr),
\]
valid for every $B$-module $M$, shows that $\Homop_A(B,I)$ is a
dualizing module for $B$.
\end{proof}

\begin{sgastatement}{Corollary 4.4}\sgareaderlabel{IV.4.4}{4.4}
Let $A$ be a noetherian local ring, let $\mathfrak a$ be an ideal of
$A$, and put $B=A/\mathfrak a$. If $I$ is a dualizing module for $A$,
then the annihilator of $\mathfrak a$ in $I$ is a dualizing module for
$B$.
\end{sgastatement}

\begin{sgastatement}{Lemma 4.5}\sgareaderlabel{IV.4.5}{4.5}
Let $A$ be a noetherian local ring and let $I$ be a locally artinian
$A$-module. There is a canonical isomorphism
\[
  I\xrightarrow{\sim}\widehat I=I\otimes_A\widehat A.
\]
\end{sgastatement}

\begin{proof}
Let $I_n$ be the annihilator of $\mathfrak m^n$ in $I$, where
$\mathfrak m$ is the maximal ideal of $A$. To say that $I$ is locally
artinian means that $I$ is the direct limit of the $I_n$ and that each
$I_n$ has finite length. Since tensor product commutes with direct
limits, it is enough to treat each $I_n$ separately. The relation
$\mathfrak m^nI_n=0$ gives, for every $p\geq n$,
\[
  I_n\xrightarrow{\sim}I_n\otimes_A A/\mathfrak m^p
      \xrightarrow{\sim}I_n\otimes_A\widehat A.
\]
The second isomorphism uses
$\widehat A/\mathfrak m^p\widehat A\simeq A/\mathfrak m^p$. Passing
to the direct limit over $n$, and again using commutation with tensor
product, gives
\[ I\xrightarrow{\sim}I\otimes_A\widehat A. \]

It follows that\srcpage{53} restriction of scalars from $\widehat A$ to
$A$ and extension of scalars induce mutually quasi-inverse
\emph{equivalences} between the category of locally artinian
$\widehat A$-modules and the category of locally artinian $A$-modules.
\end{proof}

% SOURCE COMPONENT: 09_unit_IV_4_6_9\part.tex

\begin{sgastatement}{Proposition 4.6}\sgareaderlabel{IV.4.6}{4.6}
Let\srcpage{53} $A$ be a noetherian local ring, let $\widehat A$ be its completion,
let $I$ be a dualizing module for $A$ (respectively, for $\widehat A$), and
let $J$ be the completed module\footnote{Here ``the completed module'' means
$\widehat I=I\otimes_A\widehat A$ (cf. \hyperref[IV.4.5]{Lemma 4.5}): this is $I$ endowed with
its canonical $\widehat A$-module structure, not its $\mideal$-adic
completion. For example, if $p$ is prime and
$A=\widehat A=\mathbb Z_p$ is the ring of $p$-adic integers, then the injective envelope of
$k=\mathbb Z/p\mathbb Z$ is the discrete $\mathbb Z_p$-module
$\mathbb Q_p/\mathbb Z_p$, whose $p$-adic completion is zero.}
of $I$ (respectively, the $A$-module obtained by restriction of scalars).
Then $J$ is a dualizing module for $\widehat A$ (respectively, for $A$).
Moreover, the underlying abelian groups of $I$ and $J$ are isomorphic.
\end{sgastatement}

\begin{proof}
One merely observes that the equivalence between the category of locally
artinian $A$-modules and the category of locally artinian
$\widehat A$-modules induces an isomorphism between the bifunctors
$\Homop_A(-,-)$ and $\Homop_{\widehat A}(-,-)$, and that the
characterization of a dualizing functor or a dualizing module involves only
these bifunctors.
\end{proof}

\begin{sgastatement}{Theorem 4.7}\sgareaderlabel{IV.4.7}{4.7}
Let $A$ be a noetherian local ring.
\begin{enumerate}[label=\textup{\alph*)},leftmargin=2.5em]
\item A dualizing module $I$ always exists.
\item Any two dualizing modules are isomorphic (by a noncanonical
isomorphism).
\item An $A$-module $I$ is dualizing if and only if it is an injective
envelope of the residue field $k$ of $A$.
\end{enumerate}
\end{sgastatement}

\begin{sgastatement}{Remark 4.8}\sgareaderlabel{IV.4.8}{4.8}
\hyperref[IV.4.6]{Proposition 4.6} reduces the question to the case of a complete noetherian
local ring. By a structure theorem of Cohen,\footnote{See I.~S. Cohen,
``On the structure and ideal theory of complete local rings,''
\emph{Trans. Amer. Math. Soc.} \textbf{59} (1946), 54--106.}
such a ring is a quotient of a regular ring. \hyperref[IV.4.3]{Proposition 4.3} then permits us
to assume that $A$ is regular. As we shall see later, this observation gives
an explicit calculation of the dualizing module.\footnote{This was the
method followed by Grothendieck in 1957. The method by injective envelopes
used below appears to be due to K.~Morita, ``Duality for modules and its
applications to the theory of rings with minimum conditions,''
  \emph{Sc. Rep. Tokyo Kyoiku Daigaku} \textbf{6} (1958/59), 83--142.
Morita's work is independent of Grothendieck's, considerably predates this
seminar, and is not restricted to commutative base rings.}
We shall nevertheless prove \hyperref[IV.4.7]{Theorem 4.7} by other means.
\end{sgastatement}

\subsection*{Recollections}

Before\srcpage{54} proving the theorem, we recall several facts about injective
envelopes; see Gabriel, \emph{Des Cat\'egories Ab\'eliennes}, thesis,
Paris, 1961, Chapter II, \S5.

Let $\Ccat$ be an abelian category in which direct limits exist and are
exact\footnote{More precisely, small filtered direct limits are assumed
exact; one should also assume the existence of a generator. See
  [T\^ohoku]. For the category of modules, which suffices here, one may
instead consult Chapter 10 of Bourbaki's \emph{Alg\`ebre}.}
(for example, the category of modules). Every object $M$ embeds in an
injective object, and an \emph{injective envelope} of $M$ means an injective
object containing $M$ that is minimal among such objects. The following
properties hold:
\begin{enumerate}[label=\textup{\roman*)},leftmargin=2.7em]
\item Every object $M$ has an injective envelope $I$.
\item If $I$ and $J$ are two injective envelopes of $M$, there is an
isomorphism $I\simeq J$ (generally not unique) that induces the identity on
$M$.
\item The object $I$ is an \emph{essential extension}\sgareaderlabel{pIV.13}{13} of $M$: if
$P\subset I$ and $P\cap M=\{0\}$, then $P=\{0\}$. Conversely, if $I$ is
injective and is an essential extension of $M$, then $I$ is an injective
envelope of $M$.
\end{enumerate}

Granting these facts, it plainly suffices to prove part~\emph{c)} of
\hyperref[IV.4.7]{Theorem 4.7}.

\begin{proof}
Suppose first that $I$ is a dualizing module for $A$. Then $I$ is injective
and $\Homop_A(k,I)\simeq k$. Composing this isomorphism with the inclusion
\[
  \Homop_A(k,I)\hookrightarrow\Homop_A(A,I)\simeq I
\]
gives an inclusion
\[
  k\hookrightarrow I.
\]
We show that $I$ is an injective envelope of $k$. Let $J$ be an injective
module such that
\[
  k\subset J\subset I.
\]
Since $J$ is injective, there is an injective $A$-submodule $J'$ of $I$ such
that $I=J\oplus J'$. We claim that $\Homop_A(k,J')=0$. Indeed,
\[
  \Homop_A(k,I)\simeq\Homop_A(k,J)\oplus\Homop_A(k,J').
\]
The first summand is\srcpage{55} a nonzero vector subspace of
$\Homop_A(k,I)\simeq k$ (it contains the inclusion $k\subset J$).
Consequently $\Homop_A(k,J)\simeq k$ and $\Homop_A(k,J')=0$.

Induction on length now gives $\Homop_A(M,J')=0$ for every finite-length
$A$-module $M$. The module $I$ is the direct limit of the modules
$\Homop_A(A/\mideal^n,I)$ (cf. \hyperref[IV.1.3]{Proposition 1.3}), which have finite length by
hypothesis. Hence the projection $I\to J'$ vanishes, and therefore $J'=0$.

Conversely, let $I$ be an injective envelope of $k$. By \hyperref[IV.2.1]{Propositions 2.1}
and \hyperref[IV.3.1]{3.1\emph{(ii)}}, to prove that $I$ is dualizing it suffices to show that
$V=\Homop_A(k,I)$ is isomorphic to $k$. There are inclusions
\[
  k\subset V\subset I.
\]
As a $k$-vector space, $V$ decomposes as $k\oplus V'$, where $V'$ is a
vector subspace of $I$ with $V'\cap k=0$. Since $I$ is an essential
extension of $k$, we have $V'=0$, and thus $V=k$.
\end{proof}

\begin{sgastatement}{Corollary 4.9}\sgareaderlabel{IV.4.9}{4.9}
Let $A$ be a noetherian local ring. Every dualizing module for $A$ is
locally artinian.
\end{sgastatement}

\begin{proof}
Let $I$ be a dualizing module; it is an injective envelope of $k$. Using the
notation and result of \hyperref[IV.2.2]{Corollary 2.2}, we have
\[
  k\subset H^0_{\mideal}(I)\subset I,
\]
and $H^0_{\mideal}(I)$ is injective. It follows that
$I=H^0_{\mideal}(I)$, and hence that $I$ is locally artinian.\footnote{As
already observed, one may alternatively note directly that $I$ is the
direct limit of the modules $\Homop_A(A/\mideal^n,I)$.}
\end{proof}

% SOURCE COMPONENT: 10_unit_IV_5_1\part.tex

\section*{5. Consequences of the theory of dualizing modules}\sgareaderlabel{IV.5}{5}\srcpage{56}

The functor
\[
  T=\Homop_A(-,I)\colon \Ctors\longrightarrow\Ctors
\]
is an anti-equivalence. Indeed, $T\circ T$ is isomorphic to the identity
functor, and the argument is formal from that point on.

\emph{One deduces the usual properties of the notion of orthogonality:}

Let $M^*=\Homop_A(M,I)=T(M)$, and let $N\subset M$ be a submodule. The
\emph{orthogonal of} $N$ is the submodule $N'$ of $M^*$ consisting of those
elements of $M^*$ that vanish on $N$. In this way one obtains a bijection
between the set of submodules of $M$ and the set of submodules of $M^*$,
reversing the order relation.

In particular:
\begin{itemize}[leftmargin=2.3em]
\item $\operatorname{length}_{M}N=
  \operatorname{colength}_{M^*}N'$.
\item Cyclic modules, that is, those for which $M/\mideal M$ has dimension
  $0$ or $1$, correspond under the duality to modules whose socle has length
  $0$ or $1$.
\item If $A$ is artinian, the ideals of $A$ correspond to the submodules of
  $I$.
\end{itemize}
And so on.

Let $A$ be a noetherian local ring. Write $\DA$ for the category of
$A$-modules $M$ such that, for every $n\in\mathbb N$,
\[
  M^{(n)}=M/\mideal^{n+1}M
\]
has finite length and such that
\[
  M=\varprojlim_n M^{(n)},
\]
and let $\widehat A$ be the completion of $A$. Restriction of scalars and
completion are quasi-inverse equivalences between $\DA$ and $\DhatA$; up to
isomorphism, they are compatible with taking the underlying abelian groups of
the modules in question. Write $\CA$ for the category of locally artinian
$A$-modules with finite-dimensional socle.

\begin{sgastatement}{Proposition 5.1}\sgareaderlabel{IV.5.1}{5.1}\srcpage{57}
Let $A$ be a noetherian local ring, and let $I$ be a dualizing module for
$A$. The functors
\[
  \Homop_A(-,I)\colon (\CA)^{\mathrm{op}}\longrightarrow\DA
\]
and
\[
  \Homop_{\widehat A}(-,I)\colon \DA\longrightarrow(\CA)^{\mathrm{op}}
\]
are equivalences of categories, quasi-inverse to one another.

Moreover, transporting these functors by the equivalences of categories
between $\DA$ and $\DhatA$ on the one hand, and between $\CA$ and $\ChatA$ on
the other, yields the functor $\Homop_{\widehat A}(-,I)$.
\end{sgastatement}

\begin{proof}
Let $X\in\operatorname{Ob}(\CA)$. By definition,
\[
  X=\varinjlim_{k\in\mathbb N}X_k,
  \qquad
  X_k=\Homop_A(A/\mideal^{k+1},X),
\]
and therefore
\[
  \Homop_A(X,I)=\varprojlim_k\Homop_A(X_k,I).
\]
Thus
\[
  Y=\Homop_A(X,I)=\varprojlim_k\Homop_A(X_k,I)
\]
is a finitely generated $\widehat A$-module, as follows from
\emph{EGA}~$0_{\mathrm I}$, 7.2.9. In this connection, observe that $\DhatA$
is also the category of finitely generated $\widehat A$-modules, or,
equivalently, that $\DA$ is the category of $A$-modules that are complete and
finitely generated over $\widehat A$.\footnote{The printed French and the
corrected TeX write $Y=\varprojlim_k X_k$.  But the $X_k$ were introduced as
an increasing direct system, whereas the immediately preceding display gives
$\Homop_A(X,I)=\varprojlim_k\Homop_A(X_k,I)$.  The target therefore restores
this latter module as $Y$, which is also the object required by the remainder
of the proof.}

Now let $Y$ be such a module, and let $f\colon Y\to I$ be a
$\widehat A$-homomorphism. The image of $f$ is a finitely generated
submodule, hence is annihilated by $\mideal^k$ for some $k$: indeed, every
$x\in I$ is annihilated by some power of $\mideal$. Consequently $f$ factors
through $Y/\mideal^kY$, and it follows that
\begin{align*}
  \Homop_{\widehat A}(Y,I)
    &=\varinjlim_k\Homop_{\widehat A}(Y^{(k)},I),
      &\text{where }Y^{(k)}&=Y/\mideal^{k+1}Y,\\
    &=\varinjlim_k (Y^{(k)})^*.
\end{align*}
This object belongs to $\operatorname{Ob}(\CA)$. It follows immediately that
the two functors in the statement are quasi-inverse to one another.
\end{proof}

% SOURCE COMPONENT: 11_unit_IV_5_2_4\part.tex

It follows from the preceding discussion that replacing $A$ by $\widehat A$
changes neither the categories or functors under consideration nor the
underlying abelian groups of the modules involved. \hyperref[IV.5.1]{Proposition 5.1} may
therefore be restated as follows.\srcpage{58}

The restriction of $\Homop_{\widehat A}(-,I)$ to the category of finitely
generated $\widehat A$-modules takes its values in the category of locally
artinian $\widehat A$-modules having finite-dimensional socle. It admits a
quasi-inverse, namely the restriction of $\Homop_{\widehat A}(-,I)$ in the
opposite direction. On the intersection of these two categories the two
functors coincide (of course) and define an autoduality of the category of
finite-length $\widehat A$-modules.

\begin{sgastatement}{Example 5.2 (Macaulay)}\sgareaderlabel{IV.5.2}{5.2}
Let $A$ be a local ring with residue field $k$. Let $k_0$ be a subfield of
$A$ such that $k$ is finite over $k_0$, and put $[k:k_0]=d$. Every
finite-length $A$-module may be regarded as a finite-dimensional
$k_0$-vector space of dimension
\[
  d\,\operatorname{length}(M).
\]
The functor
\[
  T(M)=\Homop_{k_0}(M,k_0)
\]
is exact and preserves length, hence is dualizing for $A$. Its associated
dualizing module is
\[
  A'=\varinjlim_n\Homop_{k_0}(A/\mideal^n,k_0),
\]
the topological dual of $A$ endowed with the $\mideal$-adic topology.
\end{sgastatement}

\begin{sgastatement}{Example 5.3}\sgareaderlabel{IV.5.3}{5.3}
Let $A$ be a regular noetherian local ring of dimension $n$, let
$\mideal$ be its maximal ideal, and let $k$ be its residue field. There is a
regular system of parameters $(x_1,x_2,\ldots,x_n)$ that generates
$\mideal$ and is an $A$-regular sequence. The Koszul complex therefore gives
\[
\begin{aligned}
  \Extop_A^i(k,A)&=0 &&\text{if }i\ne n,\\
  \Extop_A^n(k,A)&\simeq k.
\end{aligned}
\]
Since the depth\srcpage{59} of $A$ is $n$, one has
$\Extop_A^i(M,A)=0$ for $i<n$ whenever
$M$ is annihilated by a power of $\mideal$. One also has
$\Extop_A^i(M,A)=0$ for $i>n$, because the global cohomological dimension
of $A$ equals $n$. Thus $\Extop_A^n(-,A)$ is exact, and
$\Extop_A^n(k,A)\simeq k$. It follows that:
\end{sgastatement}

\begin{sgastatement}{Proposition 5.4}\sgareaderlabel{IV.5.4}{5.4}
If $A$ is a regular noetherian local ring of dimension $n$, the functor
\[
  M\longmapsto\Extop_A^n(M,A)
\]
is dualizing. The associated dualizing module is
\[
  I=\varinjlim_r\Extop_A^n(A/\mideal^r,A).
\]
It is isomorphic to $H_{\mideal}^n(A)$ (Expos\'e II, \hyperref[II.6]{Theorem 6}).\footnote{For
a ring $A$, an ideal $J\subset A$, an $A$-module $M$, and
$i\in\mathbb Z$, set $H_J^i(M)=H_Y^i(X,\widetilde M)$, where
$X=\Spec(A)$ and $Y=V(J)$.}
\end{sgastatement}

% SOURCE COMPONENT: 12_unit_IV_5_5\part.tex

\begin{sgastatement}{Remark 5.5}\sgareaderlabel{IV.5.5}{5.5}\srcpage{59}
If $A$ satisfies the hypotheses of both preceding examples, the two resulting
dualizing modules are isomorphic. Suppose, for example, that $A$ is regular
of dimension $n$, complete, and of equal characteristic. There is then a
coefficient field, say $K$. After choosing a system of parameters
$(x_1,\ldots,x_n)$ of $A$, one can construct an isomorphism
\[
  A\simeq K[[T_1,\ldots,T_n]],
\]
and hence, as we shall see, an explicit isomorphism between the two dualizing
modules
\[
  v:H_{\mideal}^n(A)\longrightarrow A'.
\]

This isomorphism has an intrinsic interpretation in terms of the module
\[
  \Omega^n=\Omega^n(A/K)
\]
of completed relative differentials of top degree. Indeed, $\Omega^n$ has a
basis consisting of
\[
  dx_1\wedge dx_2\wedge\cdots\wedge dx_n.
\]
There is therefore an isomorphism
\[
  u:H_{\mideal}^n(\Omega^n)\longrightarrow H_{\mideal}^n(A).
\]
A remarkable fact is that the composite
\[
  vu=w:H_{\mideal}^n(\Omega^n)\longrightarrow A'
\]
is\srcpage{60} independent of the chosen system of parameters and commutes
with change of the base field.

To construct $v$, compute $H_{\mideal}^n(A)$ from the Koszul complex
associated with the $x_i$. One obtains
\[
  H_{\mideal}^n(A)=
  \varinjlim_r A/(x_1^r,\ldots,x_n^r).
\]
The transition morphisms are as follows. Put
\[
  I_r=A/(x_1^r,\ldots,x_n^r),
\]
and let $e_{a_1,\ldots,a_n}^r$ be the image of
$x_1^{a_1}x_2^{a_2}\cdots x_n^{a_n}$ in $I_r$. The elements
$e_{a_1,\ldots,a_n}^r$, for $0\le a_i<r$, form a basis of $I_r$.

For an integer $s\geq0$, the transition morphism
\[
  t_{r,r+s}:I_r\longrightarrow I_{r+s}
\]
is multiplication by $x_1^s x_2^s\cdots x_n^s$; consequently,
\[
  t_{r,r+s}\bigl(e_{a_1,\ldots,a_n}^r\bigr)
    =e_{a_1+s,\ldots,a_n+s}^{r+s}.\footnote{The printed French and the
  corrected TeX write $u_{r,r+s}$ in this display, immediately after defining
  the transition map as $t_{r,r+s}$. The target emends the isolated $u$ to
  $t$; the domain, codomain, and multiplication rule make the reference
  unambiguous.}
\]

Giving an $A$-homomorphism $w:M\to A'$ is equivalent to giving a
$K$-linear form $w':M\to K$ that is continuous on every finitely generated
submodule. For $M=H_{\mideal}^n(\Omega^n)$, defining $w$ is therefore
equivalent to defining a linear form
\[
  \rho:H_{\mideal}^n(\Omega^n)\longrightarrow K,
\]
called the \emph{residue form}\sgareaderlabel{pagecithar}{46}.\footnote{\sgareaderlabel{cithar}{*}For a fuller treatment of residues,
see R.~Hartshorne, \emph{Residues and Duality}, Lecture Notes in Mathematics
20, Springer, 1966.} To construct $\rho$, it is enough to define compatible
forms $\rho_r:I_r\to K$ by
\[
  \rho_r\bigl(e_{a_1,\ldots,a_n}^r\bigr)=
  \begin{cases}
    1,&\text{if }a_i=r-1\text{ for }1\le i\le n,\\
    0,&\text{otherwise.}
  \end{cases}
\]
\end{sgastatement}

% END INCLUDED active_source/parts/04_expose_IV.tex

% BEGIN INCLUDED active_source/parts/05_expose_V.tex
% BEGIN COMPONENT SGA2-V-C01 :: unit_V_1_1_1_3_opening
\makeatletter
\providecommand{\sgareaderlabel}[2]{%
  \begingroup
  \def\@currentlabel{#2}%
  \phantomsection\label{#1}%
  \endgroup}
\makeatother

\chapter{\texorpdfstring{Local Duality and Structure of the $H^i(M)$}{Local Duality and Structure of the H\textasciicircum i(M)}}
\label{V}

\section*{1. Complexes of homomorphisms}
\sgareaderlabel{V.1}{1}

\subsubsection*{1.1}
\sgareaderlabel{V.1.1}{1.1}

Let $F^{\bullet}$ and $G^{\bullet}$ be two graded modules. Write\srcpage{61}
\begin{equation}\tag{1}\label{eq:V.1}
  \Homop^{\bullet}(F^{\bullet},G^{\bullet})
\end{equation}
for the graded module of homomorphisms of graded modules from $F^{\bullet}$
to $G^{\bullet}$. Thus
\begin{equation}\tag{2}\label{eq:V.2}
  \Homop^{s}(F^{\bullet},G^{\bullet})
    =\prod_k \Homop(F^k,G^{k+s}).
\end{equation}

Suppose that $F^{\bullet}$ (respectively $G^{\bullet}$) is a complex, and let
$d_1$ (respectively $d_2$) be its differential. For
$h\in\Homop^s(F^{\bullet},G^{\bullet})$, set\footnote{\textit{Editor's note.}
The original sign convention was different, but it is incompatible with the
convention of Expos\'e VIII. The latter seems more reasonable because in that
case the degree-zero cohomology is the set of homotopy classes of morphisms
from $F^{\bullet}$ to $G^{\bullet}$. The calculations below have been modified
accordingly.}
\begin{equation}\tag{3}\label{eq:V.3}
  d(h)=h\circ d_1+(-1)^{s+1}d_2\circ h.
\end{equation}
One checks immediately that $d\circ d=0$; hence
$\Homop^{\bullet}(F^{\bullet},G^{\bullet})$, equipped with $d$, is a complex.
The cohomology group of this complex is denoted by
\begin{equation}\tag{4}\label{eq:V.4}
  \UH^{\bullet}(F^{\bullet},G^{\bullet}).
\end{equation}
If $G^{\bullet}$ is injective in every degree, then
\[
  F^{\bullet}\longmapsto
  \UH^{\bullet}(F^{\bullet},G^{\bullet})
\]
is an exact $\partial$-functor. Likewise, for arbitrary $F^{\bullet}$,
\[
  G^{\bullet}\longmapsto
  \UH^{\bullet}(F^{\bullet},G^{\bullet})
\]
is an exact $\delta$-functor on the category of complexes $G^{\bullet}$ that
are injective in every degree.

\begin{sgaremark}{Remark 1.2}\sgareaderlabel{V.1.2}{1.2}
The cycles of $\Homop^{\bullet}(F^{\bullet},G^{\bullet})$ are the
homomorphisms from $F^{\bullet}$ to $G^{\bullet}$ that commute or
anticommute with the differentials, according to the degree. The boundaries
of $\Homop^{\bullet}(F^{\bullet},G^{\bullet})$ are the homomorphisms from
$F^{\bullet}$ to $G^{\bullet}$ that are homotopic to zero.
\end{sgaremark}

\newpage
Let $A$ be a ring, let $M$ (respectively $N$) be an $A$-module, and let
$R(M)$ (respectively $R(N)$) be an injective resolution of $M$ (respectively
$N$). There is then a canonical isomorphism\srcpage{62}\footnote{\textit{Editor's
note.} The strange original numbering has been preserved.}
\begin{equation*}\tag{1.3}\label{eq:V.1.3}
  \UH^s(R(M),R(N))\simeq\Extop^s(M,N).
\end{equation*}
Indeed, let $i\colon M\to R(M)$ be the canonical augmentation, and let
$h\in\Homop^s(R(M),R(N))$. Denote by $t_s$ the map
\[
  h\longmapsto h_0\circ i
\]
from $\Homop^s(R(M),R(N))$ to $\Homop(M,R(N)^s)$. The family
$(t_s)_{s\geq 0}$ defines a homomorphism of (ordinary) complexes\footnote{Here
$M$ also denotes the complex $M[0]$ consisting of $M$ placed in degree zero.}
\[
  t\colon \Homop^{\bullet}(R(M),R(N))
    \longrightarrow \Homop^{\bullet}(M,R(N));
\]
that is,
\[
  (dh)_0\circ i=d_2\circ h_0\circ i.
\]

One checks easily that, on passing to cohomology, $t$ induces an isomorphism.
It follows in particular that
\[
  \UH^{\bullet}(R(M),R(N))
\]
does not ``depend'' on the chosen injective resolution $R(M)$ (respectively
$R(N)$) of $M$ (respectively $N$).
% END COMPONENT SGA2-V-C01

% BEGIN COMPONENT SGA2-V-C02 :: unit_V_1_4_1_5_section_completion
To every exact sequence of $A$-modules\srcpage{62}
\begin{equation}\tag{5}\label{eq:V.5}
  0\longrightarrow M'\longrightarrow M\longrightarrow M''\longrightarrow 0
\end{equation}
one associates an exact sequence of injective resolutions
\begin{equation}\tag{6}\label{eq:V.6}
  0\longrightarrow R(M')\longrightarrow R(M)
    \longrightarrow R(M'')\longrightarrow 0.
\end{equation}
One checks that the isomorphism \emph{\hyperref[eq:V.1.3]{(1.3)}} commutes with the homomorphisms
\begin{align}
  \UH^s(R(M'),R(N))
    &\longrightarrow \UH^{s+1}(R(M''),R(N)),\tag{8}\label{eq:V.8}\\
  \Extop^s(R(M'),R(N))
    &\longrightarrow \Extop^{s+1}(R(M''),R(N)),\tag{9}\label{eq:V.9}
\end{align}
deduced from \emph{\hyperref[eq:V.6]{(6)}} and \emph{\hyperref[eq:V.5]{(5)}}.\footnote{\textit{Source note.}
The corrected French TeX and the direct compiled source page both use
resolution arguments in \emph{(9)}: $R(M')$ and $R(N)$ in the source term,
and $R(M'')$ and $R(N)$ in the target term. This bounded translation
reproduces that literal reading; it does not silently replace those arguments
by $M'$, $M''$, and $N$.}

Let $P$\srcpage{63} be a third $A$-module, and let $R(P)$ be an injective
resolution of $P$. Then composition of graded morphisms gives a pairing
\begin{equation}\tag{10}\label{eq:V.10}
  \Homop^i(R(N),R(M))\times\Homop^j(R(M),R(P))
    \longrightarrow \Homop^{i+j}(R(N),R(P)),
\end{equation}
which defines a pairing
\begin{equation}\tag{11}\label{eq:V.11}
  \UH^i(R(N),R(M))\times\UH^j(R(M),R(P))
    \longrightarrow \UH^{i+j}(R(N),R(P)),
\end{equation}
and hence a homomorphism of functors in $M$:
\begin{equation*}\tag{1.4}\label{eq:V.1.4}
  \UH^i(R(N),R(M))\longrightarrow
  \Homop\!\left(\UH^j(R(M),R(P)),\UH^{i+j}(R(N),R(P))\right).
\end{equation*}

We shall see that \emph{\hyperref[eq:V.1.4]{(1.4)}} is a homomorphism of $\delta$-functors in
$M$. The exact sequences \emph{\hyperref[eq:V.5]{(5)}} and \emph{\hyperref[eq:V.6]{(6)}} give a commutative
diagram:
\begin{center}
\resizebox{\textwidth}{!}{%
\begin{tikzcd}[ampersand replacement=\&,column sep=large,row sep=large]
\Homop^i(R(N),R(M'))
  \arrow[r] \arrow[d]
\& \Homop\!\left(\Homop^j(R(M'),R(P)),\Homop^{i+j}(R(N),R(P))\right)
  \arrow[d,"{\Homop(q,\mathrm{id})}"] \\
\Homop^i(R(N),R(M))
  \arrow[r] \arrow[d,"p"]
\& \Homop\!\left(\Homop^j(R(M),R(P)),\Homop^{i+j}(R(N),R(P))\right)
  \arrow[d] \\
\Homop^i(R(N),R(M''))
  \arrow[r]
\& \Homop\!\left(\Homop^j(R(M''),R(P)),\Homop^{i+j}(R(N),R(P))\right).
\end{tikzcd}%
}
\end{center}
Let $h\in\Homop^i(R(N),R(M''))$ (respectively
$g\in\Homop^j(R(M'),R(P))$) be a cycle, and let
$h'\in\Homop^i(R(N),R(M))$ (respectively
$g'\in\Homop^j(R(M),R(P))$) be such that $p(h')=h$ (respectively
$q(g')=g$). To say that \emph{\hyperref[eq:V.1.4]{(1.4)}} is a homomorphism of $\delta$-functors
in $M$ is to say that
\begin{equation}\tag{12}\label{eq:V.12}
  g\circ dh'-dg'\circ h
\end{equation}
is a coboundary in $\Homop^{\bullet}(R(N),R(P))$.

Now
\begin{align*}
  dh'&=h'\circ d_1+(-1)^{i+1}d_2\circ h',\\
  dg'&=g'\circ d_2+(-1)^{j+1}d_3\circ g',
\end{align*}
with the obvious notation. Thus \emph{\hyperref[eq:V.12]{(12)}} becomes
\[
  g\circ h'\circ d_1+(-1)^{i+1}g\circ d_2\circ h'
  -g'\circ d_2\circ h-(-1)^{j+1}d_3\circ g'\circ h.
\]

On the other hand,\srcpage{64} since $h$ and $g$ are cycles,
\begin{align*}
  g\circ d_2&=(-1)^j d_3\circ g,\\
  d_2\circ h&=(-1)^i h\circ d_1.
\end{align*}
Consequently \emph{\hyperref[eq:V.12]{(12)}} is finally written as
\[
  d\!\left(g\circ h'+(-1)^{i+1}g'\circ h\right),
\]
which completes the proof.

Thus \emph{\hyperref[eq:V.1.3]{(1.3)}} and \emph{\hyperref[eq:V.1.4]{(1.4)}} give a homomorphism of $\delta$-functors
in $M$:
\begin{equation*}\tag{1.5}\label{eq:V.1.5}
  \Extop^i(N,M)\longrightarrow
  \Homop\!\left(\Extop^j(M,P),\Extop^{i+j}(N,P)\right).
\end{equation*}
% END COMPONENT SGA2-V-C02

% BEGIN COMPONENT SGA2-V-C03 :: unit_V_2_local_duality_theorem
\section*{2. The local duality theorem for a regular local ring}
\sgareaderlabel{V.2}{2}

Let $A$ be a regular local ring of dimension $r$, let $\mathfrak m$ be the
maximal ideal of $A$, and let $M$ be a finitely generated $A$-module. Set
\srcpage{64}\sgareaderlabel{pV2.4}{2.4}
\[
  H^i(M)=H^i_{\mathfrak m}(M)
  \qquad\text{(so }H^i(M)=\varinjlim_k
    \Extop^i(A/\mathfrak m^k,M)\text{)}.
\]
We saw in \emph{\hyperref[IV.5.4]{IV, 5.4}} that $I=H^r(A)$ is a dualizing module for $A$;
denote the associated dualizing functor by $D$. In \emph{\hyperref[eq:V.1.5]{(1.5)}}, put
$N=A/\mathfrak m^k$ and $P=A$. We then obtain a homomorphism of
$\delta$-functors in $M$,
\begin{equation}\tag{13}\label{eq:V.13}
  \varphi_k\colon \Extop^i(A/\mathfrak m^k,M)
  \longrightarrow
  \Homop\!\left(\Extop^{r-i}(M,A),
    \Extop^r(A/\mathfrak m^k,A)\right).
\end{equation}
Passing to the direct limit over $k$ gives a homomorphism of
$\delta$-functors
\begin{equation}\tag{14}\label{eq:V.14}
  \varphi\colon H^i(M)\longrightarrow
    D\!\left(\Extop^{r-i}(M,A)\right).
\end{equation}

\medskip
\noindent\textbf{Theorem 2.1 (Local duality theorem).}\sgareaderlabel{V.2.1}{2.1}
The preceding functorial homomorphism $\varphi$ is an isomorphism.

\medskip
\noindent\textit{Proof.}
If $i>r$,\srcpage{65} the right-hand side of \emph{\hyperref[eq:V.14]{(14)}} is trivially zero,
and the left-hand side is zero because
\[
  H^i(M)=\varinjlim_k \Extop^i(A/\mathfrak m^k,M),
\]
and each $\Extop^i(A/\mathfrak m^k,M)$ is zero (syzygy theorem).

If $i=r$, the preceding discussion shows that the two functors in $M$,
$H^r(M)$ and $D(\Homop(M,A))$, are right exact. Since $A$ is noetherian
and $M$ is finitely generated, it is enough to verify the isomorphism for
$M=A$, where it is immediate.

To show that $\varphi$ is a functorial isomorphism, it now suffices to
proceed by descending induction on $i$ and observe that every finitely
generated module admits a finite presentation, while for $i<r$ both sides
of \emph{\hyperref[eq:V.14]{(14)}} vanish when $M$ is finitely generated and free. This is
clear for the right-hand side. Since $H^i$ commutes with finite direct sums,
for the left-hand side it is enough to show that $H^i(A)=0$ for $i<r$.
This follows from \emph{\hyperref[III.3.4]{III, 3.4}}, since $\prof(A)=r$. \hfill$\square$
% END COMPONENT SGA2-V-C03

% BEGIN COMPONENT SGA2-V-C04 :: unit_V_3_1_opening_draft
\section*{3. Application to the structure of the $H^i(M)$}
\sgareaderlabel{V.3}{3}

\begin{sgatheorem}{Theorem 3.1}\sgareaderlabel{V.3.1}{3.1}
Let $A$ be a Noetherian local ring, let $D$ be a dualizing functor for $A$,
and let $M\ne 0$ be a finitely generated $A$-module of dimension $n$.
Then:\srcpage{65}
\begin{enumerate}
  \item[(i)] $H^i(M)=0$ if $i<0$ or if $i>n$.

  \item[(ii)] $\widehat{D(H^i(M))}$ is a finitely generated
  $\widehat A$-module of dimension at most $i$.

  \item[(iii)] $H^n(M)\ne 0$; moreover, if $A$ is complete, then
  $D(H^n(M))$ has dimension $n$ and
  \[
    \Assop\bigl(D(H^n(M))\bigr)
      =\{\,\mathfrak p\in\Assop(M)\mid \dim(A/\mathfrak p)=n\,\}.
  \]
\end{enumerate}
\end{sgatheorem}

\par\addvspace{0.8\baselineskip}\noindent\textit{Proof.}\ 
Let $I$ be the dualizing module associated with $D$. We know that
$\widehat I$ is a dualizing module for $\widehat A$. On the other hand,
\begin{align*}
  \widehat{H^i(M)}&=H^i(\widehat M),\\
  \widehat{D(H^i(M))}&=\Homop\bigl(H^i(\widehat M),\widehat I\bigr),\\
  \dim\widehat M&=\dim M.
\end{align*}
Consequently,\srcpage{66} we may suppose that $A$ is complete. By Cohen's
theorem, every complete local ring is a quotient of a regular local ring.
To reduce to that case, we need the following lemma.
% END COMPONENT SGA2-V-C04

% BEGIN COMPONENT SGA2-V-C05 :: unit_V_3_2_lemma_draft
\begin{sgalemma}{Lemma 3.2}\sgareaderlabel{V.3.2}{3.2}\srcpage{66}
Let $X$ (respectively $Y$) be a ringed space, let $X'$ (respectively $Y'$)
be a closed subspace of $X$ (respectively $Y$), and let
$f\colon X\to Y$ be a morphism of ringed spaces such that
$f^{-1}(Y')=X'$. Let $F$ be an $\mathcal O_X$-module. Denote by $A$
(respectively $B$) the ring $\Gamma(\mathcal O_X)$ (respectively
$\Gamma(\mathcal O_Y)$), and by $\overline f\colon B\to A$ the ring
homomorphism corresponding to $f$. There is a spectral sequence of
$B$-modules with initial term
\begin{equation}\tag{15}\label{eq:V.15}
  E_2^{p,q}=H_{Y'}^p\bigl(Y,R^q f_*(F)\bigr),
\end{equation}
converging to the $B$-module
$H_{X'}^{\bullet}(X,F)_{[\overline f]}$.
\end{sgalemma}

\begin{proof}
Let $\mathcal O_{Y,Y'}$ be the sheaf
$\mathcal O_Y\vert_{Y'}$, extended by zero outside $Y'$ (see Expos\'e I).
There is an isomorphism of $B$-modules
\begin{equation}\tag{16}\label{eq:V.16}
  \Homop\bigl(\mathcal O_{Y,Y'},f_*(F)\bigr)
  \simeq
  \Homop\bigl(f^*(\mathcal O_{Y,Y'}),F\bigr)_{[\overline f]}.
\end{equation}
Now
\begin{equation}\tag{17}\label{eq:V.17}
  f^*(\mathcal O_{Y,Y'})=\mathcal O_{X,X'}.
\end{equation}
Moreover, if $G$ is an injective $\mathcal O_X$-module, then $f_*(G)$ is an
injective $\mathcal O_Y$-module, at least when $f$ is flat; one can easily
reduce to this case by replacing $\mathcal O_X$, and so forth, by the constant
sheaves of rings $\mathbb Z$. Hence the spectral sequence of the composite
functor
\[
  F\longmapsto\Homop\bigl(\mathcal O_{Y,Y'},f_*(F)\bigr),
\]
whose initial term is
\[
  E_2^{p,q}=\Extop^p\bigl(Y;\mathcal O_{Y,Y'},R^q f_*(F)\bigr),
\]
converges, in view of \textup{\hyperref[eq:V.16]{(16)}} and \textup{\hyperref[eq:V.17]{(17)}}, to
\[
  \Extop^{\bullet}\bigl(X;\mathcal O_{X,X'},F\bigr)_{[\overline f]}.
\]
The lemma then\srcpage{67} follows from Expos\'e~I,
\hyperref[I.2.3bof]{Proposition~2.3~bis}.
\end{proof}
% END COMPONENT SGA2-V-C05

% BEGIN COMPONENT SGA2-V-C06 :: unit_V_3_change_of_rings_reduction
\srcpage{67}Now let $f\colon B\to A$ be a surjective homomorphism of local rings. Let
\[
  f\colon \Specop(A)\longrightarrow\Specop(B)
\]
be the corresponding morphism of affine schemes. Set
$X=\Specop(A)$ (respectively $X'=\{\mathfrak m_A\}$) and
$Y=\Specop(B)$ (respectively $Y'=\{\mathfrak m_B\}$), and let $M$ be an
$A$-module and $\widetilde M$ the corresponding $\mathcal O_X$-module. Since
$R^qf_*(\widetilde M)=0$ for $q>0$, the spectral sequence~\textup{\hyperref[eq:V.15]{(15)}}
degenerates, and \hyperref[V.3.2]{Lemma~3.2} gives an isomorphism of $B$-modules
\begin{equation}\tag{18}\label{eq:V.18}
  H_{\{\mathfrak m_B\}}^n\bigl(Y,f_*(\widetilde M)\bigr)
  \simeq
  H_{\{\mathfrak m_A\}}^n\bigl(X,\widetilde M\bigr)_{[f]},
\end{equation}
hence an isomorphism of $B$-modules
\begin{equation}\tag{19}\label{eq:V.19}
  H_{\mathfrak m_B}^n\bigl(M_{[f]}\bigr)
  \simeq
  H_{\mathfrak m_A}^n(M)_{[f]}.
\end{equation}
On the other hand, if $D_A$ (respectively $D_B$) is the dualizing functor for
$A$ (respectively $B$), then
\begin{equation}\tag{20}\label{eq:V.20}
  D_A(M)_{[f]}\simeq D_B\bigl(M_{[f]}\bigr).
\end{equation}
Finally, since there is an isomorphism of rings
\begin{equation}\tag{21}\label{eq:V.21}
  B/\Annop M_{[f]}\simeq A/\Annop M,
\end{equation}
it follows that the change of base rings under consideration changes
nothing. We may therefore suppose that $A$ is regular of dimension $r$.
% END COMPONENT SGA2-V-C06

% BEGIN COMPONENT SGA2-V-C07 :: unit_V_3_equivalence_conditions_a_c_draft
By \textup{\hyperref[V.2.1]{(2.1)}},\srcpage{67} we have
\begin{equation}\tag{22}\label{eq:V.22}
  D\bigl(H^i(M)\bigr)=\Extop^{r-i}(M,A).
\end{equation}
We shall\srcpage{68} prove the equivalence of the following properties:
\begin{enumerate}
\item[(a)] $\dim \Extop^j(M,A)\leq r-j$;
\item[(b)] for every $\mathfrak p\in X=\operatorname{Spec}(A)$ such that
  $\dim A_{\mathfrak p}<j$, one has
  $\Extop^j(M,A)_{\mathfrak p}=0$;
\item[(c)]
  $\codimop\bigl(\Supp(\Extop^j(M,A)),X\bigr)\geq j$.
\end{enumerate}

To prove \textup{(a)} $\Longrightarrow$ \textup{(b)}, let
$\mathfrak p\in X$ with $\dim A_{\mathfrak p}<j$. Then
$\dim A/\mathfrak p>r-j$, and hence, by \textup{(a)},
\[
  \Ann\bigl(\Extop^j(M,A)\bigr)\not\subset\mathfrak p,
\]
which implies that $\Extop^j(M,A)_{\mathfrak p}=0$. Let
$\mathfrak p\in\Supp(\Extop^j(M,A))$. Then
$\Extop^j(M,A)_{\mathfrak p}\neq0$, so by \textup{(b)} we have
$\dim A_{\mathfrak p}\geq j$. Therefore
\[
  \codimop\bigl(\Supp(\Extop^j(M,A)),X\bigr)
  =\inf\bigl\{\dim A_{\mathfrak p}\mid
  \mathfrak p\in\Supp(\Extop^j(M,A))\bigr\}\geq j,
\]
that is, \textup{(b)} $\Longrightarrow$ \textup{(c)}. Finally,
\textup{(c)} implies \textup{(a)} trivially.
% END COMPONENT SGA2-V-C07

% BEGIN COMPONENT SGA2-V-C08 :: unit_V_3_1_proof_items_i_ii_draft
\noindent\srcpage{68}Let us now prove the theorem.
\begin{enumerate}
\item[(i)]
Let $x=(x_1,\dots,x_r)$ be a system of parameters for $A$ such that
$x_i\in\Ann M$ for $i=1,\dots,r-n$. Let
$K^\bullet((x^k),M)$ be the Koszul complex. It is easy to see that the map
\[
  K^i((x^k),M)\longrightarrow K^i((x^{k'}),M)
\]
for $k<k'$ is zero if $i>n$. It follows that
\[
  H^i(M)=\varinjlim H^i((x^k),M)=0
\]
if $i>n$. On the other hand, it is trivial that $H^i(M)=0$ if $i<0$;
hence \textup{(i)} is proved.

\item[(ii)]
Since $A$ is regular, $\dim A_{\mathfrak p}<j$ implies that the global
homological dimension of $A_{\mathfrak p}$ is strictly less than $j$, and
hence
\[
  \Extop^j(M,A)_{\mathfrak p}
  =\Extop^j_{A_{\mathfrak p}}(M_{\mathfrak p},A_{\mathfrak p})=0.
\]
Thus we have proved \textup{(b)}, and consequently \textup{(a)}.
Assertion \textup{(ii)} then follows from \textup{\hyperref[eq:V.22]{(22)}} and
\textup{(a)}.
\end{enumerate}
% END COMPONENT SGA2-V-C08

% BEGIN COMPONENT SGA2-V-C09 :: unit_V_3_theorem31_proof_iii
\medskip
\noindent\textup{(iii)} There exists a
$\mathfrak p\in\Suppop(M)$ such that $\dim A_{\mathfrak p}=r-n$ and such that
$\Suppop(M_{\mathfrak p})=\{\mathfrak m A_{\mathfrak p}\}$. Since
$A_{\mathfrak p}$ is regular if $A$ is, we have
$\depthop A_{\mathfrak p}=r-n$, and therefore
\begin{equation}\tag{23}\label{eq:V.23}
  \Extop_A^{r-n}(M,A)_{\mathfrak p}
  =\Extop_{A_{\mathfrak p}}^{r-n}(M_{\mathfrak p},A_{\mathfrak p})\neq 0.
\end{equation}
Taking~\textup{\hyperref[eq:V.22]{(22)}} into account, this\srcpage{69} implies, on the one hand,
\[
  H^n(M)\neq 0,
\]
and, on the other hand,
\[
  \dim D\bigl(H^n(M)\bigr)\geq n,
\]
hence, by~\textup{(ii)},
\[
  \dim D\bigl(H^n(M)\bigr)=n.
\]
% END COMPONENT SGA2-V-C09

% BEGIN COMPONENT SGA2-V-C10 :: unit_V_3_theorem31_item_iii_completion
\medskip
Now let $Y=\Suppop(M)$. By~\textup{(i)}, we know that
$D\bigl(H^n(M')\bigr)=\Extop^{r-n}(M',A)$ is a left exact functor in $M'$ on
the category $(\mathcal C_Y)^\circ$. Hence there exists an $A$-module $H$ and
an isomorphism of functors in $M'$:
\[
  \Extop^{r-n}(M',A)=\Homop(M',H).
\]

Let $Y_i$, $i=1,\ldots,k$, be the irreducible components of $Y$ of maximum
dimension. We shall see that the assertion
$\Extop^{r-n}(M',A)\neq 0$ is equivalent to the assertion that there exists an
$i$ such that $\Suppop M'\supset Y_i$. Indeed, if
$\Suppop M'\supset Y_i$, then $\dim(M')=n$, and hence
$\Extop^{r-n}(M',A)\neq 0$.

If $\Suppop M'\not\supset Y_i$ for every $i=1,\ldots,k$, then
$\dim M'<n$, and
\[
  D\bigl(H^n(M')\bigr)=\Extop^{r-n}(M',A)=0.
\]

Since
$\Assop\bigl(\Extop^{r-n}(M,A)\bigr)=\Suppop M\cap\Assop(H)$, we see that
the last assertion of~\textup{(iii)} follows from the following lemma:
% END COMPONENT SGA2-V-C10

% BEGIN COMPONENT SGA2-V-C11 :: unit_V_3_3_lemma_and_theorem31_closure
\begin{sgalemma}{Lemma 3.3}\sgareaderlabel{V.3.3}{3.3}
Let $X=\Specop(A)$, let $Y$ be a closed subset of $X$, let
\[
  T\colon(\mathcal C_Y)^\circ\longrightarrow\Ab
\]
be a left exact functor, and let $Y_i$, $i=1,\ldots,k$, be a family of
irreducible components of $Y$ such that the assertion $T(M)=0$ is equivalent
to the assertion that, for every $i$, $\Suppop M\not\supset Y_i$. Then $T$ is
representable by a module $H$ such that
\[
  \Assop(H)=\bigcup_{i=1}^k\{y_i\},
\]
where $y_i$ is the generic point of $Y_i$, for $i=1,\ldots,k$.
\end{sgalemma}

\begin{proof}
For $y\in Y$, construct\srcpage{70} an $A$-module $M(y)$ such that
\[
  \Suppop\bigl(M(y)\bigr)=\overline{\{y\}}.
\]
Suppose that $y\neq y_i$ for every $i=1,\ldots,k$. Then
$Y_i\not\subset\Suppop\bigl(M(y)\bigr)$ for every $i=1,\ldots,k$, and hence
$T\bigl(M(y)\bigr)=0$. It follows that
\[
  \Assop\bigl(T(M(y))\bigr)
  =\Suppop\bigl(M(y)\bigr)\cap\Assop(H)=\varnothing,
\]
and therefore $y\notin\Assop(H)$. If $y=y_i$, then
$Y_i\subset\Suppop\bigl(M(y)\bigr)$, so $T\bigl(M(y)\bigr)\neq0$, and hence
\[
  \Assop\bigl(T(M(y))\bigr)
  =\Suppop\bigl(M(y)\bigr)\cap\Assop(H)\neq\varnothing.
\]
By the first part of the proof, this implies $y\in\Assop(H)$, which proves
the lemma.\qedhere
\end{proof}
% END COMPONENT SGA2-V-C11

% BEGIN COMPONENT SGA2-V-C12 :: unit_V_3_4_example_affine_complement_codimension
\begin{sgaexample}{Example 3.4}\sgareaderlabel{V.3.4}{3.4}
Let $A$ be a noetherian ring, let $X=\Specop(A)$, and let $Y$ be a closed
subset of $X$ such that $X-Y$ is affine. Then, for every irreducible component
$Y_{\alpha}$ of $Y$, one has
\[
  \codimop(Y_{\alpha},X)\leq 1.
\]
\end{sgaexample}

\begin{proof}
Consider $X$ as a prescheme over $X$. Let $y_{\alpha}\in Y_{\alpha}$ be a
generic point, and consider the morphism
\[
  \Specop(\mathcal O_{X,y_{\alpha}})\longrightarrow X.
\]
The affine scheme obtained by base-changing $X$ to
$\Specop(\mathcal O_{X,y_{\alpha}})$ is canonically isomorphic to
$\Specop(\mathcal O_{X,y_{\alpha}})$.

By \textup{(EGA I 3.2.7)}, if $y_0$ is the unique closed point of
\[
  Y_0=\Specop(\mathcal O_{X,y_{\alpha}}),
\]
then $Y_0-y_0$ is affine. By \textup{(EGA III 1.3.1)}, one obtains
\[
  \HH^i(Y_0-y_0,\mathcal O_{Y_0})=0\qquad\text{if }i>0.
\]
Hence,\srcpdfpage{62 / running 54} by \textup{(\hyperref[I.2.9]{I 2.9)}},
\[
  \HH^{i-1}(\mathcal O_{X,y_{\alpha}})
  =\HH^i_{\{y_0\}}(Y_0,\mathcal O_{Y_0})=0
  \qquad\text{if }i\geq 2.
\]
By \hyperref[V.3.1]{3.1\textup{(iii)}}, it follows that
\[
  \dim\mathcal O_{X,y_{\alpha}}\leq 1,
\]
and therefore
\[
  \codimop(Y_{\alpha},X)
  =\inf_{y\in Y_{\alpha}}\dim\mathcal O_{X,y}\leq 1.
\]
\end{proof}
% END COMPONENT SGA2-V-C12

% BEGIN COMPONENT SGA2-V-C13 :: unit_V_3_5_finite_length_localization
Let\srcpage{71} $A$ be a noetherian local ring, let $\mathfrak m$ be its
maximal ideal, and let $M$ be a finitely generated $A$-module. Suppose that
$A$ is a quotient of a regular local ring. Set $X=\Specop(A)$ and, for every
$x\in X$, set $\mathfrak m_x=\mathfrak m A_x$.

\begin{sgaproposition}\sgareaderlabel{V.3.5}{3.5}
The following two conditions are equivalent:
\begin{enumerate}
  \item[(a)] $\HH^i(M)$ has finite length;
  \item[(b)] for every $x\in X-\{\mathfrak m\}$,
  \[
    \HH_{\mathfrak m_x}^{\,i-\dim\overline{\{x\}}}(M_x)=0.
  \]
\end{enumerate}
\end{sgaproposition}

\begin{proof}
\renewcommand{\qedsymbol}{}
In view of \textup{\hyperref[V.3.2]{(3.2)}}, we may suppose that $A$ is regular. By
\textup{\hyperref[V.2.1]{(2.1)}}, we have
\[
  \HH^i(M)=\DD\bigl(\Extop^{r-i}(M,A)\bigr),
\]
where $r=\dim A$. By \textup{(\hyperref[IV.4.7]{IV 4.7)}}, condition \textup{(a)} is
equivalent\footnote{\sgareaderlabel{dualnul}{*}Indeed, one must show that, if $E$ is a finitely generated
$A$-module and $\DD(E)$ has finite length, then $E$ has finite length. Let $K$
(respectively $Q$) be the kernel (respectively cokernel) of the canonical
morphism $\epsilon\colon E\to\DD\DD(E)$. The canonical morphism
$\gamma\colon\DD(E)\to\DD\DD\DD(E)$ satisfies
$\DD(\epsilon)\circ\gamma=\operatorname{id}_{\DD(E)}$. Since $\DD(E)$ has
finite length, $\gamma$ is an isomorphism, and hence so is
$\DD(\epsilon)$. Since $\DD$ is exact, it follows that $\DD(K)$ and
$\DD(Q)$ are zero. It is enough to prove that an $A$-module whose dual is zero
is itself zero, for then $E=\DD\DD(E)$ has finite length, just as $\DD(E)$
does. Indeed, let $M_0$ be a finitely generated submodule of such a module
$M$. Since $\DD$ is exact, $\DD(M_0)$ is a quotient of $\DD(M)$ and is
therefore zero. Again by exactness,
$\DD(M_0/\mathfrak m_A M_0)=0$; by biduality, the finite-length module
$M_0/\mathfrak m_A M_0$ is therefore zero. Nakayama's lemma then gives
$M_0=0$, and finally $M=0$.} to
\begin{equation}
  \Extop^{r-i}(M,A)\text{ has finite length}. \tag{24}\label{eq:V.24}
\end{equation}
Condition \textup{\hyperref[eq:V.24]{(24)}} is equivalent to
\begin{equation}
  \text{for every }x\in X-\{\mathfrak m\},\qquad
  \Extop_A^{r-i}(M,A)_x=0. \tag{25}\label{eq:V.25}
\end{equation}
On the other hand, $A_x$ is regular of dimension
$r-\dim\overline{\{x\}}$. Hence, by \textup{\hyperref[V.2.1]{(2.1)}},
\begin{align}
  \HH_{\mathfrak m_x}^{\,i-\dim\overline{\{x\}}}(M_x)
  &={\DD}\!\left(
    \Extop_{A_x}^{(r-\dim\overline{\{x\}})
      -(i-\dim\overline{\{x\}})}(M_x,A_x)
  \right) \notag\\
  &={\DD}\!\left(\Extop_{A_x}^{r-i}(M_x,A_x)\right). \tag{26}\label{eq:V.26}
\end{align}
Since $M$ is finitely generated, one has
\[
  \Extop_A^{r-i}(M,A)_x
  =\Extop_{A_x}^{r-i}(M_x,A_x),
\]
which proves the proposition.
\end{proof}
% END COMPONENT SGA2-V-C13

% BEGIN COMPONENT SGA2-V-C14 :: unit_V_3_6_corollary_exposeV_close
\begin{sgacorollary}\sgareaderlabel{V.3.6}{3.6}
For $\HH^i(M)$ to have finite length for $i\leq n$, it is necessary and
sufficient that
\[
  \depthop(M_x)>n-\dim\overline{\{x\}}
\]
for every $x\in X-\{\mathfrak m\}$.
\end{sgacorollary}

\begin{proof}
\renewcommand{\qedsymbol}{}
This follows from \textup{\hyperref[V.3.5]{(3.5)}} and \textup{(\hyperref[III.3.1]{III 3.1)}}.
\end{proof}
% END COMPONENT SGA2-V-C14

% END INCLUDED active_source/parts/05_expose_V.tex

% BEGIN INCLUDED active_source/parts/06_expose_VI.tex
\makeatletter
\providecommand{\sgareaderlabel}[2]{%
  \begingroup
  \def\@currentlabel{#2}%
  \phantomsection\label{#1}%
  \endgroup}
\makeatother

\chapter[\texorpdfstring{The Functors $\Ext_Z^{\bullet}$ and $\SheafExt_Z^{\bullet}$}{The Functors Ext over Z and sheaf Ext over Z (all degrees)}]
{The Functors $\Ext_Z^{\bullet}(X;F,G)$ and
$\SheafExt_Z^{\bullet}(F,G)$}
\label{VI}

\section*{1. Generalities}
\sgareaderlabel{VI.1}{1}
\subsection*{1.1}
\sgareaderlabel{VI.1.1}{1.1}
Let $(X,\OX)$ be a ringed space and let $Z$ be a locally closed subset of
$X$. Let $F$ and $G$ be $\OX$-modules. We shall denote by
$\Ext_Z^i(X;F,G)$ (respectively $\SheafExt_Z^i(F,G)$) the $i$th derived
functor of
\[
 G\longmapsto \Gamma_Z\bigl(\SheafHom_{\OX}(F,G)\bigr)
 \quad\text{(respectively }G\longmapsto
 \SheafGamma_Z\bigl(\SheafHom_{\OX}(F,G)\bigr)\text{)}.
\]

\begin{sgastatement}{Lemma 1.2}\sgareaderlabel{VI.1.2}{1.2}
The sheaf $\SheafExt_Z^i(F,G)$ is canonically isomorphic to the sheaf
associated with the presheaf
\[
 U\longmapsto \Ext_{Z\cap U}^i\bigl(U;F|_U,G|_U\bigr).
\]
\end{sgastatement}

This follows from \textup{([T\^ohoku], 3.7.2)} and from the fact that
\[
 \Gamma\!\left(U;\SheafGamma_Z\bigl(\SheafHom_{\OX}(F,G)\bigr)\right)
 \simeq
 \Gamma_{Z\cap U}\!\left(\SheafHom_{\OX|_U}(F|_U,G|_U)\right)
\]
canonically.

\begin{sgastatement}{\hyperref[VI.1.3]{Theorem 1.3} (Excision theorem)}
\sgareaderlabel{VI.1.3}{1.3}
Let $V$ be an open subset of $X$ containing $Z$. Then there is an
isomorphism of cohomological functors
\begin{equation}
 \Ext_Z^{\bullet}(X;F,G)
 \simeq
 \Ext_Z^{\bullet}\bigl(V;F|_V,G|_V\bigr).
 \tag{1.3.1}\label{VI.1.3.1}
\end{equation}
\end{sgastatement}

Indeed, if $G^{\bullet}$ is an injective resolution of $G$, then
$G^{\bullet}|_V$ is an injective resolution of $G|_V$. The theorem follows
immediately.

\subsection*{1.4}
\sgareaderlabel{VI.1.4}{1.4}
Let $\OoXZ$ be the $\OX$-module defined by the following conditions
\textup{([Godement], 2.9.2)}:
\[
 \OoXZ|_{X-Z}=0
 \qquad\text{and}\qquad
 \OoXZ|_Z=\OX|_Z.
\]
We have seen that, for every $\OX$-module $H$, there is a functorial
isomorphism
\[
 \Gamma_Z(H)\simeq \Homop_{\OX}(\OoXZ,H).
\]
Consequently, there are isomorphisms functorial in $F$ and $G$:
\begin{align}
 \Gamma_Z\bigl(\SheafHom_{\OX}(F,G)\bigr)
  &\simeq \Homop_{\OX}\bigl(\OoXZ,\SheafHom_{\OX}(F,G)\bigr),
   \tag{1.4.1}\label{VI.1.4.1}\\[-0.15em]
 \Gamma_Z\bigl(\SheafHom_{\OX}(F,G)\bigr)
  &\simeq \Homop_{\OX}\bigl(\OoXZ\otimes_{\OX}F,G\bigr),
   \tag{1.4.2}\label{VI.1.4.2}\\[-0.15em]
 \Gamma_Z\bigl(\SheafHom_{\OX}(F,G)\bigr)
  &\simeq \Homop_{\OX}\bigl(F,\SheafHom_{\OX}(\OoXZ,G)\bigr)
   =\Homop_{\OX}\bigl(F,\SheafGamma_Z(G)\bigr).
   \tag{1.4.3}\label{VI.1.4.3}
\end{align}
In particular, it follows from~\eqref{VI.1.4.2} that there is an
isomorphism that is $\partial$-functorial in $F$ and $\delta$-functorial in $G$,
\[
 \theta\colon
 \Ext_{\OX}^{i}\bigl(\OoXZ\otimes_{\OX}F,G\bigr)
 \xrightarrow{\ \sim\ }
 \Ext_Z^{i}(X;F,G).
\]

\subsection*{1.5}
\sgareaderlabel{VI.1.5}{1.5}
By definition, the functor
$G\longmapsto\Gamma_Z\bigl(\SheafHom_{\OX}(F,G)\bigr)$ is the composite of
$G\longmapsto\SheafHom_{\OX}(F,G)$ and $\Gamma_Z$. Since $\Gamma_Z$ is left
exact \textup{(\hyperref[I.1.9]{I 1.9)}}, since $\SheafHom_{\OX}(F,G)$ is flasque when $G$ is
injective, and since $\Gamma_Z$ is exact on flasque sheaves \textup{(\hyperref[I.2.12]{I 2.12)}},
it follows from \textup{([T\^ohoku], 2.4.1)} that there is a spectral functor
abutting to $\Ext_Z^{\bullet}(X;F,G)$ whose initial term is
\[
 H_Z^p\bigl(X,\SheafExt_{\OX}^q(F,G)\bigr).
\]

On the other hand, it follows from~\eqref{VI.1.4.3} that
$\Gamma_Z\bigl(\SheafHom_{\OX}(F,G)\bigr)$ is the composite of
$G\longmapsto\SheafGamma_Z(G)$ and
$H\longmapsto\Homop_{\OX}(F,H)$. Since $\SheafGamma_Z$ takes injectives to
injectives \textup{(\hyperref[I.1.4]{I 1.4)}}, it follows from \textup{([T\^ohoku], 2.4.1)}
that there is a spectral functor abutting to $\Ext_Z^{\bullet}(X;F,G)$ whose
initial term is
\[
 \Ext_{\OX}^p\bigl(X;F,\SheafH_Z^q(G)\bigr).
\]

Finally, it follows from~\eqref{VI.1.4.2} and the $\Ext$ spectral sequence
that there is a spectral functor abutting to $\Ext_Z^{\bullet}(X;F,G)$ whose
initial term is $H^p\bigl(X;\SheafExt_Z^{\bullet}(F,G)\bigr)$. Hence:

\begin{sgastatement}{Theorem 1.6}\sgareaderlabel{VI.1.6}{1.6}
There are three spectral functors abutting to
$\Ext_Z^{\bullet}(X;F,G)$ whose initial terms are respectively
\begin{align}
 &H_Z^p\bigl(X,\SheafExt_{\OX}^q(F,G)\bigr),
   \tag{1.6.1}\label{VI.1.6.1}\\[-0.15em]
 &H^p\bigl(X,\SheafExt_Z^q(F,G)\bigr),
   \tag{1.6.2}\label{VI.1.6.2}\\[-0.15em]
 &\Ext_{\OX}^p\bigl(X;F,\SheafH_Z^q(G)\bigr).
   \tag{1.6.3}\label{VI.1.6.3}
\end{align}
\end{sgastatement}

\subsection*{1.7}
\sgareaderlabel{VI.1.7}{1.7}
Let $Z'$ be a closed subset of $Z$, and set $Z''=Z-Z'$. There is an exact
sequence
\begin{equation}
 0\longrightarrow \OXZ{Z''}\longrightarrow \OXZ{Z}
 \longrightarrow \OXZ{Z'}\longrightarrow 0.
 \tag{1.7.1}\label{VI.1.7.1}
\end{equation}
It generalizes the exact sequence of \textup{([Godement], 2.9.3)}. This exact
sequence splits locally; hence, for every $\OX$-module $F$, there is another
exact sequence
\begin{equation}
 0\longrightarrow F\otimes_{\OX}\OXZ{Z''}
 \longrightarrow F\otimes_{\OX}\OXZ{Z}
 \longrightarrow F\otimes_{\OX}\OXZ{Z'}\longrightarrow 0.
 \tag{1.7.2}\label{VI.1.7.2}
\end{equation}
Now let $G$ be an $\OX$-module. Applying the functor
$\Homop_{\OX}(\mathbin{\bullet},G)$ to \hyperref[VI.1.7.2]{(1.7.2)}, one obtains from \hyperref[VI.1.4.2]{(1.4.2)} and
the long exact sequence of $\Ext$ the following theorem.

\begin{sgastatement}{Theorem 1.8}\sgareaderlabel{VI.1.8}{1.8}
Let $Z$ be a locally closed subset of $X$, let $Z'$ be a closed subset of
$Z$, and set $Z''=Z-Z'$. There is an exact sequence, functorial in $F$ and
$G$,
\[
\begin{aligned}
 0&\longrightarrow \Homop_{Z'}(F,G)\longrightarrow \Homop_Z(F,G)
 \longrightarrow \Homop_{Z''}(F,G)\longrightarrow \Ext_{Z'}^1(F,G)
 \longrightarrow\cdots \\
 &\cdots\longrightarrow \Ext_Z^i(F,G)\longrightarrow \Ext_{Z''}^i(F,G)
 \longrightarrow \Ext_{Z'}^{i+1}(F,G)\longrightarrow\cdots .
\end{aligned}
\]
\end{sgastatement}

\begin{sgastatement}{Corollary 1.9}\sgareaderlabel{VI.1.9}{1.9}
Let $Y$ be a closed subset of $X$, and set $U=X-Y$. There is an exact
sequence, functorial in $F$ and $G$,
\[
\begin{aligned}
 0&\longrightarrow \Homop_Y(F,G)\longrightarrow \Homop_{\OX}(F,G)
 \longrightarrow \Homop_{\OX|_U}(F|_U,G|_U)
 \longrightarrow \Ext_Y^1(F,G)\longrightarrow\cdots \\
 &\cdots\longrightarrow \Ext_{\OX}^i(F,G)
 \longrightarrow \Ext_{\OX|_U}^i(F|_U,G|_U)
 \longrightarrow \Ext_Y^{i+1}(F,G)\longrightarrow\cdots .
\end{aligned}
\]
\end{sgastatement}

This corollary is an immediate consequence of \hyperref[VI.1.3]{Theorems (1.3)} and \hyperref[VI.1.8]{(1.8)}.

\section*{2. Applications to Quasi-Coherent Sheaves on Preschemes}
\sgareaderlabel{VI.2}{2}

\begin{sgastatement}{Proposition 2.1}\sgareaderlabel{VI.2.1}{2.1}
Let $X$ be a locally Noetherian prescheme. For every locally closed subset
$Z$ of $X$, every coherent $\OX$-module $F$, and every quasi-coherent
$\OX$-module $G$, the sheaves $\SheafExt_Z^i(F,G)$ are quasi-coherent.
\end{sgastatement}

As in \hyperref[VI.1.6.3]{(1.6.3)}, one shows that the sheaves $\SheafExt_Z^i(F,G)$ are obtained
as the abutment of a spectral sequence whose initial term is
\[
 \SheafExt_{\OX}^p\bigl(F,\SheafH_Z^q(G)\bigr).
\]
By \textup{(II, \hyperref[II.3]{Corollary 3})}, the $\SheafH_Z^q(G)$ are quasi-coherent, and
therefore so are the $\SheafExt_{\OX}^p(F,\SheafH_Z^q(G))$, since $F$ is
coherent. The proposition follows immediately.

\subsection*{2.2}
\sgareaderlabel{VI.2.2}{2.2}
Let $Y$ now be a closed subprescheme of $X$, and let $\SheafI$ be an ideal
defining $Y$. Let $m$ and $n$ be integers such that $m\geq n\geq 0$. We
denote by $i_{n,m}$ the canonical map
\[
 \OY{m}=\OX/\SheafI^{m+1}\longrightarrow
 \OX/\SheafI^{n+1}=\OY{n}
\]
and by $j_n$ the map $\OXY\to\OY{n}$. The
$(\OY{n},i_{n,m})$ form a projective system, and the $j_n$ are compatible
with the $i_{n,m}$.

Applying the functor $\Ext_{\OX}^{i}(F\otimes\mathord\bullet,G)$, we obtain a
morphism
\[
 \varphi'\colon
 \varinjlim_n \Ext_{\OX}^{i}(X;F\otimes\OY{n},G)
 \longrightarrow \Ext_{\OX}^{i}(X;F\otimes\OXY,G);
\]
this is a morphism of cohomological functors in $G$. The morphism
\[
 \varphi\colon
 \varinjlim_n \Ext_{\OX}^{i}(X;F\otimes\OY{n},G)
 \longrightarrow \Ext_Y^{i}(X;F,G),
\]
which is the composite $\theta\circ\varphi'$ (cf.~\hyperref[VI.1.4]{(1.4)}), is therefore also a
morphism of cohomological functors in $G$.

In the same way one defines
\[
 \Sheafphi\colon
 \varinjlim_n \SheafExt_{\OX}^{i}(F\otimes\OY{n},G)
 \longrightarrow \SheafExt_Y^{i}(F,G).
\]

\begin{sgastatement}{Theorem 2.3}\sgareaderlabel{VI.2.3}{2.3}
Let $X$ be a locally Noetherian prescheme, let $Y$ be a closed subset of $X$
defined by a coherent ideal $\SheafI$, let $F$ be a coherent $\OX$-module,
and let $G$ be a quasi-coherent $\OX$-module. Then:
\begin{enumerate}
 \renewcommand{\labelenumi}{\alph{enumi})}
 \item $\Sheafphi$ is an isomorphism.
 \item If $X$ is Noetherian, $\varphi$ is an isomorphism.
\end{enumerate}
\end{sgastatement}

\subsection*{Proof of \hyperref[VI.2.3]{Theorem 2.3}}
The proof of b) is almost word for word the same as that of
\textup{(\hyperref[II.6]{II 6} b))}, using the spectral sequence \hyperref[VI.1.6.2]{(1.6.2)}, so we do not reproduce
it.

For the proof of a), by \hyperref[VI.2.1]{(2.1)} we may assume that $X$ is affine with ring $A$,
that $F$ (respectively $G$) is defined by an $A$-module $M$ (respectively
$N$), and that $\SheafI$ is defined by an ideal $I$. It is enough to prove
that the homomorphism
\begin{equation}
 \varinjlim_n \Ext_A^i(M/I^nM,N)
 \longrightarrow \Ext_Y^i(X,F,G)
 \tag{2.3.1}\label{VI.2.3.1}
\end{equation}
induced by $\Sheafphi$ is an isomorphism.

Indeed, for $i=0$, the two sides of \hyperref[VI.2.3.1]{(2.3.1)} may be canonically identified with
the submodule of $\Homop_A(M,N)$ consisting of the elements of
$\Homop_A(M,N)$ annihilated by a power of $I$. One then sees that the
homomorphism \hyperref[VI.2.3.1]{(2.3.1)} is simply the identity map.

The functor
\[
 N\longmapsto \varinjlim_n \Ext_A^{\bullet}(M/I^nM,N)
\]
is a universal $\delta$-functor. We shall show that the same holds for the
functor $N\longmapsto \Ext_Y^{\bullet}(M,N)$. Indeed, if $N$ is an injective
module, then by \textup{(\hyperref[II.9]{II, 9} and \hyperref[II.11]{11})}, $\SheafH_Y^q(N)=0$ for $q\ne0$; and
by \textup{(\hyperref[IV.2.2]{IV.2.2)}}, $H_Y^0(N)$ is injective.

It follows that
\[
 \Ext_{\OX}^{p}\bigl(X;M,\SheafH_Y^q(N)\bigr)=0
 \qquad\text{if }p+q\ne0.
\]
Thus, by \hyperref[VI.1.6.3]{(1.6.3)}, $\Ext_Y^i(M,N)=0$ for $i\ne0$ when $N$ is injective. This
completes the proof.

\section*{Bibliography}
The same references as those listed at the end of Expos\'e I, cited
respectively as [T\^ohoku] and [Godement].

% END INCLUDED active_source/parts/06_expose_VI.tex

% BEGIN INCLUDED active_source/parts/07_expose_VII.tex
\makeatletter
\providecommand{\sgareaderlabel}[2]{%
  \begingroup
  \def\@currentlabel{#2}%
  \phantomsection\label{#1}%
  \endgroup}
\makeatother

\chapter[Vanishing and Coherence Criteria]
{Vanishing Criteria and Coherence Conditions for the Sheaves
$\SheafExt_Y^i(F,G)$}
\label{VII}

% BEGIN COMPONENT SGA2-VII-CUM-C01
\section*{1. Study for $i<n$}
\sgareaderlabel{VII.1}{1}

We first prove a lemma.

\begin{sgastatement}{Lemma 1.1}\sgareaderlabel{VII.1.1}{1.1}
Let $X$ be a locally Noetherian prescheme, let $Y$ be a closed subset of $X$,
and let $G$ be a quasi-coherent $\OX$-module. Suppose that, for every coherent
$\OX$-module $F$ whose support is contained in $Y$, one has
\[
  \SheafExt^{n-1}(F,G)=0.
\]
Then, for every coherent $\OX$-module $F$ and every closed subset $Z$ of $X$
such that $Z\cap\Supp F=Y\cap\Supp F$, one has an isomorphism
\[
  \SheafExt_Z^n(F,G)\cong
  \SheafHom\bigl(F,\SheafH_Y^n(G)\bigr).
\]
\end{sgastatement}
% END COMPONENT SGA2-VII-CUM-C01

% BEGIN COMPONENT SGA2-VII-CUM-C02
\subsection*{Proof of \hyperref[VII.1.1]{Lemma 1.1}}

We first observe that
\[
  \SheafExt_Z^i(F,G)
  =\SheafExt_{Z\cap\Supp F}^i(F,G).
\]
(This is immediate; cf. Expos\'e VI.) We first prove the assertion for
$Z=X$, so that $\Supp F\subset Y$. On the category of coherent $\OX$-modules
whose support is contained in $Y$, the functor
\[
  F\longmapsto\SheafExt^n(F,G)
\]
is left exact. By (\hyperref[IV.1.3]{IV, 1.3)}, it is represented by
\[
  I=\varinjlim_k\SheafExt^n
    \bigl(\OX/\mathcal I^{k+1},G\bigr),
\]
where $\mathcal I$ is a defining ideal of $Y$. By (\hyperref[II.6]{II, 6)}, one has
\[
  \SheafH_Y^n(G)\cong
  \varinjlim_k\SheafExt^n
    \bigl(\OX/\mathcal I^{k+1},G\bigr).
\]
This proves the assertion when $Z=X$. Again by (\hyperref[VI.2.3]{VI, 2.3)},
\[
  \SheafExt_Z^n(F,G)\cong
  \varinjlim_k\SheafExt^n
    \bigl(F/\mathcal J^{k+1}F,G\bigr),
\]
where $\mathcal J$ is a defining ideal of $Z$. The support of
$F/\mathcal J^{k+1}F$ is contained in $Y$ whenever
$Z\cap\Supp F\subset Y$. What we have just proved therefore gives
\[
  \SheafExt_Z^n(F,G)\cong
  \varinjlim_k\SheafHom
    \bigl(F/\mathcal J^{k+1}F,\SheafH_Y^n(G)\bigr).
\]
It remains to show that the natural homomorphism
\[
  \varinjlim_k\SheafHom
    \bigl(F/\mathcal J^{k+1}F,\SheafH_Y^n(G)\bigr)
  \longrightarrow
  \SheafHom\bigl(F,\SheafH_Y^n(G)\bigr)
\]
is an isomorphism when $Z\cap\Supp F=Y\cap\Supp F$. Since $X$ admits a cover by
Noetherian affine open subsets, we reduce to the case where $X$ is affine
Noetherian. Then $F(X)$ is a finitely generated $\OX(X)$-module and
$\Supp\SheafH_Y^n(G)\subset Y$. Hence every homomorphism
\[
  u:F(X)\longrightarrow\SheafH_Y^n(G)(X)
\]
is annihilated by a power of $\mathcal I$, and therefore by a power of
$\mathcal J$. This proves the lemma.
% END COMPONENT SGA2-VII-CUM-C02

% BEGIN COMPONENT SGA2-VII-CUM-C03
\begin{sgastatement}{Proposition 1.2}\sgareaderlabel{VII.1.2}{1.2}
Let $X$ be a locally Noetherian prescheme, let $Y$ be a closed subset of $X$,
let $G$ be a quasi-coherent $\OX$-module, and let $n$ be an integer. For any
closed subsets $Z$ and $S$ of $X$ such that $Z\cap S=Y$, the following
conditions are equivalent:
\begin{enumerate}
  \renewcommand{\labelenumi}{(\roman{enumi})}
  \item $\SheafH_Y^i(G)=0$ for $i<n$;
  \item there exists a coherent $\OX$-module $F$ with support $S$ such that
  \[
    \SheafExt_Z^i(F,G)=0\qquad\text{for }i<n;
  \]
  \item for every coherent $\OX$-module $F$ whose support is contained in $S$
  (that is, $\Supp F\cap Z=\Supp F\cap Y$), one has
  \[
    \SheafExt_Z^i(F,G)=0\qquad\text{for }i<n;
  \]
  \item for every coherent $\OX$-module $F$, one has
  \[
    \SheafExt_Y^i(F,G)=0\qquad\text{for }i<n.
  \]
\end{enumerate}
Moreover, if these conditions hold, then for every coherent $\OX$-module $F$
and every closed subset $Z'$ of $X$ such that
$Z'\cap\Supp F=Y\cap\Supp F$, one has isomorphisms
\[
  \SheafExt_{Z'}^n(F,G)\cong
  \SheafExt_Y^n(F,G)\cong
  \SheafHom\bigl(F,\SheafH_Y^n(G)\bigr).
\]
\end{sgastatement}

\noindent\textit{Source note.} The quantified subset in the preceding
sentence is $Z'$. The prior maintained transcription and this English reader
dropped the prime from the first term of the final display. The primary scan
and the proof both read $Z'$. Canon correction SGA-SRC-001624 restores
the missing prime as a diplomatic transcription repair; this reader now
reproduces that scan-faithful reading.
% END COMPONENT SGA2-VII-CUM-C03

% BEGIN COMPONENT SGA2-VII-CUM-C04
\subsection*{Proof of \hyperref[VII.1.2]{Proposition 1.2}: induction}

We argue by induction. The proposition is immediate for $n<0$. Suppose that it
has been proved for $n<q$. If one of the conditions holds for $n=q$ and for two
subsets $Z$ and $S$ as in the statement, then the induction hypothesis gives,
for every closed subset $Z'$ of $X$ and every coherent $\OX$-module $F$ such
that $Z'\cap\Supp F=Y\cap\Supp F$, isomorphisms
\begin{equation}\tag{1.1}\label{eq:VII.1.1}
  \SheafExt_{Z'}^{q-1}(F,G)
  \cong\SheafHom\bigl(F,\SheafH_Y^{q-1}(G)\bigr)
  \cong\SheafExt_Y^{q-1}(F,G).
\end{equation}
Consequently:
\begin{itemize}
  \item[(i)$\Rightarrow$(iv)] take $Z'=Y$ in \hyperref[eq:VII.1.1]{(1.1)};
  \item[(iv)$\Rightarrow$(iii)] take $Z'=Z$ in \hyperref[eq:VII.1.1]{(1.1)};
  \item[(iii)$\Rightarrow$(ii)] take $F=\mathcal O_S$;
  \item[(ii)$\Rightarrow$(i)] take $Z'=Z$ in \hyperref[eq:VII.1.1]{(1.1)}. This gives
  $\SheafHom\bigl(F,\SheafH_Y^{q-1}(G)\bigr)=0$. Moreover,
  \[
    \Supp\SheafH_Y^{q-1}(G)
    \subset Y=Z\cap S\subset S=\Supp F,
  \]
  so the following lemma applies.
\end{itemize}

\begin{sgastatement}{Lemma 1.3}\sgareaderlabel{VII.1.3}{1.3}
Let $X$ be a prescheme, let $P$ be a coherent $\OX$-module, and let $H$ be a
quasi-coherent $\OX$-module such that
\[
  \SheafHom(P,H)=0
  \qquad\text{and}\qquad
  \Supp P\supset\Supp H.
\]
Then $H=0$.
\end{sgastatement}
% END COMPONENT SGA2-VII-CUM-C04

% BEGIN COMPONENT SGA2-VII-CUM-C05
\subsection*{Proof of \hyperref[VII.1.3]{Lemma 1.3}}
It suffices to prove the lemma when $X$ is affine, since the affine open
subsets form a basis for the topology of $X$ and the hypotheses are preserved
under restriction to an open subset. In that case, we are reduced to a problem
about $A$-modules, where $X=\Spec(A)$. We apply the formula (valid under the
sole hypothesis that $\sisi{M}{P}$ is of finite type):
\[
  \Assop\Homop_A(P,H)=\Supp P\cap\Assop H.
\]
We know that
\[
  \Assop H\subset\Supp H\subset\Supp P
  \qquad\text{and}\qquad
  \Assop\Homop_A(P,H)=\varnothing;
\]
hence $\Assop H=\varnothing$, and therefore $H=0$.

\noindent\textit{Source note.} The printed French text reads $M$ in the
finite-type hypothesis. The diplomatic French and source apparatus preserve
$M$; corrected French and this English edition use the accepted reading $P$,
the module in the displayed formula.

\subsection*{Close of the proof of \hyperref[VII.1.2]{Proposition 1.2}}
To complete the proof of the proposition, it remains to observe that
condition~(iv) permits the application of \hyperref[VII.1.1]{Lemma~1.1}.
% END COMPONENT SGA2-VII-CUM-C05

% BEGIN COMPONENT SGA2-VII-CUM-C06
\begin{sgastatement}{Corollary 1.4}\sgareaderlabel{VII.1.4}{1.4}
Let $G$ be a coherent Cohen--Macaulay $\OX$-module, and let $n\in\mathbb Z$.
The conditions of \hyperref[VII.1.2]{Proposition~1.2} are equivalent to:
\[
  \textup{(v)}\qquad
  \codimop(Y\cap\Supp G,\Supp G)\geq n.
\]
\end{sgastatement}

\subsection*{Proof}
Recall first that an $\OX$-module is said to be Cohen--Macaulay if, for every
$x\in X$, the stalk $G_x$ is a Cohen--Macaulay $\Oo_{X,x}$-module; that is,
for every $x\in S=\Supp G$,
\begin{equation}
  \depthop G_x=\dim G_x=\dim\Oo_{S,x}.
  \tag{1.2}\label{eq:VII.1.2}
\end{equation}
By Proposition~\hyperref[III.3.3]{III.3.3}, condition~(i) of \hyperref[VII.1.2]{Proposition~1.2} is equivalent to
\begin{equation}
  \depthop_Y G=\inf_{x\in Y}\depthop G_x\geq n,
  \tag{1.3}\label{eq:VII.1.3}
\end{equation}
and hence also to
\[
  \depthop_Y G=\inf_{x\in Y\cap S}\depthop G_x\geq n,
\]
since the depth of the zero module is infinite.

By definition,
\[
  \codimop(Y\cap S,S)=\inf_{x\in S\cap Y}\dim\Oo_{S,x};
\]
the conclusion follows by applying formula~\eqref{eq:VII.1.2}.
% END COMPONENT SGA2-VII-CUM-C06

% BEGIN COMPONENT SGA2-VII-CUM-C07
We shall now prove a result that will allow us to deduce the coherence
conditions we have in view from certain vanishing criteria.

\begin{sgastatement}{Lemma 1.5}\sgareaderlabel{VII.1.5}{1.5}
Let $X$ be a locally Noetherian prescheme. Let $T^{\ast}$ be an exact
contravariant $\partial$-functor defined on the category of coherent
$\OX$-modules and taking values in the category of $\OX$-modules. Let $Y$ be a
closed subset of $X$, and let $i\in\mathbb Z$. Suppose that, for every coherent
$\OX$-module $F$ whose support is contained in $Y$, the modules $T^iF$ and
$T^{i-1}F$ are coherent. Let $F$ be a coherent $\OX$-module. Then $T^iF$ is
coherent if and only if $T^iF''$ is coherent, where
\[
  F''=F/\sisi{\Gamma}{\SheafGamma}_Y(F).
\]
\end{sgastatement}

Indeed, $F'=\sisi{\Gamma}{\SheafGamma}_Y(F)$ is coherent because $X$ is locally Noetherian. The
long exact cohomology sequence of $T^{\ast}$ then gives
\[
  T^{i-1}F'\longrightarrow T^iF''\longrightarrow T^iF
  \longrightarrow T^iF',
\]
where the two outer terms are coherent; the conclusion follows.
% END COMPONENT SGA2-VII-CUM-C07

% BEGIN COMPONENT SGA2-VII-CUM-C08
\begin{sgastatement}{Lemma 1.6}\sgareaderlabel{VII.1.6}{1.6}
If $F$ and $G$ are coherent and $\Supp F\subset Y$, then
$\SheafExt_Y^i(F,G)$ is coherent.
\end{sgastatement}

Indeed, $\SheafExt_Y^i(F,G)$ is isomorphic to $\SheafExt^i(F,G)$. More
generally, this holds on any ringed space $X$: if $Z$ is a closed subset
containing $Y\cap\Supp F$, then $\SheafExt_Z^i(F,G)$ is isomorphic to
$\SheafExt_Y^i(F,G)$ (cf. Expos\'e~VI).

\begin{sgastatement}{Proposition 1.7}\sgareaderlabel{VII.1.7}{1.7}
Suppose that $F$ and $G$ are coherent, and put
\[
  \Supp F=S,\qquad S'=\overline{S\cap(X-Y)}.
\]
Suppose that $\depthop G_x\geq n$ for every $x\in Y\cap S'$. Then
$\SheafExt_Y^i(F,G)$ is coherent for $i<n$.
\end{sgastatement}

Indeed, \hyperref[VII.1.6]{Lemma~1.6} allows us to apply \hyperref[VII.1.5]{Lemma~1.5} to
\[
  T^{\ast}(F)=\SheafExt_Y^{\ast}(F,G).
\]
Putting $F''=F/\sisi{\Gamma}{\SheafGamma}_Y(F)$, we see that $\Supp F''=S'$. By
Proposition~\hyperref[III.3.3]{III.3.3}, the hypothesis on the depth of $G$ ensures that
$\SheafH^i_{Y\cap S'}(G)=0$ for $i<n$; \hyperref[VII.1.2]{Proposition~1.2} then gives
$T^iF''=0$ for $i<n$, and the conclusion follows from \hyperref[VII.1.5]{Lemma~1.5}.
% END COMPONENT SGA2-VII-CUM-C08

% BEGIN COMPONENT SGA2-VII-CUM-C09
\section*{2. Study for $i>n$}
\sgareaderlabel{VII.2}{2}

Let $X$ be a locally Noetherian \emph{regular} prescheme; that is, suppose that
all its local rings are regular. Let $Y$ be a closed subset of $X$, and let
$F$ and $G$ be two coherent $\OX$-modules. Put
\[
  S=\Supp F,\qquad S'=\overline{S\cap(X-Y)},
\]
and put
\begin{align*}
  m&=\sup_{x\in Y\cap S}\dim\Oo_{X,x},\\
  n&=\sup_{x\in Y\cap S'}\dim\Oo_{X,x}.
\end{align*}
Then $n\leq m$.

\begin{sgastatement}{Proposition 2.1}\sgareaderlabel{VII.2.1}{2.1}
In the situation described above:
\begin{enumerate}
\item $\SheafExt_Y^i(F,G)=0$ if $i>m$;
\item $\SheafExt_Y^i(F,G)$ is coherent if $i>n$.
\end{enumerate}
\end{sgastatement}
% END COMPONENT SGA2-VII-CUM-C09

% BEGIN COMPONENT SGA2-VII-CUM-C10
\subsection*{Proof of \hyperref[VII.2.1]{Proposition 2.1}}

First note that $\SheafExt_Y^i(F,G)$ is coherent for every $i$ when
$\Supp F\subset Y$. Moreover, putting $F''=F/\sisi{\Gamma}{\SheafGamma}_Y(F)$ as above, we see
that $\Supp F''=S'$. Thus \hyperref[VII.2.1]{assertion~(2)} follows from \hyperref[VII.2.1]{assertion~(1)} and
\hyperref[VII.1.5]{Lemma~1.5}.\footnote{The French TeX and printed page cite Lemma~1.3. The
coherence reduction used here is Lemma~1.5; Lemma~1.3 is the earlier
$\mathcal Hom$/support lemma.}

To prove \hyperref[VII.2.1]{assertion~(1)}, first observe that
\[
  \SheafExt_Y^i(F,G)\cong
  \varinjlim_k \SheafExt^i(F/\Jj^kF,G),
\]
where $\Jj$ is the ideal defining $Y$. On the other hand, Theorem~4.2.2 of
A.~Grothendieck, ``Sur quelques points d'alg\`ebre homologique,''
\emph{T\^ohoku Mathematical Journal} \textbf{9} (1957), 119--221, shows that
the $\SheafExt$ sheaves commute with the formation of stalks, at least when
$X$ is a locally Noetherian prescheme and the first argument is coherent.
Since direct limits do likewise, for every $x\in X$ there are
isomorphisms\footnote{The French TeX and printed stalk formula omit the
superscript $i$ on the local $\Extop$ term. The preceding direct-limit formula
and the following vanishing display both require $\Extop^i$.}
\[
  (\SheafExt_Y^i(F,G))_x\cong
  \varinjlim_k
  \Extop^i_{\Oo_{X,x}}((F/\Jj^kF)_x,G_x).
\]

Since
\[
  \Supp\SheafExt_Y^i(F,G)\subset S\cap Y,
\]
it is enough to note that $x\in Y\cap S$ implies
$\dim\Oo_{X,x}\leq m$, and hence
\[
  \Extop^i_{\Oo_{X,x}}((F/\Jj^kF)_x,G_x)=0
  \qquad\text{if }i>m,
\]
because the global cohomological dimension of a regular local ring is equal
to its dimension.\footnote{Cf. EGA $0_{\mathrm{IV}}$, 17.3.1.}
% END COMPONENT SGA2-VII-CUM-C10

% BEGIN COMPONENT SGA2-VII-CUM-C11
Let $X$ be a locally Noetherian prescheme. For every subset $P$ of $X$, set
\[
  D(P)=\left\{\dim\Oo_{X,p}\mid p\in P\right\}.
\]

\subsection*{\hyperref[VII.2.2]{Lemma 2.2}}
\sgareaderlabel{VII.2.2}{2.2}

If $P$ is the underlying space of a connected subprescheme of $X$, then
$D(P)$ is an interval.\footnote{The French TeX and printed witness read
``of $A$.'' The ambient prescheme introduced immediately before is $X$, and
no object $A$ occurs in this argument. The reader therefore uses $X$
transparently.}

Indeed, let $a$ and $b$ belong to $D(P)$, corresponding to points $p$ and $q$
of $P$. We show that there is a sequence of points of $P$,
\[
  p=p_1,\ldots,p_n=q,
\]
such that, for $1\leq i<n$,
\[
  \lvert\dim\Oo_{X,p_i}-\dim\Oo_{X,p_{i+1}}\rvert=1.
\]
It will follow that $D(P)$ contains the interval $[a,b]$.\footnote{The French
TeX originally introduced the numerical endpoints as ``$a$ and $B$''
and printed the interval $[p,q]$, although $p$ and $q$ are points of $P$.
The reader normalizes the two numerical endpoints to $a,b$ and the interval to
$[a,b]$. The absolute-value bars are present in the current French TeX and
printed witness and are preserved.}

For this, observe that $p$ and $q$ can be joined by a chain of irreducible
components of $P$ in which two successive components meet. We are reduced to
the case where $p$ is the generic point of an irreducible component $Q$ of
$P$, and where $q\in Q$, hence $q\supset p$ as ideals of $\Oo_q$; there the
assertion follows immediately from the definition of dimension.
% END COMPONENT SGA2-VII-CUM-C11

\clearpage
% BEGIN COMPONENT SGA2-VII-CUM-C12
\subsection*{\hyperref[VII.2.3]{Proposition 2.3}}
\sgareaderlabel{VII.2.3}{2.3}

Let $X$ be a locally Noetherian regular prescheme, let $Y$ be a closed subset
of $X$, and let $F$ be a coherent $\OX$-module. Put
\[
  P=Y\cap\overline{\Supp F\cap(X-Y)}.
\]
Let $n\in\mathbb Z$, and suppose that $n\notin D(P)$. Then
$\SheafExt_Y^n(F,\OX)$ is coherent.\footnote{The French TeX and printed
statement omit the closure and read
$P=Y\cap\Supp F\cap(X-Y)$, which is empty literally. Section~2 defines
$S'=\overline{\Supp F\cap(X-Y)}$, and Propositions~2.1 and~1.7 both use the
locus $Y\cap S'$. The reader restores that uniquely supported closure.}

\subsection*{Proof}

The conclusion is local, and the hypotheses are preserved under restriction
to an open subset. Since $P$ is closed, it is locally Noetherian and hence
locally connected. We may therefore suppose that $X$ is affine and Noetherian
and that $P$ is connected.

If $P=\varnothing$, the depth hypothesis of \hyperref[VII.1.7]{Proposition~1.7} with threshold
$n+1$ is vacuous, and the conclusion follows. Assume henceforth that
$P\neq\varnothing$. By \hyperref[VII.2.2]{Lemma~2.2}, $D(P)$ is an interval. Since
$n\notin D(P)$, exactly one of the following alternatives holds:\footnote{The
French TeX and print write $D(P)=[a,b[$ and then treat only $n>b$ and $n<a$,
leaving the allowed case $n=b$ uncovered. The print does not decide whether
the intended one-character repair was a closed upper endpoint or a non-strict
upper inequality. The reader states the order alternatives forced by
$n\notin D(P)$ and \hyperref[VII.2.2]{Lemma~2.2} and therefore does not guess that glyph.}
\[
  n>d\quad\text{for every }d\in D(P),
  \qquad\text{or}\qquad
  n<d\quad\text{for every }d\in D(P).
\]

In the first case, put $m=\sup D(P)$. Then $n>m$, and \hyperref[VII.2.1]{Proposition~2.1},
applied with $G=\OX$, shows that $\SheafExt_Y^n(F,\OX)$ is coherent.

In the second case, for every $x\in P$ regularity gives
\[
  n<\dim\Oo_{X,x}=\depthop\Oo_{X,x},
\]
and hence $\depthop\Oo_{X,x}\geq n+1$. \hyperref[VII.1.7]{Proposition~1.7} with threshold
$n+1$ therefore gives coherence in every degree $i<n+1$, in particular in
degree $i=n$.
% END COMPONENT SGA2-VII-CUM-C12

% END INCLUDED active_source/parts/07_expose_VII.tex

% BEGIN INCLUDED active_source/parts/08_expose_VIII.tex
% BEGIN COMPONENT SGA2-VIII-CUM-C01
\makeatletter
\providecommand{\sgareaderlabel}[2]{%
  \begingroup
  \def\@currentlabel{#2}%
  \phantomsection\label{#1}%
  \endgroup}
\makeatother

\chapter{The Finiteness Theorem}
\label{VIII}

\section*{1. A Biduality Spectral Sequence\footnote{A reader familiar with
Verdier's language of derived categories will recognize the spectral sequence
associated with a \emph{biduality isomorphism}. Cf. SGA~6, I.}}
\sgareaderlabel{VIII.1}{1}

Let us state the result we wish to establish:

\begin{sgastatement}{Proposition 1.1}\sgareaderlabel{VIII.1.1}{1.1}
Let $A$ be a Noetherian ring and let $I$ be an ideal of $A$. Put
$X=\operatorname{Spec}(A)$ and $Y=V(I)$. Let $M$ be a finitely generated
$A$-module of finite projective dimension. Let
$\mathcal F=\widetilde M$ be the $\OX$-module associated with $M$.
\begin{enumerate}
\item[\textup{1)}] There is a spectral sequence
\[
  \Hloc_Y(X,\mathcal F)\Longleftarrow
  \Ext_Y^p(\Ext^{-q}(M,A),A).
\]
\item[\textup{2)}] There is a spectral sequence
\[
  \SheafH_Y(X,\mathcal F)\Longleftarrow
  \SheafExt_Y^p(\SheafExt^{-q}(\mathcal F,\OX),\OX).
\]
\end{enumerate}
\end{sgastatement}
% END COMPONENT SGA2-VIII-CUM-C01

% VERIFIED SOURCE SEAM: French line 2526 is blank and excluded.

% INTEGRATION CONTROL BEFORE SGA2-VIII-CUM-C02
\let\cF\uF

% BEGIN COMPONENT SGA2-VIII-CUM-C02
Of course, 2) follows from 1) by observing that, if $M$ and $N$ are two
finitely generated $A$-modules and if we put $\mathcal F=\widetilde M$ and
$\mathcal G=\widetilde N$, then there are isomorphisms
\begin{align*}
  \SheafH_Y(\mathcal F)
    &\simeq \widetilde{\Hloc_Y(X,\mathcal F)}, \\
  \SheafExt_Y(\mathcal F,\mathcal G)
    &\simeq \widetilde{\Ext_Y(\mathcal F,\mathcal G)}, \\
  \SheafExt_{\OX}(\mathcal F,\mathcal G)
    &\simeq \widetilde{\Ext_A(M,N)}.
\end{align*}

Let $\CatC$ be the category of $A$-modules, and let $\Ab$ be the category of
abelian groups. Let $\cF$ be the functor
\[
  \begin{aligned}
    \cF\colon \CatC &\longrightarrow \Ab, \\
    M &\longmapsto \Gamma_Y(\widetilde M).
  \end{aligned}
\]
We know from Expos\'e~II that there is an isomorphism of $\partial$-functors
\[
  \Hloc_Y^*(X,\widetilde M)\simeq \R^*\cF(M).
\]

Furthermore, let $\Ext_Y^*$ denote the right derived functors, in the second
argument, of
\[
  \cF\circ\operatorname{Hom}\colon \CatC^{\circ}\times\CatC\longrightarrow\Ab.
\]
We know from Expos\'e~VI that there is an isomorphism of $\partial$-functors
\[
  \Ext_Y^*(M,N)\simeq \Ext_Y^*(\widetilde M,\widetilde N).
\]

Finally, let us record the following result from Expos\'e~VI: if $C$ is an
injective $A$-module and $N$ is a finitely generated $A$-module, the sheaf
\[
  \SheafHom(\widetilde N,\widetilde C)
    \simeq \widetilde{\operatorname{Hom}(N,C)}
\]
is flasque; hence
\[
  \R^1\cF(\operatorname{Hom}(N,C))=0.
\]
% END COMPONENT SGA2-VIII-CUM-C02

% VERIFIED SOURCE SEAM: French line 2552 is blank and excluded.

% BEGIN COMPONENT SGA2-VIII-CUM-C03
It remains for us to prove the following result.

\medskip
\sgareaderlabel{VIII.1.2}{1.2}\noindent\textbf{Lemma 1.2.} --- Let $A$ be a Noetherian ring and let $\CatC$
be the category of $A$-modules. Let
$\cF\colon\CatC\longrightarrow\Ab$ be a left exact additive functor such that,
for every finitely generated $A$-module $N$ and every injective $A$-module
$C$, one has
\[
  \R^1\cF(\operatorname{Hom}(N,C))=0.
\]
Let $M$ be a finitely generated $A$-module of finite projective dimension.
There is a spectral sequence
\[
  \R^\ast\cF(M)
    \Longleftarrow
  \Ext_{\cF}^{p}(\Ext^{-q}(M,A),A),
\]
where $\Ext_{\cF}^{p}$ denotes the $p$-th right derived functor of
$\cF\circ\operatorname{Hom}$.
% END COMPONENT SGA2-VIII-CUM-C03

% VERIFIED SOURCE SEAM: French line 2560 is blank and excluded.

% BEGIN COMPONENT SGA2-VIII-CUM-C04
\emph{We shall consider only complexes whose differential has degree $+1$.}
By the hypothesis on $M$, there exists a projective resolution of $M$ of
finite length,
\[
  u\colon L^\bullet\longrightarrow M,
\]
where, moreover, the $L^p$ are finitely generated modules and $L^p=0$ if
$p\notin[-n,0]$. On the other hand, let
$v\colon M\longrightarrow I^\bullet$ be an injective resolution of $M$. I
claim that
\begin{equation}
  v\circ u\colon L^\bullet\longrightarrow I^\bullet
  \tag{1.1}\label{eq:VIII.1.1}
\end{equation}
is an injective resolution of $L^\bullet$. We must specify what this means.

\medskip
\sgareaderlabel{VIII.1.3}{1.3}\noindent\textbf{Definition 1.3.} --- Let $X^\bullet$ be a complex of
$A$-modules. An \emph{injective resolution} of $X^\bullet$ is a homomorphism
of complexes
\[
  x\colon X^\bullet\longrightarrow CX^\bullet
\]
such that $CX^p$ is injective for every $p\in\mathbb Z$, and $x$ induces an
isomorphism in homology.
% END COMPONENT SGA2-VIII-CUM-C04

% VERIFIED SOURCE SEAM: French line 2572 is blank and excluded.

% BEGIN COMPONENT SGA2-VIII-CUM-C05
\sgareaderlabel{VIII.1.4}{1.4}\noindent\textbf{Proposition 1.4.} --- Every left-bounded complex, i.e., every
complex for which there exists $q\in\ZZ$ such that $X^p=0$ for $p<q$, admits
an injective resolution. Moreover, suppose that
$u\colon X^\bullet\longrightarrow Y^\bullet$ is a homomorphism of complexes
(both left-bounded), and that
$x\colon X^\bullet\longrightarrow CX^\bullet$ and
$y\colon Y^\bullet\longrightarrow CY^\bullet$ are injective resolutions of
$X^\bullet$ and $Y^\bullet$. Then there exists a homomorphism of complexes
\[
  Cu\colon CX^\bullet\longrightarrow CY^\bullet,
\]
unique up to homotopy, such that the diagram
\[
\begin{CD}
X^\bullet @>{x}>> CX^\bullet \\
@V{u}VV @VV{Cu}V \\
Y^\bullet @>{y}>> CY^\bullet
\end{CD}
\]
commutes up to homotopy.
% END COMPONENT SGA2-VIII-CUM-C05

% VERIFIED SOURCE SEAM: French line 2582 is blank and excluded.

% BEGIN COMPONENT SGA2-VIII-CUM-C06
The proof is left to the reader.\footnote{See also H.~Cartan and
S.~Eilenberg, \emph{Homological Algebra}, Princeton Mathematical Series,
vol.~19, Princeton University Press, 1956.}

Let us recall a notation introduced in Expos\'e V.

\medskip
\noindent\textbf{Notation.} --- Let $X^\bullet$ and $Y^\bullet$ be two
complexes. We denote by $\operatorname{Hom}^\bullet(X^\bullet,Y^\bullet)$ the
\emph{simple complex} whose component of degree $n$ is
\[
  (\operatorname{Hom}^\bullet(X^\bullet,Y^\bullet))^n
  =\prod_{-p+q=n}\operatorname{Hom}(X^p,Y^q),
\]
also denoted by $\operatorname{Hom}^n(X^\bullet,Y^\bullet)$, and whose
differential is given by
\begin{gather*}
  \partial_n\colon
  \operatorname{Hom}^n(X^\bullet,Y^\bullet)
  \longrightarrow
  \operatorname{Hom}^{n+1}(X^\bullet,Y^\bullet),\\
  \partial_n=d'+(-1)^{n+1}d'',
\end{gather*}
where $d'$ and $d''$ are the differentials (of degree $+1$) induced by those
of $X^\bullet$ and $Y^\bullet$.
% END COMPONENT SGA2-VIII-CUM-C06

% VERIFIED SOURCE SEAM: French line 2596 is blank and excluded.

% INTEGRATION CONTROL BEFORE SGA2-VIII-CUM-C07
\renewcommand{\theequation}{1.\arabic{equation}}\setcounter{equation}{1}

% BEGIN COMPONENT SGA2-VIII-CUM-C07
Let $A^\bullet$ then be the complex defined by $A^p=0$ if $p\ne 0$ and
$A^0=A$. Let
\[
  a\colon A^\bullet\longrightarrow CA^\bullet
\]
be an injective resolution of $A^\bullet$. Consider the double complex
\begin{equation}\label{eq:VIII.1.2}
  Q^{p,q}=\operatorname{Hom}(\operatorname{Hom}(L^{-q},A),CA^p).
\end{equation}
The first spectral sequence of the bicomplex $\cF Q^{\bullet\bullet}$ will
give the conclusion of \hyperref[VIII.1.2]{Lemma~1.2}.

Set
\begin{equation}\label{eq:VIII.1.3}
  {L'}^\bullet=\operatorname{Hom}^\bullet(L^\bullet,A^\bullet),
\end{equation}
and
\begin{equation}\label{eq:VIII.1.4}
  P^\bullet=\operatorname{Hom}^\bullet({L'}^\bullet,CA^\bullet).
\end{equation}
% END COMPONENT SGA2-VIII-CUM-C07

% VERIFIED SOURCE SEAM: French line 2610 is blank and excluded.

% INTEGRATION CONTROL BEFORE SGA2-VIII-CUM-C08
\renewcommand{\theequation}{1.\arabic{equation}}\setcounter{equation}{4}

% BEGIN COMPONENT SGA2-VIII-CUM-C08
One sees easily that $P^\bullet$ is the simple complex associated with
$Q^{\bullet\bullet}$. Let us compute the abutment of the spectral sequence,
that is, the homology of $\uF P^\bullet$. To this end, using the fact that
$L^\bullet$ is projective of finite type in every degree, one proves that
\[
  L^\bullet \cong \operatorname{Hom}^\bullet({L'}^\bullet,A^\bullet).
\]
From the homomorphism $a\colon A^\bullet\to CA^\bullet$, one obtains a
homomorphism
\[
  b\colon \operatorname{Hom}^\bullet({L'}^\bullet,A^\bullet)
  \longrightarrow \operatorname{Hom}^\bullet({L'}^\bullet,CA^\bullet),
\]
or equivalently a homomorphism
\begin{equation}\label{eq:VIII.1.5}
  c\colon L^\bullet\longrightarrow P^\bullet.
\end{equation}

This being said, it is easy to see, using the fact that ${L'}^\bullet$ is
projective of finite type in every degree and bounded on the left, that
\hyperref[eq:VIII.1.5]{(1.5)} is an injective resolution of $L^\bullet$. Using \hyperref[VIII.1.4]{Proposition~1.4}, one
concludes that $P^\bullet$ is homotopy equivalent to $I^\bullet$, where
$I^\bullet$ is the injective resolution of $M$ introduced above in \hyperref[eq:VIII.1.1]{(1.1)}.
It follows that the \emph{abutment} of the first spectral sequence of the
double complex $\uF Q^{\bullet\bullet}$, which is
$H^*(\uF P^\bullet)$, \emph{is isomorphic} to $R^*\uF(M)$.
% END COMPONENT SGA2-VIII-CUM-C08

% VERIFIED SOURCE SEAM: French line 2617 is blank and excluded.

% BEGIN COMPONENT SGA2-VIII-CUM-C09
The initial term of the first spectral sequence of the bicomplex
$\cF Q^{\bullet\bullet}$ is
\[
  E_2^{p,q}
  ={}'\!H^p({}''\!H^q(\cF Q^{\bullet\bullet})).
\]

For every $p\in\ZZ$, $CA^p$ is injective. By the hypothesis on $\cF$, the
functor (restricted to the category of finitely generated modules)
\[
  N\longmapsto \cF\operatorname{Hom}(N,CA^p)
\]
is exact. Hence one obtains isomorphisms
\[
  {}''\!H^q(\cF\operatorname{Hom}({L'}^\bullet,CA^p))
  \cong
  \cF\operatorname{Hom}(H^{-q}({L'}^\bullet),CA^p).
\]
By the definition of $\operatorname{Ext}^{\ast}_{\cF}$ as the right-derived functor
of $\cF\circ\operatorname{Hom}$, one therefore obtains isomorphisms
\[
  E_2^{p,q}
  \cong
  \operatorname{Ext}^{p}_{\cF}(H^{-q}({L'}^\bullet),A).
\]

\noindent\emph{Source note.} The French TeX and printed display have
$H^q({L'}^\bullet)$ in the last formula. The preceding displayed
isomorphism, the statement of Lemma~1.2, and the identity below require
$H^{-q}({L'}^\bullet)$. This target therefore uses $H^{-q}$ as a transparent
editorial emendation. It is not a restoration from an independent original
scan; the French authority remains unchanged pending editorial action.

Now
\[
  {L'}^\bullet=\operatorname{Hom}^\bullet(L^\bullet,A^\bullet),
\]
where $L^\bullet$ is a projective resolution of $M$; hence there are
isomorphisms
\[
  \operatorname{Ext}^{q}(M,A)\cong H^q({L'}^\bullet).
\]
Substituting $-q$ in this last identity gives the conclusion.\hfill$\square$
% END COMPONENT SGA2-VIII-CUM-C09

% VERIFIED SOURCE SEAM: French line 2639 is blank and excluded.

% INTEGRATION CONTROL BEFORE SGA2-VIII-CUM-C10
\let\cF\calF\renewcommand{\theequation}{\thesection.\arabic{equation}}

% BEGIN COMPONENT SGA2-VIII-CUM-C10
\setcounter{section}{1}
\section{The Finiteness Theorem}\label{VIII.2}

\begin{theorem}\label{VIII.2.1}
\markedfootnote{(1)}{\textit{Editor's note.} For an analogous statement, but
in a somewhat more general situation, see Mme Raynaud (Raynaud M., ``Th\'eor\`emes de
Lefschetz en cohomologie des faisceaux coh\'erents et en cohomologie
\'etale. Application au groupe fondamental,'' \emph{Ann. Sci. \'Ec. Norm.
Sup.} (4) \textbf{7} (1974), pp.~29--52, Proposition~II.2.1).}%
\kern-.18em\textbf{.}\enspace Let $X$ be a locally Noetherian prescheme, let $Y$ be a closed subset of $X$,
and let $\cF$ be a coherent $\OX$-module. Suppose that $X$ is locally
embeddable in a regular prescheme\markedfootnote{*}{This condition may be
generalized to the hypothesis that a ``dualizing complex,'' in the sense
defined in R.~Hartshorne, \emph{Residues and Duality} (cited in the note to
Expos\'e~IV, p.~46), exists locally on $X$.}. Let $i\in\ZZ$. Suppose that:
\begin{enumerate}[label=\textup{\alph*)},leftmargin=2.2em,itemsep=0.5\baselineskip]
\item for every $x\in U=X-Y$, one has
\begin{equation*}
  H^{i-c(x)}(\cF_x)=0,
\end{equation*}
where we have set\markedfootnote{(2)}{\textit{Editor's note.} As in
Expos\'e~V, $H^{\ast}(\cF_x)$ denotes the local cohomology
$H^{\ast}_{\mathfrak m_x}(\cF_x)$.}:
\begin{equation}\label{eq:VIII.2.1}
  c(x)=\operatorname{codim}
  (\overline{\{x\}}\cap Y,\overline{\{x\}}).
\end{equation}
Then:
\item $\mathcal H^i_Y(\cF)$ is coherent.
\end{enumerate}
\end{theorem}
% END COMPONENT SGA2-VIII-CUM-C10

% VERIFIED SOURCE SEAM: French line 2660 is blank and excluded.

% BEGIN COMPONENT SGA2-VIII-CUM-C11
\setcounter{section}{2}
\setcounter{sourcecorollary}{1}
\begin{sourcecorollary}{\textit{Editor's note.} Strictly speaking, this is a
corollary of the proof that follows, not of the statement. The implication
\textup{c)}$\Rightarrow$\textup{a)} is tautological. The converse is not, but
follows from the proof. More precisely, as below, one covers $X$ by open
subsets embeddable in regular schemes; as explained below, this permits one to
reduce to the case where $X=\operatorname{Spec}(A)$ is affine and regular and
$\cF=\widetilde M$, where $M$ is an $A$-module of finite projective dimension.
In this case it is shown that conditions \textup{a)} and \textup{c)} are
equivalent to the dual conditions \textup{a$'$)} and \textup{c$'$)}. One then
shows that \textup{c$'$)} implies condition \textup{d)} (see below), which in
turn implies \textup{a$'$)}. See the discussion following~\hyperref[VIII.2.4]{2.4}.}\label{VIII.2.2}
Under the hypotheses of the preceding theorem, condition \textup{a)} is
equivalent to:
\begin{enumerate}[label=\textup{\alph*)},start=3,leftmargin=2.2em]
\item for every $x\in U$ such that $c(x)=1$, one has
  $H^{i-1}(\cF_x)=0$.
\end{enumerate}
\end{sourcecorollary}
% END COMPONENT SGA2-VIII-CUM-C11

% VERIFIED SOURCE SEAM: French line 2669 is blank and excluded.

% BEGIN COMPONENT SGA2-VIII-CUM-C12
\setcounter{section}{2}
\setcounter{sourcecorollary}{2}
\begin{sourcecorollary}\label{VIII.2.3}
Let $X$ be a locally Noetherian prescheme that is locally embeddable in a
regular prescheme, let $Y$ be a closed subset of $X$, let $\cF$ be a coherent
$\mathcal O_X$-module, and let $n$ be an integer. The following conditions are
equivalent:
\begin{enumerate}[label=\textup{(\roman*)},leftmargin=2.8em]
\item for every $x\in U$, one has
  $\prof \cF_x>n-c(x)$;
\item for every $x\in U$ such that $c(x)=1$, one has
  $\prof \cF_x\geq n$;
\item for every $i\in\mathbb Z$, $\mathcal H^i_Y(\cF)$ is coherent if
  $i\leq n$;
\item $R^p i_*(\cF|_U)$ is coherent for $p<n$, where $i\colon U\to X$ is the canonical immersion%
  \markedfootnote{(4)}{\textit{Editor's note.} This condition appeared only
  in the body of the proof, but not in the statement of the corollary; since
  it is used in \hyperref[VIII.3]{Section~3}, it has been added.}.
\end{enumerate}
\end{sourcecorollary}
% END COMPONENT SGA2-VIII-CUM-C12

% VERIFIED SOURCE SEAM: French line 2682 is blank and excluded.

% BEGIN COMPONENT SGA2-VIII-CUM-C13
\medskip
Assume these results established when $X$ is the spectrum of a regular
Noetherian ring $A$ and when $\cF$ is the sheaf associated with an $A$-module
of finite projective dimension.

First observe that, if $(X_j)_{j\in J}$ is an open covering of $X$ by open
subsets embeddable in a regular scheme, each of the conditions above is
equivalent to the conjunction of the analogous conditions obtained by
replacing $X$ by $X_j$, $Y$ by $Y_j=Y\cap X_j$, and $\cF$ by
$\cF|_{X_j}$. Indeed, only the conditions involving $c(x)$ can present a
difficulty. Let $x\in U$. If $x\in X_j$, set
\[
  c_j(x)=\codim(X_j\cap\overline{\{x\}}\cap Y,
                X_j\cap\overline{\{x\}}).
\]
Necessarily $c_j(x)\geq c(x)$. Let $y\in\overline{\{x\}}\cap Y$ be a point
that ``realizes the codimension,'' that is, such that
$c(x)=\dim\mathcal O_{\overline{\{x\}},y}$, and let $X_j$ be a member of
the covering such that $y\in X_j$. Then $x\in X_j$, and hence
$c_j(x)=c(x)$. It follows that condition~(a) for the $X_j$ implies
condition~(a) for $X$. At this point we have only a partial converse: condition
(a) for $X$ implies condition~(a) for those $X_j$ such that
$c(x)=c_j(x)$, which suffices for our purposes.%
\markedfootnote{(5)}{\textit{Editor's note.} In fact, condition~(a) for $X$
implies condition~(a) for every $X_j$, as asserted in the original text, but
establishing this requires reading the proof below in detail. At this point
the implication does not appear formal. For a property (a), (b), or (c),
append a subscript $J$ to denote its conjunction over the $X_j$. The proof
\textit{below} establishes $(c_J)\Rightarrow(a_J)$ (this is the chain of implications
$(c')\Rightarrow(d)\Rightarrow(a')$). But tautologically
$(a)\Rightarrow(c)$, and $(c)\Longleftrightarrow(c_J)$; hence
$(a)\Rightarrow(a_J)$.}
% END COMPONENT SGA2-VIII-CUM-C13

% VERIFIED SOURCE SEAM: French line 2714 is blank and excluded.

% INTEGRATION CONTROL BEFORE SGA2-VIII-CUM-C14
\setcounter{section}{2}\setcounter{lemma}{3}

% BEGIN COMPONENT SGA2-VIII-CUM-C14
\medskip
Choose a covering of $X$ by open subsets embeddable in a regular prescheme.
Applying the preceding observation, we see that we may assume that $X$ is
closed in a regular $X'$. The reduction to $X'$ is then immediate.

We may therefore assume that $X$ is regular, and even affine, by covering
$X$ with affine open subsets. That we may assume that $\cF=\widetilde M$, where
$M$ has finite projective dimension, will follow from the next lemma.

\begin{lemma}\label{VIII.2.4}
Let $X$ be a regular Noetherian prescheme. Let $\cF$ be a coherent
$\mathcal O_X$-module. The function that assigns to each $x\in X$ the
projective dimension of $\cF_x$ is bounded above.
\end{lemma}
% END COMPONENT SGA2-VIII-CUM-C14

% VERIFIED SOURCE SEAM: French line 2722 is blank and excluded.

% BEGIN COMPONENT SGA2-VIII-CUM-C15
\setcounter{equation}{1}
\medskip
Indeed, let $x\in X$, and let $U$ be an affine open neighborhood of $x$.
Let $L^\bullet$ be a projective resolution of the module $\cF(U)$, where the
$L^i$ are finitely generated. By hypothesis, the ring $\cO_{X,x}$ is regular,
so the projective dimension of $\cF_x$ is finite; denote this integer by $d$.
Set
\[
  K=\operatorname{im}(L^{-d}\longrightarrow L^{-d+1})=\ker(L^{-d+1}\longrightarrow L^{-d+2}).
\]
The module $K_x$ is free, since $d$ is the projective dimension of $\cF_x$
([M], Ch.~VI, Prop.~2.1). By (EGA $0_{\mathrm I}$, 5.4.1 Errata), it follows
that the $\cO_U$-module $\widetilde K$ is free on a neighborhood $U'$ of $x$,
with $U'\subset U$. Choose $f\in\cO_X(U)$ such that
$x\in D(f)\subset U'$. We then have a projective resolution of $M_f$, where
$M=\cF(U)$:
\[
  0\longrightarrow K_f\longrightarrow (L^{-d+1})_f\longrightarrow
  \cdots\longrightarrow M_f\longrightarrow 0.
\]

\noindent\emph{Source note.} The printed French kernel and the earlier
maintained-English emendation both misplace the $d$-th syzygy. Exactness
gives $K=\operatorname{im}(L^{-d}\to L^{-d+1})$
$=\ker(L^{-d+1}\to L^{-d+2})$, so the truncated resolution begins with
$(L^{-d+1})_f$. The printed $M_{f'}$ is also unbound after only $f$ is
chosen and is corrected to $M_f$. Canon correction SGA-SRC-001691
preserves the printed proof diplomatically and uses these type-correct
terms in the current branch.

This proves that the function under consideration is upper semicontinuous.
Since $X$ is quasi-compact, the conclusion follows.
% END COMPONENT SGA2-VIII-CUM-C15

% VERIFIED SOURCE SEAM: French line 2732 is blank and excluded.

% BEGIN COMPONENT SGA2-VIII-CUM-C16
\medskip
We henceforth assume that $X$ is affine, Noetherian, and regular, and that
$\cF=\widetilde M$, where $M$ is a finitely generated $A$-module, necessarily
of finite projective dimension. We shall proceed in several steps. First, we
find a condition~(d), equivalent to~(a), and prove that it is also equivalent
to~(c). Then, using the spectral sequence of the preceding subsection, we
prove $(d)\Rightarrow(b)$. It remains to prove
$(\mathrm{iii})\Rightarrow(\mathrm{ii})$; indeed,
$(\mathrm{i})\Leftrightarrow(\mathrm{ii})\Rightarrow(\mathrm{iii})$ follows
immediately from $(a)\Leftrightarrow(c)\Rightarrow(b)$.

Let $x\in U$. By hypothesis, $\cO_{X,x}$ is a regular local ring. If $D$
denotes the dualizing functor for $\cO_{X,x}$, it follows from (\hyperref[V.2.1]{V, 2.1)} that
\[
  D\,H^{i-c(x)}(\cF_x)\simeq
  \operatorname{Ext}^{d(x)-i}_{\cO_{X,x}}(\cF_x,\cO_{X,x}),
\]
where
\begin{equation}\label{eq:VIII.2.2}
  d(x)=\dim\cO_{X,x}+c(x)
      =\dim\cO_{X,x}
       +\codim(\overline{\{x\}}\cap Y,\overline{\{x\}}).
  \tag{2.2}
\end{equation}
Since $X$ is Noetherian and $\cF$ is coherent,
\begin{equation}\label{eq:VIII.2.3}
  D\,H^{i-c(x)}(\cF_x)\simeq
  (\sExt^{d(x)-i}_{\cO_X}(\cF,\cO_X))_x.
  \tag{2.3}
\end{equation}
Moreover, a module vanishes if and only if its dual vanishes (cf. the
editor's note~(4) on p.~54). For every $q\in\mathbb Z$, set
\begin{equation}\label{eq:VIII.2.4}
  \begin{cases}
    S_q=\Supp\sExt^q_{\cO_X}(\cF,\cO_X),\\
    S'_q=S_q\cap U,\quad (U=X-Y),\\
    Z_q=\overline{S'_q}\cap Y.
  \end{cases}
  \tag{2.4}
\end{equation}
% END COMPONENT SGA2-VIII-CUM-C16

% BEGIN COMPONENT SGA2-VIII-CUM-C17
\medskip
It follows from \hyperref[eq:VIII.2.3]{formula~(2.3)} that conditions~(a) and~(c) are respectively
equivalent to:

\medskip
\noindent\textup{(a$'$)} for every $q\in\mathbb Z$ and every $x\in S'_q$,
one has $q+i\ne d(x)$;

\smallskip
\noindent\textup{(c$'$)} for every $q\in\mathbb Z$ and every $x\in S'_q$,
if $c(x)=1$, then $q+i\ne d(x)$.

\medskip
\noindent Here is condition~(d), promised above:

\medskip
\noindent\textup{(d)} for every $q\in\mathbb Z$ and every $y\in Z_q$,
one has $q+i\ne\dim\cO_{X,y}$.

\medskip
\noindent These conditions are equivalent.

\smallskip
\noindent\textup{(a$'$)} $\Rightarrow$ \textup{(c$'$)} is recorded for
completeness.

\smallskip
\noindent\textup{(d)} $\Rightarrow$ \textup{(a$'$)}. Indeed, let
$q\in\mathbb Z$ and $x\in S'_q$. Choose a point
$y\in\overline{\{x\}}\cap Y$ that realizes the
codimension,\markedfootnote{(6)}{\textit{Editor's note.} Throughout what
follows, closures of points are endowed with the reduced structure.} that is,
such that
\begin{equation}\label{eq:VIII.2.5}
  \dim\cO_{\overline{\{x\}},y}
  =\codim(\overline{\{x\}}\cap Y,\overline{\{x\}})=c(x).
  \tag{2.5}
\end{equation}
Since $X$ is regular at $y$, it follows that
\begin{equation}\label{eq:VIII.2.6}
  \dim\cO_{X,y}=d(x)\qquad\text{(cf.~\hyperref[eq:VIII.2.2]{(2.2)}).}
  \tag{2.6}
\end{equation}
But $y\in\overline{\{x\}}$, and hence $y\in Z_q$; the conclusion follows.

\smallskip
\noindent\textup{(c$'$)} $\Rightarrow$ \textup{(d)}. Let $q\in\mathbb Z$
and $y\in Z_q$. Assume provisionally that there exists $x\in S'_q$ such
that
\[
  y\in\overline{\{x\}}\qquad\text{and}\qquad
  \dim\cO_{\overline{\{x\}},y}=1.
\]
(One also says that $y$ is an immediate specialization of $x$.) It follows
that $c(x)=1$, since $y$ realizes the codimension of
$\overline{\{x\}}\cap Y$ in $\overline{\{x\}}$, while $x\notin Y$.
Condition~\textup{(c$'$)} therefore gives
\[
  q+i\ne d(x).
\]
Whence the conclusion, upon observing that $d(x)=\dim\cO_{X,y}$ by~\hyperref[eq:VIII.2.6]{(2.6)}.
The assertion provisionally assumed above is the content of the following
lemma.

\medskip
\sgareaderlabel{VIII.2.5}{2.5}\noindent\textbf{Lemma 2.5.} Let $X$ be a locally Noetherian prescheme and
$Y$ a closed subset of $X$. Set $U=X-Y$, and suppose that $U$ is dense in
$X$. For every $y\in Y$, there exists $x\in U$ such that $y$ is an immediate
specialization of $x$, that is,
\[
  y\in\overline{\{x\}}\qquad\text{and}\qquad
  \dim\cO_{\overline{\{x\}},y}=1.
\]

We have applied the lemma with $X$ equal to the prescheme
$\overline{S'_q}$ and $Y$ equal to the subset
$Y\cap\overline{S'_q}$.

\medskip
\noindent\textit{Proof of \hyperref[VIII.2.5]{Lemma 2.5}.} There exists $x\in U$ such that
$y\in\overline{\{x\}}$. Choose such an $x$ for which
$r=\dim\cO_{\overline{\{x\}},y}$ is minimal. We must prove that $r=1$.
By the choice of $x$, every
$z\in\Spec(\cO_{\overline{\{x\}},y})$ with $z\ne x$ lies in $Y$.
Consequently, $\{x\}$ is open in
$\Spec(\cO_{\overline{\{x\}},y})$, and the conclusion follows.\hfill$\square$
% END COMPONENT SGA2-VIII-CUM-C17

% VERIFIED SOURCE SEAM: French line 2795 is blank and excluded.

% BEGIN COMPONENT SGA2-VIII-CUM-C18
\medskip
\noindent\textit{The second step is to deduce \textup{(b)} from
\textup{(d)}.}

\medskip
Set
\[
  D(Z_q)=\{\dim\cO_{X,y}\mid y\in Z_q\}.
\]
By~\textup{(d)}, for every $q\in\mathbb Z$ one has
$q+i\notin D(Z_q)$. Applying Proposition~\hyperref[VII.2.3]{VII.2.3}, we see that
\[
  \sExt_Y^{q+i}(\sExt^q(\cF,\cO_X),\cO_X)
  \quad\text{is coherent.}
\]
The initial term of the spectral sequence of the preceding subsection is
\[
  E_2^{p,q}
  =\sExt_Y^p(\sExt^{-q}(\cF,\cO_X),\cO_X).
\]
It follows that $E_2^{p,q}$ is coherent for every $p,q\in\mathbb Z$ such
that $p+q=i$. Now \sisi{there are only finitely many pairs $(p,q)$ such that
$p+q=i$}{only finitely many pairs $(p,q)$ with $p+q=i$ contribute nonzero $E_2$ terms},\markedfootnote{[S]}{\textit{Source note.} Taken literally for
arbitrary integers, this sentence is false. In context it must be read as
concerning the terms that can contribute, in view of the finite-projective-dimension
setup above. The French authority is unchanged.} and this spectral
sequence converges to $\sH_Y^\ast(\cF)$; whence the conclusion.
% END COMPONENT SGA2-VIII-CUM-C18

% VERIFIED SOURCE SEAM: French line 2806 is blank and excluded.

% BEGIN COMPONENT SGA2-VIII-CUM-C19
\medskip
\noindent\textit{It remains to prove that \textup{(iii)} implies
\textup{(ii)}.} Let
\[
  i\colon U\longrightarrow X
\]
denote the canonical immersion of $U$ into $X$. In view of the exact homology
sequence associated with the closed subset $Y$ (\hyperref[I.2.11]{I, 2.11)}, one sees that
\textup{(iii)} is \textit{equivalent to}:

\medskip
\noindent\textup{(iv)} $R^i i_*(\cF|_U)$ is coherent for $i<n$.
\refstepcounter{sourcecondition}\label{condition}

\medskip
\noindent{\scriptsize\color{gray}French printed p.~94}\quad Indeed, there is
an exact sequence
\[
  0\longrightarrow \sH_Y^0(\cF)\longrightarrow \cF
  \longrightarrow i_*(\cF|_U)\longrightarrow \sH_Y^1(\cF)
  \longrightarrow 0.
\]
Now $\sH_Y^0(\cF)$ is a quasi-coherent subsheaf of the coherent sheaf $\cF$,
and is therefore coherent. Hence $\sH_Y^1(\cF)$ is coherent if and only if
$i_*(\cF|_U)$ is coherent. Moreover, for $p>0$, the exact cohomology sequence
associated with the closed subset $Y$ reduces to isomorphisms
\[
  R^p i_*(\cF|_U)\overset{\approx}{\longrightarrow}\sH_Y^{p+1}(\cF).
\]
% END COMPONENT SGA2-VIII-CUM-C19

% BEGIN COMPONENT SGA2-VIII-CUM-C20
\medskip
\noindent\textit{We shall prove that \textup{(iv)} implies
\textup{(ii)}.} To this end, recall \textup{(ii)}:

\medskip
\noindent\textup{(ii)} For every $x\in U$ such that $c(x)=1$, one has
\[
  \operatorname{prof}\cF_x\geq n.
\]

\noindent We argue by induction on $n$.

\medskip
\noindent If $n=0$, both conditions are vacuous.
% END COMPONENT SGA2-VIII-CUM-C20

% VERIFIED SOURCE SEAM: French line 2827 is blank and excluded.

% BEGIN COMPONENT SGA2-VIII-CUM-C21
\medskip
\noindent If $n=1$, assume that $i_*(\cF|_U)$ is coherent. We argue by
contradiction and suppose that there exists $x\in U$ such that $c(x)=1$ and
\[
  \prof\cF_x=0,
\]
i.e. $x\in\Ass\cF_x$. Let $y\in\overline{\{x\}}\cap Y$ be such that
\[
  \dim\cO_{\overline{\{x\}},y}=1.
\]
Set
\[
  A=\cO_{X,y}\qquad\text{and}\qquad X'=\Spec(A).
\]

\setcounter{equation}{6}
\noindent Perform the base change $v\colon X'\longrightarrow X$, which is
flat:
\begin{equation}
\begin{CD}
U'=X'\times_X U @>{v'}>> U\\
@V{i'}VV @VV{i}V\\
X' @>{v}>> X .
\end{CD}
\tag{2.7}\label{eq:VIII.2.7}
\end{equation}
The morphism $i$ is separated, since it is an immersion, and is of finite
type, since it is an open immersion and $X$ is locally Noetherian. Since the
base change is flat, EGA~III, 1.4.15 gives an isomorphism
\begin{equation}
  v^*(i_*(\cF|_U))
  \approx
  i'_*(v'^*(\cF|_U)).
\tag{2.8}\label{eq:VIII.2.8}
\end{equation}

\noindent{\scriptsize\color{gray}French printed p.~95}\quad Denote by
$\underline{x}$ (respectively, $\underline{y}$) the ideal of $A$
corresponding to $x$ (respectively, $y$).

\medskip
\noindent Set $\cG=v'^*(\cF|_U)$. Then $\cG$ is coherent and
$\underline{x}\in\Ass\cG$, so there exists a monomorphism
\[
  \cO_{\overline{\{x\}}}\longrightarrow\cG;
\]
consequently, $i'_*(\cO_{\overline{\{x\}}}|_{U'})$ is coherent. By the choice
of $y$, $\dim(A/\underline{x})=1$, and therefore the support of
$\cO_{\overline{\{x\}}}$ is precisely
\[
  \overline{\{\underline{x}\}}
  =\{\underline{x}\}\cup\{\underline{y}\},
\]
since $\overline{\{\underline{x}\}}=\Spec(A/\underline{x})$ as a scheme. It
follows that
\[
  (\cO_{\overline{\{x\}}}|_{U'})(U')
  =\Frac(A/\underline{x}),
\]
the field of fractions of $A/\underline{x}$, and
\[
  i'_*(\cO_{\overline{\{x\}}}|_{U'})(X')
  =\Frac(A/\underline{x}).
\]
But $\Frac(A/\underline{x})$ is not a finitely generated $A$-module, because
$\underline{x}$ is different from the maximal ideal of $A$. This is a
contradiction.
% END COMPONENT SGA2-VIII-CUM-C21

% VERIFIED SOURCE SEAM: French line 2855 is blank and excluded.

% BEGIN COMPONENT SGA2-VIII-CUM-C22
\medskip
\noindent Suppose that $n>1$, and assume the result established for every
$n'<n$. By the induction hypothesis, for every $x\in U$ such that $c(x)=1$,
one has $x\notin\Ass\cF_x$. Let $x$ be such a point, and let
$y\in\overline{\{x\}}\cap Y$ be an immediate specialization of $x$, i.e.
\[
  \dim\cO_{\overline{\{x\}},y}=1.
\]
Perform the base change
\[
  v\colon\Spec(\cO_{X,y})\longrightarrow X,
\]
retaining the notation of diagram~\textup{\hyperref[eq:VIII.2.7]{(2.7)}}. Applying EGA~III, 1.4.15,
one obtains isomorphisms
\[
  v^*(R^p i_*(\cF|_U))
  \simeq
  R^p i'_*(v'^*(\cF|_U)),
  \qquad p\in\mathbb Z.
\]

\noindent\emph{Thus we are reduced to the case in which $X$ is the spectrum
of a local ring $A$ in which $\underline{x}$ is a prime ideal of dimension
$1$}, i.e. $\dim(A/\underline{x})=1$. Set
\[
  \cF'=\SheafGamma_Y(\cF)
  \qquad\text{and}\qquad
  \cF''=\cF/\cF'.
\]
We have $\cF_x\simeq\cF''_x$ and $y\notin\Ass\cF''$. Moreover,
$\cF'|_U=0$; hence the long exact sequence of the $R^p i_*$ gives
isomorphisms
\[
  R^p i_*(\cF|_U)
  \simeq
  R^p i_*(\cF''|_U),
  \qquad p\in\mathbb Z.
\]

\noindent{\scriptsize\color{gray}French printed p.~96}\quad Since $n>1$,
neither $x$ nor $y$ belongs to $\Ass\cF''$. Now $\underline{x}$ and
$\underline{y}$ are the only prime ideals of $A$ containing $\underline{x}$.
It follows from Expos\'e~\hyperref[III.2.1]{III, 2.1} that there exists an element
$g\in\underline{x}$ that is $M$-regular, where we have set
\[
  \cF''=\widetilde M,
  \qquad M=\cF''(X).
\]
Thus there is an exact sequence
\[
  0\longrightarrow M\xrightarrow{g'}M\longrightarrow N\longrightarrow0,
\]
in which $g'$ denotes multiplication by $g$ on $M$. From the resulting long
exact sequence one sees that
\[
  R^p i_*(\widetilde N|_U)
  \quad\text{is coherent for }p<n-1.
\]
Therefore, by the induction hypothesis,
\[
  \prof(\widetilde N)_x\geq n-1,
  \qquad\text{and hence}\qquad
  \prof\cF_x\geq n.
\]
\hfill$\square$
% END COMPONENT SGA2-VIII-CUM-C22

% VERIFIED SOURCE SEAM: French line 2873 is blank and excluded.

% BEGIN COMPONENT SGA2-VIII-CUM-C23
\section*{3. Applications}
\sgareaderlabel{VIII.3}{3}

\noindent From these results one deduces a coherence condition for the higher
direct images of a coherent sheaf under a \emph{morphism that is not proper}.

\medskip
\sgareaderlabel{VIII.3.1}{3.1}\noindent\textbf{Theorem 3.1.}---Let $f\colon X\to Y$ be a morphism of
preschemes. Suppose that $Y$ is \emph{locally Noetherian} and that $f$ is
\emph{proper}. Suppose that $X$ is \emph{locally embeddable in a regular
prescheme}. Let $n\in\mathbb Z$. Let $U$ be an open subset of $X$, and let
$\cF$ be a \emph{coherent} $\cO_U$-module. Suppose that, for every $x\in U$
such that
\[
  \codimn(\overline{\{x\}}\cap(X\setminus U),
  \overline{\{x\}})=1,
\]
one has $\prof\cF_x\geq n$. \emph{Then} the $\cO_Y$-modules
\[
  R^p(f\circ g)_*(\cF)
\]
are coherent for $p<n$, where $g$ is the canonical immersion of $U$ into
$X$.

\medskip
\noindent Indeed, there is a Leray spectral sequence whose abutment is
$R^*(f\circ g)_*(\cF)$ and whose initial term is given by
\[
  E_2^{p,q}=R^p f_*(R^q g_*(\cF)).
\]
Moreover, there exists a \emph{coherent} $\cO_X$-module $\cG$ such that
$\cG|_U\simeq\cF$ (EGA~I, 9.4.3). It then follows from the preceding section
that condition~\textup{(iv)} on p.~74 is satisfied
{\scriptsize\color{gray}[French printed p.~97]}, i.e. that
$R^q g_*(\cG|_U)$ is coherent for $q<n$. Applying the finiteness theorem of
EGA~III, 3.2.1 to $f$ and the sheaves $R^q g_*(\cF)$, one finds that
$E_2^{p,q}$ is coherent for $q<n$, whence the conclusion.
% END COMPONENT SGA2-VIII-CUM-C23

% VERIFIED SOURCE SEAM: French line 2887 is blank and excluded.

% BEGIN COMPONENT SGA2-VIII-CUM-C24
\medskip
\sgareaderlabel{VIII.3.2}{3.2}\noindent\textbf{Proposition 3.2.}---Let $X$ be a locally Noetherian
prescheme, locally embeddable in a regular prescheme. Let $U$ be an open
subset of $X$, and let $i\colon U\to X$ be the canonical immersion. Let
$n\in\mathbb Z$. Finally, let $\cF$ be a coherent Cohen--Macaulay
$\cO_U$-module (on $U$!). The following conditions are equivalent:
\begin{enumerate}
\item[(a)] $R^p i_*(\cF)$ is coherent for $p<n$;
\item[(b)] for every irreducible component $S'$ of the closure $\overline S$
of the support $S$ of $\cF$, one has
\[
  \codimn(S'\cap(X\setminus U),S')>n.
\]
\end{enumerate}
% END COMPONENT SGA2-VIII-CUM-C24

% VERIFIED SOURCE SEAM: French line 2900 is blank and excluded.

% BEGIN COMPONENT SGA2-VIII-CUM-C25
\medskip
\noindent\textit{Proof, opening.}---Let $\cG$ be a coherent $\cO_X$-module
such that $\cG|_U\simeq\cF$ (EGA~I, 9.4.3). Applying \hyperref[VIII.2.3]{Corollary~2.3} to
$\cG$, we find that condition \textup{(a)} is equivalent to \textup{(c)}:
\begin{enumerate}
\item[(c)] for every $x\in S$, one has
\[
  \prof\cF_x>n-c(x),
  \qquad
  c(x)=\codimn(\overline{\{x\}}\cap (X-U),\overline{\{x\}}).
\]
\end{enumerate}
% END COMPONENT SGA2-VIII-CUM-C25

% VERIFIED SOURCE SEAM: French line 2908 is blank and excluded.

% BEGIN COMPONENT SGA2-VIII-CUM-C26
\medskip
\noindent\textit{Proof, \textup{(a)}$\Rightarrow$\textup{(b)}.}---Indeed,
let $S'$ be an irreducible component of $\overline S$, and let $s$ be its
generic point. Since $\cF$ is Cohen--Macaulay, one has
$\prof \cF_s=\dim\cO_{S,s}=0$. Moreover, $\overline{\{s\}}=S'$, and the
conclusion follows.
% END COMPONENT SGA2-VIII-CUM-C26

% VERIFIED SOURCE SEAM: French line 2910 is blank and excluded.

% BEGIN COMPONENT SGA2-VIII-CUM-C27
\medskip
\noindent\textit{Proof, \textup{(b)}$\Rightarrow$\textup{(a)}, opening.}---Let
$x\in S$, and let $S'$ be an irreducible component of $\overline S$ such that
$x\in S'$. Let $s$ be the generic point of $S'$. Since $\cF$ is
Cohen--Macaulay, we know that
\[
  \prof \cF_x=\dim\cO_{\overline{\{s\}},x}.
\]
% END COMPONENT SGA2-VIII-CUM-C27

% BEGIN COMPONENT SGA2-VIII-CUM-C28
\medskip
\noindent\textit{Proof, \textup{(b)}$\Rightarrow$\textup{(a)}, case
split.}---If $c(x)=+\infty$, there is nothing to prove. Otherwise, there exists
$y\in (X-U)\cap\overline{\{x\}}$ such that
\[
  c(x)=\dim\cO_{\overline{\{x\}},y}.
\]
% END COMPONENT SGA2-VIII-CUM-C28

% BEGIN COMPONENT SGA2-VIII-CUM-C29
\medskip
\noindent\textit{Proof, \textup{(b)}$\Rightarrow$\textup{(a)},
conclusion.}---Now $\cO_{X,y}$ is a quotient of a regular local ring by
hypothesis; hence
\[
  \dim\cO_{\overline{\{s\}},y}
  =\dim\cO_{\overline{\{s\}},x}
   +\dim\cO_{\overline{\{x\}},y}>n.
\]
\hfill$\square$
% END COMPONENT SGA2-VIII-CUM-C29

% VERIFIED SOURCE SEAM: French line 2923 is blank and excluded.

% INTEGRATION CONTROL BEFORE SGA2-VIII-CUM-C30
\begingroup

% BEGIN COMPONENT SGA2-VIII-CUM-C30
\renewcommand{\refname}{Bibliography}
\begin{thebibliography}{M}
\bibitem[M]{M}
\textsc{H.~Cartan \& S.~Eilenberg} -- \emph{Homological Algebra},
Princeton Math. Series vol.~19, Princeton University Press, 1956.
\end{thebibliography}
\endgroup
% END COMPONENT SGA2-VIII-CUM-C30

% END INCLUDED active_source/parts/08_expose_VIII.tex

% BEGIN INCLUDED active_source/parts/09_expose_IX.tex
\makeatletter
\providecommand{\sgareaderlabel}[2]{%
  \begingroup
  \def\@currentlabel{#2}%
  \phantomsection\label{#1}%
  \endgroup}
\makeatother

\chapterspace{-1}
\chapter{Algebraic Geometry and Formal Geometry}
\label{IX}

\medskip
The purpose of this expos\'e is to generalize Theorems 3.4.2 and 4.1.5 of
\emph{EGA III} to the case of a morphism that is not proper.

\section*{1. The Comparison Theorem}
\sgareaderlabel{IX.1}{1}

Let $f\colon X\to X'$ be a morphism of preschemes that is
\emph{separated} and of \emph{finite type}. Suppose that $X'$ is locally
noetherian. Let $Y'$ be a closed subset of $X'$, and let
$Y=f^{-1}(Y')$. Let $\widehat X$ and $\widehat{X'}$ be the formal
completions of $X$ and $X'$ along $Y$ and $Y'$, respectively. Let
$\widehat f$ be the morphism induced by $f$ upon passing to the completions.

\begin{equation}
\tag{1.1}\label{eq:IX.1.1}
\begin{tikzcd}[column sep=3.1em,row sep=2.6em]
X \arrow[d,"f"'] & Y \arrow[l] \arrow[d] &&
\widehat{X} \arrow[d,"\widehat f"'] \arrow[r,"j"] & X \arrow[d,"f"] \\
X' & Y' \arrow[l] & {,} &
\widehat{X'} \arrow[r,"i"] & X'.
\end{tikzcd}
\end{equation}

We denote by $j$ (resp. $i$) the morphism from $\widehat X$ to $X$
(resp. from $\widehat{X'}$ to $X'$). It is known that $i$ and $j$ are
\emph{flat}.

Let $\mathcal I'$ be an ideal of definition of $Y'$, and let
$\mathcal J=f^*(\mathcal I')\cdot\mathcal O_X$; this is an ideal of
definition of $Y$. Hence
\begin{equation}
\tag{1.2}\label{eq:IX.1.2}
\widehat{X'}=
\left(Y',\varprojlim_{k\in\mathbb N}
\mathcal O_{X'}/{\mathcal I'}^{k+1}\right),
\qquad
\widehat X=
\left(Y,\varprojlim_{k\in\mathbb N}
\mathcal O_X/\mathcal J^{k+1}\right).
\end{equation}

For every $k\in\mathbb N$, set
\begin{equation}
\tag{1.3}\label{eq:IX.1.3}
Y'_k=\left(Y',\mathcal O_{X'}/(\mathcal I')^{k+1}\right),
\qquad
Y_k=\left(Y,\mathcal O_X/\mathcal J^{k+1}\right).
\end{equation}

Let $F$ be a coherent $\mathcal O_X$-module. For every $k\in\mathbb N$, set
\begin{equation}
\tag{1.4}\label{eq:IX.1.4}
F_k=F/\mathcal J^{k+1}F,
\qquad
\widehat F=j^*(F)\simeq\varprojlim F_k.
\end{equation}

If we set
\begin{equation}
\tag{1.5}\label{eq:IX.1.5}
R^i f_*(F)^{\wedge}
=\varprojlim_{k\in\mathbb N}
\left(R^i f_*(F)\otimes_{\mathcal O_{X'}}\mathcal O_{Y'_k}\right),
\qquad i\in\mathbb Z,
\end{equation}
there is a natural homomorphism
\begin{equation}
\tag{1.6}\label{eq:IX.1.6}
r_i:i^*(R^i f_*(F))\longrightarrow R^i f_*(F)^{\wedge},
\end{equation}
which is an \emph{isomorphism} when $R^i f_*(F)$ is \emph{coherent}.

As explained in \emph{EGA III}, 4.1.1, there is a commutative diagram
\begin{equation}
\tag{1.7}\label{eq:IX.1.7}
\begin{tikzcd}[column sep=5.2em,row sep=3.2em]
i^*(R^i f_*(F))
  \arrow[r,"\rho_i"]
  \arrow[d,"r_i"]
& R^i\widehat f_*(\widehat F)
  \arrow[d,"\psi_i"] \\
R^i f_*(F)^{\wedge}
  \arrow[r,"\varphi_i"]
& \displaystyle\varprojlim_{k\in\mathbb N} R^i f_*(F_k).
\end{tikzcd}
\end{equation}

In the cited passage one obtains a commutative diagram because
$R^i f_*(F)$ is known to be coherent, and the source and target of~\hyperref[eq:IX.1.6]{(1.6)}
are identified. In our case, $R^i f_*(F)$ will be coherent only for certain
values of $i$, for which we shall study~\hyperref[eq:IX.1.7]{(1.7)}.

Consider the graded ring
\begin{equation}
\tag{1.8}\label{eq:IX.1.8}
\mathcal S=\bigoplus_{k\in\mathbb N}{\mathcal I'}^k,
\end{equation}
and the graded $\mathcal S$-module
\begin{equation}
\tag{1.9}\label{eq:IX.1.9}
\mathcal H^i=\bigoplus_{k\in\mathbb N}R^i f_*(\mathcal J^kF),
\qquad i\in\mathbb Z,
\end{equation}
whose $\mathcal S$-module structure is defined as follows.

The sheaf $R^i f_*(\mathcal J^kF)$ is associated with the presheaf that
assigns to each affine open subset $U'$ of $X'$
\begin{equation}
\tag{1.10}\label{eq:IX.1.10}
H^i\!\left(f^{-1}(U'),
  \left.\mathcal J^kF\right|_{f^{-1}(U')}\right).
\end{equation}
Thus let $U'$ be an affine open subset of $X'$, set
\[
U=f^{-1}(U'),
\]
and let $x'\in{\mathcal I'}^m(U')$. Let $x$ be the image of $x'$ in
$\mathcal J^m(U)$. Multiplication by $x$ on $F|_U$ maps
$\mathcal J^kF|_U$ into $\mathcal J^{k+m}F|_U$ and hence, by
functoriality, gives a morphism
\begin{equation}
\tag{1.11}\label{eq:IX.1.11}
\mu^i_{x',k}(U'):
H^i(U,\mathcal J^kF|_U)
\longrightarrow
H^i(U,\mathcal J^{k+m}F|_U),
\end{equation}
defined for every $i\in\mathbb Z$ and every $k\in\mathbb N$. Upon passage
to the associated sheaf, these morphisms give $\mathcal H^i$ its graded
$\mathcal S$-module structure.

\begin{theoremIXoneone}\sgareaderlabel{IX.1.1}{1.1}
Let $n$ be an integer. Suppose that the graded $\mathcal S$-module
$\mathcal H^i$ is finitely generated for $i=n-1$ and $i=n$. Then:
\begin{enumerate}
\renewcommand{\labelenumi}{(\arabic{enumi})}
\setcounter{enumi}{-1}
\item
$r_n$ and $r_{n-1}$ are isomorphisms, and
$R^i\widehat f_*(\widehat F)$ is coherent for $i=n-1$;

\item
for $i=n-1$, the maps $\rho_i$, $\varphi_i$, and $\psi_i$ are topological
isomorphisms. In particular, the filtration on $R^{n-1}f_*(F)$ defined by
the kernels of the homomorphisms
\begin{equation}
\tag{1.12}\label{eq:IX.1.12}
R^{n-1}f_*(F)\longrightarrow R^{n-1}f_*(F_k)
\end{equation}
is $\mathcal I'$-good;

\item
for $i=n$, the maps $\rho_i$, $\varphi_i$, and $\psi_i$ are monomorphisms;
moreover, the filtration on $R^n f_*(F)$ is $\mathcal I'$-good, and $\psi_n$
is an isomorphism;

\item
for $i=n-2,n-1$, the projective system $R^i f_*(F_k)$ satisfies the uniform
Mittag--Leffler condition; that is, there is a \emph{fixed} integer $k\geq0$
such that, for every $p\geq0$ and every $p'\geq p+k$,
\[
\operatorname{Im}\!\left[R^i f_*(F_{p'})\longrightarrow R^i f_*(F_p)\right]
=
\operatorname{Im}\!\left[R^i f_*(F_{p+k})\longrightarrow R^i f_*(F_p)\right].
\]
\end{enumerate}
\end{theoremIXoneone}

Proceeding as in \emph{EGA III}, 4.1.8, it is easy to \emph{reduce to the case
where $X'$ is the spectrum of a noetherian ring} $A$. In this case, it is
known that
\begin{equation}
\tag{1.13}\label{eq:IX.1.13}
R^i f_*(F)=\widetilde{H^i(X,F)}
\qquad\text{(cf.~\hyperref[eq:IX.1.10]{(1.10)}).}
\end{equation}
Let $I$ be the ideal of $A$ such that $\widetilde I=\mathcal I'$, and set
\begin{equation}
\tag{1.14}\label{eq:IX.1.14}
S=\bigoplus_{k\in\mathbb N}I^k,
\end{equation}
\begin{equation}
\tag{1.15}\label{eq:IX.1.15}
H^i=\bigoplus_{k\in\mathbb N}H^i(X,\mathcal J^kF),
\qquad i\in\mathbb Z,
\end{equation}
where $H^i$ carries the graded $S$-module structure defined by~\hyperref[eq:IX.1.11]{(1.11)}, with
$U'=X'$.
\footnote{The corrected French TeX and same-edition PDF print $U=X'$ here.
In the construction preceding~(1.11), however, $U=f^{-1}(U')$; specializing
the affine open to the whole base therefore requires $U'=X'$. The French
authority is left unchanged.}

The proof is modeled on that of \emph{EGA III}, 4.1.5; let us give a summary.

We now work with $\varphi_i$ and $\psi_i$, to which the following module
homomorphisms correspond:
\begin{equation}
\tag{1.16}\label{eq:IX.1.16}
\begin{tikzcd}[column sep=5.2em,row sep=3.0em]
& H^i(\widehat X,\widehat F)
    \arrow[d,"\psi_i"] \\
H^i(X,F)^{\wedge}
    \arrow[r,"\varphi_i"]
& \displaystyle\varprojlim_{k\in\mathbb N}H^i(X,F_k).
\end{tikzcd}
\end{equation}

\noindent\textup{(a)} Suppose only that $H^i$ is a graded $S$-module of finite
type. It follows that \emph{the filtration on $H^i(X,F)$ defined by the
modules}
\begin{equation}
\tag{1.17}\label{eq:IX.1.17}
R_k^i=\ker\bigl(H^i(X,F)\longrightarrow H^i(X,F_k)\bigr)
\end{equation}
\emph{is $I$-good}. Indeed, the exact cohomology sequence
\begin{equation}
\tag{1.18}\label{eq:IX.1.18}
H^i(X,\mathcal J^{k+1}F)\longrightarrow H^i(X,F)
  \longrightarrow H^i(X,F_k)
\end{equation}
shows that the \emph{graded $S$-module}
$\bigoplus_{k\in\mathbb N}R_k^i$ \emph{is a quotient} of the graded
$S$-submodule
\[
\bigoplus_{k\in\mathbb N}H^i(X,\mathcal J^{k+1}F)
\]
of $H^i$, and hence \emph{is of finite type}, since $S$ is noetherian. This
proves the first assertion.

Set
\begin{equation}
\tag{1.19}\label{eq:IX.1.19}
M^i=H^i(X,F),\qquad H_k^i=H^i(X,F_k).
\end{equation}
There is a commutative diagram
\begin{equation}
\tag{1.20}\label{eq:IX.1.20}
\begin{tikzcd}[column sep=5.5em,row sep=3.2em]
H^i(X,F)^{\wedge}
  \arrow[r,"s_i"]
  \arrow[dr,"\varphi_i"']
& \displaystyle\varprojlim_k(M^i/R_k^i)
  \arrow[d,"t_i"] \\
& \displaystyle\varprojlim_k H_k^i,
\end{tikzcd}
\end{equation}
in which $s_i$ is an \emph{isomorphism}, since the filtration of $H^i(X,F)$
is $I$-good. Moreover, $t_i$ is a \emph{monomorphism}: the functor
$\varprojlim$ is left exact and, for every $k\geq 0$, the natural morphism
$M^i/R_k^i\longrightarrow H_k^i$ is a monomorphism by the definition of
$R_k^i$.

To study the surjectivity of $t_i$, introduce
\begin{equation}
\tag{1.21}\label{eq:IX.1.21}
Q_k^i=\operatorname{coker}\bigl(H^i(X,F)\longrightarrow H^i(X,F_k)\bigr),
\end{equation}
which gives a projective system of exact sequences
\begin{equation}
\tag{1.22}\label{eq:IX.1.22}
0\longrightarrow R_k^i\longrightarrow M^i\longrightarrow H_k^i
 \longrightarrow Q_k^i\longrightarrow 0.
\end{equation}
Using the exact cohomology sequence
\begin{equation}
\tag{1.23}\label{eq:IX.1.23}
H^i(X,F)\longrightarrow H^i(X,F_k)
 \longrightarrow H^{i+1}(X,\mathcal J^{k+1}F),
\end{equation}
we see that the \emph{graded $S$-module}
\begin{equation}
\tag{1.24}\label{eq:IX.1.24}
Q^i=\bigoplus_{k\in\mathbb N}Q_k^i
\end{equation}
\emph{is a graded $S$-submodule of $H^{i+1}$}. Moreover, for every
$k\geq0$,
\begin{equation}
\tag{1.25}\label{eq:IX.1.25}
I^{k+1}Q_k^i=0,
\end{equation}
because $Q_k^i$ is the image of $H_k^i$.

\sgareaderlabel{stepb}{b}\noindent\textup{(b)} Suppose only that $H^{i+1}$ is of finite type, and
consider $t_i$ alone, disregarding $s_i$. Since $S$ is noetherian, $Q^i$ is of
finite type; as $I^{k+1}Q_k^i=0$, there exist integers $r\geq0$ and $k_0\geq0$
such that
\begin{equation}
\tag{1.26}\label{eq:IX.1.26}
I^rQ_k^i=0\qquad\text{for }k\geq k_0.
\end{equation}
It follows that the \emph{projective system $(Q_k^i)_{k\in\mathbb N}$ is
essentially zero}, and hence that the \emph{projective system
$(H_k^i)_{k\in\mathbb N}$ satisfies the uniform Mittag--Leffler condition}.
From the exact sequence~\hyperref[eq:IX.1.22]{(1.22)} we obtain the exact sequence
\begin{equation}
\tag{1.27}\label{eq:IX.1.27}
0\longrightarrow M^i/R_k^i\longrightarrow H_k^i
 \longrightarrow Q_k^i\longrightarrow0,
\end{equation}
and therefore the exact sequence
\begin{equation}
\tag{1.28}\label{eq:IX.1.28}
0\longrightarrow\varprojlim_k M^i/R_k^i
 \xrightarrow{t_i}\varprojlim_k H_k^i
 \longrightarrow\varprojlim_k Q_k^i.
\end{equation}
But the projective system $(Q_k^i)_{k\in\mathbb N}$ is essentially zero, so
$t_i$ is an \emph{isomorphism}.

\noindent\textup{(c)} Let us prove that if $H^i$ is of finite type, then
$\psi_i$ is an isomorphism. It suffices to apply \emph{EGA}~$0_{\mathrm{III}}$,
13.3.1, taking the affine open subsets as a basis of open subsets of $X$.
This is legitimate: by~(b), the projective system
$(H_k^{i-1})_{k\in\mathbb N}$ satisfies the Mittag--Leffler condition.

The theorem follows formally from~(a), (b), and~(c). Observe that, in fact, at
each step the proof uses finite generation of $H^i$ for only one value of $i$.

Let us give some examples in which the hypothesis of \hyperref[IX.1.1]{Theorem~1.1} is satisfied.

\medskip
\sgareaderlabel{IX.1.2}{1.2}\noindent\textbf{Corollary 1.2.}\quad
\emph{Suppose that $\mathcal I'$ is generated by a section $t'$ of
$\mathcal O_{X'}$, and denote by $t$ the corresponding section of
$\mathcal O_X$. Let $F$ be a coherent $\mathcal O_X$-module, and let $n$ be
an integer. Assume that:}
\begin{enumerate}
\item[(i)] \emph{$t$ is $F$-regular (that is, multiplication by $t$ on $F$
is a monomorphism);}
\item[(ii)] \emph{$R^i f_*(F)$ is coherent for $i=n-1$ and $i=n$.}
\end{enumerate}
\emph{Then the hypothesis of \hyperref[IX.1.1]{Theorem~1.1} is satisfied.}

Indeed, multiplication by $t^k$ defines an isomorphism
$F\xrightarrow{\sim}\mathcal J^kF$, from which it follows that
\begin{equation}
\tag{1.29}\label{eq:IX.1.29}
H^i\simeq R^i f_*(F)\otimes_{\mathcal O_{X'}}\mathcal O_{X'}[T],
\end{equation}
where $T$ is an indeterminate. This proves the result.

\sgareaderlabel{IX.1.3}{1.3}\noindent\textbf{Corollary 1.3.}\quad
\emph{Suppose that $X'=\operatorname{Spec}(A)$, where $A$ is a noetherian ring
that is separated and complete for the $I$-adic topology. Suppose that the
graded $S$-module $H^i$ is of finite type for $i=n-1$ and $i=n$
(cf.~\hyperref[eq:IX.1.14]{(1.14)} and~\hyperref[eq:IX.1.15]{(1.15)}). Then the hypotheses of \hyperref[IX.1.1]{Theorem~1.1} are satisfied, and
there is a commutative diagram of isomorphisms}
\begin{equation}
\tag{1.30}\label{eq:IX.1.30}
\begin{tikzcd}[column sep=4.6em,row sep=3.0em]
H^i(X,F)
  \arrow[rr,"\rho_i'"]
  \arrow[dr,"\varphi_i'"']
&& H^i(\widehat X,\widehat F)
  \arrow[dl,"\psi_i"] \\
& \displaystyle\varprojlim_k H^i(X,F_k)
&& \text{for }i=n-1.
\end{tikzcd}
\end{equation}
\noindent{\footnotesize\emph{Source-numbering note
(SGA2-IX-EQ130-RENDERDEF-001).} The corrected French TeX
labels this diagram \texttt{eq:IX.1.30}, the following proof cites (1.30), and
the next numbered formula is (1.31). The same-edition rendered PDF nevertheless
prints (1.1), apparently because its subsection step resets the equation
counter. The French files are left unchanged; the intended tag (1.30) is used
here.}

\section*{Editor's note (1)}
In the same vein, see the article by W.-L. Chow, ``Formal functions on
homogeneous spaces,'' \emph{Invent. Math.} \textbf{86} (1986), no.~1,
pp.~115--130. The author proves the following result. Let $X$ be an algebraic
variety over a field, homogeneous under an algebraic group $G$, and let $Z$ be
a complete subvariety of $X$ of dimension $>0$. Suppose that $Z$ generates $X$
in the following sense: given $p\in Z$, let $\Gamma_p$ be the set of elements
of $G$ that send $p$ into $Z$. One says that $Z$ generates $X$ if the group
generated by the identity component of $\Gamma_p$ is all of $G$. In that case,
every formal rational function on $X$ along $Z$ is algebraic; compare the
results of Hironaka and Matsumura cited in editor's note~(3), p.~138.

Along the lines of the techniques introduced by those authors, we also mention
the elegant algebraization result due to D.~Gieseker, ``On two theorems of
Griffiths about embeddings with ample normal bundle,'' \emph{Amer. J. Math.}
\textbf{99} (1977), no.~6, pp.~1137--1150, Theorems~4.1 and~4.2. Let $X$ be a
connected projective variety of dimension $>0$, locally a complete intersection
over an algebraically closed field. Suppose there are two embeddings of $X$
into smooth projective varieties $Y$ and $W$. If the formal completions of $X$
in $Y$ and $W$ are equivalent, then there exists a scheme $U$ containing $X$
as a closed subscheme which embeds into both $Y$ and $W$ as an étale
neighborhood of $X$. In other words, formal equivalence implies étale
equivalence.

See also G.~Faltings, ``Formale Geometrie und homogene Räume,''
\emph{Invent. Math.} \textbf{64} (1981), pp.~123--165.

\medskip
\noindent{\footnotesize\emph{Source-defect note
(SGA2-IX-L3156-PUNCT-SRCDEF-001;
SGA2-IX-L3157-FALTINGS-TITLE-SRCDEF-001).} The French TeX and its
same-edition PDF omit the sentence stop before ``Autrement dit.'' They also
print the Faltings title in the mixed and misspelled form ``Formale Geometrie
and homogene Raüme.'' The target restores the sentence boundary and the
published German title; the French authority remains unchanged.}

\noindent\emph{Proof.}\quad
We simply note that $H^i(X,F)$ is of finite type and hence isomorphic to its
completion. One obtains~\hyperref[eq:IX.1.30]{(1.30)} by transcribing the diagram of modules~\hyperref[eq:IX.1.7]{(1.7)}
into the category of $A$-modules and replacing the left-hand vertical by
$H^i(X,F)$.

\medskip
\noindent{\footnotesize\emph{Source-numbering note
(SGA2-IX-EQ130-RENDERDEF-001).} The corrected French
TeX cites its label \texttt{eq:IX.1.30}; the next numbered formula is (1.31).
The same-edition rendered PDF prints the referenced diagram as (1.1) because
of an apparent counter reset. The French authority remains unchanged; the
intended reference (1.30) is retained here.}

\sgareaderlabel{IX.1.4}{1.4}\noindent\textbf{Proposition 1.4.}\quad
{\itshape Let $A$ be a noetherian ring. Let $t\in A$, and suppose that $A$ is
separated and complete for the $(tA)$-adic topology. Set
\begin{equation}
\tag{1.31}\label{eq:IX.1.31}
X'=\operatorname{Spec}(A),\qquad Y'=V(t),\qquad I=(tA).
\end{equation}
Let $T$ be a closed subset of $X'$, and set
\begin{equation}
\tag{1.32}\label{eq:IX.1.32}
X=X'\setminus T,\qquad
Y=Y'\cap X=Y'\setminus(Y'\cap T).
\end{equation}
Let $F$ be a coherent $\mathcal O_X$-module. Finally, set
\begin{equation}
\tag{1.33}\label{eq:IX.1.33}
T'=\bigl\{x\in X'\mid
\codimop\bigl(\overline{\{x\}}\cap T,\overline{\{x\}}\bigr)=1\bigr\}.
\end{equation}
Suppose that
\begin{enumerate}
\renewcommand{\labelenumi}{\alph{enumi})}
\item $t$ is $F$-regular;
\item $\prof_{T'}(F)\ge n+1$;
\item $A$ is a quotient of a regular noetherian ring.
\end{enumerate}
Then, in diagram~\hyperref[eq:IX.1.30]{(1.30)}, the morphisms $\rho'_i$, $\varphi'_i$, and $\psi_i$
are isomorphisms for $i<n$ and monomorphisms for $i=n$. Moreover, $\psi_n$
is an isomorphism.}

\medskip
\noindent{\footnotesize\emph{Source-notation and numbering note
(SGA2-IX-L3177-CLOSURE-BRACING-SRCDEF-001;
SGA2-IX-EQ130-RENDERDEF-001).} In (1.33)
the French TeX and its same-edition PDF put braces around $\overline{x}$,
whereas the defining codimension formula in Expos\'e VIII, Theorem~2.1, uses
$\overline{\{x\}}$. The target makes that intended closure-of-the-point
notation explicit; French remains unchanged. French line 3189 cites
\texttt{eq:IX.1.30}; the target retains (1.30), rather than the same-edition
PDF's accidental counter-reset display (1.1).}

\noindent\emph{Proof.}\quad
By~\hyperref[eq:IX.1.3]{(1.3)} and~\hyperref[eq:IX.1.2]{(1.2)}, it suffices to prove that $R^i f_*(F)$ is coherent for
$i\le n$; this follows from the finiteness theorem, Expos\'e~VIII,
\hyperref[IX.2.1]{Theorem~2.1}.

\medskip
\noindent In particular:

\sgareaderlabel{IX.1.5}{1.5}\noindent\textbf{Example 1.5.}\quad
\hyperref[IX.1.4]{Proposition~1.4} applies when $A$ is a local ring and $t$ belongs to the
radical $\mathfrak r(A)$ of $A$. We then take
$T=\{\mathfrak r(A)\}$. In this case, for $n=1$, we obtain the following
statement:

\medskip
{\itshape If $A$ is noetherian, separated and complete for the $t$-adic
topology, and is a quotient of a regular ring (for example, if $A$ is
complete), if moreover $t$ is $F$-regular and $\prof F_x\ge2$ for every
$x\in\operatorname{Spec}(A)$ such that $\dim(A/x)=1$, then the natural
homomorphism
\[
  \Gamma(X,F)\longrightarrow\Gamma(\widehat X,\widehat F)
\]
is an isomorphism.}

\medskip
Indeed, retaining the notation of \hyperref[IX.1.4]{Proposition~1.4}, we have
$T=\{\mathfrak r(A)\}$, and \hyperref[eq:IX.1.33]{formula~(1.33)} says that
\[
  T'=\{x\in\operatorname{Spec}(A)\mid\dim(A/x)=1\}.
\]

\medskip
\noindent{\footnotesize\emph{Source-cross-reference note
(SGA2-IX-EXAMPLE15-SRCDEF-001).} French lines 3196
and 3202 use the key \texttt{eq:IX.1.4}, which formally names equation~(1.4),
although the grammar and the retained hypotheses/notation unambiguously refer
to Proposition~1.4. The same-edition PDF displays only the number `1.4' and
does not resolve the key. The target makes the semantic reference explicit;
the French authority remains unchanged.}

\setcounter{section}{1}
\section{The Existence Theorem}\label{IX.2}

Let us first state EGA~III, 3.4.2 in a slightly more general form.

Let $f:\mathfrak X\to\mathfrak X'$ be an adic
morphism\markedfootnote{(*)}{This hypothesis is not essential; cf.
Expos\'e~XII, p.~118.} of formal preschemes, with $\mathfrak X'$ noetherian.
Let $\mathcal I'$ be an ideal of definition of $\mathfrak X'$; since $f$ is
adic,
$f^*\mathcal I'=\mathcal J$ is\markedfootnote{(2)}{\textit{Editor's note.}
By the very definition; cf. EGA~I, 10.12.1.} an ideal of definition of
$\mathfrak X$.

For every $n\in\mathbb N$, set
\begin{equation}\label{eq:IX.2.1}
  \mathfrak X_n=
  \bigl(\mathfrak X,\mathcal O_{\mathfrak X}/\mathcal J^{n+1}\bigr);
\end{equation}
this is an ordinary prescheme with the same underlying topological space as
$\mathfrak X$.

Let $\mathfrak F$ be a coherent $\mathcal O_{\mathfrak X}$-module. For every
$k\in\mathbb N$, the $\mathcal O_{\mathfrak X_k}$-modules
\begin{equation}
\tag{2.2}\label{eq:IX.2.2}
F_k=\mathfrak F/\mathcal J^{k+1}\mathfrak F
\end{equation}
are coherent. For every $i$, there is a homomorphism
\begin{equation}
\tag{2.3}\label{eq:IX.2.3}
\psi_i:R^i f_*(\mathfrak F)
  \longrightarrow\varprojlim_k R^i f_*(F_k),
\end{equation}
induced by functoriality from the natural homomorphism
\begin{equation}
\tag{2.4}\label{eq:IX.2.4}
\mathfrak F\longrightarrow F_k.
\end{equation}

Set
\begin{equation}
\tag{2.5}\label{eq:IX.2.5}
\mathbf S
  =\operatorname{gr}_{\mathcal I'}\mathcal O_{\mathfrak X'}
  =\bigoplus_{k\in\mathbb N}\mathcal I'^{k}/\mathcal I'^{k+1},
\end{equation}
\begin{equation}
\tag{2.6}\label{eq:IX.2.6}
\operatorname{gr}_{\mathcal J}(\mathfrak F)
  =\bigoplus_{k\in\mathbb N}
    \mathcal J^{k}\mathfrak F/\mathcal J^{k+1}\mathfrak F,
\end{equation}
\begin{equation}
\tag{2.7}\label{eq:IX.2.7}
\mathbf K^{i}
  =R^{i}f_*\bigl(\operatorname{gr}_{\mathcal J}(\mathfrak F)\bigr)
  =\bigoplus_{k\in\mathbb N}
    R^{i}f_*\bigl(\mathcal J^{k}\mathfrak F/
      \mathcal J^{k+1}\mathfrak F\bigr).
\end{equation}
It is clear that $\mathbf K^{i}$ is endowed with the structure of a graded
$\mathbf S$-module.

\medskip
\noindent\textbf{Theorem 2.1.}\sgareaderlabel{IX.2.1}{2.1}\quad
{\itshape Suppose that $\mathbf K^i$ is a graded $\mathbf S$-module of finite
type for $i=n-1$, $i=n$, and $i=n+1$. Then:
\begin{enumerate}
\renewcommand{\labelenumi}{(\roman{enumi})}
\setlength{\itemsep}{0.35em}
\item $R^n f_*(\mathfrak F)$ is coherent.
\item The homomorphism $\psi_n$ in~\hyperref[eq:IX.2.3]{(2.3)} is a topological isomorphism. The
natural filtration on the right-hand side of~\hyperref[eq:IX.2.3]{(2.3)} is $\mathcal I'$-good.
\item The projective system of the $R^n f_*(F_k)$ satisfies the uniform
Mittag-Leffler condition.
\end{enumerate}}

\medskip
\noindent\textit{Proof.}\quad
The proof follows very easily from EGA~$0_{\mathrm{III}}$, 13.7.7
(cf. EGA~III, 3.4.2), provided that the text on page~78 is corrected as
indicated\sgareaderlabel{pagerectif}{11} in (EGA~III~2, $\mathrm{Err}_{\mathrm{III}}\,24$).

\medskip
\sgareaderlabel{IX.2.2}{2.2}\noindent\textbf{Theorem 2.2}\textsuperscript{(3)}.\quad
{\itshape Let $A$ be a noetherian adic ring, and let $I$ be an ideal of
definition of $A$. Let $T$ be a closed subset of
$X'=\operatorname{Spec}(A)$. Suppose that $I$ is generated by an element
$t\in A$. Retain the notation of~\hyperref[eq:IX.1.31]{(1.31)}, \hyperref[eq:IX.1.32]{(1.32)}, and~\hyperref[eq:IX.1.33]{(1.33)}. Let
$\mathfrak F$ be a coherent $\mathcal O_{\widehat X}$-module. Set
\begin{equation}
\tag{2.8}\label{eq:IX.2.8}
F_0=\mathfrak F/\mathcal J\mathfrak F,
\end{equation}
\par\smallskip
\noindent{\normalfont\footnotesize\color{notegray}
\textbf{Source-render note
[SGA2-IX-T22-EQ28-COUNTER-SRCDEF-001].}
The admitted French counter order renders this display as~(2.1), although
its label and the surrounding sequence require~(2.8). The English target
uses~(2.8); the French authority remains unchanged.}
\par\smallskip
where $\mathcal J=t\mathcal O_{\widehat X}$ is an ideal of definition of
$\widehat X$. \emph{Suppose that} $A$ is a quotient of a regular
noetherian ring and that:
\begin{enumerate}
\renewcommand{\labelenumi}{(\arabic{enumi})}
\setlength{\itemsep}{0.35em}
\item $t$ is $\mathfrak F$-regular;
\item $\operatorname{prof}_{T'}F_0\geq 2$.
\end{enumerate}
Then there exists a coherent $\mathcal O_X$-module $F$ such that
$\widehat F\simeq\mathfrak F$.}

\medskip
\noindent\colorbox{softgray}{\parbox{0.95\linewidth}{\small
\color{notegray}\textbf{Editor's note (3).}
Numerous algebraization statements have been obtained since, without
mentioning those cited below; see the articles by Faltings or by Mme Raynaud
cited, respectively, in editor's note~(22) on p.~155 and editor's note~(7)
on p.~203.

\smallskip
{\footnotesize\textbf{Source-defect note
[SGA2-IX-L3265-ONT-GRAMMAR-SRCDEF-001].}
The admitted French reads \emph{on \'et\'e obtenus}; the plural subject
requires \emph{ont \'et\'e obtenus}. The English phrase ``have been
obtained'' gives the intended plural sense; the French authority remains
unchanged.}

\smallskip
One has in mind in particular Artin's results (see especially
Artin M., ``Algebraization of formal moduli. I,'' in \emph{Global Analysis
(Papers in Honor of K.~Kodaira)}, Univ. Tokyo Press, Tokyo, 1969,
pp.~21--71), but also recent results on the algebraicity of leaves of
foliations; see especially Bost J.-B., ``Algebraic leaves of algebraic
foliations over number fields,'' \emph{Publ. Math. Inst. Hautes
\'Etudes Sci.} \textbf{93} (2001), pp.~161--221, and Chambert-Loir A.,
``Th\'eor\`emes d'alg\'ebricit\'e en g\'eom\'etrie diophantienne (d'apr\`es
J.-B. Bost, Y.~Andr\'e, D.~\& G.~Chudnovsky),'' in \emph{S\'eminaire
Bourbaki, Vol. 2000/2001}, \emph{Ast\'erisque}, vol.~282, Soci\'et\'e
math\'ematique de France, Paris, 2002, Exp.~886, pp.~175--209, and the
references cited there. In particular, the latter two articles contain
discussions of the connection between algebraization questions and the
theory of Diophantine approximation.}}

\medskip
\noindent\textit{Proof.}\quad
It suffices to prove that $\widehat f_*(\mathfrak F)$ is a coherent
$\mathcal O_{\widehat{X'}}$-module, where
$\widehat f\colon\widehat X\to\widehat{X'}$ is the morphism of formal
preschemes obtained from the inclusion of $X$ into $X'$ by completion with
respect to $t$. Indeed, $A$ is separated and complete for the $t$-adic
topology, so there exists an $A$-module $F'$ whose completion is isomorphic
to $\widehat f_*(\mathfrak F)$. Since $X$ is open in $X'$, we may take
\[
F=\widetilde{F'}\big|_X.
\]

\noindent\colorbox{softgray}{\parbox{0.95\linewidth}{\footnotesize
\color{notegray}\textbf{Source-defect note
[SGA2-IX-L3280-PUSHFORWARD-BASE-SRCDEF-001].}
The French TeX and the same-edition reader print
$\mathcal O_{\widehat X}$ in the first sentence. Since the pushforward along
$\widehat f\colon\widehat X\to\widehat{X'}$ is a sheaf on the target, the
English target uses the type-correct $\mathcal O_{\widehat{X'}}$. The French
authority remains unchanged under manager decision
EG-SGA2-IX-T22-SOURCE-DEFECT-ADJUDICATION-20260719-0001.}}

\medskip
It remains to show that \hyperref[IX.2.1]{Theorem~2.1} applies to the morphism of formal
preschemes $\widehat f$ and to $\mathfrak F$. By hypothesis~\hyperref[IX.2.2]{(1)}, for every
$k\in\mathbb N$ there is an isomorphism
\[
\mathcal J^k\mathfrak F/\mathcal J^{k+1}\mathfrak F
  \to \mathfrak F/\mathcal J\mathfrak F,
\]
and hence the hypothesis of \hyperref[IX.2.1]{Theorem~2.1} is satisfied provided that
\[
R^i f_*(F_0)\text{ is coherent for }i\leq 1.
\]
This follows from~\hyperref[IX.2.2]{(2)} and the finiteness theorem of Expos\'e~VIII,
\hyperref[IX.2.1]{Theorem~2.1}. This proves the result.

\medskip
\noindent\colorbox{softgray}{\parbox{0.95\linewidth}{\footnotesize
\color{notegray}\textbf{Source-notation observation
[SGA2-IX-L3283-PLAIN-ARROW-OBS-001].}
The corrected TeX and its same-edition reader use a plain $\to$ in the
display while the preceding prose explicitly calls the map an isomorphism.
This is an observation, not a confirmed defect. The target retains both
features; no $\sim$ has been added over the arrow.}}

\medskip
It remains to specialize this statement under the assumption that $A$ is a
local ring.

\medskip
\sgareaderlabel{IX.2.3}{2.3}\noindent\textbf{Corollary 2.3.}\quad
\emph{Let $A$ be a noetherian local ring, and let
$t\in\mathfrak r(A)$ be an element of the radical of $A$. Suppose that $A$
is separated and complete for the $t$-adic topology and, moreover, is a
quotient of a regular ring (for example, suppose that $A$ is a complete
noetherian local ring). Set}
\begin{equation}
\tag{2.9}\label{eq:IX.2.9}
X'=\operatorname{Spec}A,\qquad T=\{\mathfrak r(A)\},
\end{equation}
\emph{and retain the notation of \hyperref[eq:IX.1.31]{(1.31)}, \hyperref[eq:IX.1.32]{(1.32)}, and \hyperref[eq:IX.1.33]{(1.33)}. Let
$\mathfrak F$ be a coherent $\mathcal O_{\widehat X}$-module. Suppose that}
\begin{enumerate}
\item $t$ is $\mathfrak F$-regular;
\item $\prof_{T'}F_0\geq 2$, where
$F_0=\mathfrak F/\mathcal J\mathfrak F$ and
$\mathcal J=t\mathcal O_{\widehat X}$.
\end{enumerate}
\emph{Then there exists a coherent $\mathcal O_X$-module $F$ such that
$\widehat F\simeq\mathfrak F$.}

\medskip
\noindent\colorbox{softgray}{\parbox{0.95\linewidth}{\footnotesize
\color{notegray}\textbf{Source-structure note
(SGA2-IX-L3304-ENUM-STRUCTURE-SRCDEF-001).}
The French TeX closes its \texttt{enumerate} environment only after the
conclusion, syntactically nesting ``Alors il existe'' inside item~(2). The
same-edition PDF preserves the source order but visually masks this nesting:
the conclusion begins at the corollary margin, not at the item-text margin.
The conclusion is consequent on both hypotheses, not part of hypothesis~(2);
the English therefore closes the list before ``Then.'' The French authority
remains byte-identical; this disclosed repair follows the manager's final
adjudication.}}

\medskip
\noindent Note that here $T'$ is the set of prime ideals $\mathfrak p$ of
$A$ such that
\[
  \dim(A/\mathfrak p)=1.
\]

% END INCLUDED active_source/parts/09_expose_IX.tex

% BEGIN INCLUDED active_source/parts/10_expose_X.tex
\makeatletter
\providecommand{\sgareaderlabel}[2]{%
  \begingroup
  \def\@currentlabel{#2}%
  \phantomsection\label{#1}%
  \endgroup}
\makeatother

\chapter{Application to the Fundamental Group}
\label{X}

\medskip
Throughout this expos\'e, $X$ will denote a locally noetherian prescheme,
$Y$ a closed subset of $X$, $U$ a variable open neighborhood of $Y$ in $X$,
and $\widehat X$ the formal completion of $X$ along $Y$
(\textup{EGA}~I, 10.8). For every prescheme $Z$, let $\Et(Z)$ denote the
category of \'etale coverings of $Z$, and let $\LLcat(Z)$\sgareaderlabel{pX.1}{1} denote the category
of coherent locally free modules on $Z$.

\section*{1. Comparison of $\Et(\widehat X)$ and $\Et(Y)$}
\sgareaderlabel{X.1}{1}

Let $\mathcal I$ be an ideal of definition of $Y$ in $X$. For every
$n\in\mathbb N$, set
\[
  Y_n=\left(Y,\left.(\mathcal O_X/\mathcal I^{\,n+1})\right|_Y\right).
\]
The $Y_n$ form a direct system of ordinary preschemes, or also of formal
preschemes when the structure sheaves are endowed with the discrete topology.
We know (EGA~I, 10.6.2) that $\widehat X$ is the direct limit, in the category
of formal preschemes, of the direct system $(Y_n)$. We also know (EGA~I,
10.13) that to give a formal $\widehat X$-prescheme $R$ of \emph{finite type}
is to give a direct system of ordinary $Y_n$-preschemes $R_n$ of
\emph{finite type}, such that
\[
  R_n\simeq R_{n+1}\times_{Y_{n+1}}Y_n.
\]
Moreover, for $R$ to be an étale covering of $\widehat X$, it is necessary
and sufficient that, for every $n$, $R_n$ be an étale covering of $Y_n$.

That said, it is easy to see that nilpotent elements make no difference to
étale coverings (SGA~1, 8.3); in other words, the base-change functor
\[
  \Et(Y_{n+1})\longrightarrow\Et(Y_n)
\]
is an equivalence of categories for every $n\in\mathbb N$. Hence:

\medskip
\sgareaderlabel{X.1.1}{1.1}\noindent\textbf{Proposition 1.1.}\quad
\emph{With the notation introduced above, the natural functor
\[
  \Et(\widehat X)\longrightarrow\Et(Y)
\]
is an equivalence of categories (cf.\ SGA~1, 8.4).}

\section*{2. Comparison of $\Et(Y)$ and $\Et(U)$, for varying $U$}
\sgareaderlabel{X.2}{2}

We shall introduce two conditions from which the comparison theorem announced
above will follow easily. Let $X$ be a locally noetherian prescheme, and let
$Y$ be a closed subset of $X$. We say that the pair $(X,Y)$ satisfies the
\emph{Lefschetz condition}, written $\Lefc(X,Y)$,\sgareaderlabel{pX.2}{2} if, for every open subset
$U$ of $X$ containing $Y$ and every locally free coherent sheaf $E$ on $U$,
the natural homomorphism
\[
  \Gamma(U,E)\longrightarrow\Gamma(\widehat X,\widehat E)
\]
is an isomorphism.

We say that the pair $(X,Y)$ satisfies the \emph{effective Lefschetz
condition}, written $\Leffc(X,Y)$, if $\Lefc(X,Y)$ holds and, in addition,
for every locally free coherent sheaf $\mathcal E$ on $\widehat X$, there
exist an open neighborhood $U$ of $Y$, a locally free coherent sheaf $E$ on
$U$, and an isomorphism
\[
  \widehat E\simeq\mathcal E.
\]

These conditions hold in two important examples:

\medskip
\sgareaderlabel{X.2.1}{2.1}\noindent\textbf{\textit{Example 2.1}}\footnote{\sgareaderlabel{noteX.2.1}{*}\textit{Editor's note.}
One can slightly improve (i): see Mme Raynaud (Raynaud~M., ``Théorèmes de
Lefschetz en cohomologie des faisceaux cohérents et en cohomologie étale.
Application au groupe fondamental,'' \emph{Ann. Sci. Éc. Norm. Sup. (4)}
\textbf{7} (1974), pp.~29--52, Corollaries I.1.4 and I.5). Condition (ii) can
be improved so as to dispense with the depth conditions along $Y$ (see
Theorem~3.3 of \emph{loc. cit.} for a precise statement). The proof of this
latter point is very technical; moreover, the article cited above gives only
indications of proof and refers to an earlier, detailed version published in
the \emph{Bulletin de la Société mathématique de France}.}.\ ---\ 
Let $A$ be a noetherian ring, and let $t\in\mathfrak r(A)$ be an
$A$-regular element belonging to the radical $\mathfrak r(A)$ of $A$.
Assume that $A$ is a quotient of a regular local ring and that $A$ is
complete for the $t$-adic topology (for example, that $A$ is complete for
the $\mathfrak r(A)$-adic topology). Set
\[
  X'=\operatorname{Spec}(A),\qquad Y'=V(t),
\]
and, further, set
\[
  x=\mathfrak r(A),\qquad X=X'\setminus\{x\},\qquad
  Y=Y'\setminus\{x\}.
\]
Thus $X$ is open in $X'$ and $Y=X\cap Y'$. Then:

\begin{enumerate}[label=(\roman*),leftmargin=2.5em]
\item If, for every prime ideal $\mathfrak p$ of $A$ such that
$\dim A/\mathfrak p=1$ (i.e. for every closed point of $X$), we have
$\depthop A_{\mathfrak p}\geq 2$, then $\Lefc(X,Y)$ holds;

\item if, moreover, for every prime ideal $\mathfrak p$ of $A$ such that
$t\in\mathfrak p$ and $\dim A/\mathfrak p=1$ (i.e. for every closed point
of $Y$), we have $\depthop A_{\mathfrak p}\geq 3$, then $\Leffc(X,Y)$
holds.
\end{enumerate}

Let us first prove that, for every open neighborhood $U$ of $Y$ in $X$, the
complement of $U$ in $X$ is a union of finitely many closed points (in $X$).
Observe that $U$ is open in $X$, hence in $X'$, and therefore
$Z'=X'\setminus U$ is closed. Let $I$ be an ideal defining $Z'$; it is
enough to prove that $A/I$ has dimension $1$. Since
\[
  Z'\cap Y'=\{x\},
\]
the ring $A/(I+(t))$ is Artinian, and the conclusion follows from the
``Hauptidealsatz.''

The \emph{first hypothesis is equivalent to}: ``for every prime ideal
$\mathfrak p$ of $A$ with $\mathfrak p\ne\mathfrak r(A)$, one has
\[
  \operatorname{depth} A_{\mathfrak p}
  \geq 3-\dim(A/\mathfrak p).
\]
Indeed, $A$ is a quotient of a regular ring, so we may apply \hyperref[VIII.2.3]{VIII, 2.3} to
the prescheme $X'$, the closed subset $\{x\}$, and the coherent sheaf
$\mathcal O_{X'}$, observing that
$c(\mathfrak p)=\dim(A/\mathfrak p)$ for
$\mathfrak p\in X=X'\setminus\{x\}$ (since $x$ is the closed point of
$X'$).

\medskip
\noindent\colorbox{noteblue}{\parbox{0.95\linewidth}{\footnotesize
\textbf{Source oddity SGA2-X-L3357-U-VS-X-SRCDEF-001.}
The French authority prints $\mathfrak p\in U=X'-\{x\}$, although $U$ was
just quantified as an arbitrary open neighborhood and line 3345 defines
$X=X'-\{x\}$. The English therefore uses $X$ here. The French file remains
unchanged; the same-edition PDF confirms only the printed $U$, not the
repair.}}

Let $U$ be an open neighborhood of $Y$ in $X$, and let $E$ be a locally
free $\mathcal O_U$-module. Set $Z=X\setminus U$, and let
$u\colon U\to X$ be the canonical immersion. We shall first prove that
$u_*(E)$ is a coherent $\mathcal O_X$-module, or, equivalently, that
$\mathcal H^i_Z(E')$ is coherent for $i=0,1$, where $E'$ is a coherent
extension of $E$ to $X$. For this, we apply Theorem \hyperref[VIII.2.1]{VIII, 2.1} to the
prescheme $X$, the closed subset $Z$, and the coherent sheaf $E'$. It is
enough to verify that, for every point $p\in U$ such that $c(p)=1$, one has
$\operatorname{depth}E'_p\geq1$, where we have set
\[
  c(p)=\operatorname{codim}
  \bigl(\overline{\{p\}}\cap Z,\overline{\{p\}}\bigr).
\]
Now suppose $p\in U$ and $c(p)=1$, and let $\mathfrak p$ denote the ideal
of $A$ corresponding to $p$. Then $\dim(A/\mathfrak p)=2$, since the
complement of $U$ is a union of finitely many closed points and $A$ is a
quotient of a regular ring. Moreover, $E$ is locally free, so for every
$p\in\operatorname{Supp}E$ one has
$\operatorname{depth}E_p=\operatorname{depth}\mathcal O_{U,p}$. Finally,
if $p\in U$ and $c(p)=1$, then
\[
  \operatorname{depth}E'_p
  =\operatorname{depth}E_p
  =\operatorname{depth}\mathcal O_{U,p}
  =\operatorname{depth}A_{\mathfrak p}
  \geq3-2=1.
\]

We must now prove that the natural homomorphism
\begin{equation}\label{eq:X.1}
  \Gamma(U,E)\longrightarrow\Gamma(\widehat X,\widehat E)
\end{equation}
is an isomorphism. Set $\overline{\!E}=u_*(E)$. The sheaf
$\overline{\!E}$ is coherent and has depth at least $2$ at every closed
point of $X$. It follows that $R^if_*(\overline{\!E})$ is coherent for
$i=0,1$, where $f\colon X\to X'$ denotes the canonical immersion of $X$
into $X'=\operatorname{Spec}(A)$ (Exposé VIII). Applying \hyperref[IX.1.5]{IX, 1.5}, we
conclude that
\begin{equation}\label{eq:X.2}
  \Gamma(U,\overline{\!E})
  \longrightarrow
  \Gamma(\widehat X,\widehat{\overline{\!E}})
\end{equation}
is an isomorphism, since $A$ is complete for the $t$-adic topology.

There is a commutative diagram
\[
\xymatrix@C=18mm@R=9mm{
  \Gamma(X,\overline{\!E})
    \ar[rr]^{\simeq}
    \ar[dr]_{\simeq}
  && \Gamma(U,E)\ar[dl] \\
  & \Gamma(\widehat X,\widehat E) &
}
\]
and the conclusion follows.

\medskip
Now let $\mathcal E$ be a locally free coherent sheaf on $\widehat X$. If
one has proved that $\mathcal E$ is algebraizable, that is, isomorphic to the
formal completion of a coherent $\mathcal O_X$-module $E$, it is easy to see
that $E$ is locally free in a neighborhood of $Y$, and hence to prove
$\Leffc(X,Y)$. Let $\widehat{X'}$ be the formal spectrum of $A$ for the
$t$-adic topology; it is identified with the formal completion of $X'$ along
$Y'$. Denote by $f$ the canonical immersion of $X$ into $X'$, by $f'$ the
canonical immersion of $Y$ into $Y'$, and by $\widehat f$ the morphism
obtained by passing to the completions. For $\mathcal E$ to be algebraizable,
it suffices that $\widehat f_*(\mathcal E)$ be a coherent
$\mathcal O_{\widehat{X'}}$-module, since $A$ is complete for the $t$-adic
topology.
\sourceoddity{SGA2-X-L3384-PUSHFORWARD-BASE-SRCDEF-001}{The French authority
prints $\mathcal O_{\widehat X}$ as the coefficient ring of the pushforward.
Because $\widehat f\colon\widehat X\to\widehat{X'}$, the target uses
$\mathcal O_{\widehat{X'}}$. Only this coefficient ring is corrected.}
Set $\mathcal I=t\mathcal O_{\widehat X}$; this is an ideal of definition of
$\widehat X$.

For every $n\geq0$, set
\[
  \mathcal E_n=\mathcal E/\mathcal I^{n+1}\mathcal E.
\]
At every closed point $y\in Y$, the depth of $\mathcal E_0$ is at least two;
indeed, $t$ is an $A$-regular element, and therefore
\[
  \depthop\mathcal O_{Y,y}
  =\depthop\mathcal O_{X,y}-1\geq2.
\]
\sourceoddity{SGA2-X-L3386-Y0-VS-Y-SRCDEF-001}{The French authority prints
$\mathcal O_{Y_0,y}$, but $Y_0$ is undefined in this argument. Since
$Y=V(t)\cap X$, the regular-element quotient gives $\mathcal O_{Y,y}$.}
It follows that $\widehat f_*(\mathcal E)$ is coherent (\hyperref[IX.2.3]{IX~2.3)}.
\hfill\textsc{q.e.d.}

\medskip
\sgareaderlabel{X.2.2}{2.2}\noindent\textbf{\textit{Example 2.2}}\quad
\textup{(This will make it possible to compare the fundamental group of a
projective variety with that of a hyperplane section).}

Let $K$ be a field, and let $X$ be a proper $K$-prescheme. Let $\mathcal L$
be an ample invertible $\mathcal O_X$-module. Let
$t\in\Gamma(X,\mathcal L)$ be an $\mathcal O_X$-regular element; this means
that, for every open subset $U$ and every isomorphism
\[
  u\colon \mathcal L|_U \xrightarrow{\ \sim\ } \mathcal O_U,
\]
the section $u(t)$ is a non-zero-divisor in $\mathcal O_U$ (a condition that
does not depend on $u$). Let $Y=V(t)$ be the subscheme of $X$ defined by the
equation $t=0$\footnote{\textit{Editor's note.} Condition (ii) is
superfluous; see the editor's note to \hyperref[X.2.1]{Example~2.1} (original running p.~90).}.
Then:

\begin{enumerate}[label=(\roman*),leftmargin=2.5em]
\item If, for every closed point $x$ of $X$, one has
$\depthop\mathcal O_{X,x}\geq2$, then $\Lefc(X,Y)$ holds;

\item if, moreover, for every closed point $y\in Y$, one has
$\depthop\mathcal O_{X,y}\geq3$, then $\Leffc(X,Y)$ holds.
\end{enumerate}

\medskip
This example will be treated in detail in Exposé~XII.

\medskip
Let $S$ be a prescheme. It is known (EGA~II, 6.1.2) that the functor which
assigns to every finite flat covering $r\colon R\to S$ the
$\mathcal O_S$-algebra $r_*(\mathcal O_R)$ induces an equivalence between
the category of finite flat coverings of $S$ and the category of coherent
locally free $\mathcal O_S$-algebras.
\sourceoddity{SGA2-X-L3402-OS-VS-OX-SRCDEF-001}{The French authority prints
$\mathcal O_X$ in both coefficient-ring positions. Since $S$ is arbitrary
and $r\colon R\to S$, the direct image and the equivalent algebra category
are over $\mathcal O_S$.}
Let $U$ be an open neighborhood of $Y$, and let $r\colon R\to U$ be a
finite flat covering of $U$. Let $\widehat R$ be the finite flat covering of
$\widehat X$ obtained from it by base change. Then
\[
  \widehat r_*(\mathcal O_{\widehat R})
  \simeq \widehat{r_*(\mathcal O_R)}.
\]

Assume now that $\Lefc(X,Y)$ holds. This implies that, for every $U$, the
inverse-image functor
\[
  \mathsf L(U)\longrightarrow\mathsf L(\widehat X)
\]
is fully faithful. Indeed, let $E$ and $F$ be two coherent locally free
$\mathcal O_U$-modules; $\cHom(E,F)$ is likewise coherent and locally free.
\sourceoddity{SGA2-X-L3408-DUPLICATE-HOM-SRCDEF-001}{The French authority
prints the preceding sentence twice consecutively. The target states the
identical assertion once.}
By hypothesis, the natural map
\[
  \Gamma\bigl(U,\cHom(E,F)\bigr)
  \longrightarrow
  \Gamma\bigl(\widehat X,\widehat{\cHom(E,F)}\bigr)
\]
is an isomorphism, which proves the assertion, since $\cHom$ commutes with
formal completion when all the modules involved are locally free. Formal
completion also commutes with tensor products. It follows that the functor
which assigns to every coherent locally free $\mathcal O_U$-algebra
$\mathcal A$ the $\mathcal O_{\widehat X}$-algebra $\widehat{\mathcal A}$
is fully faithful. More precisely, if $E$ is a coherent locally free
$\mathcal O_U$-module, commutative $\mathcal O_U$-algebra structures on $E$
correspond bijectively to commutative $\mathcal O_{\widehat X}$-algebra
structures on $\widehat E$.
\sourceoddity{SGA2-X-L3412-BIJECTION-MISSING-FIRST-SIDE-SRCDEF-001}{The
French authority introduces a bijective correspondence but prints only its
$\mathcal O_{\widehat X}$ side. The target restores the omitted
$\mathcal O_U$-algebra structures on $E$, as required by the preceding
fully faithful tensor-compatible completion functor.}

\medskip
\sgareaderlabel{X.2.3}{2.3}\noindent\textbf{\textit{Proposition 2.3.}}
Let $X$ be a locally noetherian prescheme, and let $Y$ be a closed subset of
$X$. Let $\widehat X$ be the formal completion of $X$ along $Y$. For every
open subset $U$ of $X$ containing $Y$, let $L_U$ (respectively $P_U$,
respectively $E_U$) denote the functor that assigns to every locally free
coherent $\mathcal O_U$-module (respectively every finite flat covering of
$U$, respectively every \'{e}tale covering of $U$) its pullback along
$\widehat X\to X$.

\begin{enumerate}[label=(\roman*),leftmargin=2.5em]
\item If $\Lefc(X,Y)$ holds, then, for every open neighborhood $U$ of $Y$,
the functors $L_U$, $P_U$, and $E_U$ are fully faithful.

\item If $\Leffc(X,Y)$ holds, then, for every locally free coherent
$\mathcal O_{\widehat X}$-module $\mathcal E$, there are an open
neighborhood $U$ of $Y$ and a locally free coherent $\mathcal O_U$-module
$E$ such that
\[
  L_U(E)\simeq\mathcal E.
\]
The analogous assertions hold for finite flat coverings and for \'{e}tale
coverings, with $P_U$ and $E_U$, respectively.
\end{enumerate}

\medskip
\noindent\textit{Proof.}
Part~(i) has already been proved.

Part~(ii) follows from (i) and the hypothesis, at least for $L_U$ and $P_U$.
Moreover, if $R$ is an \'{e}tale covering of $\widehat X$, there are an open
neighborhood $U$ of $Y$ in $X$ and a finite flat covering $R'$ of $U$ such
that
\[
  \widehat{R'}\simeq R.
\]
By restriction this gives a covering $R''$ of $Y$, which is \'{e}tale by
\hyperref[X.1.1]{Proposition~1.1}. Hence the restriction of $R'$ to some neighborhood $U'$ of
$Y$ is \'{e}tale.
\hfill$\square$

\medskip
\sgareaderlabel{X.2.4}{2.4}\noindent\textbf{\textit{Corollary 2.4.}}
Suppose that $\Lefc(X,Y)$ holds. A finite flat covering $R$ of an open
neighborhood $U$ of $Y$ is connected if and only if
$R\times_U\widehat X$ is connected. In particular, $Y$ is connected if and
only if $U$ is connected, and this is also equivalent to $X$ being
connected.

\medskip
\noindent\textit{Proof.}
Indeed, a locally ringed space $(X,\mathcal O_X)$ is connected if and only if
$\Gamma(X,\mathcal O_X)$ is not a direct product of two nonzero rings. But,
by $\Lefc(X,Y)$, one has
\[
  \Gamma\bigl(U,r_*\mathcal O_R\bigr)
  \simeq
  \Gamma\bigl(\widehat X,\widehat r_*\mathcal O_{\widehat R}\bigr).
\]

\medskip
\sgareaderlabel{X.2.5}{2.5}\noindent\textbf{\textit{Corollary 2.5.}}
Suppose that $\Lefc(X,Y)$ holds. Then, for every open neighborhood $U$ of
$Y$, the functor
\[
  \Etc(U)\longrightarrow\Etc(Y)
\]
is fully faithful. Suppose that $\Leffc(X,Y)$ holds. Then, for every
\'{e}tale covering $R$ of $Y$, there are an open neighborhood $U$ of $Y$
and a covering $R'$ of $U$ such that
$R'\times_U Y\simeq R$.

\medskip
\sgareaderlabel{X.2.6}{2.6}\noindent\textbf{\textit{Corollary 2.6\textsuperscript{(3)}.}}
Suppose that $\Lefc(X,Y)$ holds and that $Y$ is connected. Then every open
neighborhood $U$ of $Y$ is connected, and the natural homomorphism
$\pi_1(Y)\longrightarrow\pi_1(U)$ is surjective. Suppose, moreover, that
$\Leffc(X,Y)$ holds. Then the natural homomorphism
\[
  \pi_1(Y)\longrightarrow\varprojlim_{U}\,\pi_1(U)
\]
is an isomorphism. (N.B. A ``base point'' in $Y$ is understood to have been
chosen and is also taken as a base point in $X$ for the definition of the
fundamental groups.)

\medskip
\noindent All of this follows immediately from \hyperref[X.1.1]{Propositions~1.1} and~\hyperref[X.2.3]{2.3}.

\section*{Editor's note (3)}
Combined with \hyperref[X.3.3]{Proposition~3.3} and criteria \hyperref[XII.2.4]{XII.2.4} and \hyperref[XII.3.4]{XII.3.4}, this yields the
following relative Lefschetz theorem. Let $f\colon X\to S$ be a projective and
flat morphism of connected noetherian schemes, and let $D$ be an effective
relative Cartier divisor on $X$ that is relatively ample. If, for every
$s\in S$, the depth of $X_s$ at every closed point is $\geq 2$, then $D$ is
connected and, for every open subset $U$ of $X$ containing $D$, the map
\[
  i_U\colon \pi_1(D)\longrightarrow\pi_1(U)
\]
is surjective. If, in addition, the depth of $X_s$ at every closed point of
$D_s$ is $\geq 3$ and the local rings of $X$ at its closed points are
pure---for example, are complete intersections (cf. \hyperref[X.3.4]{X, 3.4)}---then $i_X$ is an
isomorphism.

Cf. J.-B. Bost, ``Lefschetz theorem for Arithmetic Surfaces,''
\emph{Ann. Sci. \'Ec. Norm. Sup. (4)} \textbf{32} (1999), pp.~241--312,
Theorems~1.1 and~2.1.

When $X$ is merely a smooth projective surface that is geometrically connected
over a field, $D$ is always connected and
\[
  \pi_1(D)\longrightarrow\pi_1(U)
\]
is surjective (where $U$ is an open subset containing $D$), provided only that
$D$ is nef and has self-intersection $>0$; cf. \emph{loc. cit.}, Theorem~2.3,
and also Theorem~2.4 for surfaces assumed only normal and complete.

In the case of an arithmetic surface (normal and quasi-projective) $X$ over a
ring of integers $\mathcal O_K$, Bost, improving results of Y.~Ihara
(Y.~Ihara, ``Horizontal divisors on arithmetic surfaces associated with
Bely\u{\i} uniformizations,'' in \emph{The Grothendieck theory of dessins
d'enfants (Luminy, 1993)}, London Math. Soc. Lect. Note Series, vol.~200,
Cambridge Univ. Press, Cambridge, 1994, pp.~245--254, or \emph{loc. cit.},
Corollary~7.2), proved that if a point
$P\in X(\mathcal O_K)$, which plays the role of the divisor $D$ in the
geometric situation, satisfies certain positivity conditions, then the map
\[
  \pi_1(X)\longrightarrow\pi_1(\operatorname{Spec}\mathcal O_K)
\]
induced by the projection is an isomorphism whose inverse is the map
\[
  \pi_1(\operatorname{Spec}\mathcal O_K)\longrightarrow\pi_1(X)
\]
induced by $P$; cf. \emph{loc. cit.}, Theorem~1.2.

\medskip
\noindent{\footnotesize\emph{Accepted source-normalization note
(SGA2-X-L3456-AU-DESSUS-NORMALIZATION-001@1).} The corrected French control
selects \textup{au-dessus} from
\texttt{\string\sisi\{au dessus\}\{au-dessus\}}; the target renders this as
``over.'' This is an explicit source normalization, not a source defect.}

\section*{3. Comparison of $\pi_1(X)$ and $\pi_1(U)$}
\sgareaderlabel{X.3}{3}

\medskip
\sgareaderlabel{X.3.1}{3.1}\noindent\textbf{Definition 3.1.}\quad
Let $X$ be a prescheme and $Z$ a closed subset of $X$. Set
$U=X\setminus Z$. We say that the pair $(X,Z)$ is \emph{pure} if, for every
open subset $V$ of $X$, the functor
\[
\begin{aligned}
  \Et(V)&\longrightarrow\Et(V\cap U),\\
  V'&\longmapsto V'\times_V(V\cap U)
\end{aligned}
\]
is an equivalence of categories\renewcommand{\thefootnote}{*}\footnote{For a
notion that is more satisfactory in certain respects, cf. the commentary
\hyperref[XIV.1.6]{XIV, 1.6(d)}.}.

\medskip
\noindent{\footnotesize\emph{Accepted source-normalization note
(SGA2-X-L3464-MAPSTO-NORMALIZATION-001@1).} The corrected French control
selects \texttt{\string\mto} from
\texttt{\string\sisi\{\string\rightsquigarrow\}\{\string\mto\}}; the target
renders the object map with $\longmapsto$. This is an explicit source
normalization, not a source defect.}

\medskip
\sgareaderlabel{X.3.2}{3.2}\noindent\textbf{Definition 3.2.}\quad
Let $A$ be a noetherian local ring. Set $X=\Spec A$. Let $\rr(A)$ be the
radical of $A$, and let $x=\rr(A)$ be the closed point of $X$. We say that
$A$ is \emph{pure} if the pair $(X,\{x\})$ is pure.

\medskip
\noindent We leave to the reader the task of not proving the following
proposition:

\medskip
\sgareaderlabel{X.3.3}{3.3}\noindent\textbf{Proposition 3.3.}\quad
Let $X$ be a locally noetherian prescheme and let $Z$ be a closed subset of
$X$. The pair $(X,Z)$ is pure if and only if, for every $z\in Z$, the local
ring $\Oo_{X,z}$ is pure\renewcommand{\thefootnote}{**}\footnote{Compare the
noncommutative case of \hyperref[XIV.1.8]{XIV, 1.8}, whose proof is essentially the same as that
of \hyperref[X.3.3]{Proposition~3.3}.}.

\medskip
\noindent{\footnotesize\color{notebluetext}\emph{Source-defect note
(SGA2-X-L3477-FERME-VS-FERMEE-SRCDEF-001).} French line~3477 reads
\emph{une partie ferm\'e de $X$}; the masculine adjective \emph{ferm\'e}
disagrees with the feminine noun \emph{partie}. The target renders the
intended sense as ``a closed subset of $X$.'' The admitted French control is
unchanged.}

\medskip
\noindent That said, the following theorem is the main result of this
section:

\medskip
\sgareaderlabel{X.3.4}{3.4}\noindent\textbf{Theorem 3.4 (Purity theorem)}%
\renewcommand{\thefootnote}{(4)}\footnote{\textsc{N.D.E.}: For the history of
the methods employed, see Grothendieck's letter of October~1, 1961 to Serre,
\emph{Correspondance Grothendieck-Serre}, edited by Pierre Colmez and
Jean-Pierre Serre, Documents Math\'ematiques, vol.~2, Soci\'et\'e
Math\'ematique de France, Paris, 2001.}\textbf{.}

\begin{enumerate}
  \renewcommand{\labelenumi}{(\roman{enumi})}
  \item A regular noetherian local ring of dimension $\geq 2$ is pure
  (Zariski--Nagata purity theorem).
  \item A noetherian local ring of dimension $\geq 3$ that is a complete
  intersection is pure.
\end{enumerate}

\medskip
\noindent Recall that a local ring $A$ is said to be a \emph{complete
intersection} if and only if there exist a \emph{regular} noetherian local
ring $B$ and a $B$-regular sequence $(t_1,\ldots,t_k)$ of elements of the
radical $\mathfrak r(B)$ of $B$ such that
\[
  A \simeq B/(t_1,\ldots,t_k).
\]

\medskip
\noindent In this connection, let us note that it would be less ambiguous to
say that $A$ is an \emph{absolute} complete intersection, in contrast to the
situation already encountered, in which $X$ is a locally noetherian prescheme
(not necessarily regular) and $Y$ is a closed subset of $X$ that is said to
be ``locally a set-theoretic complete intersection \emph{in $X$}.''

\medskip
\noindent Let us first prove a few lemmas.

\medskip
{\itshape
\sgareaderlabel{X.3.5}{3.5}\noindent\textbf{Lemma 3.5.} --- Let $X$ be a locally noetherian prescheme,
and let $U$ be an open subset of $X$. Set $Z=X\setminus U$. Let
$i\colon U\to X$ be the canonical immersion of $U$ into $X$. The following
conditions are equivalent:
\begin{enumerate}[label=\textup{(\roman*)},leftmargin=2.2em,itemsep=0.45em,topsep=0.45em]
\item For every open subset $V$ of $X$, if $V'=V\cap U$, the restriction
functor $F\longmapsto F|_{V'}$ from the category of locally free coherent
$\mathcal O_V$-modules to the category of locally free coherent
$\mathcal O_{V'}$-modules is fully faithful;
\item the natural homomorphism $\mathcal O_X\to i_*\mathcal O_U$ is an
isomorphism;
\item for every $z\in Z$, one has $\prof\mathcal O_{X,z}\geq 2$.
\end{enumerate}
}

We have already seen (\Ref{III}~\Ref{III.3.3}) the equivalence of (ii)
and (iii). Let us show that (ii) implies (i). Let $F$ and $G$ be two
locally free coherent $\Oo_V$-Modules. Then $\SheafHom(F,G)$ is also
locally free and coherent, so
$$
\SheafHom(F,G)\longrightarrow
i_\ast\bigl(\SheafHom(\sisi{F|{|V'}}{F_{|V'}},
\sisi{G|{|V'}}{G_{|V'}})\bigr)
$$
is an isomorphism, and therefore
$$
\Hom(F,G)\simeq
\Hom(\sisi{F|{|V'}}{F_{|V'}},\sisi{G|{|V'}}{G_{|V'}}).
$$
Conversely, take $F=G=\OX$ and apply (i) to every open subset $V$ of
$X$.

\medskip
\noindent Here is a useful ``descent lemma'':

\medskip
{\itshape
\sgareaderlabel{X.3.6}{3.6}\noindent\textbf{Lemma 3.6.} --- Let $X$ be a locally noetherian prescheme,
and let $Z$ be a closed subset of $X$. Set $U=X\setminus Z$. Suppose that
the homomorphism $\mathcal O_X\to i_*\mathcal O_U$ is an isomorphism. Let
$f\colon X_1\to X$ be a faithfully flat and quasi-compact morphism. Set
$Z_1=f^{-1}(Z)$. If the pair $(X_1,Z_1)$ is pure, then so is $(X,Z)$.
}

\medskip
\noindent\textit{Proof of \hyperref[X.3.6]{Lemma 3.6}.}\quad
Observe that the hypothesis
$\Oo_X\simeq i_*\Oo_U$ is preserved under flat base extension, since $i$ is
quasi-compact and, in this case, formation of the direct image commutes with
inverse image. This hypothesis implies that the functor
\[
  \Et(V)\longrightarrow\Et(U\cap V)
\]
defined by
\[
  V'\longmapsto V'\times_V(V\cap U)
\]
is fully faithful, as follows from the interpretation of an étale covering
in terms of a coherent locally free algebra. It remains to prove effectivity.
One may, for example, form the fiber square $X_2$ and the fiber cube $X_3$ of
$X_1$ over $X$, and observe that a faithfully flat and quasi-compact
morphism is a morphism of universal effective descent for the fibered
category of étale coverings over the category of preschemes. The conclusion
then follows formally\renewcommand{\thefootnote}{(*)}\footnote{Cf. J.~Giraud,
\emph{Méthode de la descente}, Mémoire no.~2 du \emph{Bulletin de la Société
Mathématique de France} (1964).}.

\medskip
{\itshape
\sgareaderlabel{X.3.7}{3.7}\noindent\textbf{Remark 3.7.} --- We proved along the way that if
$\mathcal O_X\to i_*\mathcal O_U$ is an isomorphism, then $X$ is connected if
and only if $U$ is connected, and then $\pi_1(U)\to\pi_1(X)$ is surjective.
}

\medskip
{\itshape
\sgareaderlabel{X.3.8}{3.8}\noindent\textbf{Corollary 3.8.} --- Let $A$ be a noetherian local ring.
Assume that $\depthop A\geq 2$. If $\widehat{A}$ is pure, then $A$ is pure.
}

\medskip
\noindent\emph{Proof.}\quad This follows from \hyperref[X.3.5]{Lemma~3.5} and \hyperref[X.3.6]{Lemma~3.6}.

\medskip
\noindent The following lemma is the essential point in the proof of the purity theorem:

\medskip
{\itshape
\sgareaderlabel{X.3.9}{3.9}\noindent\textbf{Lemma 3.9.} --- Let $A$ be a noetherian local ring, and
let $t\in\mathfrak r(A)$ be an $A$-regular element. Suppose that $A$ is
complete for the $t$-adic topology and is moreover a quotient of a regular
local ring (for example, if $A$ is complete). Set $B=A/tA$.
\begin{enumerate}[label=\textup{(\roman*)},leftmargin=2.2em,itemsep=0.6em,topsep=0.5em]
\item If, for every prime ideal $\mathfrak p$ of $A$ such that
$\dim A/\mathfrak p=1$, one has $\depthop A_{\mathfrak p}\geq 2$, then
the purity of $B$ implies the purity of $A$.

\item If, for every prime ideal $\mathfrak p$ of $A$ such that
$\dim A/\mathfrak p=1$, one has $\depthop A_{\mathfrak p}\geq 2$, if
$A_{\mathfrak p}$ is pure whenever $t\notin\mathfrak p$, and
if\renewcommand{\thefootnote}{(5)}\footnote{\textsc{N.D.E.}: This last
condition can be improved; cf. the editor's note~(1) on p.~90.}
$\depthop A_{\mathfrak p}\geq 3$ whenever $t\in\mathfrak p$, then the
purity of $A$ implies the purity of $B$.
\sourceoddity{SGA2-X-L3551-MISSING-EST-SRCDEF-001}{The French authority
prints \emph{si $A_{\mathfrak p}$ pur lorsque $t\notin\mathfrak p$},
omitting the finite copula \emph{est}. The English supplies ``is'' in the
unambiguous intended condition. The French authority remains unchanged.}
\end{enumerate}
}

\medskip
Let $X'=\operatorname{Spec}(A)$ and let $Y'=V(t)$, which we identify with
the spectrum of $B$. Let $x=\mathfrak r(A)$, and set
$X=X'\setminus\{x\}$ and $Y=Y'\setminus\{x\}=X\cap Y'$. Denote by
$\widehat{X'}$ the formal spectrum of $A$ for the $t$-adic topology; it is
identified with the formal completion of $X'$ along $Y'$.

Since $A$ is complete for the $t$-adic topology, observe that
$\Et(X')\longrightarrow\Et(\widehat{X'})$ is an equivalence of categories.
Likewise, $\Et(\widehat{X'})\longrightarrow\Et(Y')$ is an equivalence by
\hyperref[X.1.1]{Proposition~1.1}; hence $\Et(X')\longrightarrow\Et(Y')$ is an equivalence of
categories.

We prove \textup{(i)}. Consider the diagrams
\[
\xymatrix@C=12mm@R=8mm{
  X' & X\ar[l] &&
  \Et(X')\ar[r]^{a}\ar[d]_{c} & \Et(X)\ar[d]^{b}\\
  Y'\ar[u] & Y\ar[l]\ar[u] &&
  \Et(Y')\ar[r]^{d} & \Et(Y)
}
\]
We have just seen that $c$ is an equivalence; $d$ is also an equivalence by
the hypothesis that $B$ is pure; and finally $b$ is fully faithful, as was
seen in \hyperref[X.2.1]{Example~2.1}; cf. \hyperref[X.2.3]{Proposition~2.3}\textup{(i)}.

We prove \textup{(ii)}. This time suppose that $A$ is pure, so $a$ is an
equivalence; likewise $c$. Let us show that $b$ is an equivalence. By
\hyperref[X.2.1]{Example~2.1}, we know that $\Leffc(X,Y)$ holds, so $b$ is already fully
faithful; let us prove that it is essentially surjective. Apply
\hyperref[X.2.3]{Proposition~2.3}\textup{(ii)}, noting that if $U$ is an open neighborhood of
$Y$ in $X$, then $X\setminus U$ is a union of finitely many closed points.
The pair $(X,X\setminus U)$ is therefore pure by \hyperref[X.3.3]{Proposition~3.3}, since at
any such point $\mathfrak p$ one has
$\mathcal O_{X,\mathfrak p}=A_{\mathfrak p}$, which is pure by hypothesis.
This proves the assertion.

\medskip
\textit{Proof of the purity theorem.}
We first prove \textup{(i)} by induction on the dimension. Let $A$ be a
noetherian local ring of dimension~$2$. Set
$X'=\Spec(A)$, $x=\mathfrak r(A)$, and $X=X'\setminus\{x\}$. We have
$\depthop A=2$. We may therefore apply \hyperref[X.3.5]{Lemma~3.5} to the pair
$(X',\{x\})$, and hence
$\Et(X')\longrightarrow\Et(X)$ is fully faithful.

Now let $r\colon R\longrightarrow X$ be an \'{e}tale covering defined by
the coherent, locally free, \'{e}tale $\mathcal O_X$-algebra
$\mathcal A=r_*(\mathcal O_R)$. Let $i\colon X\longrightarrow X'$ denote
the canonical immersion. I claim that
$i_*(\mathcal A)=\mathcal B$ is a coherent $\mathcal O_{X'}$-algebra.
Indeed, it suffices to apply the ``finiteness theorem,'' \hyperref[VIII.2.3]{VIII, 2.3}. I also
claim that this algebra has depth at least~$2$ at $x$: it is the direct image
of an $\mathcal O_X$-module, with $X=X'\setminus\{x\}$. Since $A$ is a
regular ring of dimension~$2$, we have
\[
  \pdop \mathcal B+\depthop \mathcal B=\dim A=2,
\]
where $\pdop \mathcal B$ denotes the projective dimension of $\mathcal B$.
Thus $\pdop \mathcal B=0$, so $\mathcal B$ is projective and therefore
free. Consequently, $\mathcal B$ defines a finite flat covering of
$X'=\Spec(A)$. The set of points of $X'$ at which this covering is not
\'{e}tale is a closed subset of $X'$ whose defining ideal is principal,
namely the discriminant ideal of $\mathcal B/A$. By construction this closed
subset is contained in $x=\mathfrak r(A)$; it is therefore empty because
$\dim A=2$.

Let $A$ now be a regular noetherian local ring with $\dim A=n\geq 3$.
Suppose that \textup{(i)} has been proved for rings of dimension less than
$n$. To prove that $A$ is pure, we may assume by \hyperref[X.3.8]{Corollary~3.8} that $A$ is
complete. Choose $t\in\mathfrak r(A)$ whose image in
$\mathfrak r(A)/\mathfrak r(A)^2$ is nonzero. Then $B=A/tA$ is a regular
noetherian local ring of dimension $n-1$, and is therefore pure because
$n-1\geq 2$. The conclusion follows from \hyperref[X.3.9]{Lemma~3.9}\textup{(i)}, which is
applicable because $A$ is complete.

\medskip
\textit{Let us prove \textup{(ii)}.}
Let $A$ be a noetherian local ring of dimension at least~$3$. Suppose that
there exist a regular noetherian local ring $B$ and a $B$-sequence
$(t_1,\ldots,t_k)$ such that
\[
  A\simeq B/(t_1,\ldots,t_k).
\]
We prove that $A$ is pure by induction on $k$. If $k=0$, this follows from
\textup{(i)}. Suppose that $k\geq1$ and that the result has been proved for
$k'<k$. By \hyperref[X.3.8]{Corollary~3.8}, we may assume that $A$, and hence also $B$, is
complete. Set
\[
  C=B/(t_1,\ldots,t_{k-1}).
\]
Then $A\simeq C/t_kC$, and $t_k$ is $C$-regular. By the induction hypothesis,
$C$ is pure; it therefore suffices to verify that \hyperref[X.3.9]{Lemma~3.9}\textup{(ii)} is
applicable. In the notation of that lemma, its $A$ and $B$ are here $C$ and
$A$, respectively. We have $\dim C\geq4$; hence, for every prime ideal
$\mathfrak p$ of $C$ such that $\dim(C/\mathfrak p)=1$, one has
$\depthop C_{\mathfrak p}\geq3$. Moreover, $C_{\mathfrak p}$ is a complete
intersection with $k'\leq k-1$, and is therefore pure by the induction
hypothesis.

\medskip
{\itshape
\sgareaderlabel{X.3.10}{3.10}\noindent\textbf{Theorem 3.10.} --- Let $X$ be a locally noetherian
prescheme, and let $Y$ be a closed subset of $X$. Suppose that
$\Leff(X,Y)$ holds (cf. \hyperref[X.2.1]{Examples~2.1} and~\hyperref[X.2.2]{2.2}). Suppose further that, for
every open neighborhood $U$ of $Y$ and every $x\in X\setminus U$, the
local ring $\mathcal O_{X,x}$ is either regular of dimension $\geq 2$ or
a complete intersection of dimension $\geq 3$. Then
\[
  \pi_0(Y)\longrightarrow\pi_0(X)
\]
is a bijection, and, if $X$ is connected,
\[
  \pi_1(Y)\longrightarrow\pi_1(X)
\]
is an isomorphism.
}

\medskip

There is nothing more to prove. Observe that in the two cited
\hyperref[X.2.1]{Examples~2.1} and~\hyperref[X.2.2]{2.2}, the complement of $U$ is a union of finitely many
closed points; it follows that the hypothesis on the dimension of
$\mathcal O_{X,x}$ is not farcical.

% END INCLUDED active_source/parts/10_expose_X.tex

% BEGIN INCLUDED active_source/parts/11_expose_XI.tex
\chapter{Application to the Picard Group}\label{XI}
\medskip
This expos\'e follows the pattern of the preceding one, but this time the result of no.~1 is weaker.

\medskip
Throughout this expos\'e, $X$ will denote a locally noetherian prescheme;
$\Ii$, a quasi-coherent ideal of $\OX$ (so $Y=V(\Ii)$ is a closed subset of
$X$); $U$, a variable open neighborhood of $Y$ in $X$; and $\widehat X$,
the formal completion of $X$ along $Y$.

For every ringed space $(Z,\mathcal O_Z)$, let $\PP(Z)$\refstepcounter{toto}\label{pXI.1}
denote the category of invertible $\mathcal O_Z$-modules, that is, modules locally
free of rank~$1$, and let $\Pic(Z)$ denote the group of isomorphism classes of
invertible modules on $Z$.

\section{\texorpdfstring{Comparison of $\Pic(\widehat X)$ and $\Pic(Y)$}{Comparison of Pic(X-hat) and Pic(Y)}}\label{XI.1}

For every $n\in\NN$, set
\[
X_n=(Y,\OX/\Ii^{n+1})
\qquad\text{and}\qquad
P_n=\Ii^{n+1}/\Ii^{n+2}.
\]
The sequence of sheaves of abelian groups on $Y$
\begin{equation}\label{eq:XI.1.1}
0\longrightarrow P_n\xrightarrow{u}\mathcal O_{X_{n+1}}^*
\xrightarrow{v}\mathcal O_{X_n}^*\longrightarrow 1
\end{equation}
is \emph{exact}. More precisely, the group structure on $P_n$ is the additive
structure, $u(x)=1+x$ for every $x\in P_n$, and $v$ is the homomorphism induced
by the injection $\Ii^{n+2}\longrightarrow\Ii^{n+1}$. To see that $v$ is
surjective, observe that, for every $y\in Y$, $\mathcal O_{X_n,y}$ is a local
ring, the quotient of $\mathcal O_{X_{n+1},y}$ by a nilpotent ideal; the rest is
equally trivial. From~\eqref{eq:XI.1.1} we deduce an exact cohomology sequence:
\begin{equation*}\tag{$*$}\label{eq:XI.1.*}
\mathrm H^1(Y,P_n)\xrightarrow{u^1}\mathrm H^1(Y,\mathcal O_{X_{n+1}}^*)
\xrightarrow{v^1}\mathrm H^1(Y,\mathcal O_{X_n}^*)
\xrightarrow{d}\mathrm H^2(Y,P_n).
\end{equation*}

Moreover, for every $n\in\NN$, we know how to identify
$\Pic(X_n)$ with $\mathrm H^1(Y,\mathcal O_{X_n}^*)$. Furthermore, if $E$
is an invertible $\mathcal O_{X_{n+1}}$-module corresponding to a cohomology
class $c(E)$\pageoriginale, the cohomology class corresponding to the pullback
of $E$ to $X_n$ is equal to $v^1(c(E))$. Hence the following proposition:

\begin{proposition}\label{XI.1.1}
Retain the notation introduced above. Let $p\in\NN$. The map
\[
  \Pic(\widehat X)\longrightarrow\Pic(X_n)
\]
has the following properties:

\sourceoddity{SGA2-XI-L3609-PIC-YN-VS-XN-SRCDEF-001}{The French prints
$Y_n$. Expos\'e~XI locally defines and uses $X_n$, while Expos\'e~X used
$Y_n$ for the same infinitesimal thickening. We therefore read $X_n$ here.}

\begin{enumerate}
\renewcommand{\labelenumi}{(\roman{enumi})}
\item it is injective for $n\geq p$ if $\mathrm H^1(Y,P_n)=0$ for every
$n\geq p$;
\item it is an isomorphism for $n\geq p$ if $\mathrm H^i(Y,P_n)=0$ for
every $n\geq p$ and $i=1,2$.
\end{enumerate}
\end{proposition}

Of course, the exact sequence~\textup{($\ast$)} contains more
information than the proposition above. The reader will have noticed that
nothing has been said about the functor
$\PP(\widehat X)\longrightarrow\PP(Y)$. Given two invertible
$\mathcal O_{\widehat X}$-modules $E,F$,
$H=\SheafHom(E,F)$ is also invertible. If reduction modulo
$\Ii^{n+1}$ is denoted by a subscript $n$, we obtain an exact sequence:
\[
0\longrightarrow H_0\otimes P_n
\longrightarrow \SheafHom(E_{n+1},F_{n+1})
\longrightarrow \SheafHom(E_n,F_n)
\longrightarrow 0.
\]
Hence there is an exact cohomology sequence that we shall not write down and
whose interpretation is evident; this observation can be used to study the
functor~$\PP$.

\section{\texorpdfstring{Comparison of $\Pic(X)$ and $\Pic(\widehat X)$}{Comparison of Pic(X) and Pic(X-hat)}}\label{XI.2}

The reader will find in Expos\'e~X, no.~2, the proof of the following:






\begin{proposition} \label{XI.2.1} Suppose that $\Lef(X, Y)$ holds; then, for every open neighborhood $U$ of $Y$ in $X$, the functor
\begin{equation} \label{eq:XI.2.1} \PP(U)\to\PP(\hat X)
\end{equation}
is fully faithful, and the map
\begin{equation} \label{eq:XI.2.2} \Pic(U)\to\Pic(\hat X)
\end{equation}
is\pageoriginale therefore injective. If $\Leff(X, Y)$ holds, the map~(\ref{eq:XI.2.3}) is an isomorphism:
\begin{equation} \label{eq:XI.2.3}
\varinjlim_U\Pic(U)\to\Pic(\hat X){\quoi}.
\end{equation}
\end{proposition}

\begin{corollaire} \label{XI.2.2} Suppose that $\Lef(X, Y)$ holds and that, for every integer $n\geq p$, one has $\H^1(Y, P_n)=0$; then, for every open subset $U\supset Y$, the maps
\[
\Pic(X)\to\Pic(U)\to\Pic(X_n)
\]
are injective if $n\geq p$. If $\Leff(X, Y)$ holds and if, moreover, for every integer $n\geq p$, one has $\H^i(Y, P_n)=0$ for $i=1$ and $i=2$, then the map
\[
\varinjlim_U\Pic(U)\to\Pic(X_n)
\]
is an isomorphism for $n\geq p$.
\end{corollaire}

\section{Comparison of $\PP(X)$ and $\PP(U)$} \label{XI.3}

A definition:
\refstepcounter{subsection}\label{XI.3.1}
\begin{enonce*}[remark]{Definition 3.1\sfootnotemark}\sfootnotetext{For a more detailed study of the notion of parafactoriality, and for the proof of~\Ref{XI.3.3}, \cf \EGA IV 21.13, 21.14.}
Let $X$ be a prescheme and let $Z$ be a closed subset of~$X$. Put $U=X-Z$. We say that $X$ is \emph{parafactorial} at the points of $Z$ if, for every open subset $V$ of~$X$, the functor $\PP(V)\to \PP(V\cap U)$ is an equivalence of categories. We also say that the pair $(X, Z)$ is parafactorial.
\end{enonce*}

Recall that $\PP(Z)$ denotes the category of locally free modules of rank~$1$ on~$Z$.

\begin{definition} \label{XI.3.2}
A noetherian local ring is said to be \emph{parafactorial} if the pair $(\Spec(A), \{\rr(A)\})$ is parafactorial.
\end{definition}

One proves the following proposition, which shows that the notion is \og pointwise\fg:

\begin{proposition} \label{XI.3.3}
Suppose\pageoriginale that $X$ is locally noetherian. The pair $(X, Z)$ is parafactorial if and only if, for every $z\in Z$, the local ring $\Oo_{X, z}$ is parafactorial.
\end{proposition}

Note that \og parafactorial\fg includes \og fully faithful\fg. As in Lemma~\Ref{X.3.5} of Expos\'e~\Ref{X}, one proves the following:

\begin{lemme} \label{XI.3.4}
If $X$ is a locally noetherian prescheme and if $Z=X-U$ is a closed subset of $X$, the following conditions are equivalent:
\begin{enumeratei}
\item
for every open subset $V$ of $X$, the functor $\PP(V)\to \PP(V\cap U)$ is fully faithful;
\item
the homomorphism $\OX \to i_*(\Oo_U)$ is an isomorphism;
\item
for every $z\in Z$, one has $\prof(\Oo_{X, z})\geq 2$.
\end{enumeratei}
\end{lemme}

Thus \og parafactorial\fg means that the conditions of~\Ref{XI.3.4} are satisfied and that, for every open subset $V$ of $X$, the homomorphism $\Pic(V)\to\Pic(V\cap U)$ is surjective. In particular, if $X$ is the spectrum of a noetherian local ring, one obtains:

\begin{proposition} \label{XI.3.5}
Let $A$ be a noetherian local ring; it is parafactorial if and only if $\prof A\geq 2$ and $\Pic(X'-\{x\})=0$, where $X'=\Spec(A)$ and $x$ is the unique closed point of $X'$.
\end{proposition}

Notice that a local ring of dimension $\leq 1$ is never parafactorial, since its depth is $\leq 1$. Thus \emph{factorial does not imply parafactorial}; the implication does, however, hold for noetherian local rings of dimension $\geq 2$, as we shall see below.

\begin{lemme} \label{XI.3.6} Let $X$ be a locally noetherian prescheme and let $Z$ be a closed subset of $X$. Let $f\colon X_1\to X$ be a faithfully flat and quasi-compact morphism. Put $Z_1=f^{-1}(Z)$. If $(X_1, Z_1)$ is parafactorial, then $(X, Z)$ is parafactorial.
\end{lemme}

First note\pageoriginale that, if $i\colon(X-Z)\to X$ denotes the canonical immersion of $U=X-Z$ into $X$, the formation of the direct image under $i$ of a quasi-coherent $\Oo_U$-module commutes with base change by $f$, since $f$ is flat. It is therefore equivalent to assume the equivalent conditions of Lemma~\Ref{XI.3.4} for $(X, Z)$ or for $(X_1, Z_1)$, since $f$ is a morphism of descent for the category of quasi-coherent sheaves. It remains to prove that, for every open subset $V$ of $X$, $\Pic(V)\to\Pic(V\cap U)$ is surjective. We make the base change $V\to X$, which changes nothing (\sisi{sic}{\textit{sic}}), and we are reduced to $V=X$. We then note that, if ${\Lcal}$ is an invertible $\Oo_U$-module and if ${\Lcal}$ admits a locally free extension, this extension is isomorphic to $i_*({\Lcal})$, by what we have just seen. It remains to prove that $i_*({\Lcal})$ is invertible. Using once again the fact that direct image under $i$ commutes with flat base change, and that \og locally free of rank~$1$\fg is a property that descends under faithfully flat and quasi-compact morphisms, we are done.

\begin{corollaire} \label{XI.3.7} Let $A$ be a noetherian local ring; if $\hat A$ is parafactorial, then $A$ is parafactorial.
\end{corollaire}

One should not believe that if $A$ is parafactorial, then $\hat A$ is parafactorial\nde{For a precise study of the relationship between the factoriality of $A$ and that of its completion, see (Heitmann~R., {\og Characterization of completions of unique factorization domains\fg}, \emph{Trans. Amer. Math. Soc.} \textbf{337} (1993), \numero 1, pp\ptbl 379--387).}.

\danger

Before turning to the statement of the main theorem of this \no, let us make the connection with divisor theory and the notion of a factorial ring\sfootnote{For the generalities that follow, \cf also \EGA IV~21.}.

Let $X$ be a \emph{noetherian} and \emph{normal} prescheme. Let $Z^1(X)$ be the free abelian group generated by the $x\in X$ such that $\dim\OXx=1$. The local ring at such a point is a discrete valuation ring. We denote the corresponding normalized valuation by $v_x$. Let $K(X)$ be the ring of rational functions on $X$, and let
\[
p\colon K(X)^*\to Z^1(X)
\]
be the map that associates to every $f\in K(X)^*$ the codimension-one cycle:
\[
(f) = \sum_{x\in X, \;\dim\OXx=1} v_x(f)\cdot x{\quoi}.
\]
The image\pageoriginale of $p$ is denoted by $P(X)$, and its elements are called \emph{principal divisors}\sfootnote{In accordance with the terminology of \EGA IV~21, we now prefer to reserve the term \og divisors\fg for \og locally principal divisors\fg or \og Cartier divisors\fg.}. We set
\[
\Cl(X) = Z^1(X)/P(X){\quoi}.
\]
Let $\sisi{Z}{Z^{\prime1}(X)}$ be the subgroup of $Z^1(X)$ whose elements are the \emph{locally principal divisors}. We know that
\[
\Pic(X)\simeq Z^{\prime1}(X)/P(X){\quoi},
\]
and consequently $\Pic(X)$ is identified with a subgroup of $\Cl(X)$.

Note that if $U$ is a dense open subset of $X$, then $K(X)\to K(U)$ is an isomorphism, and that if $\codim(X-U, X)\geq 2$, \ie if every $x\in X$ such that $\dim\OXx\leq 1$ belongs to $U$, the homomorphism $Z^1(X)\to Z^1(U)$, and consequently $\Cl(X)\to\Cl(U)$, is also an isomorphism. Finally, if every $x\in U$ is factorial, \ie if $\OXx$ is factorial, then $Z^1(U) = Z^{\prime1}(U)$, and hence $\Pic(U)\simeq\Cl(U)$.

\begin{propositionblah} \label{XI.3.7.1}
Let $X$ be a noetherian and normal prescheme. Let $(U_i)_{i\in I}$ be a family of open subsets of $X$ such that:
\begin{enumeratea}
\item
the $U_i$ form a filter base\nde{\ie a decreasing filtered family.};
\item
if we set $Y_i=X-U_i$, then $\codim(Y_i, X)\geq 2$ for every $i$,
\item
if $x\in U_i$ for every $i\in I$, then $\OXx$ is factorial.
\end{enumeratea}
\noindent
Then there is an isomorphism:
\[
\varinjlim_{i\in I}\Pic(U_i)\lto{\approx}\Cl(X){\quoi}.
\]
\end{propositionblah} Note that b) implies that every $x\in X$ such that $\dim\OXx\leq 1$ belongs to $U_i$ for every $i$. Thus the $U_i$ are dense, and moreover the homomorphism $Z^1(U_i)\to Z^1(X)$ is an isomorphism, as is $K(U_i)\to K(X)$. Thus $\Pic(U_i)\subset \Cl(U_i) \simeq \Cl(X)$. To prove the desired assertion, it is therefore enough to show that every $D\in Z^1(X)$ belongs to $Z^{\prime1}(U_i)$ for a suitable $i$. It is enough to do this for irreducible positive \og divisors\fg. Thus let $x\in X$ be such that\pageoriginaled $\dim\OXx=1$. It is enough to prove that there exists an $i\in I$ such that $\overline{\{x\}}$ is locally principal at the points of $U_i$. Let $\Ii$ be the largest ideal of definition of the closed subset $\overline{\{x\}}$. The set of points in a neighborhood of which $\Ii$ is free is an open subset $U$. Now $U\supset \bigcap_{i\in I}U_i$ by~c). If we set $Y=X-U$, then $Y\subset \bigcup_{i\in I}Y_i$\sisi{}{ with $Y_i=X-U_i$}; now $Y$ is closed, and therefore has only \emph{finitely} many generic points, so it is contained in the union of finitely many $Y_i$, and hence in $Y_j$ for some $j\in I$, since the $U_i$ form a filter base. Thus $U\supset U_j$.\qed

\begin{corollaire} \label{XI.3.8}
Let $X$ be a noetherian and normal prescheme, and let $Y$ be a closed subset of codimension $\geq 2$. Suppose that, for every $p\in X-Y$, $\Oo_{X, p}$ is factorial. Then
\[
\Pic(X-Y)\to\Cl(X-Y)\to\Cl(X)
\]
are isomorphisms.
\end{corollaire}

\begin{corollaire} \label{XI.3.9}
Let $A$ be a noetherian normal local ring. Set $X'=\Spec(A)$ and $x = \rr(A)$. For $A$ to be \emph{factorial}, it is necessary and sufficient that $\Pic(X'-\{x\})=0$ and that $\pp\in X'-\{x\}$ imply that $A_\pp$ is factorial.
\end{corollaire}
Indeed, for $A$ to be factorial, it is necessary and sufficient that $\Cl(X')=0$\,\nde{see Bourbaki, \emph{Alg\`ebre commutative} VII.1.4, corollary to Theorem~2, and VII.3.2, Theorem~1}.

\begin{corollaire} \label{XI.3.10}
Let $A$ be a noetherian local ring of dimension $\geq 2$. Set $X'=\Spec(A)$, and let $x = \rr(A)$. Set $X=X'-\{x\}$. The following conditions are equivalent:
\begin{enumeratei}
\item
$A$ is factorial;
\item
\textup{a)}\enspace for every $y\in X$, $\Oo_{X, y}$ is factorial, and
\par\hspace*{6.3mm}\textup{b)}\enspace $A$ is parafactorial, \ie $\prof A\geq 2$ and $\Pic(X)=0$.
\end{enumeratei}
\end{corollaire}

Before\pageoriginale proving this corollary, let us state the
\refstepcounter{subsection}\label{XI.3.11}
\begin{enonce*}{Serre's normality criterion 3.11\sfootnotemark}\sfootnotetext{\Cf \EGA IV~5.8.6.} Let $A$ be a noetherian local ring. For $A$ to be normal, it is necessary and sufficient that
\begin{enumeratei}
\item
for every prime ideal ${\pp}$ of $A$ such that $\dim A_{{\pp}}\leq 1$, $A_{{\pp}}$ is normal,
\item
for every prime ideal ${\pp}$ of $A$ such that $\dim A_{{\pp}}\geq 2$, one has $\prof A_{{\pp}}\geq 2$.
\end{enumeratei}
\end{enonce*}

Let us prove~\Ref{XI.3.10}.

(i)\ALORS (ii). Since a localization of a factorial ring is again factorial, we have (ii)~a). Moreover, $A$ is normal, and hence $\prof A\geq 2$, since $\dim A\geq 2$ (\Ref{XI.3.11}~(ii)). Finally, $A$ is parafactorial; indeed, $\Pic(X)\simeq\Cl(X')=0$ (\cf \Ref{XI.3.9}).

(ii)\ALORS (i). Let us first prove that $A$ is \emph{normal} by applying \sisi{SERRE}{Serre}'s criterion. Since $\dim A\geq 2$, condition~(i) of the criterion is among the hypotheses. Moreover, for every ${\pp}\in X$, $A_{{\pp}}$ is factorial, hence normal, and therefore of depth $\geq 2$, at least if $\dim A_{{\pp}}\geq 2$. Finally, $\prof A\geq 2$ by (ii)~b). It remains to apply~\Ref{XI.3.9}.

\smallskip
Let us summarize the foregoing:

\begin{proposition} \label{XI.3.12}
Let $X$ be a locally noetherian prescheme, and let $\Ii$ be a quasi-coherent ideal of $X$. Let $Y=V(\Ii)$. Let $p\in\NN$. Suppose that:
\begin{enumerate}
\item[1)] one \sisi{has}{has} $\Leff(X, Y)$ (\textup{Expos\'e} \Ref{X});
\item[2)] $\H^i(X, \Ii^{n+1}/\Ii^{n+2})=0$ if $i=1$ or $2$ and if $n\geq p$;
\item[3)] for every open neighborhood $U$ of $Y$ in $X$ and every $x\in X-U$, the ring $\OXx$ is parafactorial.
\end{enumerate}
\noindent
Then, for every $n\geq p$, and every open neighborhood $U$ of $Y$, the homomorphisms
\[
\xymatrix{\Pic(X)\ar[rr]\ar[dr]&&\Pic(X_n)\\&\Pic(U)\ar[ur]&}
\]
are isomorphisms.
\end{proposition}

We\pageoriginale know some parafactorial rings:

\skpt
\begin{theoreme} \label{XI.3.13}
\begin{enumeratei}
\item
(Auslander--Buchsbaum)\nde{compare with the following purity result, due to Gabber. Let $X$ be the spectrum of a regular local ring $A$ of dimension $3$, let $a$ have nonzero differential, \ie $a\in\mm-\mm^2$, and let $U$ be the complement of $V(a)$. Then a vector bundle on $U$ is free (for a simple proof, see (Swan~R.G., {\og A simple proof of Gabber's theorem on projective modules over a localized local ring\fg} \emph{Proc. Amer. Math. Soc.} \textbf{103} (1988), \numero 4, p\ptbl 1025--1030). The rank-$1$ case is a special case of Theorem~\Ref{XI.3.13}. For purity results concerning vector bundles of arbitrary rank, whether in the analytic or algebraic setting, see (Gabber~O., {\og On purity theorems for vector bundles\fg}, \emph{Internat. Math. Res. Notices} (2002), \numero 15, p\ptbl 783--788).} A regular noetherian local ring is factorial (and hence parafactorial if its dimension is $\geq 2$).
\item
A noetherian local ring of dimension $\geq 4$ that is a \emph{complete intersection} is parafactorial.
\end{enumeratei}
\end{theoreme}

\begin{corollaire}[Conjecture of\label{XI.3.14} \sisi{SAMUEL}{Samuel}\ndemark]\ndetext{for a proof in the same spirit, but more elementary, see Call~F.\ \&\ Lyubeznik~G., {\og A simple proof of Grothendieck's theorem on the parafactoriality of local rings\fg}, in \emph{Commutative algebra: syzygies, multiplicities, and birational algebra (South Hadley, MA, 1992)}, Contemp. Math., vol.~159, American Mathematical Society, Providence, RI, 1994, p\ptbl 15--18.} A noetherian local ring $A$ that is a complete intersection and is factorial in codimension $\geq 3$ (\ie $\dim A_{{\pp}}\leq 3$ implies that $A_{{\pp}}$ is factorial) is factorial.
\end{corollaire}

\emph{Proof of the corollary}

We argue by induction on the dimension of $A$.

If $\dim A\leq 3$, then $A$ is factorial by hypothesis.

If $\dim A>3$, then by the induction hypothesis, noting that a localization of a complete intersection is again a complete intersection, all localizations of $A$ other than $A$ are factorial. By Theorem~\Ref{XI.3.13}~(ii), $A$ is parafactorial, and therefore factorial by~\Ref{XI.3.10}.

Let us prove~\Ref{XI.3.13}~(i) (following \sisi{KAPLANSKY}{Kaplansky})\sfootnote{This is the proof reproduced in \EGA IV~21.11.1.}.

Let $A$ be a regular noetherian local ring, and set $\dim A=n$. If $n=0$ or $1$, the result is known. Suppose $n\geq 2$; we argue by induction on $n$, and suppose $n\geq 2$ and the theorem proved for rings of dimension $<n$. Set $X'=\Spec(A)$ and $X=X'-\{x\}$, where $x=\rr(A)$. The localizations of $A$ other than $A$ are regular and of dimension $<n$, and hence factorial. Moreover, $\prof A = \dim A\geq 2$. It is therefore enough to prove that $\Pic(X)=0$ (cor\ptbl\Ref{XI.3.10}). Thus let ${\Lcal}$ be an invertible $\OX$-Module; we know that it can be extended to a coherent $\Oo_{X'}$-Module ${\Lcal}'$. There exists a resolution of ${\Lcal}'$ by free $\Oo_{X'}$-Modules:
\[
0\from {\Lcal}'\from {\Lcal}'_1\from\cdots\from{\Lcal}'_n\from 0{\quoi},
\]
because\pageoriginale the cohomological dimension of $A$ is finite. By restriction to $X$ we obtain a finite free resolution. It is therefore enough to prove the following lemma:

\begin{lemme} \label{XI.3.15} Let $(X, \OX)$ be a ringed space, and let ${\Lcal}$ be a locally free $\OX$-Module that admits a finite resolution by \emph{finite-type free} modules. Then $\det({\Lcal}) \simeq \OX$.
\end{lemme}

Recall that $\det({\Lcal})$ is defined as the highest exterior power of ${\Lcal}$. In the case under consideration, $\det({\Lcal}) \simeq{\Lcal}$, because ${\Lcal}$ is invertible, and hence the lemma gives the conclusion. Let us prove this lemma. Let
\[
0\from {\Lcal}_0\from {\Lcal}_1\from {\Lcal}_2\from\cdots\from{\Lcal}_n\from 0{\quoi},
\]
be the asserted exact sequence, where ${\Lcal}_0={\Lcal}$. Since everything is locally free, we have:
\[
\tbigotimes_{0\leq i\leq n}(\det({\Lcal}_i))^{(-1)^i} \simeq \OX
\]
and all the ${\Lcal}_i$, $i>0$, are free, and hence so are their determinants, and consequently so is the determinant of ${\Lcal}_0 = {\Lcal}$.\qed

\medskip
It remains for us to prove part~(ii) of the theorem. First, let us prove a lemma that will allow us to proceed by induction:

\begin{lemme} \label{XI.3.16}
Let $A$ be a noetherian local ring that is a quotient of a regular ring. Let $t\in\rr(A)$ be an $A$-regular element. Suppose that $A$ is complete for the $t$-adic topology. Set $X'=\Spec(A)$, $Y'=V(t)\simeq \Spec(B)$, $B=A/tA$, $X=X'-\{x\}$, $Y=Y'-\{x\}$, and $x=\rr(A)$. Suppose that:
\begin{enumeratea}
\item
for every $y\in X$ closed in $X$, one has $\prof\Oo_{X, y}\geq 2$,
\item
$\prof A/tA\geq 3$,
\end{enumeratea}
\noindent
then\pageoriginale the map $\Pic(X)\to\Pic(Y)$ is injective. In particular, if $B$ is parafactorial, then $A$ is also parafactorial.
\end{lemme}

We know that a) implies $\Lef(X, Y)$ by \Ref{X}~\Ref{X.2.1}. If we prove that $\H^1(Y, P_n)=0$ for every $n\geq 0$, then we will know \sisi{\ignorespaces}{by} (\Ref{XI.2.2}) that $\Pic(X)\to\Pic(Y)$ is injective. If, moreover, $B$ is parafactorial, then we will know that $\Pic(Y)=0$ (\Ref{XI.3.5}), and hence $\Pic(X)=0$; moreover, $\prof(A)=\prof(B)+1\geq 3+1\geq 2$, since $t$ is $A$-regular, and thus $A$ will be parafactorial by~\Ref{XI.3.5}.

Let $\Ii=\widetilde{(tA)}$ be the $\Oo_{X'}$-Module associated with the ideal $tA$. In \numero \Ref{XI.1}, we set $P_n = \sisi{(\Ii^{n+1}/\Ii^{n+2}){|Y}}{(\Ii^{n+1}/\Ii^{n+2})_{|Y}}$ for every $n\geq 0$. Since $t$ is $A$-regular, $P_n\simeq \Oo_Y$. \emph{It therefore remains for us to prove that $\H^1(Y, \Oo_Y)=0$.} Now $Y = Y'-\{x\}$ is an open subset of $Y'$, and hence we have an exact sequence~(\Ref{I}~(\Ref{eq:I.27})):
\[
\H^1(Y', \Oo_{Y'})\to\H^1(Y, \Oo_{Y})\to\H^2_x(Y', \Oo_{Y'}){\quoi},
\]
whose right-hand term is zero by hypothesis~b), and whose left-hand term is zero because $Y'$ is affine.\qed

\begin{lemme} \label{XI.3.17}
Keeping the hypotheses of~\Ref{XI.3.16}, suppose in addition that:
\begin{itemize}
\item[c)] for every $y$ closed in $Y$, one has $\prof\Oo_{X, y}\geq 3$,
\item[d)] $\prof A/tA\geq 4$ (stronger than b)),
\item[e)] for every $y$ closed in $X$ with $y\in Y$, the ring $\Oo_{X, y}$ is parafactorial.
\end{itemize}
Then the map $\Pic(X)\to\Pic(Y)$ is an isomorphism; in particular, for $A$ to be parafactorial, it is necessary and sufficient that $B$ be parafactorial.
\end{lemme}

\enlargethispage{5mm}%
We know~(\Ref{X}~\Ref{X.2.1}) that a)~and~c) imply $\Leff(X, Y)$. Moreover, by the argument just given, d) implies that $\H^i(Y, P_n)=0$ for every $n\geq 0$ and $i=1$ or $i=2$. Moreover, for every open neighborhood $U$ of $Y$ in $X$, the complement of $U$ in $X$ consists of finitely many closed points. By~e) and Theorem~\Ref{XI.3.12}, we deduce that $\Pic(X)\to \Pic(Y)$ is an isomorphism. On the other hand, $\prof A\geq \prof B\geq 2$; by the criterion~\Ref{XI.3.5}, we deduce\pageoriginale that $A$ is parafactorial if and only if $B$ is.

Let us now prove~\Ref{XI.3.13}~(ii). Let $R$ be a regular noetherian local ring. Let $(t_1, \ldots, t_k)$ be an $R$-sequence. Set $B=R/(t_1, \ldots, t_k)$ and suppose that $\dim B\geq 4$. We must prove that $B$ is parafactorial. We argue by induction on~$k$. If $k=0$, then $B$ is regular, and hence factorial by \Ref{XI.3.13}~(i), and therefore parafactorial by~\Ref{XI.3.10}. Suppose $k\geq 1$ and the theorem proved for $k'<k$. Set $A=R/(t_1, \ldots, t_{k-1})$, so that $B=A/t_k A$. We may suppose that $B$ is complete by~\Ref{XI.3.7}. By the induction hypothesis, $A$ is parafactorial. Let us prove that Lemma~\Ref{XI.3.17} applies. We have assumed that $B$ is complete, and hence so is $A$, and therefore $A$ is complete for the $t_k$-adic topology. If $x\in X$ and $x$ is closed in $X$, then $A_x$ is a complete intersection of dimension $\geq 4$, with $k'<k$. By the induction hypothesis, $A_x$ is parafactorial, and moreover has depth $\geq 4$. This gives a), c), and~e). Moreover, $\dim A\geq 5$, which gives~d).\qed

\begin{theoreme} \label{XI.3.18}
Let $X$ be a locally noetherian prescheme, and let $\Ii$ be a coherent sheaf of ideals of $X$. Set $Y=V(\Ii)$. Let $n$ be an integer. Suppose that:
\begin{enumeratei}
\item
one has $\sisi{\leff}{\Leff}(X, Y)$ (\cf examples \Ref{X}~\Ref{X.2.1}~and~\Ref{X}~\Ref{X.2.2}),
\item
for every $p\geq n$, one has $\H^i(Y, \Ii^{p+1}/\Ii^{p+2})=0$ for $i=1$ and $i=2$,
\item
for every open subset $U\supset Y$ and every $x\in X-U$, the ring $\OXx$ is regular of dimension $\geq 2$ or a complete intersection of dimension $\geq 4$.
\end{enumeratei}
\noindent
Then, for every open subset $U\supset Y$ and every integer $p\geq n$, the homomorphisms
\[
\Pic(X)\to\Pic(U)\to\Pic(Y_p)
\]
are isomorphisms, where $Y_p$ denotes the prescheme $(Y, \OX/\Ii^{p+1})$.
\end{theoreme}

It is enough to combine \Ref{XI.3.12}~and~\Ref{XI.3.13}.

% END INCLUDED active_source/parts/11_expose_XI.tex

% BEGIN INCLUDED active_source/parts/12_expose_XII.tex
\chapter[Applications to projective algebraic schemes]
{Applications\\ to projective algebraic schemes} \label{XII}

\renewcommand{\theequation}{\arabic{equation}}

\section[Projective duality theorem and finiteness theorem]
{Projective duality theorem and finiteness theorem\protect\sfootnotemark}\label{XII.1}

The\pageoriginale following theorem,\sfootnotetext{The present expos\'e, written in January~1963, is considerably more detailed than was the oral presentation in June~1962.} essentially contained in FAC\sfootnote{J.-P.~Serre, \sisi{Faisceaux alg\'ebriques Coh\'erents, Ann. of Math. t.61 (1955) 197-278}{{\og Faisceaux alg\'ebriques coh\'erents\fg}, \emph{Ann. of Math.} \textbf{61} (1955), p\ptbl197-278}.} (except that, at the time, Serre did not have at his disposal the language of Ext for sheaves of modules\nde{Readers interested in the history of mathematics will profitably consult the letter that Grothendieck wrote to Serre on 15 December 1955 and Serre's reply of 22 December of the same year; see \emph{Correspondance Grothendieck-Serre}, edited by Pierre Colmez and Jean-Pierre Serre, Documents Math\'ematiques, vol.~2, Soci\'et\'e Math\'ematique de France, Paris, 2001.}), is the global analogue of the local duality theorem (Expos\'e \Ref{IV}), which was modeled on it.

\begin{theoreme} \label{XII.1.1}
Let $k$ be a field, $X=\PP_k^r$ \sisi{the standard projective scheme}{projective space} of dimension $r$ over $k$, and $F$~a variable coherent module on $X$; then there is an isomorphism of $\partial$-functors in~$F$:
\steco
\begin{equation} \label{eq:XII.1} {}\H^i(X, F)' \isomfrom \Ext^{r-i}(X;F, \omx^r),
\end{equation}
where we set
\steco
\begin{equation} \label{eq:XII.2} {}\Omega_{X/k}^r= \Oo_{\PP_k^r}(-r-1).
\end{equation}
\end{theoreme}

\begin{remarquestar}
This module is of course also the module of relative differentials of degree $r$ of $X$ over $k$. In this form, the theorem remains valid if $X$ is a proper smooth scheme over $k$ (for the projective case, see A. Grothendieck, \og \sisi{Th\'eor\`emes de dualit\'e pour les faisceaux alg\'ebriques coh\'erents}{\emph{Th\'eor\`emes de dualit\'e pour les faisceaux alg\'ebriques coh\'erents}}\fg, S\'eminaire Bourbaki, May 1957)\sfootnote{For a more general duality theorem, \cf the Hartshorne seminar cited at the end of \Exp \Ref{IV}.}. When $F$ is locally free, one recovers Serre's duality theorem $\H^i(X, F)' \isomto \H^{r-i}(\SheafHom_{\OX}(F, \omx^r))$. Theorem \Ref{XII.1.1} (which also recovers the case of a projective $X$ over $k$, as in {\loccit}) will suffice for our purposes.
\end{remarquestar} The homomorphism\pageoriginale (\Ref{eq:XII.1}) is deduced from the Yoneda pairing
\steco
\begin{equation} \label{eq:XII.3} {}\H^i(X, F) \times \Ext^{r-i}(X;F, \omx^r) \to \H^r(X, \omx^r),
\end{equation}
and from a well-known isomorphism (\cf FAC, or \EGA III 2.1.12):
\steco
\begin{equation} \label{eq:XII.4} {}\H^r(X, \omx^r) = \H^r(\PP_k^r, \Oo_{\PP_k^r}(-r-1)) \isomto k.
\end{equation}
To show that (\Ref{eq:XII.1}) is an isomorphism, one proceeds as in the case of the local duality theorem, noting that $\H^r(X, F)$, as a functor of $F$, is right exact (since $\H^n(X, F)=0$ for $n>r$), and that every coherent module is isomorphic to a quotient of a direct sum of modules of the form $\Oo(-m)$, with $m$ large. By descending induction on $i$, this reduces us to checking the assertion for a sheaf of the form $\Oo(-m)$, where it is contained in the well-known explicit computations (FAC, or \EGA III 2.1.12). We may moreover assume $-m \leq -r-1$, in which case $\H^i(X, \Oo(-m))=0$ for $i\neq r$.

\begin{corollaire} \label{XII.1.2} For a given coherent $F$, and for $m$ sufficiently large, there is a canonical isomorphism
\steco
\begin{equation} \label{eq:XII.5} {}\H^i(X, F(-m))'\isomto\H^0(X, \SheafExt_{\OX}^{r-i}(F, \omx^r)(m))
\end{equation}
(where $'$ denotes the vector-space dual).
\end{corollaire}
Indeed, on \sisi{a projective scheme}{projective space} $X$, for every pair of coherent sheaves $F,G$ and for $m$ sufficiently large, there is a canonical isomorphism:
\steco
\begin{equation} \label{eq:XII.6} {}\Ext^n(X;F(-m), G)\simeq \Ext^n(X;F, G(m)) \isomto \H^0(X, \SheafExt_{\OX}^{n}(F, G)(m)),
\end{equation}
(the isomorphism between the first two terms being trivially valid for every $m$), as follows from the global $\Ext$ spectral sequence
\[
\H^p(X, \SheafExt^q_{\OX}(F, G(m))) \To \Ext^{\boule}(X;F, G(m)),
\]
which\pageoriginale degenerates for $m$ sufficiently large, by virtue of the fact that
\[
\SheafExt^q_{\OX}(F, G(m)) \simeq \SheafExt^q_{\OX}(F, G)(m),
\]
and that the $\SheafExt^q_{\OX}(F, G)$ are coherent sheaves. Thus (\Ref{eq:XII.5}) follows from (\Ref{eq:XII.6}) and (\Ref{eq:XII.1}).

\begin{corollaire} \label{XII.1.3}
\sisi{Equivalent conditions for given $i,F$}{For given $i,F$, the following conditions are equivalent}:
\begin{enumeratei}
\item
$\H^i(X, F(-m))=0$ for $m$ large.
\item[\textup{(i bis)}] $\H^i(X, F(\cdot))= \sisi{\coprod}{\bigoplus}_{m\in\ZZ}\H^i(X, F(-m))$ is a finite-type $S$-module, where $S=k[t_0, \ldots, t_r]$.
\item
$\SheafExt^{r-i}_{\OX}(F, \omx^r)=0$.
\item[\textup{(ii bis)}] $\SheafExt^{r-i}_{\OX}(F, \OX)=0$.
\item
$\H^i_x(F_x)=0$ for every closed point $x$ of $X$.
\item
$\H^{i+1}_x(\widetilde{F}_x)=0$ for every closed point $x$ of the punctured affine cone $\widetilde{X}=\Spec(S)-\Spec(k)$ over $X$, where $\widetilde{F}$ denotes the inverse image of $F$ under the canonical morphism $\widetilde{X}\to X$.
\end{enumeratei}
\end{corollaire}

\begin{proof}
(i) \SSI (i bis), since the submodule of $\H^i(X, F(\cdot))$ that is the sum of the homogeneous components of degree $\geq \nu$ is of finite type over $S$ (in fact, for $i\neq 0$, it is even of finite type over $k$), (\cf FAC or \EGA III 2.2.1 and 2.3.2).

(i) \SSI (ii) by Corollary \Ref{XII.1.2}.

(ii) \SSI (ii bis), because $\Omega_{X/k}^r$ is locally isomorphic to $\OX$.

(ii bis) \SSI (iii) by the local duality theorem for $\OXx$ (which is a regular\pageoriginale local ring of dimension $r$), according to which the \og dual\fg of $\SheafExt^{r-i}_{\OX}(F, \OX)_x$ is identified with $\H^i_x(F_x)$ (\hyperref[V.2.1]{V 2.1)}.

(ii bis) is equivalent to the analogous relation
\[
\SheafExt^{r-i}_{\Oo_{\widetilde{X}}}(\widetilde{F}, \Oo_{\widetilde{X}})=0
\]
(because $\widetilde{X}\to X$ is faithfully flat, so the inverse image of $\SheafExt^{r-i}_{\OX}(F, \OX)$ is isomorphic to $\SheafExt^{r-i}_{\Oo_{\widetilde{X}}}(\widetilde{F}, \Oo_{\widetilde{X}})$), and this last relation is equivalent to (iv) by the local duality theorem for the local ring $\Oo_{\widetilde{X},x}$, which is regular of dimension $r+1$.
\skipqed
\end{proof}

In particular, applying this to all $i\leq n$, we find:

\begin{corollaire} \label{XII.1.4}
Equivalent conditions for given $n,F$:
\begin{enumeratei}
\item
$\H^i(X, F(-m))=0$ for $i\leq n$ and $m$ large.
\item[\textup{(i bis)}] $\H^i(X, F(\cdot))$ is a finite-type $S$-module for $i\leq n$.
\item
$\prof(F_x)>n$ for every closed point $x$ of $X$.
\item
$\prof(\widetilde{F}_x)>n+1$ for every closed point $x$ of $\widetilde{X}$.
\end{enumeratei}
\end{corollaire}

The point of Corollaries \Ref{XII.1.3} and \Ref{XII.1.4} is to express a global condition (i) or~(i~bis) in terms of local conditions, \sisi{\ignorespaces}{namely} the vanishing of local invariants such as $\H_x^i(X, F_x)$ or $\H_x^i(X, \widetilde{F}_x)$, or an inequality involving depth. In this form, these results remain trivially valid for an arbitrary projective scheme $X$, and a very ample invertible sheaf $\OX(1)$ on $X$, as one sees by means of a suitable projective immersion $X\to\PP_k^r$. (Of course, conditions \Ref{XII.1.3}~(ii) and \Ref{XII.1.3}~(ii bis) are no longer equivalent to the others in this general case, unless, for example, one assumes that $X$ is regular.) One can moreover generalize to the case of a projective morphism $X\to S$ as follows:

\begin{proposition} \label{XII.1.5}
Let\pageoriginaled $f:X\to S$ be a projective morphism, with $S$ noetherian, let $\OX(1)$ be an invertible module on $X$ that is very ample relative to $S$, let $F$ be a coherent module on $X$, flat over $S$, let $s$ be a point of $S$, let $X_s$ be the fiber of $X$ at $s$ (regarded as a projective scheme over $k(s)$), let $F_s$ be the sheaf induced on $X_s$ by $F$, and finally let $i$ be an integer. Suppose that, for every closed point $x$ of $X_s$, one has $\H^i_x(F_{s, x})=0$ (for example, $\prof(F_{s, x})>i$). Then there exists an open neighborhood $U$ of $s$ such that the same condition holds for every $s'\in U$. Moreover, for such a $U$, one has
\[
\R^if_\ast(F(-m))= 0 \; \text{for}\; m \; \text{large, }
\]
and if $\Ss$ is a quasi-coherent graded algebra on $S$, generated by $\Ss^1$, and defining~$X$ together with $\OX(1)$ as $X=\Proj(\Ss)$, $\OX(1)=\Proj(\Ss(1))$, then the $S$-module
\[
\R^if_\ast(F(\cdot)) = \sisi{\coprod}{\tbigoplus}_{m\in\ZZ}\R^if_\ast(F(m))
\]
is of finite type over $U$.
\end{proposition}

Embed $X$ in an $X'=\PP_S^r$ in such a way that $\OX(1)$ is induced by $\Oo_{\sisi{P}{{X'}}}(1)$ (this is possible provided that $S$ is replaced by an affine neighborhood of $s$). Put\refstepcounter{toto}\label{P5}\sfootnote{The first part of \Ref{XII.1.5} can also be obtained immediately by applying the purely local statement \EGA IV 12.3.4 to the preceding $\underline{E}^j$, thereby bypassing most of the proof that follows.} for every integer $j$, and every $t\in S$:
\steco
\begin{equation} \label{eq:XII.7} {}\underline{E}^j(t)= \SheafExt^j_{\Oo_{X'_t}} (F_t, \Oo_{X'_t}(-r-1))
\end{equation}
Thus $\underline{E}^j(t)$ is a coherent module on $X_t$; I claim that, as $t$ varies, the family of these modules is \og constructible\fg in the following sense: for every $t\in S$ there exists a nonempty open subset $V$ of the closure $\overline{t}$, endowed with the induced reduced structure, and a coherent module $\underline{E}^j(V)$ on $X_V=X\times_SV$, \emph{flat} over $V$, such that, for every $t'\in V$, $\underline{E}^j(t')$ is isomorphic to the module induced by $\underline{E}^j(V)$ on $X_{t'}$. To verify this assertion, putting $Z=\overline{t}$ with its induced structure, we consider the coherent modules\pageoriginale
\[
\underline{E}^j(Z)= \SheafExt^j_{\Oo_{X'_Z}} (F_Z, \Oo_{X'_Z}(-r-1))
\]
(where the subscript $Z$ still means that one takes the induced object \sisi{over}{over} $Z$), and we take for $V$ a nonempty open subset of $Z$ such that the modules $\underline{E}^j(Z)$ are flat \sisi{over}{over} $V$: this is possible, since one verifies immediately that $\underline{E}^j(Z)=0$ for $j$ outside the interval $[0, r]$, and one can apply \SGA 1 IV 6.11. We then take $\underline{E}^j(V)=\underline{E}^j(Z)|_{X_V}$, and one readily verifies that it has the required property.

It follows from the preceding remark that there exists a finite partition of $S$ formed by sets $V_\alpha$ of the form $V=V(t)$ as above (noetherian induction), and, applying Serre's theorem \EGA III 2.2.1 to the $\underline{E}^j(V_\alpha)$, one sees that there exists an integer $m_0$ such that
\[
\R^if_{V_\alpha\ast}(\underline{E}^j(V_\alpha)(m))=0\textup{ for } i\neq 0, m\geq m_0, \textup{ for all }j,
\]
from which it follows, using the flatness of $\underline{E}^j(V_\alpha)$ over $V_\alpha$ and straightforward K\"unneth-type relations (\cf \EGA III par\ptbl7), that
\[
\H^i(X_t, \underline{E}^j(t)(m))=0\textup{ for } i\neq 0, m\geq m_0, \textup{ for all }j,
\]
for every $t\in V_\alpha$, and hence for every $t$, since the $ V_\alpha$ cover $S$. From this and the spectral sequence for global Ext it follows, by virtue of \Ref{XII.1.1}, as in the proof of \Ref{XII.1.2}, that there is an isomorphism
\steco
\begin{equation} \label{eq:XII.8} {}\H^i(X_t, F_t(-m))' \isomfrom \H^0(X_t, \underline{E}^{r-i}(t)(m))\textup{ for } m\geq m_0,
\end{equation}
for every integer $i$, and every $t\in S$.

Now use the hypothesis on $F_s$, which can be written
\steco
\begin{equation} \label{eq:XII.9} {}\underline{E}^{r-i}(s)=0,
\end{equation}
and which, by virtue of (\Ref{eq:XII.8}), is equivalent to
\steco
\begin{equation} \label{eq:XII.10} {}\H^i(X_s, F_s(-m))=0 \; \text{for}\; m\geq m_0.
\end{equation}
Since $F$, and hence $F(-m)$, is flat over $S$, it follows from
the K\"unneth-type relations already invoked that (for a given $m$) the
same relation (\Ref{eq:XII.10}) holds when\pageoriginale $s$ is replaced by a
$t$ near $s$, in particular by every generization $t$ of
$s$. By virtue of (\Ref{eq:XII.8}), we shall therefore have, for such a
generization, \steco
\begin{equation} \label{eq:XII.11} {}\underline{E}^{r-i}(t)=0,
\end{equation}
and the set of $t\in \sisi{s}{S}$ for which this relation holds is plainly constructible (since it induces an open subset on each $V_\alpha$); as it contains the generizations of $s$, it contains an open neighborhood $U$ of $s$. This proves the first assertion of \Ref{XII.1.5}. Moreover, for $t\in U$, it follows from (\Ref{eq:XII.11}) and (\Ref{eq:XII.8}) that
\steco
\begin{equation} \label{eq:XII.12} {}\H^i(X_t, \sisi{F_t(-m)}{F_t(-m))}=0 \; \text{for}\; m\geq m_0, \; t\in U,
\end{equation}
which, by virtue of the K\"unneth-type relations, implies (and in fact is much stronger than)
\steco
\begin{equation} \label{eq:XII.13} {}\R^if_\ast(F(-m))=0 \; \text{on}\; U, \; \text{for} \; m\geq m_0.
\end{equation}
This proves the second assertion of \Ref{XII.1.5}. Finally, the last assertion follows immediately by proceeding as at the beginning of the proof of \Ref{XII.1.3}.

\refstepcounter{subsection}\label{XII.1.6}
\begin{enonce*}[remark]{Remark 1.6\sfootnotemark} \sfootnotetext{This remark is made more precise by footnote \sisi{5.}{\pageref{P5}.}} The proof simplifies considerably (eliminating all consideration of constructibility) when one assumes from the outset that the hypothesis made for $s\in S$ holds for every $s'\in S$. In fact, when one assumes that $F_s$ has depth $>i$ at the closed points of $X$, there is a general statement, \emph{local in nature on} $X$, which says that the same condition holds for all the $X_t$, provided that $X$ is replaced by a suitable open neighborhood of the fiber $X_s$ (in other words, a certain subset of $X$, defined by conditions on the modules induced by $F$ on the fibers, is open, \cf \EGA IV). Since $f$ is proper here, this neighborhood can therefore be taken to be of the form $f^{-1}(U)$, which recovers the first assertion of \Ref{XII.1.5} without laborious d\'evissage. In this general case, one can still prove by the method of \loccit that the first assertion of \Ref{XII.1.5} (proved here globally, using the fact that $X$ is projective over $S$) again follows from a purely local statement on $X$ (which the reader may formulate explicitly if useful).
\end{enonce*}

\section[Lefschetz theory for a projective morphism: comparison]
{Lefschetz theory for a projective morphism: Grauert comparison theorem} \label{XII.2}

\setcounter{equation}{13}
It is\pageoriginale the following theorem:

\begin{theoreme} \label{XII.2.1}
Let $f:X\to S$ be a projective morphism, with $S$ Noetherian, let $\OX(1)$ be an invertible module on $X$, relatively ample over $S$, let $Y$ be the prescheme of zeros of a section $t$ of $\OX(1)$, let $J$ be the ideal defining $Y$, let $X_n$ be the closed subprescheme of $X$ defined by $J^{n+1}$, let $\widehat{X}$ be the formal completion of $X$ along $Y$, let $\widehat{f}:\widehat{X}\to S$ be the composite morphism $\widehat{X}\to X\to S$, and let $F$ be a coherent module on $X$, flat over $S$. Assume in addition that, for every $s\in S$, the module $F_s$ induced on the fiber $X_s$ has depth $>n$ at the points of that fiber, and that $t$ is $F$-regular. Under these conditions:
\begin{enumeratei}
\item
The canonical homomorphism
\[
\R^if_\ast(F) \to \R^i\widehat{f}_\ast(\widehat{F})
\]
is an isomorphism for $i<n$, and a monomorphism for $i=n$.
\item
The canonical homomorphism
\[
\R^i\widehat{f}_\ast(\widehat{F}) \to \varprojlim_m \R^if_\ast(F_m)
\]
is an isomorphism for $i\leq n$.
\end{enumeratei}
\end{theoreme}

\begin{proof}
We immediately reduce to the case where $S$ is affine, and then to proving the

\begin{corollaire} \label{XII.2.2}
Under the conditions of \Ref{XII.2.1}, assume in addition that $S$ is affine. Then:
\begin{enumeratei}
\item
The canonical homomorphism
\[
\H^i(X, F) \to \H^i(\widehat{X}, \widehat{F})
\]
is\pageoriginale an isomorphism for $i<n$, and a monomorphism for $i=n$.
\item
The canonical homomorphism
\[
\H^i(\widehat{X}, \widehat{F}) \to \varprojlim_m \H^i(X_m, F_m)
\]
is an isomorphism for $i\leq n$.
\end{enumeratei}
\end{corollaire}

After replacing $\OX(1)$ by a tensor power, and $t$ by a power of~$t$, we may assume that $\OX(1)$ is very ample relative to $S$. On the other hand, since $t$, and hence $t^m$, is $F$-regular, multiplication by $t^m$, regarded as a homomorphism from $F(-m)$ to $F$, is injective; therefore, for every $m\geq 0$, there is an exact sequence:
\steco
\begin{equation} \label{eq:XII.14}
0\to F(-m) \lto{t^m} F \to F_m \to0,
\end{equation}
and hence an exact cohomology sequence
\[
\H^i(X, F(-m)) \to \H^i(X, F) \to \H^i(X, F_m) \to \H^{i+1}(X, F(-m)).
\]
But by \Ref{XII.1.5}, we have $\H^i(X, F(-m))=0$ for $i\leq n$ and $m$ sufficiently large, which proves the

\begin{lemme} \label{XII.2.3}
For $m$ large, the canonical homomorphism
\[
\H^i(X, F) \to \H^i(X, F_m)
\]
is bijective if $i<n$, and injective if $i=n$.
\end{lemme}

\enlargethispage{-2\baselineskip}%
This shows that, for $i<n$, the inverse system $(\H^i(X_m, F_m))_{m\geq 0}$ is essentially constant, and \afortiori satisfies the Mittag-Leffler condition; hence (\sisi{given that}{in view of} $\widehat{F}= \varprojlim F_m$), we obtain (ii) from $\EGA 0_{\textup{III}}$ 13.3. On the other hand, (i) follows trivially from this, taking \Ref{XII.2.3} into account.

\refstepcounter{subsection}\label{XII.2.4}
\begin{enonce*}{Corollary 2.4\ndemark}\ndetext{See Corollary I.1.4 of the article by Mme Raynaud (Raynaud~M., {\og Th\'eor\`emes de Lefschetz en cohomologie des faisceaux coh\'erents et en cohomologie \'etale. Application au groupe fondamental\fg}, \emph{Ann. Sci. \'Ec. Norm. Sup. (4)} \textbf{7} (1974), p\ptbl 29--52).}
Let\pageoriginale $f:X\to S$ be a projective flat morphism, with $S$ locally Noetherian, let $\OX(1)$ be an invertible module on $X$, relatively ample over $S$, let $t$ be a section of this module that is $\OX$-regular, let $Y$ be the closed subprescheme of zeros of $t$, and let $\widehat{X}$ be the formal completion of $X$ along $Y$. Assume that, for every $s\in S$, $X_s$ has depth $\geq 1$ (\resp depth $\geq 2$) at its closed points. Then, for every open neighborhood $U$ of $Y$, the functor
\[
F \mto \widehat{F}
\]
from the category of locally free coherent modules on $U$ to the category of locally free coherent modules on $\widehat{X}$ is faithful (\resp fully faithful, \ie the Lefschetz condition ($\Lef$) of \Ref{X}~\Ref{X.2} holds).
\end{enonce*}

For two locally free Modules $F$ and $G$ on $U$, introduce the Module
\[
H=\SheafHom_{\Oo_U}(F, G)
\]
we are reduced to proving that the canonical homomorphism
\steco
\begin{equation} \label{eq:XII.15} {}\H^0(U, H) \to \H^0(\widehat{U}, \widehat{H})
\end{equation}
is injective (\resp bijective). Now the Modules $H_s$ have depth $\geq 1$ (\resp $\geq 2$) at the closed points of $X_s$, and hence we can apply \Ref{XII.2.1}, which gives the conclusion \Ref{XII.2.4} in the case where $U=X$. In the case of an arbitrary $U$, note that the question is local on $S$, so we may suppose that $S$ is affine. Then every coherent Module on $X$ is a quotient of a coherent locally free Module (since $\OX(1)$ is an invertible \sisi{absolutely}{relatively} ample Module on $X$). Since the dual Module $H'=\SheafHom(H, \Oo_U)$ extends to a coherent Module on $X$, which is therefore isomorphic to a cokernel of a homomorphism of locally free Modules on $X$, it follows by transposition that one can find a homomorphism
\[
u':{L'}^0 \to {L'}^1
\]
of\pageoriginale locally free Modules on $X$, inducing a homomorphism
\[
u:L^0 \to L^1
\]
of locally free Modules on $U$, such that there is an exact sequence
\[
0\to H \to L^0 \lto{u} L^1.
\]
Using the five lemma (which becomes the three lemma), and the left exactness of the functor $H^0$, we are reduced to proving that (\Ref{eq:XII.15}) is injective (\resp bijective) when $H$ is replaced there by $L^0$, $L^1$, which brings us back to the case where $H$ is induced by a locally free Module $H'$ on $X$. Moreover, in the non-parenthetical case this reduction is even unnecessary, since the kernel of (\Ref{eq:XII.15}) consists in any event of the sections of $H$ on~$U$ that vanish in some suitable open neighborhood $V$ of $Y$; but the restriction homomorphism $\H^0(U, H)\to\H^0(V, H)$ is injective, since $H$ has depth $\geq 1$ at the points of every closed subset $Z$ of $X$ not meeting $Y$ (\cf the lemma below). In the parenthetical case, we are reduced to proving that
\[
H^0(X, H')\to\H^0(U, H')
\]
is bijective, which follows from the fact that $H'$ has depth $\geq 2$ at all the points of a closed subset $Z=X- U$ of $X$ not meeting $Y$. It therefore remains only to prove the

\begin{lemme} \label{XII.2.5} Let $F$ be a coherent Module on $X$, flat over $S$, such that for every $s\in S$, $F_s$ has depth $\geq n$ at every closed point of $X_s$. Then, for every closed subset $Z$ of $X$ not meeting $Y$, $F$ has depth $\geq n$ at all the points of $Z$.
\end{lemme}

Indeed, for every $x\in X$, setting $s=f(x)$, we have
\steco
\begin{equation} \label{eq:XII.16} {}\prof(F_x)\geq \prof(F_{s, x}),
\end{equation}
as can be seen by lifting arbitrarily a maximal $F_{s, x}$-regular sequence of elements of\pageoriginale $\rr(\Oo_{X_s, x})$, which gives an $F_x$-regular sequence by \SGA 1 IV 5.7. Now if $x$ belongs to a $Z$ as in Lemma~\Ref{XII.2.5}, then $x$ is necessarily closed in $X_s$; in other words, $Z$ is \emph{finite} over $S$. Indeed, $Z$ (endowed with a structure induced from $X$) is projective over $S$ as a closed subprescheme of $X$, which is projective, and $Z$ is affine over $S$ as a closed subprescheme of $X- Y$, which is affine.
\skipqed
\end{proof}

\begin{remarque} \label{XII.2.6}
Suppose that for every $s\in S$, the section $t_s$ of $\Oo_{X_s}(1)$ induced by $t$ is $\Oo_{X_s}$-regular (which implies by \SGA 1 IV 5.7 that $t$ is $\OX$-regular). Then the hypotheses made are stable under extension of the base $S'\to S$ ($S'$ locally noetherian). Thus the conclusion remains valid after every base change.
\end{remarque}

\section[Lefschetz theory for a projective morphism: existence]{Lefschetz theory for a projective morphism: existence theorem}\label{XII.3}%

\setcounter{equation}{16}%
\refstepcounter{subsection}\label{XII.3.1}%
\begin{enonce*}{Theorem 3.1\ndemark}
\ndetext{For a version without the flatness hypothesis, see (Raynaud~M., {\og Th\'eor\`emes de Lefschetz en cohomologie des faisceaux coh\'erents et en cohomologie \'etale. Application au groupe fondamental\fg}, \emph{Ann. Sci. \'Ec. Norm. Sup. (4)} \textbf{7} (1974), p\ptbl 29--52, Theorem II.3.3).}
Let $f:X\to S$ be a projective morphism, with $S$ noetherian, let $\OX(1)$ be an invertible Module on $X$ that is ample relative to $S$, let $X_0$ be the closed subprescheme of zeros of a section $t$ of $\OX(1)$, let $\widehat{X}$ be the formal completion of $X$ along $X_0$, let $\formelF$ be a coherent Module on $\widehat{X}$, and let $F_0$ be the Module that it induces on $X_0$. Assume moreover:
\begin{enumeratea}
\item\label{XII.3.1a}
$\formelF$ is flat over $S$.
\item\label{XII.3.1b}
For every $s\in S$, the section $t_s$ induced by $t$ on $X_s$ is $\formelF_s$-regular (which implies that $F_0$ is also flat over $S$, \cf \textup{\SGA 1 IV 5.7}).
\item\label{XII.3.1c}
For every $s\in S$, $F_{0s}$ has depth $\geq 2$ at the closed points of $X_{0s}$.
\end{enumeratea}
Assume moreover that $S$ admits an ample invertible sheaf. Under these conditions, there exists a coherent Module $F$ on $X$, and an isomorphism from its formal completion $\widehat{F}$ to $\formelF$.
\end{enonce*}

This statement will follow from the following one:

\begin{corollaire} \label{XII.3.2}
Under conditions \Ref{XII.3.1a}, \Ref{XII.3.1b}, \Ref{XII.3.1c} above, the following hold:
\begin{enumeratei}
\item
The\pageoriginale Module $\widehat{f}_\ast(\formelF)$ on $S$ is coherent; hence, for every $n$, the Module $\widehat{f}_\ast(\formelF(n))$ on~$S$ is coherent.
\item
For $n$ large, the canonical homomorphism $\widehat{f}^\ast\widehat{f}_\ast(\formelF(n))\to \formelF(n)$ is surjective.
\end{enumeratei}
\end{corollaire}

Assume the corollary for the moment, and let us prove \Ref{XII.3.1}. By virtue of the last hypothesis made in \Ref{XII.3.1}, we may reduce to the case $X=\PP^r_S$, after replacing $\OX(1)$ and $t$ by a suitable power. I claim that we may further suppose that $t_s\neq 0$ for every~$s$. Otherwise, indeed, $\formelF_s=0$ by b), or equivalently, by Nakayama, $F_{0s}=0$, \ie $s$ does not belong to the image of $\Supp F_0$ under the morphism $f_0:X_0\to S$ induced by $f$. Now this image $S'$ is open by virtue of a), b), since $F_0$ is flat over $S$, and it is clear that it suffices to prove the conclusion of \Ref{XII.3.1} in the situation obtained by restricting over $S'$, since the coherent Module $F'$ on $X|_{S'}$ thereby obtained will be the restriction of a coherent Module $F$ on $X$, which will have the required property. We may therefore suppose that, in addition to hypotheses a), b), c), the following hypotheses also hold:
\begin{enumerate}
\item[\textup{a$'$)}] $\OX$ is flat over $S$.
\item[\textup{b$'$)}] For every $s\in S$, the section $t_s$ is $\Oo_{X_s}$-regular.
\item[\textup{c$'$)}] For every $s\in S$, $\Oo_{X_{0s}}$ has depth $\geq 2$ at the closed points of $X_{0s}$. (It~suffices to choose $X=\PP_S^r$ with $r\geq 3$, which is permissible.)
\end{enumerate}
Now \Ref{XII.3.2} implies that one can find an epimorphism
\steco
\begin{equation} \label{eq:XII.17}
\widehat{L} \to \formelF \to 0,
\end{equation}
where $L$ is a Module on $X$ of the form $\sisi{f_\ast(G)}{f^\ast(G)}(-n)$, with $G$ a coherent locally free Module on $S$: indeed, for $n$ large, it suffices to express the coherent Module $\widehat{f}_\ast(\formelF)$ on $S$ as a quotient of such a $G$. On the other hand, hypotheses a), b), c) on $f$, $t$ imply that $\widehat{L}$ satisfies the same conditions a), b), c) as $\formelF$. It follows easily that the same is true of the kernel of (\Ref{eq:XII.17}), to which the same argument may therefore be applied, so that $\formelF$ is represented as the cokernel of a homomorphism
\steco
\begin{equation} \label{eq:XII.18}
\widehat{L'} \to \widehat{L},
\end{equation}
where $L$, $L'$ are locally free Modules on $X$. Now, by virtue of a$'$), the second part of c$'$), and either \Ref{XII.2.1} or \Ref{XII.2.4}, the homomorphism (\Ref{eq:XII.18}) comes from a homomorphism $L'\to L$ of Modules on $X$. It now suffices to take for $F$ the cokernel of $L'\to L$, and we are done.

It remains to prove \Ref{XII.3.2}. This had been done in the
seminar by a somewhat laborious expedient, consisting in interpreting
everything in terms of cohomology on the punctured projective cone
of $X$ relative to $S$, in order to reduce to
Theorem \Ref{IX.2.1}. A more direct and more
satisfactory method (although substantially identical) now seems to me
to be the following. It consists in noting that in \Ref{IX},
\numero \Ref{IX.2}\refstepcounter{toto}\label{XII.14} (and with the
notation of that expos\'e), the hypothesis that the morphism
$f:\Xx\to\Xx'$ is adic is not used anywhere in the
proof of \Ref{IX.2.1}, via $\EGA 0_{\textup{III}}$ 13.7.7;
it suffices instead to suppose that $\Xx$ is also adic,
and to choose two ideals of definition $\Jj$ for $\Xx'$, $\Ii$
for $\Xx$, such that $\Jj\Oo_\Xx\subset\Ii$, and to define
$\Ss=\gr_\Jj(\Oo_{\Xx'})$, and consider $\gr_\Ii(\formelF)$. In
any event, \Ref{IX.2.1} can be applied directly to the
morphism $\widehat{f}:\widehat{X}\to S$ considered in the present
subsection, where one simply takes $\Jj=0$. Thus, to verify that
$\widehat{f}_\ast(\formelF)$ is coherent, it suffices, by virtue of
\loccit, to verify that $\R^if_{0\ast}(\gr_\Ii(\formelF))$ is
coherent on $S$ for $i=0, 1$; for this, note that by virtue of
a) and b), the Module in question is none other than
$\sisi{\coprod}{\bigoplus}_{m\geq 0}\R^if_{0\ast}(F_0(-m))$, which
is indeed coherent by virtue of hypothesis c) and
\Ref{XII.1.5}.

This proves \Ref{XII.3.2}~(i). For \Ref{XII.3.2}~(ii), we shall need the

\begin{lemme} \label{XII.3.3}
Under\pageoriginaled the conditions of \Ref{XII.3.1a}, \Ref{XII.3.1b}, \Ref{XII.3.1c} of \Ref{XII.3.1}, set
\[
G_m= \widehat{f}_\ast(\formelF(\cdot)_m) = \sisi{\coprod}{\tbigoplus}_n\widehat{f}_\ast(\formelF_m(n))
\]
Then the projective system $(G_m)$ satisfies the Mittag-Leffler condition.
\end{lemme}

We may suppose that $S$ is affine, with ring $A$. Let $\Ss$ be a finitely generated graded $A$-algebra in positive degrees, and let $t'\in\Ss_1$, such that $X$ embeds in $\Proj(\Ss)$, with $\OX(1)$ induced by \sisi{$\SheafProj(\Ss(1))$}{$\Oo(1)$} and the section $t$ the image of $t'$. Endow $\Ss$ with the $\Jj$-adic filtration, where $\Jj=t'\Ss$, and consider the projective system of the $\formelF(\cdot)_m$ in the category of abelian sheaves on $X_0$. We are again under the preliminary conditions of $\EGA 0_{\textup{III}}$ 13.7.7\sfootnote{Corrected as indicated in \Ref{IX} p\ptbl\sisi{11}{\pageref{pagerectif}.}}, and moreover $\H^i(X_0, \gr(\formelF(\cdot))\sisi{}{)}$ is a finite-type Module over $\gr_\Jj(\Ss)$ for $i=0, 1$. Indeed, since $t'$ is $\formelF$-regular, one sees at once that, as a Module over $(\Ss /t'\Ss)[T]$ (of which $\gr_\Jj(\Ss)$ is a quotient), the Module in question identifies with $\H^i(X_0, F_0(\cdot))\otimes_{\Ss /t'\Ss}(\Ss /t'\Ss)[T]$; but by \Ref{XII.1.5}, $\H^i(X_0, F_0(\cdot))$ is of finite type over $\Ss$, and hence over $\Ss /t'\Ss$, for $i=0, 1$, which proves our assertion. Consequently, we are under the conditions for applying $0_{\text{III}}$ 13.7.7 with $n=1$, which proves \Ref{XII.3.3}.

This point being established (and still supposing $S$ affine, as we may for proving \Ref{XII.3.2}~(ii)), let $m_0$ be such that $m\geq m_0$ implies $\Im(G_m\to G_0)= \Im(G_{m_0}\to G_0)$, so that both sides are also equal to $\Im(\varprojlim G_m\to G_0) = \Im(\widehat{f}_\ast(\formelF(\cdot))\to f_\ast\formelF_0(\cdot))$. Now note that for large $n$, $\formelF_{m_0}(n)$ is generated by its sections, and therefore $\formelF_0(n)$ is generated by sections that lift to $\formelF_{m_0}$, and hence (by the choice of~$m_0$) that lift to $\formelF$. Thus the sections of $\formelF$ generate $\formelF_0$, and hence also~$\formelF$ by Nakayama. This proves \Ref{XII.3.2}~(ii), and therefore \Ref{XII.3.1}.

\begin{corollaire} \label{XII.3.4}
Let $f:X\to S$ be a projective flat morphism, with $S$ locally noetherian, let $\OX(1)$ be an invertible Module on $X$, relatively ample over $S$, and let $t$ be a section of this Module such that for every $s\in S$, the induced section $t_s$ on\pageoriginale the fiber $X_s$ is $\Oo_{X_s}$-regular; let $X_0$ be the closed subprescheme of zeros of $t$, and let $\widehat{X}$ be the formal completion of $X$ along~$X_0$. Suppose that for every $s\in S$, $X_{0s}$ has depth $\geq 2$ at its closed points (\ie $X_s$ has depth $\geq 3$ at the closed points of $X_{0s}$), and $X_s$ has depth $\geq 2$ at its closed points. Under these conditions, the pair $(X, X_0)$ satisfies the effective Lefschetz condition ($\Leff$) \sisi{of \Ref{X}~\Ref{X.2}}{of paragraph~\Ref{X.2} of Expos\'e~\emph{\Ref{X}}}, \ie:
\begin{enumeratea}
\item
For every open neighborhood $U$ of $X_0$ in $X$, the functor
\[
F \sisi{\rightsquigarrow}{\mto} \widehat{F}
\]
from the category of coherent locally free Modules on $U$ to the category of coherent locally free Modules on $\widehat{X}$ is fully faithful.
\item
For every coherent locally free Module $\formelF$ on $\widehat{X}$, there exist an open neighborhood~$U$ of $X_0$ and a coherent locally free Module $F$ on $U$ such that $\formelF$ is isomorphic to~$\widehat{F}$.
\end{enumeratea}
\end{corollaire}

Indeed, a) has already been noted in \Ref{XII.2.4} under weaker conditions. For b), we apply \Ref{XII.3.1}, which gives the conclusion, at least if $S$ is noetherian and admits an absolutely ample invertible Module, in particular if $S$ is affine. Indeed, if $F$ is a coherent module on $X$ such that $\widehat{F}$ is isomorphic to $\formelF$, and hence locally free, it follows that $F$ is locally free on a neighborhood $U$ of $X_0$, and $\sisi{F|U}{F|_U}$ will satisfy the required condition. But now note that by \Ref{XII.2.5}, for such an $F$, its direct image under the immersion $U\to X$ is coherent, and is moreover independent of the chosen solution $(U, F)$ (taking account \sisi{\ignorespaces}{of the fact} that two solutions coincide in a neighborhood of $X_0$, by a)). More precisely, one can find a coherent Module $F$ on $X$ and an isomorphism $\widehat{F}\isomto \formelF$, such that $F$ has depth $\geq 2$ at all the points of $X$ that are closed in their fiber, and this determines $F$ up to unique isomorphism. By this uniqueness property, the solutions of the problem obtained over the affine open subsets of $S$ glue together, giving a coherent $F$ on all of $X$ and an isomorphism $\widehat{F}\simeq \formelF$. Restricting $F$ to the open subset $U$ of points at which it is free gives what was required.

Thanks\pageoriginale to \Ref{XII.2.4} and \Ref{XII.3.4}, in the setting of a projective algebraic scheme and a \og hyperplane section\fg thereof, one can apply the general facts established in Expos\'es \Ref{X} and \Ref{XI} concerning conditions (Lef) and (Leff). Thus:

\begin{corollaire} \label{XII.3.5} Let $X$ be a projective algebraic scheme equipped with an ample invertible module $\OX(1)$, let $t$ be a section of this module that is $\OX$-regular, and let $X_0$ be the closed subscheme of zeros of $t$. Assume that $X$ has depth $\geq 2$ at its closed points (\resp and depth $\geq 3$ at the closed points of $X_0$). Then $\pi_0(X_0)\to \pi_0(X)$ is bijective; in particular, $X$ is connected if and only if $X_0$ is connected, and, after choosing a geometric base point in $X_0$, $\pi_1(X_0)\to \pi_1(X)$ is surjective. More generally, for every open $U\supset X_0$, the homomorphism $\pi_1(X_0)\to \pi_1(U)$ is surjective (\resp the homomorphism $\pi_1(X_0)\to \varprojlim_U\pi_1(U)$ is bijective). In the parenthesized case, if one assumes in addition that the local ring at every closed point of $X$ not in $X_0$ is pure (\Ref{X.3.2}) (for example, regular, or merely a complete intersection), then $\pi_1(X_0)\to \pi_1(X)$ is an isomorphism.
\end{corollaire}

Apply \Ref{X.2.4} and \Ref{XII.3.4}. Note that, in the parenthesized case, the hypothesis that $X$ have depth $\geq 3$ at the closed points of $X_0$ implies that every irreducible component of $X$ having dimension $\neq 0$ has dimension $\geq 3$ (as one sees by noting that such a component necessarily meets $X_0$, and by considering a closed point of the intersection).

\begin{remarquestar}
When $X$ is normal and has dimension $\geq 2$ at all its points, it has depth $\geq 2$ at its closed points, and the non-parenthesized hypotheses of \Ref{XII.3.5} hold. In this case there is a more elementary proof of the surjectivity of $\pi_1(X_0)\to \pi_1(X)$ using Bertini's theorem (\cf \SGA 1 X.2.10). If one assumes in addition that $X_0$ is normal and that $X$ has dimension $\geq 3$ at all its points, then the parenthesized hypotheses of \Ref{XII.3.5} hold. In this case, \Ref{XII.3.5} was established by Grauert (indeed, thanks to the normality hypothesis, one can then dispense with the existence theorem \Ref{XII.3.1} by various devices)\pageoriginale. This proof of Grauert's was the starting point of the \og Lefschetz theory\fg that is the subject of the present seminar.
\end{remarquestar}

\begin{corollaire} \label{XII.3.6} Let $X$, $\OX(1)$, $t$, and $X_0$ be as in \Ref{XII.3.5}. Assume that $X$ has depth $\geq 2$ at its closed points, and that
\[
\H^i(X_0, \Oo_{X_0}(-n))=0
\]
for $n>0$ and for $i=1$ (\resp for $i=1$ and $i=2$), which, by \Ref{XII.1.4}, implies that $X_0$ has depth $\geq 2$ (\resp $\geq 3$) at its closed points, \ie that $X$ has depth $\geq 3$ (\resp $\geq 4$) at the closed points of $X_0$. Under these conditions, for every open neighborhood $U$ of $X_0$, $\Pic(U)\to\Pic(X_0)$ is injective; in particular, $\Pic(X)\to\Pic(X_0)$ is injective (\resp $\varprojlim_U \Pic(U)\to\Pic(X_0)$ is bijective). In the parenthesized case, if one assumes in addition that the local ring of $X$ at every closed point not in $X_0$ is parafactorial (\Ref{XI.3.1}) (for example, regular, or more generally a complete intersection), then $\Pic(X)\to\Pic(X_0)$ is bijective.
\end{corollaire}
Apply \Ref{XI}~\Ref{XI.3.12} and \Ref{XI.3.13}, noting that the parenthesized hypothesis implies that the irreducible components of $X$ having dimension $\neq 0$ have dimension $\geq 4$. In particular, applying this when $X$ is a global complete intersection of dimension $\geq 4$ in projective space, one obtains:

\begin{corollaire} \label{XII.3.7} Let $X$ be an algebraic scheme of dimension $\geq 3$ that is a complete intersection in the scheme $\PP_k^r$. Then $\Pic(X)$ is the free group generated by the class of the sheaf $\OX(1)$.
\end{corollaire}
Proceed by induction on the number of hypersurfaces whose intersection is $X$, applying \Ref{XII.3.6} and noting that, for a complete intersection $X$ of dimension $\geq 3$, one has $\H^i(X, \OX(n))=0$ for $i=1,2$ and every $n$.

\begin{remarque} \label{XII.3.8}
In\pageoriginale the case where $X$ is a nonsingular hypersurface, \Ref{XII.3.7} is due to Andreotti. The result \Ref{XII.3.7} can also be expressed (when $X$ is nonsingular) by saying that the homogeneous coordinate ring of $X$ is factorial, and in this form it is contained in \Ref{XI}~\Ref{XI.3.13} (ii). We also note that Serre had given a proof of \Ref{XII.3.7} in the nonsingular case by transcendental methods, using a specialization argument to reduce to characteristic 0, where the Lefschetz theorem is available in its classical form. Of course, the fact that the purely algebraic proof given here makes it possible to dispense with nonsingularity hypotheses in the statement of the Lefschetz theorem invites one to reconsider it in the classical case as well. \Cf the following expos\'e, which proposes conjectures in this direction.

In Corollaries \Ref{XII.3.5} and \Ref{XII.3.7}, we have worked over a base field, whereas the key theorems \Ref{XII.2.4} and \Ref{XII.3.4} are valid over an arbitrary base. To generalize Corollaries \Ref{XII.3.5} and \Ref{XII.3.6} to a general $S$, we must give workable criteria for a point of $X$ (flat over $S$) to have a \og pure\fg \resp parafactorial local ring. This will be the subject of the following n\textsuperscript{o}.
\end{remarque}

\section{Formal completion and normal flatness} \label{XII.4}

\setcounter{equation}{18}

\begin{theoreme} \label{XII.4.1} Let $X$ be a locally Noetherian prescheme, locally embeddable in a regular scheme, $Y$ a closed subset of $X$, $U=X- Y$, $X_0$ a closed subprescheme of $X$ defined by an ideal $\Jj$, $\widehat{X}$ the formal completion of $X$ along $X_0$, $U_0$ the trace of $X_0$ on $U$, $\widehat{U}$ the formal completion of $U$ along $U_0$, $i:U\to X$ and $\widehat{i}:\widehat{U}\to \widehat{X}$ the canonical immersions, and $n$ an integer. Assume that
\begin{enumeratea}
\item
$X$ is normally flat along $X_0$ at the points of $Y\cap X_0$, \ie at these points the \sisi{m}{M}odules $\Jj^n/\Jj^{n+1}$ over $X_0$ are flat, \ie \sisi{\ignorespaces}{locally} free.
\item
For every $x\in Y\cap X_0$, one has $\prof \Oo_{X_0, x}\geq n+2$.
\end{enumeratea}
Under\pageoriginale these conditions, the following assertions hold:
\begin{enumerate}
\sisi{\item[$1^\circ)$]}{\item} Let $F$ be a coherent Module on $U$, and suppose that:
\begin{enumerate}
\item[\textup{c)}] For every $x\in Y- Y\cap X_0$, one has $\prof \Oo_{X, x}\geq n+2$.

\item[\textup{d)}] $F$ is free at the points of $U_0$, and has depth $\geq n+1$ at all the points of $U$ where it is not free.
\end{enumerate}
Then the graded Module
\[
\sisi{\coprod}{\tbigoplus}_{m\geq 0}\R^pi_\ast(\Jj^mF)
\]
over $\sisi{\coprod}{\bigoplus}_{m\geq 0}\Jj^m$ is of finite type for $p\leq n$. \sisi{\item[$2^\circ$]}{\item} Let $\formelF$ be a coherent \sisi{}{locally free }Module on $\widehat{U}$; then the graded Module
\[
\sisi{\coprod}{\tbigoplus}_{m\geq 0}\R^pi_\ast(\Jj^m\formelF/\Jj^{m+1}\formelF)
\]
over $\sisi{\coprod}{\bigoplus}_{m\geq 0}\Jj^m/\Jj^{m+1}=\gr_{\Jj}(\OX)$ is of finite type for $p\leq n$.
\end{enumerate}
\end{theoreme}

\begin{proof}
1) Let $X'\!=\!\Spec(\sisi{\coprod}{\bigoplus}_{m\geq 0}\Jj^m)$; the base change \hbox{$f:X'\!\to\! X$} then defines $U'=X'- Y'$, $X_0'$, $U_0'=X_0'\cap U'$, and immersions $i':U'\to X'$, \hbox{$i_0':U_0'\to X_0'$}. Thus we have a Cartesian square
\[
\xymatrix{
X'\ar[d]_f&\ar[l]_-{i'} U'\ar[d]^g \\
X&\ar[l]_-{i} U
}
\]
and we have
\steco
\begin{equation} \label{eq:XII.19}
\sisi{\coprod}{\tbigoplus}_{m\geq 0}\R^pi_\ast(\Jj^mF) = \R^pi_\ast \Big(\sisi{\coprod}{\tbigoplus}_{m\geq 0} \Jj^mF \Big) =\R^pi_\ast (g_\ast(F')),
\end{equation}
where\pageoriginale $F'=g^\ast F$, so that there is indeed a canonical isomorphism
\[
g_\ast(F') \isomto \sisi{\coprod}{\tbigoplus}_{m\geq 0}\Jj^mF,
\]
since this is true at the points of $U_0$, because $F$ is free there by d), and also at the points of $\sisi{U_0}{U-U_0}$, because there one has $\Jj^m=\Oo_U$, (so that in both cases, $\Jj^m\otimes_{\Oo_U}F\to \Jj^mF$ is an isomorphism).

On the other hand, since $f$, and hence $g$, are affine, one has
\steco
\begin{equation} \label{eq:XII.20} {}\R^pi_\ast(g_\ast(F')) = \R^p(ig)_\ast(F') = \R^p(fi')_\ast(F') = f_\ast(R^pi'_\ast(F'))
\end{equation}
and therefore, comparing (\Ref{eq:XII.19}) and (\Ref{eq:XII.20}), we see that \hyperref[XII.4.1]{assertion 1}\sisi{${}^\circ$}{} is equivalent to the following: $R^pi'_\ast(F')$ is a Module of finite type, \ie coherent, on $X'$, for every $p\leq n$. Now, since $X$ is locally embeddable in a regular scheme, the same is true of $X'$, which is of finite type over $X$, and we may apply the coherence criterion \sisi{\ignorespaces}{\Ref{VIII}}~\Ref{VIII.2.3} to a coherent extension $F^{''}$ of $F'$: we wish to express that $\SheafH^p\sisi{_Y}{_{Y'}}(F^{''})$ is coherent for $p\leq n+1$, and this is also equivalent to saying that, for every $x'\in U'$ such that
\begin{equation*} \label{eq:XII.20bis} \tag{20 bis} {}\codim (\overline{x'}\cap Y', \overline{x'})=1,
\end{equation*}
one has
\steco
\begin{equation} \label{eq:XII.21} {} \prof F'_{x'}\geq n+1.
\end{equation}
Now this condition is satisfied at the points $x'$ where $F'$ is not free, because for such an $x'$ one has $x'\notin U_0'$ by d), and therefore $g$ is an isomorphism there; again by d), $F$ has depth $\geq n+1$ at $g(x')$, and hence $F'$ has depth $\geq n+1$ at $x'$. It is therefore enough to verify condition (\Ref{eq:XII.21}) at those $x'\in U'$ satisfying \sisi{(\Ref{eq:XII.21})}{(\Ref{eq:XII.20bis})} at which $F'$ is free. For this, it is enough to prove that one has\pageoriginale
\begin{equation*} \label{eq:XII.21bis} \tag{21 bis}
\prof \Oo_{X', x'}\geq n+1
\end{equation*}
\enlargethispage{\baselineskip}%
at these points; \afortiori it is enough to establish this relation at all the points~$x$ of $U'$ satisfying (\Ref{eq:XII.20}\sisi{}{ bis}). Now, \sisi{by the criterion cited}{again by criterion~\Ref{VIII.2.3}} in Expos\'e \Ref{VIII}, this is equivalent to the assertion that the Modules
\[
\SheafH^p_{Y'}(\Oo_{X'}) \hspace{5mm} \text{for} \; p\leq n+1
\]
are coherent. In fact, we shall prove that they are even zero, or, equivalently by Expos\'e~\Ref{III}, that one has
\steco
\begin{equation} \label{eq:XII.22} {} \prof \Oo_{X', x'}\geq n+2 \hspace{5mm} \text{for every }\; x'\in Y'.
\end{equation}
For this, we distinguish two cases. If $x'\notin X'_0$, then $f$ is an immersion at $x'$, and we must verify that $\sisi{F}{\Oo_X}$ has depth \sisi{at}{$\geq n+2$ at } the image $x=f(x')$, which is precisely condition c). If, on the other hand, $x'\in X'_0$, \ie $x=f(x')\in X_0$ and hence $x\in Y\cap X_0$, we apply conditions a) and b) by means of the

\begin{lemme} \label{XII.4.2}
Let $X$ be a locally Noetherian prescheme, $X_0$ a closed subprescheme of $X$ defined by an ideal $\Jj$, $X'=\Spec(\sisi{\coprod}{\bigoplus}_{m\geq 0} \Jj^m)$, $X_0'=\Spec(\sisi{\coprod}{\bigoplus}_{m\geq 0} \Jj^m/\Jj^{m+1}) = X'\times_XX_0$, and $x$ a point of $X_0$ at which $X$ is normally flat along $X_0$, \ie such that $\gr_{\Jj}(\OX)= \sisi{\coprod}{\bigoplus}_{m\geq 0}\Jj^m/\Jj^{m+1}$ is flat there as a Module over $X_0$. Then, for every sequence of elements $f_i$ ($1\leq i\leq m$) of $\Oo_{X, x}$ whose images in $\Oo_{X_0, x}$ form an $\Oo_{X_0, x}$-regular sequence, and every $x'\in X'$ above $x$, the images of the $f_i$ in $\Oo_{X', x'}$ (respectively in $\Oo_{X'_0, x'}$) likewise form an $\Oo_{X', x'}$-regular sequence (respectively an $\Oo_{X'_0, x'}$-regular sequence); in particular, one has
\steco
\begin{equation} \label{eq:XII.23}
\prof \Oo_{X', x'} \geq \prof \Oo_{X_0, x}, \quad \prof \Oo_{X'_0, x'} \geq \prof \Oo_{X_0, x}.
\end{equation}
\end{lemme}

To\pageoriginale prove it, we may suppose that $X$ is local with closed point $x$, and hence affine with ring $A=\Oo_{X, x}$, with $\Jj$ defined by the Ideal $J$; it is enough to prove that, for every sequence $f_i$ ($1\leq i\leq m$) of elements of $A$ whose images in $A/J$ form an $A/J$-regular sequence, the $f_i$ likewise form a $(\sisi{\coprod}{\bigoplus}_{m\geq 0} J^m)$-regular sequence and a $(\sisi{\coprod}{\bigoplus}_{m\geq 0} J^m/J^{m+1})$-regular sequence, \ie for every $m$, it forms a $J^m$-regular sequence and a $J^m/J^{m+1}$-regular sequence. The second assertion is trivial, since $J^m/J^{m+1}$ is a free module over $A/J$. The first follows by considering the $J$-adic filtration of $J^m$ and noting that the sequence of the~$f_i$ is regular on the graded module associated with $J^m$ for that filtration.

This proves \Ref{XII.4.2}, and consequently \Ref{XII.4.1}~1\sisi{\noo}{}).

Let us prove \Ref{XII.4.1} 2\sisi{\noo}{}). For this, use the Cartesian square
\[
\xymatrix{
X'_0\ar[d]_{f_0}&\ar[l]_-{i'_0} U'_0\ar[d]^{g_0} \\
X_0&\ar[l]_-{i_0} U_0
}
\]
and, proceeding as at the beginning of the proof of 1\sisi{\noo}{}), we find that
\steco
\begin{equation} \label{eq:XII.24}
\sisi{\coprod}{\tbigoplus}_{m\geq 0} \R^pi_\ast(\Jj^m\formelF/\Jj^{m+1}\formelF) \simeq {f_0}_\ast (\R^p{\sisi{i}{i'}_0}_\ast(\formelF_0')),
\end{equation}
where $\formelF_0=\formelF/\Jj\formelF$ and $\formelF_0'=g_0^\ast(\formelF_0)$, (using the fact that $\formelF$ is locally free). Thus the conclusion of 2\sisi{\noo}{}) is equivalent to saying that, for $p\leq n$, $\R^p{i'_0}_\ast(\formelF_0')$ is a coherent Module.

\enlargethispage{\baselineskip}%
Here again, taking into account that $\formelF_0'$ is locally free, criterion \Ref{VIII}~\Ref{VIII.2.3} allows us to reduce to proving the same assertion after replacing $\formelF_0'$ by $\Oo_{U'_0}$, \ie to proving that the Modules\pageoriginale
\[
\SheafH^p_{Y_0'}(\Oo_{X_0'}) \hspace{5mm} \text{for} \; p\leq n+1 \hspace{5mm} \text{(where} \; Y'_0=Y'\cap X_0'=X_0'- U'_0 \text{)}
\]
are coherent. Once again, we prove that they are in fact zero, i.e., that one has
\steco
\begin{equation} \label{eq:XII.25} {}\prof \Oo_{X'_0, x'} \geq n+2 \hspace{5mm} \text{for every} \; x'\in Y'_0.
\end{equation}
This does indeed follow from conditions a) and b), in view of \Ref{XII.4.2}. This completes the proof of \Ref{XII.4.1}.
\skipqed
\end{proof}

\begin{remarque} \label{XII.4.3}
It follows immediately, by descent, that the hypothesis that $X$ be locally embeddable in a regular scheme can be replaced by the following weaker hypothesis: there exists a faithfully flat and quasi-compact morphism $\overline{X}\to X$ such that $\overline{X}$ is locally embeddable in a regular scheme.
\end{remarque}

Theorem \Ref{XII.4.1} enables us to apply the results of Expos\'e \Ref{IX} (the comparison and existence theorems). We shall be interested in particular in the

\begin{corollaire} \label{XII.4.4}
Assume that conditions a), b), and c) of Theorem \Ref{XII.4.1} hold, with $n=1$ and $X=\Spec(A)$, where $A$ is separated and complete for the $J$-adic topology. Then
\begin{enumerate}
\sisi{\item[\textup{1\noo})]}{\item} The functor $F\mto \widehat{F}$ from the category of locally free coherent modules on~$U$ to the category of locally free coherent modules on $\widehat{U}$ is fully faithful. \sisi{\item[\textup{2\noo})]}{\item} For every locally free coherent module $\formelF$ on $\widehat{U}$, there exists a coherent module $F$ on $U$ and an isomorphism $\widehat{F}\isomto \formelF$.
\end{enumerate}
In particular, if, for every $x\in U$ whose closure in $U$ does not meet $U_0$, \ie such that $\overline{x}\cap X_0\subset Y$, one has $\prof \Oo_{U, x}\geq 2$, then the pair $(U, U_0)$ satisfies the effective Lefschetz condition ($\Leff$) of Expos\'e \Ref{X}.
\end{corollaire}
\noindent
(For the last assertion, one proceeds as in \Ref{X}~\Ref{X.2.1}.)

\pagebreak[2]
Special case of \Ref{XII.4.4}:

\begin{corollaire} \label{XII.4.5}
Let\pageoriginaled $A$ be a Noetherian ring, let $J$ be an ideal of $A$ contained in the radical, and let $A_0=A/J$. Assume
\begin{enumeratei}
\item
$\prof A_0\geq 3$.
\item
$\gr_J(A)$ is a free $A_0$-module.
\item
$A$ is complete for the $J$-adic topology.
\end{enumeratei}

Let $X=\Spec(A)$, $X_0=\Spec(A_0)=V(J)$, let $a$ be the closed point of $X$, let $U=X-\{a\}$, $U_0=X_0-\{a\}$, and let $\widehat{U}$ be the formal completion of $U$ along $U_0$. Then the functor $F\mto \widehat{F}$ from the category of locally free coherent modules on $U$ to the category of locally free coherent modules on $\widehat{U}$ is fully faithful. Moreover, for every locally free coherent module $\formelF$ on $\widehat{U}$, there exists a coherent module (not necessarily locally free!) $F$ on $U$, and an isomorphism $\widehat{F}\simeq\formelF$.
\end{corollaire}

Note that, thanks to \Ref{XII.4.3}, we did not have to assume that $A$ is a quotient of a regular ring, because the completion of $A$ for the $\rr(A)$-adic topology satisfies this condition in any case.

Proceeding as in Expos\'es \Ref{X} and \Ref{XI}, we deduce from \Ref{XII.4.5}

\begin{corollaire} \label{XII.4.6}
Under the conditions of \Ref{XII.4.5}, one has the following:
\begin{enumeratea}
\item
$U$ and $U_0$ are connected (\Ref{III}~\Ref{III.3.1}).

Choosing a geometric base point in $U_0$, the homomorphism
\[
\pi_1(U_0)\to \pi_1(U)
\]
is surjective.
\item
The homomorphism
\[
\Pic(U)\to \Pic(U_0)
\]
is injective.
\end{enumeratea}
\end{corollaire}

To\pageoriginale prove b), in view of \Ref{XII.4.5}, this amounts to checking that every isomorphism $L_0'\isomto L_0$ lifts to an isomorphism $\widehat{L'}\isomto \widehat{L}$. For this, one lifts step by step to isomorphisms $L_n'\isomto L_n$; the obstructions lie in $\H^1(U_0, \Jj^n/\Jj^{n+1})$, and these modules are zero because $J^n/J^{n+1}$ is free and $\prof A_0\geq 3$.

We are now in a position to prove the
\enlargethispage{-1\baselineskip}%
\begin{theoreme} \label{XII.4.7}
Let $A$ be a noetherian local ring, let $J$ be an ideal of $A$ contained in its radical, and let $A_0=A/J$. Suppose that
\begin{enumeratei}
\item
$\prof A_0\geq 3$.
\item
$\gr_J(A)$ is a free module over $A_0$.
\end{enumeratei}
\noindent
Then, if $A_0$ is \og pure\fg (\Ref{X}~\Ref{X.3.1}) (\resp parafactorial (\Ref{XI}~\Ref{XI.3.1})), the same holds for~$A$.
\end{theoreme}

\begin{proof}
By descent, we may suppose that we also have
\begin{enumeratei}\setcounter{enumi}{2}
\item
$A$ is complete for the $J$-adic topology.
\end{enumeratei}

Indeed, by (i) and (ii), we have $\prof(A)\geq 3$, and hence $\prof(\widehat{A})\geq 3$, where $\widehat{A}$ is the completion of $A$ for the $J$-adic topology, and we apply \Ref{X}~\Ref{X.3.6} and \Ref{XI}~\Ref{XI.3.6}. We are therefore under the conditions of \Ref{XII.4.5}. Since $\prof(A)\geq 3\geq 2$, to say that $A$ is parafactorial means simply that $\Pic(U)=0$, and by \Ref{XII.4.6}~b), it is enough for this that $\Pic(U_0)=0$, \ie that $A_0$ be parafactorial. To prove that $A$ is \og pure\fg if $A_0$ is, we must prove that if $V$ is an \'etale covering of $U$, defined by an algebra $\Bb$ on~$U$, then $\H^0(U, \Bb)$ is a finite \'etale algebra over $A$. Now, since $A_0$ is pure, the same holds for the $A_n$ (which differ from it only by nilpotent elements), and therefore, for every $n$, $\sisi{}{B_n=}\H^0(U, \Bb_n)$ is an \'etale algebra over $A/J^{n+1}$; these algebras of course fit together, so that $\varprojlim B_n$ is an \'etale algebra over $A$. By \Ref{XII.4.5}, this algebra is none other than $\H^0(U, \Bb)$, which proves our assertion.
\skipqed
\end{proof}

\begin{corollaire} \label{XII.4.8}
Let\pageoriginale $f:X\to Y$ be a flat morphism of locally noetherian preschemes, let $x\in X$, and let $y=f(x)$. Suppose that $\Oo_{X_y, x}$ is a \og pure\fg (\resp parafactorial) local ring of depth $\geq 3$. Then the same holds for $\Oo_{X, x}$.
\end{corollaire}

This is the result of the kind promised at the end of the preceding \sisi{N\textsuperscript{o}}{\numero}, for generalizing Corollaries \Ref{XII.3.5} and following. Thus, using \Ref{XII.3.4}, we obtain the

\begin{corollaire} \label{XII.4.9}
Let $f:X\to S$ be a projective flat morphism, with $S$ locally Noetherian, let $\OX(1)$ be an invertible module on $X$, ample relative to $S$, let $t$ be a section of $\OX(1)$ such that, for every $s\in S$, the section $t_s$ induced on $X_s$ is $\Oo_{X_s}$-regular, let $X_0$ be the closed subscheme of zeros of $t$, and let $X_m$ be the closed sub-\sisi{pre}{}scheme of zeros of $t^{m+1}$. Assume that, for every $s\in S$, $X_s$ has depth $\geq 3$ at all its closed points. Then:
\begin{enumeratea}
\item
If the local rings at the closed points of $X_s-X_{0,s}$ ($s\in S$) are \og pure\fg\,---for example, if they are complete intersections---then the functor $X'\mto X_0'=X'\times_X X_0$ from the category of finite \`etale coverings of $X$ to the category of finite \`etale coverings of $X_0$ is an equivalence of categories; in particular, after choosing a geometric base point in $X_0$, the homomorphism
\[
\pi_1(X_0)\to \pi_1(X)
\]
is an isomorphism\refstepcounter{toto}\nde{\label{FH}Let us point out
the striking connectedness result obtained subsequently by Fulton and
Hansen, in the case where $S=\Spec(k)$ ($k$ an algebraically closed
field). Let $g:X\to {\PP^m_k}\times{\PP^m_k}$ be such that $\dim
g(X)>m$; then the inverse image of the diagonal is connected. Among
other things, this makes it possible to generalize
Corollary~\Ref{XII.4.9} when $f$ is the structural morphism
of projective space $\PP^m_k$ over $S=\Spec(k)$: more precisely, an
irreducible subvariety $X$ of ${\PP^m_k}$ of dimension $>m/2$ has
trivial fundamental group! (\cf Fulton~W. \& Hansen J., {\og
A connectedness theorem for projective varieties, with
applications to intersections and singularities of mappings\fg},
\emph{Ann. of Math. (2)} \textbf{110} (1979), \numero 1, p\ptbl
159-166). For generalizations to Grassmannians or
abelian varieties, see {Debarre~O.}, {\og Th\'eor\`emes de
connexit\'e pour les produits d'espaces projectifs et les
grassmanniennes\fg}, \emph{Amer. J.~Math.} \textbf{118} (1996),
\numero 6, p\ptbl 1347-1367 and {\og Th\'eor\`emes de connexit\'e et
vari\'et\'es ab\'eliennes\fg}, \emph{Amer. J.~Math.} \textbf{117}
(1995), \numero 3, p\ptbl 787--805. The triviality result for the
fundamental group of $X$ above was obtained independently
by Faltings, who moreover proves that
$\text{Pic}(X)$ has no torsion prime to the characteristic
of $k$, using methods of algebraization of formal vector bundles,
more in the spirit of Grothendieck's techniques, \cf
(Faltings~G., \og{Algebraization of some formal vector bundles\fg}
\emph{Ann. of Math. (2)} \textbf{110} (1979), \numero 3, p\ptbl
501--514).}

\item
If the local rings at the closed points of $X_s-X_{0,s}$ ($s\in S$) are \og parafactorial\fg---for example, regular, or complete intersections of dimension $\geq 4$---then, for every integer $m$ such that $\R^if_{0\ast}(\Oo_{X_0}(-n))=0$ for $n>m$ and $i=1,2$, the map $\Pic(X)\to \Pic(X_m)$ is bijective.

Moreover, if $S$ is Noetherian and the $X_{0,s}$ have depth $\geq 3$ at their closed points, such $m$ exist (\cf \Ref{XII.1.5}).
\end{enumeratea}
\end{corollaire}

\begin{remarque} \label{XII.4.10}
Under\pageoriginale the conditions of the last assertion of \Ref{XII.4.9} b), we saw in \Ref{XII.1.5} that there exists an $m$ such that $n>m$ even implies $\H^i(X_0, \Oo_{X_0}(-n))=0$ for $i=1,2$, and indeed for $i\leq 2$\sisi{)}{}. This condition is stronger than $\R^if_{\ast}(\Oo_{X_0}(-n))=0$ for $i=1,2$, and it has the further advantage of being stable under base change. The same is true of the depth hypotheses made in \Ref{XII.4.9}, and also of a hypothesis of the type \og the $X_s$ are locally complete intersections\fg. It follows under these conditions that \Ref{XII.4.9} b) also implies that the morphism of functors
\[
\SheafPic_{X/S}\to \SheafPic_{X_n/S}
\]
in $\catSch_{/S}$ is an isomorphism, and hence so is the morphism of relative Picard schemes, when these exist:
\[
\Pic_{X/S}\to \Pic_{X_n/S}.
\]
Even when $S$ is the spectrum of an algebraically closed field, this statement is considerably more precise than the statement that merely says that $\Pic(X)\to \Pic(X_n)$ is bijective.

One may ask whether one can always take $n=0$ in the preceding conclusions (thus assuming that the $X_{0,s}$ have depth $\geq 3$ at their closed points). When $X_0$ is smooth over $S$ and the residual characteristics of $S$ are zero, this is indeed so, by the Kodaira \og vanishing theorem\fg (proved by transcendental methods, using a K\"ahler metric), which implies that, for every connected smooth projective scheme of dimension $n$ over a field $k$ of characteristic zero and every ample invertible module $L$ on $X$, one has $\H^i(X,L^{-1})=0$ for $i\neq n$. It is not known\nde{As Raynaud observed, the decomposition result for the de Rham complex of Deligne and Illusie readily implies the vanishing of the group $H^i(X,L^{-1})$ (with $L$ ample on $X$, projective and smooth over $k$ of characteristic $p>0$) for $i<\inf(p,\dim(X))$, provided that $X$ can be lifted to a flat scheme over $W_2(k)$ (\cf Deligne~P. \& Illusie L., {\og Rel\`evements modulo~$p^2$ et d\'ecomposition du complexe de de Rham\fg}, \emph{Invent. Math.} \textbf{89} (1987), \numero 2, p\ptbl 247--270); this gives a purely algebraic proof of Kodaira's result for projective varieties in characteristic zero. If $X$ cannot be lifted, it is well known that the \og vanishing theorem\fg is false; \cf the example in (Raynaud~M., {\og Contre-exemple au "vanishing theorem" en caract\'eristique $p>0$\fg}, in \emph{C.P.~Ramanujam---a tribute}, Tata Inst. Fund. Res. Studies in Math., vol.~8, Springer, Berlin-New York, 1978, p\ptbl 273--278); see also the very elegant examples in (Haboush~W. \& Lauritzen~N., {\og Varieties of unseparated flags\fg}, in \emph{Linear algebraic groups and their representations (Los Angeles, CA, 1992)}, Contemp. Math., vol.~153, American Mathematical Society, Providence, RI, 1993, p\ptbl 35--57), simplified in (Lauritzen~N. \& Rao~A.P., {\og Elementary counterexamples to Kodaira vanishing in prime characteristic\fg}, \emph{Proc. Indian Acad. Sci. Math. Sci.} \textbf{107} (1997), \numero 1, p\ptbl 21--25). On the other hand, I know of no example in which the map $\Pic(X_{n+1})\to\Pic(X_n)$ is not surjective for $n>1$ in positive characteristic, where $X_n$ denotes a thickened hyperplane section of a smooth projective $X$ as above.} at present whether this theorem can be replaced by a generalization in characteristic $p>0$, and whether the smoothness hypothesis can be replaced by a hypothesis of a more general nature (involving depth, or of \og complete intersection\fg type, and so on).
\end{remarque}

\section[Universal finiteness conditions]
{Universal finiteness conditions for a nonproper morphism} \label{XII.5} Recall\pageoriginale for reference the \setcounter{equation}{25}

\begin{proposition} \label{XII.5.1}
Let $f:X\to S$ be a proper morphism of preschemes, with $S$ locally noetherian, let $U$ be an open subset of $X$, let $g:U\to X$ be the canonical immersion, let $h=fg:U\to S$, and let $F$ be a Module on $U$. Suppose that the Modules $\R^ig_\ast(F)$ are coherent for $i\leq n$ (a hypothesis local in nature on $X$, which in practice is checked using the criterion \Ref{VIII}~\sisi{\Ref{VII.2.3}}{\Ref{VIII.2.3}}). Then $\R^ih_\ast(F)$ is coherent for $i\leq n$.
\end{proposition}

This follows at once from the Leray spectral sequence
\steco
\begin{equation} \label{eq:XII.26} {} \E_2^{p, q}=\R^pf_\ast(\R^qg_\ast(F)) \To \R^\ast h_\ast(F),
\end{equation}
and from the fact that the higher direct images under $f$ of a coherent Module on $X$ are coherent (\EGA III 3.2.1).

\begin{proposition} \label{XII.5.2}
Let $S$ be a locally noetherian prescheme, let $\Ss$ be a quasi-coherent graded Algebra of finite type over $S$, generated by $\Ss_1$, let $X$ be a subprescheme of $\Proj(\Ss)$, let $\OX(1)$ be the invertible Module on $X$ induced by $\Proj(\Ss(1))$ and very ample relative to $S$, let $U$ be an open subset of $X$, let $g:U\to X$ be the canonical immersion, let $h=fg:U\to S$, and let $F$ be a quasi-coherent Module on $U$, giving twisted Modules $F(m)=F\otimes \OX(m)$\sisi{)}{} ($m\in\ZZ$); let $n$ be an integer and $m_0$ an integer. The following conditions are equivalent:
\begin{enumeratei}
\item
$\R^ig_\ast(F)$ is coherent for $i\leq n$.
\item
$\sisi{\coprod}{\bigoplus}_{m>m_0}\R^ih_\ast(F(m))$ is a finite-type $\Ss$-Module for $i\leq n$.
\end{enumeratei}
\end{proposition}

\begin{proof}
Replacing $F$ by $F(m)$ in the spectral sequence above, we
obtain a spectral sequence of graded $\Ss$-Modules\pageoriginale
\[
\E_2^{p, q}= \sisi{\coprod}{\tbigoplus}_{m\geq m_0} \R^pf_\ast(\R^qg_\ast(F(m))\sisi{}{)} \To \sisi{\coprod}{\tbigoplus}_{m\geq m_0} \R^\ast h_\ast(F(m)).
\]
Since we have
\[
\R^qg_\ast(F(m)) \simeq \R^qg_\ast(F) (m),
\]
we see that if the $\R^ig_\ast(F)$ are coherent, then $\E_2^{p, q}$ is of finite type over $\Ss$ for $q\leq n$, by part~a) of Lemma~\Ref{XII.5.3} below, which implies that the abutment is of finite type over $\Ss$ in degree $i\leq n$. This proves (i) \ALORS (ii). Moreover, reasoning in the abelian category of graded $\Ss$-Modules modulo the thick subcategory $\mathcal{C}$ of those that are quasi-coherent of finite type, we obtain from the preceding spectral sequence
\[
\sisi{\coprod}{\tbigoplus}_{m\geq m_0} \R^{n+1} h_\ast(F(m)) \simeq \sisi{\coprod}{\tbigoplus}_{m\geq m_0} f_\ast (\R^{n+1} g_\ast(F)(m)) \mod \mathcal{C},
\]
which proves that if the left-hand side is a finite-type $\Ss$-Module, then $\R^{n+1} g_\ast(F)$ is coherent, by part~b) of Lemma~\Ref{XII.5.3}. This proves the implication (ii) \ALORS (i) by induction on $n$. It remains to prove:

\begin{lemme} \label{XII.5.3} Let $S$, $\Ss$, $X$, $f$ be as in \Ref{XII.5.2}, let $G$ be a quasi-coherent Module on $X$, and let $m_0$ be an integer. Then
\begin{enumeratea}
\item
If $G$ is coherent, then for every integer $i$, the graded Module
\[
\sisi{\coprod}{\tbigoplus}_{m\geq m_0} \R^i f_\ast(G(m))
\]
over $\Ss$ is of finite type.
\item
Conversely, suppose that the Module $\sisi{\coprod}{\bigoplus}_{m\geq m_0} \sisi{\R^i }{}f_\ast(G(m))$ over $\Ss$ is of finite type; then $G$ is coherent.
\end{enumeratea}
\end{lemme}

\begin{proof}[Proof of \Ref{XII.5.3}]
For a), the case $i=0$ is given in \EGA III 2.3.2, and the case $i>0$ in \EGA III 2.2.1 (i)(ii), which says that the $ \R^i f_\ast G(m)$ are coherent\pageoriginale and zero for large $m$ (if we suppose $S$ noetherian, as we may). For b), note that $G$ is isomorphic to $\SheafProj (\sisi{\coprod}{\bigoplus}_{m\geq m_0}f_\ast(G(m))\sisi{}{)}$ (\EGA II 3.4.4 and 3.4.2), which proves that $G$ is coherent if $\sisi{\coprod}{\bigoplus}_{m\geq m_0}f_\ast(G(m))$ is of finite type over $\Ss$, by \loccit 3.4.4.
\skipqed
\end{proof}
\skipqed
\end{proof}

\begin{corollaire}[{{\normalfont$S$ noetherian}}] \label{XII.5.4}
Suppose that $\R^ig_\ast(F)$ is coherent for $i\leq n$. Then, for $i\leq n+1$ and large $m$, there is a canonical isomorphism:
\[
\R^ih_\ast(F(m)) \simeq f_\ast(\R^ig_\ast(\sisi{G}{F})(m)).
\]
\end{corollaire}

Indeed, the spectral sequence (\Ref{eq:XII.26}) for $F(m)$ then degenerates in degree $\leq \sisi{n}{n+1}$, by \EGA III 2.2.1~(ii), which gives the result at once (and in fact recovers the implication (ii) \ALORS (i) of~\hyperref[XII.5.2]{5.2}).

\enlargethispage{1.2\baselineskip}%
\begin{corollaire} \label{XII.5.5}
Under the preliminary conditions of \Ref{XII.5.2}, with $S$ noetherian, the following conditions are equivalent:
\begin{enumeratei}
\item
$\sisi{\coprod}{\bigoplus}_{m\geq m_0}h_\ast(F(m)\sisi{}{)}$ is of finite type over $\sisi{S}{\Ss}$, and $\R^ih_\ast(F(m))=0$ for $0<i\leq n$ and large $m$.
\item
$g_\ast(F)$ is coherent, and $\R^ig_\ast(F)=0$ for $0<i\leq n$.
\item[\textup{(ii bis)}] $g_\ast(F)$ is coherent, and $\prof_Y g_\ast(F)>n+1$.
\end{enumeratei}
\end{corollaire}

The equivalence of (ii) and (ii bis) is contained in \Ref{III}~\Ref{III.3.3}. Moreover, by \Ref{XII.5.2}, conditions (i) and (ii) both imply that the $\R^ig_\ast(F)$ ($i\leq n$) are coherent. The equivalence of (i) and (ii) then follows from \Ref{XII.5.4}, taking into account the fact that for every coherent Module $G$ on $X$, one has $G=0$ if and only if $f_\ast(G(m))=0$ for large $m$, for example by \EGA III 2.2.1 (iii).

\begin{remarque} \label{XII.5.6}
The\pageoriginale criteria \Ref{XII.5.2} and \Ref{XII.5.5} can be interpreted by saying that the \og simultaneous finiteness condition\fg \Ref{XII.5.2}~(ii) is expressed by local regularity properties (in terms of depth by \Ref{VIII}~\Ref{VIII.2.1}) of $F$ at the points of $U$ near $Y=X- U$, whereas the \og asymptotic vanishing condition\fg \Ref{XII.5.5}~(i) is distinctly stronger in nature and is expressed by local regularity conditions on $g_\ast(F)$ at the points of~$Y$ itself. It would be interesting, in order to generalize Lefschetz-type theorems for projective morphisms to quasi-projective morphisms, to find local criteria on $X$ that are necessary and sufficient for the $\Ss$-Modules $\sisi{\coprod}{\bigoplus}_{m\geq 0} \R^ih_\ast(F(m))$ for $i\leq n$ to be of finite type. When $S$ is the spectrum of a field (and doubtless more generally if it is the spectrum of an artinian ring) and $Y=X- U$ is finite, one can show that it is necessary and sufficient that the following conditions hold:
\begin{enumerate}
\item[\textup{1\noo)}] $\prof F_x >n$ for every closed point $x$ of $U$ (compare \Ref{XII.1.4}).
\item[\textup{2\noo)}] $\R^ig_\ast(F)$ is coherent for $i\leq n$, or equivalently, there exists an open neighborhood $V$ of $Y$ such that for every closed point $x$ of $U\cap V$, one has $\prof F_x >n+1$.
\end{enumerate}
\end{remarque}

\begin{proposition} \label{XII.5.7}
Let $S$ be a locally Noetherian prescheme, let $g:U\to X$ be a morphism of preschemes of finite type over $S$\sfootnote{It is in fact enough that $g$ be quasi-compact and quasi-separated (\textup{\EGA IV1.2.1}), with no condition on $U,X$.}, with structural morphisms $h$ and $f$, let $F$ be a quasi-coherent module on $U$, and let $n$ be an integer. The following conditions are equivalent:
\begin{enumeratei}
\item
For every base change $S'\to S$, with $S'$ Noetherian, the module $\R^ng'_\ast(F')$ on $X'$ is coherent.
\item
For every base change as above and every coherent ideal $\Jj$ on~$S'$, denoting by $\Ii$ the ideal $\Jj\Oo_{X'}$ on $X'$, the graded module
\[
\sisi{\coprod}{\tbigoplus}_{m\geq 0} \R^ng'_\ast(\Ii^mF')
\]
over $\sisi{\coprod}{\bigoplus}_{m\geq 0} \Ii^m$ is of finite type.
\item
For\pageoriginale every base change $S'\to S$, and $\Jj$ as above, the graded module
\[
\sisi{\coprod}{\tbigoplus}_{m\geq 0} \R^ng'_\ast(\Ii^mF'/\Ii^{m+1}F')
\]
over $\gr_\sisi{I}{\Ii}(\Oo_{X'})=\sisi{\coprod}{\bigoplus}_{m\geq 0} \Ii^m/\Ii^{m+1}$ is of finite type.
\end{enumeratei}
\end{proposition}

Obviously, (ii) \ALORS (i) and (iii) \ALORS (i), as one sees by taking $\Ii=0$ in conditions (ii) and (iii). The reverse implications are obtained by applying (i) to the composite base change $S^{''}\to S'\to S$, where $S^{''}$ is equal to $\Spec \sisi{\coprod}{\bigoplus}_{m\geq 0} \Jj^m$ \resp $\Spec \sisi{\coprod}{\bigoplus}_{m\geq 0} \Jj^m/\Jj^{m+1}$.

The point of this proposition is that conditions of the form (ii) are those that occur in the \og algebraic-formal comparison theorems\fg, whereas conditions of the form (iii) occur in the \og existence theorems\fg that complement them; \cf Expos\'e \Ref{IX}. A first interesting case is that in which $f:X\to S$ is the identity, so that one is dealing with conditions on a morphism $h:U\to S$ locally of finite type and a quasi-coherent module $F$ on $U$, flat over $S$. To obtain sufficient conditions, we shall assume that $U$ embeds by $g:U\to X$ as an open subprescheme of an $X$ that is \emph{proper} over $S$. Applying \Ref{XII.5.1}, we thus see:

\begin{corollaire} \label{XII.5.8}
Let $f:X\to S$ be a proper morphism, with $S$ locally Noetherian, let $U$ be an open subset of $X$, let $g:U\to X$ be the canonical immersion, let $h=fg:U\to S$, and let $F$ be a quasi-coherent module on $\sisi{X}{U}$, flat over $S$. Assume that, for every base change $S'\to S$, with $S'$ locally Noetherian, $\R^ig'_\ast(F')$ is coherent on $X'$ for $i\leq n$. Then the following hold:
\begin{enumeratei}
\item
For every base change $S'\to S$, with $S'$ locally Noetherian, $\R^ih'_\ast(F')$ is coherent on $S'$ for $i\leq n$.
\item
For every $S'\to S$ as above, and every coherent ideal $\Jj$ on $S'$, the graded modules\pageoriginale
\[
\sisi{\coprod}{\tbigoplus}_{m\geq 0} \R^ih'_\ast(\Jj^mF')
\]
over $\sisi{\coprod}{\bigoplus}_{m\geq 0}\Jj^m$ are of finite type for $i\leq n$.
\item
For every $S'\to S$ and $\Jj$ as above, the graded modules
\[
\sisi{\coprod}{\tbigoplus}_{m\geq 0} \R^ih'_\ast(\Jj^mF'/\Jj^{m+1}F')
\]
over $\gr_\Jj(\Oo_{S'})=\sisi{\coprod}{\bigoplus}_{m\geq 0}\Jj^m/\Jj^{m+1}$ are of finite type for $i\leq n$.

Moreover, under the conditions of (ii), and by the comparison theorem \Ref{IX.1.1}, denoting by $\widehat{S'}$ the formal completion of $S'$ with respect to $\Jj$, and by $\widehat{U'}$ that of $U'$ with respect to $\Jj\Oo_{U'}$, the canonical homomorphisms
\[
\widehat{\R^ih'_\ast(F')} \to \R^i\widehat{h'}_\ast(\widehat{F'}) \to \varprojlim_k\R^ih'_\ast(F'_k)
\]
are isomorphisms for $i\leq n-1$.
\end{enumeratei}
\end{corollaire}

\begin{remarque} \label{XII.5.9}
Assume in addition that the conditions of \Ref{XII.5.8} hold with $F$ \emph{coherent}, and consider a base change $S'\to S$ as in \sisi{\Ref{XII.5.9}}{\Ref{XII.5.8}} (i). Assume further that $S'$ is locally embeddable in a regular scheme, or, more generally, that there exists a faithfully flat and quasi-compact morphism $S^{''}\to \sisi{S}{S'}$ such that $S^{''}$ is locally embeddable in a regular scheme; this condition holds in particular if $S'$ is local. Then the conclusion of \Ref{XII.5.8} (i) and (ii) remains valid when $F'$ is replaced there by a module $G'$ on $U'$ such that every point of $U'$ has an open neighborhood on which $G'$ is isomorphic to a module of the form ${F'}^n$. Indeed, one reduces to the case where $S'$ itself is locally embeddable in a regular scheme, and hence so are $S^{''}=\Spec\big(\sisi{\coprod}{\bigoplus}_{m\geq 0}\Jj^m\big)$ and $X^{''}=X'\times_{S'}S^{''}=X\times_SS^{''}$, which are of finite type over it. One then applies the finiteness criterion \Ref{VIII}~\Ref{VIII.2.3} to the direct images for $i\leq n$ of $G^{''}$ under the immersion $U^{''}\to X^{''}$, noting that its conditions hold by hypothesis for $F^{''}$\pageoriginaled, and therefore also for $G^{''}$, since they are expressed in terms of depth and $G^{''}$ is locally isomorphic to an ${F^{''}}^n$. The same argument shows that, if $\Gg'$ is a coherent module on $\widehat{U'}$ (the completion of $U'$ with respect to the ideal $\Jj\Oo_{U'}$), such that $\Gg'_0=\Gg'/\Jj\Gg'$ is locally of the form $\sisi{F_0^{'n}}{{F'_0}^n}$, then the conclusion of (iii) remains valid when $F'$ is replaced there by $\Gg'$. Using the results of Expos\'e \Ref{IX}, one thus obtains the following result:
\end{remarque}

\begin{corollaire} \label{XII.5.10}
Let $f:X\to S$ be a proper morphism, with $S$ locally Noetherian, and let $U$ be an open subset of $X$; assume that $U$ is flat over $S$, and that, for every base change $S'\to S$, with $S'$ locally Noetherian, $\sisi{R^i}{\R^i}g'_\ast(\Oo_{U'})$ is coherent on $X'$ for $i=0,1$. Assume now that $S'$ is of the form $\Spec(A)$, where $A$ is a Noetherian ring equipped with an ideal $J$ such that $A$ is separated and complete for the $J$-adic topology. Under these conditions:
\begin{enumeratei}
\item
The functor $F\mto\widehat{F}$ from the category of locally free modules on $U'$ to the category of locally free modules on $\widehat{U'}$ is fully faithful.
\item
For every locally free module $\formelF$ on $\widehat{U'}$, there exists a coherent module $F$ on~$U'$ (not necessarily locally free, alas), and an isomorphism $\widehat{F}\simeq\formelF$.
\end{enumeratei}
\end{corollaire}

It remains only to prove (ii), by \Ref{XII.5.9}. But it follows from this remark and \Ref{IX.2.1} that $\formelF$ is induced by a coherent module $\formelG$ on $\widehat{X'}$. By the existence theorem \EGA III 5.1.4, $\formelG$ is of the form $\widehat{F}$, where $F$ is coherent on $\sisi{X}{X'}$, which proves the conclusion.

\begin{remarques} \label{XII.5.11}
1\textsuperscript{o}) Using \Ref{XII.5.10}, \Ref{XII.4.7}, and a suitable hypothesis saying that certain local rings of the geometric fibers of $X'\to S'$ are \og pure\fg \resp parafactorial, one should be able to obtain statements saying that the functor $Z'\mto\widehat{Z'}$ from the category of finite \`etale coverings of $X'$ to the category of finite \`etale coverings\pageoriginale of $\widehat{X'}$ (or, equivalently, of $X'_0$) is an equivalence of categories, \resp that the functor $L\mto\widehat{L}$ from the category of invertible modules on $X'$ to the category of invertible modules on $\widehat{X'}$ is an equivalence. Using recent results of Murre, it is likely that one should be able to deduce existence theorems for Picard schemes of certain \emph{nonproper} algebraic schemes\nde{Of course, in the projective case one should refer to Grothendieck's existence theorems in FGA; \cf Grothendieck~A., {\og Technique de descente et th\'eor\`emes d'existence en g\'eom\'etrie alg\'ebrique. VI. Les sch\'emas de Picard: propri\'et\'es g\'en\'erales\fg}, in \emph{S\'eminaire Bourbaki}, vol.~7, Soci\'et\'e math\'ematique de France, Paris, 1995, \Exp 236, p\ptbl 221--243 and {\og Technique de descente et th\'eor\`emes d'existence en g\'eom\'etrie alg\'ebrique. V. Les sch\'emas de Picard: th\'eor\`emes d'existence\fg}, in \emph{S\'eminaire Bourbaki}, vol.~7, Soci\'et\'e math\'ematique de France, Paris, 1995, \Exp 232, 143--161. The nine finiteness conjectures found there are proved in Expos\'es \Ref{XII} and \Ref{XIII} by Mme Raynaud and Kleiman in \SGA 6. For an excellent elementary account of the subject, see Kleiman's expository article (Kleiman~S., {\og The Picard scheme\fg}, forthcoming in Contemp. Math.). For an application of these techniques to the global generalized Jacobians of a relative smooth curve, see (Contou-Carr\`ere~C., {\og La jacobienne g\'en\'eralis\'ee d'une courbe relative; construction et propri\'et\'e universelle de factorisation\fg}, \textit{C.~R. Acad. Sci. Paris S\'er.~A-B} \textbf{289} (1979), \numero3, A203--A206 and {\og Jacobiennes g\'en\'eralis\'ees globales relatives\fg}, in \emph{The Grothendieck Festschrift, Vol. II}, Progr. Math., vol.~87, Birkh\"auser, Boston, 1990, p\ptbl 69--109). See also, by the same author and in the purely local setting, the construction and study of the \og local generalized Jacobian\fg functor ({\og Jacobienne locale, groupe de bivecteurs de Witt universel, et symbole mod\'er\'e\fg}, \textit{C.~R. Acad. Sci. Paris S\'er.~I Math.} \textbf{318} (1994), \numero 8, p\ptbl 743--746). Moreover, although in the case of a projective smooth morphism the connected components of the Picard scheme are proper, this is no longer so in the singular case. The problem of compactifying Picard schemes then arises naturally; this problem has been studied in detail, notably in (Altman~A.B. \& Kleiman~S., {\og Compactifying the Picard scheme\fg}, \emph{Adv. in Math.} \textbf{35} (1980), \numero 1, p\ptbl 50--112, and {\og Compactifying the Picard scheme. II\fg}, \emph{Amer. J.~Math.} \textbf{101} (1979), \numero 1, p\ptbl 10--41). The case of curves had been studied earlier (D'Souza~C., {\og Compactification of generalised Jacobians\fg}, \emph{Proc. Indian Acad. Sci. Sect. A Math. Sci.} \textbf{88} (1979), \numero 5, p\ptbl 419--457). One even knows exactly when the compactified Jacobian of a curve is irreducible (Rego~C.J., {\og The compactified Jacobian\fg}, \emph{Ann. Sci. \'Ec. Norm. Sup. (4)} \textbf{13} (1980), \numero 2, p\ptbl 211--223), the latter being the closure of the ordinary Jacobian when the curve is geometrically integral and locally planar; for a construction in families of compactified Jacobians, see (Esteves~E., {\og Compactifying the relative Jacobian over families of reduced curves\fg}, \emph{Trans. Amer. Math. Soc.} \textbf{353} (2001), \numero 8, p\ptbl 3045--3095). Since then, existence results for the Picard scheme in the proper case have advanced beyond the original edition of \SGA 2; \cf (Murre~J.P., {\og On contravariant functors from the category of pre-schemes over a field into the category of abelian groups (with an application to the Picard functor)\fg}, \emph{Publ. Math. Inst. Hautes \'Etudes Sci.} \textbf{23} (1964), p\ptbl 5--43), and especially (Artin~M., {\og Algebraization of formal moduli. I\fg}, in \emph{Global Analysis (Papers in Honor of K\ptbl Kodaira)}, Univ. Tokyo Press, Tokyo, 1969, p\ptbl 21--71). See also (Raynaud~M., {\og Sp\'ecialisation du foncteur de Picard\fg}, \emph{Publ. Math. Inst. Hautes \'Etudes Sci.} \textbf{38} (1970), p\ptbl 27--76) for a scheme proper over a discrete valuation ring but not necessarily flat. For a discussion of more recent results, in particular Artin's, for the Picard functor of proper flat schemes, especially in the cohomologically flat case in dimension $0$, see Chapter VIII of (Bosch~S., L\"utkebohmert~W. \& Raynaud~M., \emph{N\'eron models}, Ergebnisse der Mathematik und ihrer Grenzgebiete (3), vol.~21, Springer-Verlag, Berlin, 1990) and the references cited there. Much more recently, very precise results have been obtained for relative curves $f:X\to S$ \sisi{over}{over} the spectrum $S$ of a discrete valuation ring with perfect residue field. More precisely, one assumes that $f$ is proper and flat, that $X$ is regular, and that $f_*\OX=\Oo_S$. On the other hand, one does not assume that $f$ is cohomologically flat in dimension~$0$, \ie one does not assume that $H^1(X,\Oo)$ is torsion-free. The Picard functor is then not representable, either by a scheme or by an algebraic space, as soon as the torsion in question is nonzero. Let~$J$ be the N\'eron model of the generic fiber of $f$; it is a quotient of the Picard functor $P$. Raynaud then showed that the kernel of the tangent map $H^1(X,\Oo)=\mathrm{Lie}(P)\to\mathrm{Lie}(J)$ coincides with the torsion subgroup of $H^1$, and that the cokernel has the same length (see Theorem 3.1 of (Liu~Q., Lorenzini~D. \& Raynaud~M., {\og N\'eron models, Lie algebras, and reduction of curves of genus one\fg}, \emph{Invent. Math.} \textbf{157} (2004), p\ptbl 455--518)). This result enables those authors to study the relation between the Birch--Swinnerton-Dyer and Artin--Tate conjectures (see Th.\ptbl 6.6 of {\loccit}). As for the local Picard scheme, see Boutot's thesis, cited in editor's note~\eqref{Boutot} on page~\pageref{Boutot}.}. In general, the removal of purity hypotheses in various existence theorems, especially representability theorems for functors such as the Hilbert or Picard functors, by means of the techniques developed in this seminar, deserves systematic study.

2\textsuperscript{o}) One may seek manageable necessary and sufficient conditions, in terms of depth, for the universal finiteness condition considered in \Ref{XII.5.10} to hold. When $S$ is the spectrum of a field, it follows easily from \EGA III 1.4.15 that it is necessary and sufficient for the $\R^ig_\ast(F)$ ($i\leq n$) to be coherent, which is indeed expressed in terms of depth by \Ref{VIII}~\Ref{VIII.2.3}. In the general case, however, note that it is not enough to require the preceding condition to hold for all the fibers $U_s\subset X_s$ ($s\in S$), even when $n=0$. For example, take $X=S$, let $S$ be the spectrum of a discrete valuation ring, let $U$ be the open subset consisting only of the generic point, and let $F=\Oo_U$.

3\textsuperscript{o}) Here, however, is a \emph{sufficient} condition ensuring that the hypotheses of \Ref{XII.5.10} hold: it is enough that $f$ be flat and that, for every $s\in S$ and every $x\in Y_s=X_s-U_s$, one have
\[
\prof \Oo_{X_s,x}\geq n+2.
\]
Indeed, taking Lemma \Ref{XII.2.5} into account (\cf relation (\Ref{eq:XII.16}) following \Ref{XII.2.5}), it follows that $g_*(\Oo_U)\simeq\OX$ and $\R^ig_*(\Oo_U)=0$ for $0<i\leq n$, and the same relations will clearly hold after every base change $S'\to S$.
\end{remarques}

% END INCLUDED active_source/parts/12_expose_XII.tex

% BEGIN INCLUDED active_source/parts/13_expose_XIII.tex
\chapter{Problems and conjectures} \label{XIII}

\section[Relations between global and local results]
{Relations between global and local results. Affine problems connected with duality} \label{XIII.1}

It\pageoriginale is well known that many statements concerning a projective scheme
$X$ can be formulated in terms of statements concerning
a certain graded ring, or better, a complete local ring,
\sisi{\ignorespaces}{namely} the homogeneous coordinate ring
of $X$ (\ie the affine ring of the cone $\tilde{X}$ over
$X$), or its completion (\ie the completion of the local ring at the vertex
of $\tilde{X}$). The value of this reformulation is that,
starting from known global results, it often makes it possible to
conjecture, or even prove, analogous results for complete
Noetherian local rings more general than those which
actually occur in the global statement, for example for
local rings that are not necessarily of equal
characteristic. Thus, Serre's duality theorem for
projective space \Ref{XII}~\Ref{XII.1.1} suggested the useful
local duality theorem \sisi{V.5.6}{\Ref{V}~\Ref{V.2.1}}. Serre's
fundamental theorem on the cohomology of coherent
\sisi{m}{M}odules on projective space (finiteness,
a\sisi{s}{}symptotic behavior for large $n$ of $\H^i(X,
F(n))$, {\cf \EGA III~2.2.1}) generalizes to a
structure theorem for the local invariants $\H^i_{\mm}(M)$; see
\sisi{V.5.7}{\Ref{V}~\Ref{V.3}}. Likewise, the Lefschetz theorems
for the fundamental group and the Picard group (\og
equivalence criteria\fg), well known in the
classical case and subsequently extended to an arbitrary base field,
suggested the \og local\fg Lefschetz theorems of Exposés
\Ref{X} and \Ref{XI}. Of course, the local theorems are in
turn valuable tools for obtaining global statements.
For example, local duality allows a global asymptotic property
\Ref{XII}~\Ref{XII.1.3} (i) to be formulated by the vanishing
of certain local invariants $\H_x^i(F_x)$. More
substantially, the local Lefschetz theorems, implying for
example the \og purity\fg or parafactoriality of certain
complete-intersection local rings (\Ref{X}~\Ref{X.3.4} and
\Ref{XI}~\Ref{XI.3.13}), make it possible in the global
Lefschetz theorems to dispense with certain nonsingularity
hypotheses, as in \Ref{X}~\Ref{X.3.5}, \Ref{X.3.6},
\Ref{X.3.7}.

Another\pageoriginale useful generalization of theorems concerning projective schemes over a field $k$ consists in replacing $k$ by a general base scheme. Thus, the sequel to {\EGA III} will give\nde{in fact, this generalization is not found there; see the note below and the editor's note~\eqref{noteconrad} on page~\pageref{noteconrad}} a generalization in this direction of Serre duality\sfootnote{\cf the Hartshorne Seminar, cited at the end of {\Exp \Ref{IV}}.}; the finiteness and asymptotic-behavior theorems for $\H^i(X, F(n))$ were stated in \EGA III~2.2.1 over a general base scheme; finally, the Lefschetz theorems can likewise be developed for a projective morphism, as we saw in \Ref{XII}~\Ref{XII.4.9}, thanks to the local theorem \Ref{XII}~\Ref{XII.4.7}. Of course, working over a general base scheme also leads to essentially new statements, such as the \og comparison theorem\fg \EGA III~4.1.5 and the existence theorem for sheaves \EGA III~5.1.4 (which, as we saw elsewhere in \Ref{IX}, arise from the same key theorems of a cohomological nature as the Lefschetz theorems for $\pi_1$ and $\Pic$).

It is therefore necessary to isolate theorems that simultaneously encompass the two generalizations just mentioned of statements concerning projective schemes over a field. The natural objects for such a common generalization are the \emph{Noetherian rings separated and complete for an $I$-adic topology}. From this point of view, their study has not yet been seriously undertaken and seems to me at present the most interesting subject in the local theory of coherent sheaves. Here is a typical problem in this direction:

\refstepcounter{subsection}\label{XIII.1.1}
\begin{enonce*}{Conjecture 1.1}
 [\og Second affine finiteness theorem\fg
 \markedfootnote{**}{This conjecture, and Conjecture \hyperref[XIII.1.2]{1.2} below, are false, as was shown by \sisi{R.~HARTSHORNE (unpublished)}{R\ptbl Hartshorne, {\og Affine duality and cofinite modules\fg}, \emph{Invent. Math.} \textbf{9} (1969/70), p.~145-164, Section 3.}}
 \sisi{}{\setcounter{footnoteNDE}{1}}\ndemark]
\ndetext{However, if $A$ is complete local (respectively regular of positive characteristic) and $J$ is the maximal ideal, the statement is true for $M$ of finite type (respectively $M=A$), \cf (Hartshorne~R., {\og Affine duality and cofinite modules\fg}, \emph{Invent. Math.} \textbf{9} (1969/70), p\ptbl 145-164, Corollary 1.4) (respectively (Huneke~C. \& Sharp~R., {\og Bass numbers of local cohomology modules\fg}, \emph{Trans. Amer. Math. Soc.} \textbf{339} (1993), \numero 2, p\ptbl 765--779), which moreover contains much stronger results). For completely different methods ($D$-modules) for treating characteristic zero, see (Lyubeznik~G., {\og Finiteness properties of local cohomology modules (an application of $D$-modules to commutative algebra)\fg}, \emph{Invent. Math.} \textbf{113} (1993), \numero 1, p\ptbl 41--55); see also, by the same author, {\og Finiteness properties of local cohomology modules for regular local rings of mixed characteristic: the unramified case\fg}, \emph{Comm. Algebra} \textbf{28} (2000), \numero 12, p\ptbl 5867--5882, Special issue in honor of Robin Hartshorne, and {\og Finiteness properties of local cohomology modules: a characteristic-free approach\fg}, \emph{J.~Pure Appl. Algebra} \textbf{151} (2000), \numero 1, p\ptbl 43--50. The notion of a cofinite module has since evolved under Hartshorne's leadership. One says that $M$ is $J$-cofinite if its support is contained in $V(J)$ and if all the $\Ext^i_A(A/J, M)$ are of finite type. On this subject, see for example (Delfino~D. \& Marley~Th., {\og Cofinite modules and local cohomology\fg}, \emph{J. Pure Appl. Algebra} \textbf{121} (1997), \numero 1, p\ptbl 45--52).}
Let $M$ be a module of finite type over a Noetherian ring $A$ (which may, if necessary, be assumed to be a quotient of a regular ring), and let $J$ be an ideal of $A$. Prove that the modules $\H^i_J(M)$ are \og $J$-cofinite\fg, \ie that the modules
$$
\Hom_A(A/J, \H^i_J(M))$$
are of finite type.
\end{enonce*}

Recall that $\H^i_J (M)$ denotes the module $\H^i_Y (X, \tilde{M})$ (where $X=\Spec (A)$, $Y=V(J)$) of \Exp \Ref{I}, interpreted in \Ref{II} in terms of an inductive limit of the cohomology\pageoriginale of Koszul complexes, or, for $i\ge 2$, the module $\H^{i-1} (X-Y, \tilde{M})$. \sisi{T}{T}o tell the truth, \Ref{XIII.1.1} should be a consequence of a more precise statement, implying that the $\H^i_J ({M})$ belong to a suitable abelian subcategory ${\DDJ}$ of the category ${\CJ}$ of $A$-modules with support $\subset Y=V(J)$, such that $H\in \Ob {\DDJ}$ implies that $H$ is $J$-cofinite. (\sisi{NB}{N.B.} The category of modules $H$ with support contained in $V(J)$ and which are $J$-cofinite is unfortunately not stable under passage to quotients!). The essential problem would then consist in defining ${\DDJ}$. More precisely, the solution of problem \Ref{XIII.1.1} should emerge (at least if $A$ is a quotient of a regular ring) from a duality theory generalizing both local duality and the duality theory for projective morphisms alluded to above, and which would be of the following kind:

\refstepcounter{subsection} \label{XIII.1.2}
\begin{enonce*}{Conjecture 1.2}[\og Affine duality\fg\sfootnotemark]
\sfootnotetext{This conjecture, false as stated, was nevertheless established in a fairly similar form by \sisi{R.~HARTSHORNE, \emph{Affine duality and cofinite modules} (forthcoming).}{R.~Hartshorne, {\og Affine duality and cofinite modules\fg}, \emph{Invent. Math.} \textbf{9} (1969/70), p\ptbl 145-164.}} Suppose that $A$ is regular, separated, and complete for the $J$-adic topology. Let $C^{\boule}(A)$ be an injective resolution of $A$.
\begin{enumeratei}
\item
Prove that the functor
$$D_J: L_{\boule}\mto \Hom_J(L_{\boule}, C^{\boule}(A))$$
from the category of complexes of $A$-modules that are free of finite type in every dimension and bounded above (\sisi{where}{where} the morphisms are homomorphisms of complexes up to homotopy) to the category of complexes of $A$-modules $K^{\boule}$ that are injective in every dimension and bounded above (\sisi{where}{where} the homomorphisms are defined in the same way), is fully faithful.

\item
Prove that, for every $K^{\boule}$ of the form $D_J(L_{\boule})$, the $\H^i(K^{\boule})(= \Ext_Y^i (X; L_{\boule}, \OX))$ are $J$-cofinite.

\item
More precisely, prove that the $K^{\boule}$ which are homotopic to a complex of the form $D_J(L_{\boule})$ can be characterized by finiteness properties of the $\H^i(K^{\boule})$ stronger than the one considered in (ii), for example by\pageoriginale the property $\H^i(K^{\boule})\in \Ob {\DDJ}$, where ${\DDJ}$ is a suitable abelian category, as contemplated above.
\end{enumeratei}
\end{enonce*}

Note that the problem has an affirmative solution when $A$ is local and $J$ is an ideal of definition (\cf \Exp \Ref{IV}), and also when $J$ is the zero ideal. In these two exceptional cases, one may take for ${\DDJ}$ simply the category of \sisi{m}{M}odules with support $V(J)$ which are $J$-cofinite (which in the second case means simply that one takes the category of \sisi{m}{M}odules of finite type over $A$). An affirmative solution of Conjecture \Ref{XIII.1.2} in general would give one for \Ref{XIII.1.1}, by taking for $L_{\boule}$ the dual of a finite free resolution of $M$. On the other hand, an affirmative solution of \Ref{XIII.1.1} would give an affirmative answer to the first part of the following conjecture, which we formulate in \og global\fg form:

\begin{conjecture} \label{XIII.1.3}
Let $X\subset \PP^r_k$ be a closed subscheme of the standard projective scheme which is locally a complete intersection and each of whose irreducible components has codimension $\ge s$. Let $U=\PP^r_k -X$.
\begin{enumeratei}
\item
Prove that, for every coherent Module $F$ on $U$, one has
$$
\dim \H^i(U, F)<+\infty{\rm \ for\ } i\ge s.{}\sfootnote{Part\sisi{}{ie} (i) of this conjecture was proved by R.~\sisi{HARTSHORNE}{Hartshorne} when $\sisi{Y}{X}$ is smooth \sisi{in \emph{Ample Vector Bundles}, Pub. Math. \numero 29, (1966)}{\cf \emph{Ample subvarieties of algebraic varieties}, Notes written in collaboration with C\ptbl Musili, Lect. Notes in Math, vol.~156, Springer-Verlag, Berlin-New York, 1970, Theorem III.5.2}. This author also found an example for (ii), \cf \sisi{R.~HARTSHORNE, \emph{Cohomological dimension of algebraic varieties}, forthcoming in the Ann. of Math.}{R.~Hartshorne, {\og Cohomological dimension of algebraic varieties\fg}, \emph{Ann. of Math. (2)} \textbf{88} (1968), p\ptbl 403--450, example on page 449.}}
$$

\item
Give an example, with $X$ connected and regular, for which one has
$$
\H^s(U, F)\neq 0.
$$
\end{enumeratei}
\end{conjecture}

To see that (i) is a special case of \Ref{XIII.1.1}, consider
$$
\H^i(U, F(\cdot))=\sisi{\coprod}{\tbigoplus}_n \H^i (U, F(n))= \H^i(E^{r+1} -\sisi{X}{\tilde X}, \tilde{F})$$
as a module over the affine ring $k[t_0, \ldots, t_r]$ of the cone $E^{r+1}$ over $\PP^r$. This module is none other than $\H_J^{i+1}(M)$, where $J$ is the ideal of the cone $\tilde X$ over $X$ in $E^{r+1}$. On the other hand, the hypothesis on $X$, which implies that $\sisi{X}{\tilde X}$ is also a complete intersection of co\sisi{-}{}dimension $\ge s$ at every\pageoriginale point of $E^{r+1}$ distinct from the origin, implies that $\H^{i+1}_J(M)$ is zero away from the origin for $i\ge s$. Thus, if it is $J$-cofinite as required by \Ref{XIII.1.1}, it is \afortiori $\mm$-cofinite, which readily implies that it is finite-dimensional in every degree\refstepcounter{toto}\nde{\label{noteHM}Hartshorne proved (Hartshorne~R., {\og Cohomological dimension of algebraic varieties\fg}, \emph{Ann. of Math. (2)} \textbf{88} (1968), p\ptbl403--450) that the cohomology $H^{n-1}(\PP^n_k-X,F)$ is zero for $F$ coherent and $X$ of positive dimension ($k$ algebraically closed). In fact, thanks essentially to Serre duality and the Lichtenbaum theorem --- the vanishing of the cohomology of coherent sheaves in the maximal dimension of irreducible noncomplete quasi-projective varieties ---, one is reduced to proving that the formal completion $\hat \PP^n_k$ and $X$ have the same field of rational functions. This is the difficult point (Theorem 7.2 of {\loccit}); in other words, $\PP^n_k$ is $G3$ in the terminology of (Hironaka~H. \& Matsumura~H., {\og Formal functions and formal embeddings\fg}, \emph{J.~Math. Soc. Japan} \textbf{20} (1968), p\ptbl 52--82). These authors independently proved the preceding results, and in fact much more. They proved that $X$ is universally $G3$ and calculated the field of rational functions of the formal completion of an abelian variety along a positive-dimensional subvariety. It is in this article that the now-classical conditions $G1,G2,G3$ appear for the first time.} Note moreover that Conjecture \Ref{XIII.1.3} already arises for a nonsingular irreducible curve $X$ in $\PP^3$: it is unknown whether in this case the $\H^2(\PP^3-X, \OX(n))$ are finite-dimensional, or whether they are necessarily zero\sfootnote{The question has just been answered affirmatively by R\ptbl \sisi{HARTSHORNE}{Hartshorne}\sisi{}{ (Hartshorne~R., {\og Ample vector bundles\fg}, \emph{Publ. Math. Inst. Hautes \'Etudes Sci.} \textbf{29} (1966), p\ptbl63--94, Theorem 8.1)} and H. \sisi{HIRONAKA}{Hironaka}.}. It is not even known whether there exists an irreducible curve in $\PP^3$ which is not set-theoretically the intersection of two hypersurfaces\label{courbesurface}\nde{see Conjecture~\Ref{XII.3.5} and the corresponding note.}.

\begin{probleme} \label{XIII.1.4}
Give an affine variant of the \og comparison theorem\fg \textup{\EGA III~4.1.5} as a theorem on the commutation of the functors $\H^i_J$ with certain inverse limits.
\end{probleme}

Finally, in the present line of thought, I had posed the following problem: \sisi{L}{l}et $A$ be a \emph{complete} regular Noetherian local ring, and $K$ its field of fractions; prove that $\Ext^i_A(K, A)=0$ for every $i$. An affirmative answer was given on the spot by \sisi{Mr AUSLANDER}{M.~Auslander}; the regularity hypothesis may be replaced by the hypothesis that $A$ is integral, and in fact $\Ext^i_A(K, M)=0$ for every $i$ whenever $M$ is of finite type over $A$. This follows immediately from the following statement, due to \sisi{AUSLANDER}{Auslander}: if $A$ is a complete Noetherian local ring, then for every module of finite type $M$ over $A$, the functors $\Ext_A^i(., M)$ transform inductive limits into inverse limits\nde{write $K=\varinjlim_{a\neq 0}A[1/a]$ and observe that $a\neq 0$ is $A$-regular.}.

\section{\texorpdfstring{Problems related to $\pi_0$: local Bertini theorems}{Problems related to pi\textunderscore 0: local Bertini theorems}} \label{XIII.2}

Let $A$ be a complete noetherian local ring, $f$ an element of its
maximal ideal, $X=\Spec (A)$, $Y=\Spec(A/fA)$. The use of the local
\og Lefschetz\fg technique makes it possible to give criteria
for $Y'=X'\cap Y$ (where $X'=X-\{\mm\}$) to be connected, in terms
of hypotheses on $X'$. Thus, it suffices that: a) $X'$\pageoriginale
is connected\sisi{}{,} b) $\prof \Oo_{X', x} \ge 2$ for every closed point $x$
of $X'$\sisi{}{,} c) $f$ is $A$-regular. Note, however, that
hypotheses b) and c) are not purely
topological in nature; for example, they are not invariant\sisi{}{}
under replacement of $X$ by $X_{\text{red}}$. In the analogous situation
for a projective scheme $X'$ over a field and a hyperplane
section $Y'$ of $X'$, the use of the \og Bertini theorem\fg
and Zariski's \og connectedness theorem\fg in fact gives
results of a distinctly more satisfactory form,
which had led me, in the oral seminar, to state
a conjecture that I have since solved affirmatively.
Let us therefore state here:

\begin{theoreme} \label{XIII.2.1}
Let $A$ be a complete noetherian local ring, $X$ its spectrum, $\aaa$ the closed point of $X$, $X'=X-\{ \aaa\}$. Suppose that $X$ satisfies the following conditions (where $k$ denotes an integer $\ge 1$):
\begin{itemize}
\item[a$_k$)] The irreducible components of $X'$ have dimension $\ge k+1$.

\item[b$_k$)] $X'$ is connected in dimension $\ge k$, \ie it cannot be disconnected by a closed subset of dimension $<k$ (\cf \Ref{III}~\Ref{III.3.8}).
\end{itemize}

Let $m$ be an integer, $0\le m\le k$, and let $f_1, \dotsc, f_m\in \rr(A)$; set $B=A/\sum_i f_iA$, $Y=\Spec(B)=V(f_1)\cap\dotsb\cap V(f_m)$, $Y'=X'\cap Y=Y-\{\aaa\}$. Then $Y$ satisfies conditions a$_{k-m}$), b$_{k-m}$). In particular, for every sequence of $m\le k$ elements $f_1, \dotsc, f_m$ of $\rr(A)$, $Y'=X'\cap V(f_1)\cap\dotsb\cap V(f_m)$ is connected.
\end{theoreme}

It is moreover easy to see that if the last conclusion holds (it plainly suffices to take $m=k$ there), and if the case where $X$ is irreducible of dimension~$0$ or $1$ is excluded, it follows that the irreducible components of $X'$ have dimension $\ge k+1$, and $X'$ is connected in dimension $\geq k$, so that, in a sense, \Ref{XIII.2.1} is a \og best possible\fg result.

\bigskip

Let us give the principle of the proof of~\Ref{XIII.2.1}. Only condition b$_{k-m}$) for $Y$ presents a problem. For a given $k$, one is easily reduced to the case where $X$ is integral, and even (by passing to the normalization, which is finite over $X$) to\pageoriginale the case where $X$ is \emph{normal}. If $k=1$, hence $\dim X' \ge 2$, then $X'$ has depth $\ge 2$ at its closed points and one may apply the result recalled at the beginning of \sisi{N${}^\circ$}{n$^\circ$}, which shows that \hbox{$Y'=X'\cap V(f)$} is connected. In the case $k\ge 1$, suppose the theorem proved for \hbox{$k'\!<\!k$}. By induction on $m$, one is reduced to the case $m=1$, \ie to verifying that, for $f_1\in\nobreak \rr(A)$, $X'\cap V(f_1)$ is connected in dimension $\ge k-1$. If it were not, \ie if it were disconnected by a $Z'$ of dimension $< k-1$, there would exist a sequence $f_2, \dotsc, f_k$ such that \hbox{$X'\cap V(f_1)\cap\dotsb\cap V(f_k)$} is disconnected, and in this sequence one may choose $f_2\in \rr(A)$ arbitrarily, subject only to the condition that it not vanish at any point of a certain finite subset $F$ of $X'$ (\sisi{\ignorespaces}{namely }the set of maximal points of $Z'$). Moreover, one easily verifies, using the fact that $X'$ is normal and therefore satisfies Serre's condition ({$S_2$})\sfootnote{\Cf \EGA IV 5.7.2.}, that there exists a finite subset $F'$ of $X'$ such that $f\in \rr(A)$, $V(f)\cap F'=\emptyset$ implies that $V(f)\cap X'$ also satisfies condition ({$S_2$}). One may then choose $f_2$ so that $f_2$ vanishes neither on $F$ nor on $F'$, and hence so that $X'\cap V(f_2)$ satisfies ({$S_2$}). But then\sisi{}{,} by the theorem of \sisi{HARTSHORNE}{Hartshorne} \Ref{III}~\Ref{III.3.6}\sisi{}{,} $X'\cap V(f_2)$ is connected in codimension~$1$ and thus (since every component of $X'\cap V(f_2)$ has dimension $\ge k$) it is connected in dimension $\ge k-1$. Applying the induction hypothesis to $V(f_2)=\Spec(A/f_2A)$, it follows that $X'\cap V(f_2)\cap V(f_1)\cap V(f_3)\cap \dotsb\cap V(f_k)$ is connected, whereas it had been constructed disconnected, a contradiction.

Let us mention some interesting corollaries:

\begin{corollaire} \label{XIII.2.2}
Let $f\colon X\to Y$ be a proper morphism of locally noetherian preschemes, with $Y$ integral, $y_0\in Y$, and $y_1$ the generic point of $Y$. Assume that
\begin{enumeratea}
\item
$Y$ is unibranch at $y_0$, and every irreducible component of $X$ dominates $Y$.
\item
The irreducible components of $X_1$ have dimension $\ge k+1$, and $X_1$ is connected in dimension $\ge k$.
\end{enumeratea}

Then\pageoriginale the irreducible components of $X_0$ have dimension $\ge k+1$, and $X_0$ is connected in dimension $\ge k$.
\end{corollaire}

Indeed, Zariski's connectedness theorem (\cf\ \EGA III 4.3.1) implies that $X_0$ is connected; to show that it is not disconnected by a closed subset of dimension $<k$, one is reduced to showing that the local rings at points $x\in X_0$ such that $\dim \overline{x}<k$ have a spectrum that is not disconnected by $x$. But this is true without assuming either that $f$ is proper or that $Y$ is unibranch at $y_0$. To see this, one is reduced to the case where $X$ is integral and dominates $Y$, and, if desired, where $Y$ is affine of finite type over $\ZZ$, so that one is under the hypotheses of the dimension formula for $\Oo_{X, x}$ over $\Oo_{Y, y_0}$. Using the finiteness of the normalization in this case, one may even suppose $X$ \emph{normal}; hence, by a theorem of \sisi{NAGATA}{Nagata}\sfootnote{\Cf \EGA IV 7.8.3 (i) (ii) (v).}, the completion of a local ring $\Oo_{X, x}$ of $X$ is still normal, and therefore (if $\Oo_{X, x}$ has dimension $N$) $\Spec (\hat{\Oo}_{X, x})$ is connected in dimension $\ge N-1$. Let $n=\dim\Oo_{\sisi{X}{Y}, y_0}$; then $\degtr~k(x)/k(y_0)<k$ implies $\dim \Oo_{X, x}>n+(k+1)-k=n+1$, in view \sisi{that }{of }$\dim X_1\ge k+1$, and, taking a system $f_1, \dotsc, f_n$ of parameters of $\Oo_{Y, y_0}$ and lifting them to elements of $\Oo_{X, x}$, one sees from \Ref{XIII.2.1} that $\Spec (\hat{\Oo}_{X, x}/\sum f_i\hat{\Oo}_{X, x})$ is connected in dimension $\ge 1$, \ie is not disconnected by its closed point, or, equivalently, $\Spec(\hat{\Oo}_{X_0, x})$ is not disconnected by its closed point; \afortiori the same holds for $\Spec (\Oo_{X_0, x})$.

As in the case of the ordinary connectedness theorem, one may vary \Ref{XIII.2.2} by taking geometric fibers (over algebraic closures of the residue fields), provided that one assumes $Y$ to be \emph{geometrically} unibranch at $y_0$, or (with no hypothesis other than that $Y$ is noetherian) that $f$ is universally open. Applying this to the case where $Y$ is the dual scheme of a projective scheme $\PP_k^r$ over a field, one recovers a strengthened form of the global result that had inspired \Ref{XIII.2.1}, namely:

\refstepcounter{subsection}\label{XIII.2.3}
\begin{enonce*}{Corollary 2.3\ndemark}\ndetext{for a very beautiful direct proof, see
(Fulton~W. \& Lazarsfeld.~R., {\og Connectivity and its applications in algebraic geometry\fg},
in \emph{Algebraic geometry (Chicago, Ill., 1980)}, Lect. Notes in Math., vol.~862,
Springer, Berlin-New York, 1981, p\ptbl 26--92, Theorem 2.1). \Cf also \cite{HL}, cited in the editor's note~\eqref{HL}~on page~\pageref{HL}.}
Let $X$ be a closed subscheme of $\PP^r_k$ ($k$ a field); suppose that the irreducible components of $X$ have dimension $\ge \sisi{k}{l}+1$, and that $X$ is geometrically connected\pageoriginaled in dimension $\ge \sisi{k}{l}$. Then for every sequence $H_1, \dotsc, H_m$ of $m$ hyperplanes of $\PP^r_k$ ($0\le m\le \sisi{k}{l}-1$), $X\cap H_1\cap \dotsb \cap H_m$ satisfies the same condition with $\sisi{k}{l}-m$; in particular, it is geometrically connected in dimension $\ge \sisi{k}{l}-1$.
\end{enonce*}

One may moreover modify this statement in the
obvious way for the case where one is given a proper morphism
$X \to \PP^r_k$, which is not necessarily an
immersion, and an analogous extension is possible for
\Ref{XIII.2.1} (by considering a scheme proper over~$X'$). These
statements moreover follow formally from the statements given
here, in view of the ordinary connectedness theorem, which
reduces us to the case of a finite morphism.

\begin{corollaire} \label{XIII.2.4}
Let $A$ be a complete normal noetherian local ring of dimension $\ge k+2$. Let $X=\Spec (A)$, $X'=X-\{ \aaa\}$, and let $f_1, \dotsc, f_k$ be elements of $\rr(A)$; then $Y'=X'\cap V(f_1) \cap \dotsb \cap V(f_k)$ is connected, and $\pi_1(Y')\to \pi_1(X')$ is surjective.
\end{corollaire}

One proceeds as in \SGA 1~X~2.11.

Throughout the preceding discussion, only questions of \emph{connectedness} were involved. In the global case, however, well-known theorems assert that for an irreducible projective variety $X\subset \PP^r_k$, with $k$ algebraically closed, its intersection with a sufficiently \og general\fg hyperplane $H$ is irreducible (and not merely connected): this is the \emph{\sisi{BERTINI}{Bertini} theorem}, proved by \sisi{ZARISKI}{Zariski}, which in turn implies, by Zariski's connectedness theorem, that for \emph{every} $H$, $H\cap X$ is connected (although not necessarily irreducible). One may moreover proceed in the opposite direction, by proving the latter result using a Lefschetz-type technique and deducing the Bertini theorem by reducing to the case where $X$ is normal and using the following result: for a \og sufficiently general\fg $H$, $X\cap H$ is also normal. This suggests the

\refstepcounter{subsection}\label{XIII.2.5}
\begin{enonce*}{Conjecture 2.5\ndemark}\ndetext{see the following editor's note.}
Let $A$ be a complete normal noetherian local ring. Show that there exists a nonzero $f\in\rr(A)$ such that $Y'=X'\cap V(f)=Y-\{ \aaa\}$ (where\pageoriginale $Y=\Spec (A/fA)$) is normal (and hence irreducible by \Ref{XIII.2.1} if $\dim A\ge 3$).
\end{enonce*}

Ideally, one should show that, in a suitable sense, there are even \og many\fg elements $f$ having the property in question; for example, that $f$ may be chosen in an arbitrary power of the maximal ideal. Using Serre's normality criterion and the remark made above concerning Serre's property $(S_2)$, one sees that \Ref{XIII.2.5} would have an affirmative answer if the same were true for the

\refstepcounter{subsection}\label{XIII.2.6}
\begin{enonce*}{Conjecture 2.6\ndemark}
\ndetext{a proof of this conjecture can now be found in the literature, and the preceding conjecture must therefore also be regarded as proved, as indicated above. One can also find two earlier published but unfortunately unsuccessful proof attempts, by Flenner and Trivedi. See Trivedi V., {\og Erratum: ``A local Bertini theorem in mixed characteristic"\fg}, \emph{Comm. Algebra} \textbf{25} (1997), \numero 5, p\ptbl1685-1686. However, the editor has not verified that the proof is now complete.}
Let $A$ be a complete noetherian local ring, $U$ an open subset of its spectrum $X$, and $F$ a finite subset of $X'=X-\{ \aaa\}$. Suppose that $U$ is regular. Prove that there exists an $f \in\rr(A)$ such that $V(f)\cap U$ is regular, and $V(f)\cap F=\emptyset$.
\end{enonce*}

For a result of \og local Bertini\fg type, see \sisi{CHOW}{Chow}~\cite{XIII.2}

\section{\texorpdfstring{Problems related to $\pi_1$}{Problems related to pi\textunderscore 1}} \label{XIII.3}
Here again there are many questions, suggested by the global results or by
the transcendental results.

\refstepcounter{subsection}\label{XIII.3.1}
\begin{enonce*}{Conjecture 3.1\ndemark}
\ndetext{The analogous statement is true for (connected) schemes $X$ of
finite type over a separably closed field $k$, under the hypothesis of
strong desingularization for all $\bar k$-schemes of finite type; in
particular, this holds if $k$ has characteristic zero or if $X$ has
dimension $\leq 2$. One reduces to the case of quasi-projective surfaces
by Lefschetz-type techniques developed by Mme Raynaud; \cf the notes
\textit{supra}; see \SGA 7.1, Theorem II.2.3.1.}
Let $A$ be a complete noetherian local ring with algebraically closed
residue field, let $X=\Spec(A)$, let $X'=X-\{\aaa\}$, where $\aaa$ is the
closed point. Suppose that the irreducible components of $X$ have
dimension $\geq 2$ and that $X'$ is connected.
\begin{enumeratei}
\item
Prove that $\pi_1(X')$ is topologically finitely generated.
\item
If $p$ is the characteristic exponent of the residue field $k$ of $A$,
prove that the largest topological quotient group of $\pi_1(X')$ that is
\og of order prime to $p$\fg is finitely presented.
\end{enumeratei}
\end{enonce*}

For part (i), using the descent theory of \SGA 1~IX~5.2 and
Theorem\pageoriginale~\hyperref[XIII.2.4]{2.4}, one is reduced to the case in which $A$ is normal of
dimension $2$. In this case, a systematic method for studying the
fundamental group of $X'$, introduced by \sisi{MUMFORD}{Mumford}~\cite{XIII.5}
in the transcendental setting, consists in \emph{desingularizing} $X$,
\ie considering a projective birational morphism $Z\to X$, with $Z$
integral and regular, which induces an isomorphism
$Z'=Z|X'\to X'$. It is plausible that such a $Z$ always exists; in
\sisi{all}{any} case this is what \sisi{ABHYANKAR}{Abhyankar}'s method~\cite{XIII.1}
proves in the case of \og equal characteristic\fg\sfootnote{The
possibility of \og resolving\fg $A$ has now been proved in full
generality by Abhyankar~\cite{XIII.8}.}. Let $C$ be the fiber over the
closed point of $X$ under $Z\to X$; it is an algebraic curve over the
residue field $k$, connected by the connectedness theorem. The solution
of \Ref{XIII.3.1} then seems to be related to the following problem.

\begin{probleme} \label{XIII.3.2}
With the preceding notation, relate $\pi_1(X')$ to the topological
invariants of $C$, in particular to its fundamental group, so as to
exhibit the topological finite generation of $\pi_1(X')$, using for
example \SGA \sisi{}{\textup{I}, Theorem~}\textup{X~2.6}.
\end{probleme}

Another method would be to regard $A$ as a finite algebra over a
complete \emph{regular} local ring $B$ of dimension $2$, ramified along
a curve $C$ contained in $\Spec(B)=Y$. This leads to the following
problem.

\begin{probleme} \label{XIII.3.3}
Let $A$ be a complete regular local ring of dimension $2$, with
algebra\sisi{.}{ically} closed residue field $k$; let $X$ be its spectrum,
and let $C$ be a closed subset of $X$ of dimension~$1$. Define local
invariants of the embedded curve $C$ that have a meaning independent of
the residue characteristic and whose knowledge makes it possible to
calculate the fundamental group of $X-C$ by generators and relations
when $k$ has characteristic zero. Prove that when $k$ has characteristic
$p>0$, the \og tame\fg fundamental group of $X-C$ is a quotient of the
preceding one, and that the two fundamental groups (in
\sisi{char.}{characteristic}~$0$ and in
\sisi{char.}{characteristic} $p>0$) have the same maximal quotient of
order prime to $p$.
\end{probleme}

Of course, \Ref{XIII.3.3} shows that in \Ref{XIII.3.1} one should also
replace $X'$ by a scheme of the form $X-Y$, where $Y$ is a closed subset
of $X$ that\pageoriginale has codimension $\geq 2$ in every component of $X$ that
contains it. When this restriction on $Y$ is dropped, there should still
be an analogous finiteness result, provided that \og tame\fg-type
restrictions are imposed on the ramification at the maximal points of
the irreducible components of $Y$ that have codimension~$1$.

\begin{probleme} \label{XIII.3.4}
Let $A$ be a complete noetherian local ring of dimension $2$, with
algebraically closed residue field. Again let $X=\Spec(A)$ and
$X'=X-\{\aaa\}$. Find special structural properties of $\pi_1(X')$ in
the case where $A$ is a complete intersection.
\end{probleme}

A satisfactory solution of this problem might make it possible to solve
the following old problem:

\refstepcounter{subsection}\label{XIII.3.5}
\begin{enonce*}{Conjecture 3.5\ndemark}
\ndetext{As of autumn 2004, this problem is still open.}
Find an irreducible curve in $\PP^3_k$ (where $k$ is an algebraically
closed field), preferably nonsingular, which is not set-theoretically the
intersection of two hypersurfaces.
\sisi{(KNESER~\cite{XIII.4} shows that it can always be obtained as the
intersection of three hypersurfaces).}{\ignorespaces}
\end{enonce*}
\sisi{\ignorespaces}{\noindent(Kneser~\cite{XIII.4} shows that it can
always be obtained as the intersection of three hypersurfaces).}

\section[\texorpdfstring{Problems related to higher $\pi_i$}{Problems related to higher pi\textunderscore i}]{Problems related to higher
$\pi_i$: local and global Lefschetz theorems for complex analytic
spaces\protect\ndemark} \label{XIII.4}

Let $X$ be a scheme\ndetext{The statements are made more precise in the
Comments (Section~\Ref{XIII.6}). The conjectures appearing there have
become theorems; \cf the footnotes to Section~\Ref{XIII.6}.} locally of
finite type over the field of complex numbers $\CC$. One can associate
with it a complex analytic space $X^h$, and hence homotopy and homology
invariants $\pi_i(X^h)$, $\H_i(X^h)$, $\H^i(X^h)$, etc.
\sisi{\dotsc}{} Moreover, $X$ is known to be connected if and only if
$X^h$ is connected, so there is a bijection
$$
\pi_0(X^h)\to \pi_0(X).
$$
Similarly, since every \sisi{\'etale covering}{\'etale cover} $X'$ of
$X$ defines an \sisi{\'etale covering}{\'etale cover} $X^{\prime h}$ of
$X^h$\pageoriginale, there is a canonical homomorphism
$$
\pi_1(X^h)\to \pi_1(X),
$$
which, by a theorem of Grauert--Remmert, is known to identify the second
group with the completion of the first for the topology defined by
subgroups of finite index. This simply expresses the fact that
$X'\sisi{\rightsquigarrow}{\mto}X^{\prime h}$ is an equivalence from the
category of \sisi{\'etale coverings}{\'etale covers} of $X$ to the
category of finite \sisi{\'etale coverings}{\'etale covers} of $X^h$.
It follows that the results of this seminar, obtained by purely
algebraic methods, concerning $\pi_0(X)$ and $\pi_1(X)$ imply results
for $\pi_0(X^h)$ and $\pi_1(X^h)$, which are transcendental in nature.
Moreover, if $X$ is proper, the well-known exact sequence
$0\to\ZZ\to\CC\to\CC^*\to0$ shows that the N\'eron--Severi group of $X$
(the quotient of its Picard group by the connected component of the
identity) is isomorphic to a subgroup of $\H^2(X^h,\ZZ)$; in the
nonsingular K\"ahler case, it is the subgroup denoted
$\H^{(1,1)}(X^h,\ZZ)$ (classes of type $(1,1)$):
$$
\Pic(X)/\Pic^0(X)\subset \H^2(X^h,\ZZ).
$$

Consequently, the information that we have obtained about Picard groups
implies information, admittedly very partial, about the groups
$\H^2(X^h,\ZZ)$. It is tempting to complete all these fragmentary
results with conjectures.

Very precise indications pointing in the same direction as those just
mentioned are supplied by a classical theorem of
\sisi{LEFSCHETZ}{Lefschetz}~\cite{XIII.7}. It asserts that if $X$ is a
nonsingular irreducible projective analytic space of dimension $n$, and
if $Y$ is a nonsingular hyperplane section, then the inclusion
$$
Y^{n-1}\to X^n
$$
induces\pageoriginale a homomorphism
$$
\pi_i(Y^{n-1})\to \pi_i(X^n)
$$
that is an \emph{isomorphism for} $i\leq n-2$ and an
\emph{epimorphism for} $i=n-1$. The analogous statement follows for the
homomorphisms
$$
\H_i(Y^{n-1})\to \H_i(X^n)
$$
on homology (integral homology, to be definite), while in cohomology
$$
\H^i(X^n)\to\H^i(Y^{n-1})
$$
is an isomorphism in degree $i\leq n-2$ and a monomorphism in degree
$i=n-1$. We have obtained variants of these results in the setting of
schemes, for $\pi_0$, $\pi_1$, and $\Pic$, valid moreover, to a large
extent, without nonsingularity hypotheses; \cf \Exp \Ref{XII}.
Furthermore, in removing the nonsingularity hypotheses, we made
essential use of \og local\fg variants of these global Lefschetz
theorems. All this suggests the following problems, which will doubtless
have to be attacked simultaneously\sfootnote{The formulations
\Ref{XIII.4.1} to \Ref{XIII.4.3} below are provisional. See Conjectures
\ref{XIII.6.A} to \ref{XIII.6.D} below, in the \og Comments on
\Exp \Ref{XIII}\fg, for more satisfactory formulations, and also
\Exp \Ref{XIV}.}.

\begin{probleme} \label{XIII.4.1}
Let $X$ be an analytic space and $Y$ a closed analytic subset of $X$
(or merely a closed subset?)\nde{The meaning of this question is not
clear; indeed, the statement of the problem itself does not seem to make
sense in this case, since the codimension of $Y$ in $X$ is not defined
when $Y$ is no longer assumed analytic.} such that for every $x\in Y$,
the local ring $\OXx$ is a complete intersection. Let $n$ be the complex
codimension of $Y$ in $X$. Is the canonical homomorphism
$$
\pi_i(X-Y)\to \pi_i(X)
$$
an isomorphism for $i\leq n-2$ and an epimorphism for $i=n-1$?
\end{probleme}

In this problem and the \sisi{preceding}{following} ones, a base point
is of course implicitly understood to have been chosen in order to
define the homotopy groups. To state the following problem, one must
define, for an analytic space $X$ (more generally, for a locally
path-connected space) and a point $x\in X$, local invariants\pageoriginale
$\pi_i^x(X)$\sfootnote{If $i\geq2$. For the case $i\leq1$, \cf the
\emph{Comments} in \numero 6 below, page~\pageref{p25y}.}.
\refstepcounter{toto}\label{pagexpo} To do this, choose a nonconstant map
$f$ from the interval $[0,1]$ to $X$ such that $f(0)=x$ and $f(t)\neq x$
for $t\neq0$; such a map exists if $x$ is not an isolated point. Then,
for every neighborhood $U$ of $x$, there is an $\varepsilon>0$ such that
$0<t<\varepsilon$ implies $f(t)\in U$, and the homotopy groups
$\pi_i(U-x,f(t))$ are essentially independent of $t$: as $t$ varies,
they are connected by a transitive system of isomorphisms. They may be
denoted by $\pi_i(U-x,f)$. Set\refstepcounter{toto}\label{pip15}\label{pXIII.15}
$$
\pi_i^x(X)=\varprojlim_U\pi_{i-1}(U-x,f),
$$
where the inverse limit is taken over the system of open neighborhoods
$U$ of $x$. Strictly speaking, this limit depends on $f$ and should be
denoted $\pi_i^x(X,f)$, but one verifies that as $f$ varies these groups
are isomorphic to one another\sfootnote{At least if $x$ does not locally
disconnect $X$; \cf the \emph{Comments} below, page~\pageref{p25y}.};
more precisely, they form a local system on the space of paths of the
kind under consideration issuing from $x$. These invariants are the
homotopical counterpart of the local cohomology invariants $\H_x^i(F)$
for a sheaf $F$ on $X$, introduced in \Ref{I}, and should play the role
of \emph{relative local homotopy groups} of $X$ modulo $X-x$. Their
vanishing for $i\leq n$ and every $x\in Y$, where $Y$ is a closed subset
of $X$ of topological dimension $\leq d$, should imply that the
homomorphisms
$$
\pi_i(X-Y)\to\pi_i(X)
$$
are bijective for $i<n-d$ and surjective for
$i=n-d$\sfootnote{For a corrected formulation, \cf the \emph{Comments}
below, page~\pageref{p25y}.}. From this point of view,
\Ref{XIII.4.1} would imply, when $Y$ is reduced to a point, a conjecture
of purely local nature, expressed by\refstepcounter{toto}
$$
\pi_i^x(X)=0\ \rm{for}\ i\leq n-1\label{p15l}
$$
when $X$ is a complete intersection of dimension $n$ at $x$.

\sisi{As}{As} an example of the local invariants $\pi_i^x(X)$, note that
if $x$ is a nonsingular point of \sisi{\ignorespaces}{$X$ of} complex
dimension $n$, then\pageoriginale
$$
\pi_i^x(X)=\pi_{i-1}(S^{2n-1}),
$$
where $S^{2n-1}$ denotes the sphere of dimension $2n-1$. In particular,
in this case $\pi_i^x(X)=0$ for $i\leq2n-1$, which corresponds to the
fact that if a closed subset $Y$ of codimension $\geq m$ is removed from
a topological \emph{manifold} $X$, then
$\pi_i(X-Y)\to\pi_i(X)$ is an isomorphism for $i\leq m-2$ and an
epimorphism for $i=m-1$.

This being said:

\begin{probleme} \label{XIII.4.2}
Let $X$ be an analytic space, let $x\in X$, let $t$ be a section of
$\OX$ vanishing at $x$, and let $Y$ be the zero set of $t$. Suppose that
the following conditions are satisfied:
\begin{enumeratea}
\item
$t$ is regular at $x$ (\ie it is not a zero divisor at $x$; this
hypothesis may moreover be superfluous).
\item
At the points $x'$ of $X-Y$ near $x$, $\Oo_{X,x'}$ is a complete
intersection. If \Ref{XIII.4.1} is true, this hypothesis should be
replaceable by the following more general one: for $x'$ as above,
$\pi_i^{x'}(X)=0$ for $i\leq n-1$.
\item
At the points $y$ of $Y-\{x\}$ near $x$, one has
$$
\prof\Oo_{X,y}\geq n
$$
(it suffices, for example, that $\prof\Oo_{X,x}\geq n$).
\end{enumeratea}

Under these conditions, is the canonical homomorphism
$$
\pi_i^x(Y)\to\pi_i^x(X)
$$
an isomorphism for $i\leq n-2$ and an epimorphism for $i=n-1$?
\end{probleme}

Here, finally, is a global variant of \Ref{XIII.4.2}, which should
follow from it by considering the affine cone at its vertex and
which would generalize the classical Lefschetz theorems:

\begin{probleme} \label{XIII.4.3}
Let\pageoriginale $X$ be a projective analytic space, equipped with an
ample invertible module $L$; let $t$ be a section of $L$, and let $Y$ be
the zero set of $t$. Suppose that:
\begin{enumeratea}
\item
$t$ is a regular section (\sisi{\ignorespaces}{a hypothesis} that may
be superfluous).
\item
For every $x\in X-Y$, $\Oo_{X,x}$ is a complete intersection. This
should be replaceable by $\pi_i^x(X)=0$ for $i\leq n-1$.
\item
For every $x\in Y$, $\prof\Oo_{X,x}\geq n$.
\end{enumeratea}

Under these conditions, is the homomorphism
$$
\pi_i(Y)\to\pi_i(X)
$$
an isomorphism for $i\leq n-2$ and an epimorphism for $i=n-1$?
\end{probleme}

We leave it to the reader to state analogous conjectures of a
\emph{cohomological} nature\sfootnote{\Cf \Exp \Ref{XIV} for the
corresponding results in the theory of schemes.}, with hypotheses and
conclusions concerning local cohomological invariants (with coefficients
in a given group). In any event, the key result seems to be
\Ref{XIII.4.2}, with hypothesis b) taken in the more general form just
indicated, whether one adopts the point of view of homology or that of
homotopy.

We have stated these conjectures in the transcendental setting in the
hope of interesting topologists in them and convincing them that
\og Lefschetz\fg-type questions are far from closed. Of course, now that
we are about to possess a good theory of the cohomology of schemes (with
finite coefficients), thanks to the recent work of M.~\sisi{ARTIN}{Artin},
the same questions arise in the setting of schemes, and it is difficult
to doubt that they will receive an affirmative answer in the near
future\sfootnote{See the preceding note.}.

\section{Problems related to local Picard groups} \label{XIII.5}

A\pageoriginale first fundamental problem, pointed out for the first time by \sisi{MUMFORD}{Mumford}~\cite{XIII.5} in a particular case, is the following. Let $A$ be a complete local ring with residue field $k$, let $X=\Spec(A)$, and let $U=\Spec(A)-\{a\}$, where $a$ is the maximal ideal of $A$, \ie the closed point of $\Spec(A)$. The problem is to construct a strict projective system $G$ of locally algebraic groups $G_i$ over $k$, and a natural isomorphism
\begin{equation*} \label{eq:XIII.5.+} \tag{$+$} {\Pic(U)\simeq G(k)} \end{equation*}
where, of course, we set $G(k)=\varprojlim G_i(k)$. Heuristically, the problem is to \og put an algebraic-group structure\fg (or at least a pro-algebraic one, in a suitable sense) on the group $\Pic(\sisi{X'}{U})$.

It is clear that, as stated, the problem is not sufficiently precise, since the datum of an isomorphism~(\Ref{eq:XIII.5.+}) is far from characterizing the pro-object $G$. If $A$ contains a subfield, again denoted by $k$, that is a coefficient field, the problem can be made more precise by requiring that, for a varying extension $k'$ of $k$, there be an isomorphism, functorial in $k'$:
\begin{equation*} \label{eq:XIII.5.+'} \tag{$+'$} {\Pic(U')\simeq G(k')} \end{equation*}
where $U'$ is the open subset analogous to $U$ in $\Spec(A')$, with $A'=A\hat{\otimes}_k k'$. One can proceed analogously even if $A$ has no coefficient field, provided that $k$ is perfect, since one can then construct an $A'$ functorially \og by residue-field extension $k'/k$\fg. Moreover, when $A$ admits a coefficient field, the algebraic structure obtained on $\Pic(U)$ will depend essentially on the choice of that coefficient field (as one sees already for the projective cone over an elliptic curve); it therefore seems necessary to start from a \og pro-algebraic ring over $k$\fg in the sense of \sisi{GREENBERG}{Greenberg}~\cite{XIII.3}, in order to define the pro-object $G$. It is moreover conceivable that, when no coefficient field is given, one obtains only a projective system of\pageoriginaled \emph{quasi}-algebraic groups in the sense of Serre, or rather quasi-locally algebraic groups (the groups $G_i$ obtained will not in general be of finite type over $k$, but only locally of finite type over $k$). It is even possible that in general one obtains only a still weaker structure on $\Pic(U)$, of the kind encountered by \sisi{NERON}{N\'eron}~\cite{XIII.6} in his theory of degeneration of abelian varieties defined over local fields.

One method of attacking the problem, likewise introduced by \sisi{MUMFORD}{Mumford}, consists in resolving the singularities of $X$, \ie in considering a projective birational morphism $Y\to X$ with $Y$ regular. When $U$ is regular (\ie $a$ is an isolated singular point), one can often find $Y$ such that $Y\sisi{|U}{_|U}=V\to U$ is an isomorphism. In that case one therefore has
$$
\Pic(U)\simeq\Pic(V)\simeq\Pic(Y)/\Im \ZZ^I,
$$
where $I$ is the set of irreducible components of the fiber $Y_{a}$\sisi{, }{} (each of which defines an element of $\Pic(Y)$, since it is a locally principal divisor, by the regularity of $Y$). On the other hand, using the techniques of Formal Geometry in \EGA III~4~and~5, in particular the existence theorem, one finds
$$
\Pic(Y)\simeq\varprojlim\Pic(Y_n),
$$
where $Y_n=Y\otimes_A A_n$ and $A_n=A/\mm^{n+1}$. When $A$ admits a coefficient field $k$, the theory of Picard schemes of the projective $k$-schemes $Y_n$ is available, and hence
$$
\Pic(Y_n)\simeq \mathbf{Pic}_{Y_n/k}(k).
$$

This therefore gives a construction of a projective system of locally algebraic groups $\mathbf{Pic}_{Y_n/k}/\Im \ZZ^I$, which is the desired system.\refstepcounter{toto}\nde{\label{Boutot}The question was greatly clarified by the results of {Boutot} (Boutot~J.-F., \emph{Sch\'ema de {P}icard local}, Lect. Notes in Math., vol.~632, Springer, Berlin, 1978). In particular, if $A$ is a complete (noetherian) local $k$-algebra of depth $\geq 2$ such that $H^2_\m(A)$ has finite dimension over $k$, the local Picard group is a group scheme locally of finite type over $k$, whose tangent space at the origin is $H^2_\m(A)$. If $A$ is moreover normal of dimension $\geq 3$, Serre's normality criterion~\Ref{XI}~\Ref{XI.3.11}, together with Corollary~\Ref{V}~\Ref{V.3.6}, ensures the required finiteness and hence the existence of the local Picard scheme. See also (Lipman~J., {\og The Picard group of a scheme over an Artin ring\fg}, \emph{Publ. Math. Inst. Hautes \'Etudes Sci.} \textbf{46} (1976), p\ptbl 15--86) for an approach closer to Grothendieck's outline above.} In the case considered here, one can moreover see (using the fact that $a$ is an isolated singular point) that the connected components of the universal-image subgroups\pageoriginale in this projective system form an \emph{essentially constant} projective system; thus, in this case, one obtains a locally algebraic group $G$ as a solution of the problem. If one further assumes that $A$ is normal of dimension $2$, then a remark of \sisi{MUMFORD}{Mumford} (stating that the intersection matrix of the components of $Y_{a}$ in $Y$ is negative definite\nde{Mumford~D., {\og The topology of normal singularities of an algebraic surface and a criterion for simplicity\fg}, \emph{Publ. Math. Inst. Hautes \'Etudes Sci.} \textbf{9} (1961), p\ptbl 5--22.}) implies that $G$ is even an algebraic group, \ie of finite type over $k$ (the number of its connected components being equal, moreover, to the determinant of the intersection matrix just considered).

If, on the contrary, $a$ is not an isolated singularity, examples (with $A$ of dimension $2$) show that one obtains a projective system of algebraic groups which does not reduce to a single algebraic group.

\refstepcounter{toto}\label{pXIII20}%
Once a good notion of \og \emph{local Picard scheme}\fg were available, one should strengthen the notion of parafactoriality by saying that $A$ is \og \emph{geometrically parafactorial}\fg when not only $A$, and even $\widehat{A}$, is parafactorial, but the local Picard scheme $G(\widehat{A})$ is the trivial group (which, when the residue field is not algebraically closed, is stronger than saying that $G$ has no rational point over $k$ other than the identity). The necessity of a strengthened notion of parafactoriality becomes apparent on recalling that there exist complete normal local rings of dimension $2$ which are factorial but admit finite \'etale algebras which are not\nde{Factorial rings with nonfactorial henselization arise naturally in the study of moduli spaces of vector bundles. See, for example, (Dr\'ezet~J.-M., {\og Groupe de Picard des vari\'et\'es de modules de faisceaux semi-stables sur $\mathbf{P}\sb 2$\fg}, in \emph{Singularities, representation of algebras, and vector bundles (Lambrecht, 1985)},
Lect. Notes in Math., vol.~1273, Springer, Berlin, 1987, p\ptbl 337--362). \textit{Strictly speaking}, Dr\'ezet shows that the completion is not factorial, but in fact the proof gives the result for the henselization: the point is that Luna's \'etale slice theorem (Luna~D., {\og Slices \'etales\fg}, in \emph{Sur les groupes alg\'ebriques}, M\'em. Soc. math. France, vol.~33, Soci\'et\'e math\'ematique de France, Paris, 1973, p\ptbl 81--105) describes the local ring of a quotient in the sense of invariant theory near a semistable point, locally for the \'etale topology.}. A \og \emph{geometrically factorial}\fg local ring would then be a normal ring~$A$ such that all \sisi{t}{th}e localizations of dimension $\geq 2$ are geometrically parafactorial, or better, such that the localizations of $\widehat{A}$ are parafactorial\sfootnote{For a more flexible notion of \og geometrically factorial\fg local ring, \cf \emph{Comments}, page~\pageref{commpara}.}. Of course, it would be interesting to find a \og good\fg definition of these notions, independent of the still undeveloped theory of local Picard schemes.\nde{See page~\pageref{commpara}: a local ring is geometrically factorial (\resp parafactorial) if its strict henselization is factorial (\resp parafactorial).}

It is in any event plausible that these notions will be needed if one wishes to obtain statements of the following type: Let $A$ be a \og good ring\fg (for example an algebra of finite type over $\ZZ$, or over a complete local ring, for example over a field). Let~$U$ be the set of $x\in X=\Spec(A)$ such that $\OXx$ is \og geometrically\pageoriginale factorial\fg; is $U$ open\sisi{?.}{?} Or again: Let $f\colon X\to Y$ be a flat morphism of finite type with $Y$ locally noetherian, and let $U$ be the set of $x\in X$ such that $\Oo_{X_{f(x)}, x}$ is \og geometrically factorial\fg; is $U$ open, at least under suitable additional hypotheses on $f$\sisi{?.}{?} I doubt that true statements of this type exist with the usual notion of factorial ring.

We have raised here, in a particular case, the question of studying the geometric properties of \og varying\fg local rings, for example the $\OXx$ as $x$ ranges over a prescheme $X$. When $X$ is, for example, a scheme of finite type over a field, it is known\nde{See \EGA IV.16.} that there exists on $X$ a projective system of finite algebras $P^n_{X/k}$ (obtained by completing $X\times_k X$ along the diagonal), whose fiber at every point $x\in X$ rational over $k$ is isomorphic to the projective system of the $\OXx/\mm_X^{n+1}$. It is then natural to relate the study of the completions of the local rings $\OXx$, for varying $x$, to that of the \og algebraic family of complete local rings\fg given by the $P^n$, noting that for every $x\in X$ (rational over $k$ or not), one obtains a complete local ring
$$
P^{\infty}(x)=\varprojlim P^n(x)
$$
(where $P^n(x)=$~reduced~fiber~$P^n\otimes_{\OXx}k(x)$). Of particular interest, for example, will be the complete ring thus associated with the generic point, and one expects its algebro-geometric properties (expressed, for example, by its local Picard, homotopy, or homology groups) to be essentially those of the completions $\widehat{\Oo}_{X,x}$ for $x$ in a suitable dense open subset $U$.

More generally, one may undertake the simultaneous study of the complete local rings obtained in this way from an adic projective system $(P_n)$ of finite algebras over a given scheme $X$. It is plausible that, under certain regularity conditions (such as flatness of the $P_n$),\pageoriginale the local homotopy groups arise from a projective system of finite group schemes over $X$, and that analogous results hold for the local Picard groups. As regards the \sisi{latter}{latter}, a first interesting case deserving investigation is that in which one starts with an algebraic surface $X$ having singular curves, and seeks to study the local Picard schemes at varying points of these curves in terms of a suitable pro-group scheme defined over the singular locus.

\section[Comments]{Comments\protect\sfootnotemark} \label{XIII.6}

The\pageoriginale viewpoint of\sfootnotetext{Written in March~1963.} the \og \'etale cohomology\fg of schemes, and recent progress in this theory, lead\sisi{}{s} us both to make some of the problems posed more precise and to broaden them. For the notion of \og topology\fg and of the \og \'etale topology of a scheme\fg, I refer to M\ptbl\sisi{ARTIN}{Artin}, \sisi{Grothendieck Topologies}{\emph{Grothendieck Topologies}}, Harvard University 1962 (mimeographed notes)\sfootnote{Or preferably to \SGA 4.}.

By providing a finer notion of localization than that furnished by the traditional \og Zariski topology\fg, this theory gives particular importance to \emph{strictly local rings}\refstepcounter{toto}\label{pXIII.6}, \ie henselian local rings with separably closed residue field. For every local ring $A$ with residue field $k$, and every separable closure $k'$ of $k$, one can find a local homomorphism from $A$ to a strictly local ring $A'$, the \emph{strict local closure} of $A$, with residue field $k'$, having an evident universal property. The ring $A'$ is henselian and flat over $A$, and $A'\otimes_A k\simeq k'$; it is noetherian if and only if $A$ is. (\Cf \loccit Chap\ptbl III, section 4)\sfootnote{Or \EGA IV~18.8.}. If $X$ is a prescheme, $x$ a point of $X$, and $x'$ a point above $x$ which is the spectrum of a separable closure $k'$ of $k=k(x)$, one is led to define the \emph{strictly local ring} of $X$ at $x'$, $\Oo'_{X,x'}$, as the strict local closure of the usual local ring $\OXx$, relative to the residue-field extension $\sisi{K'/K}{k'/k}$. It is the \emph{strictly} local rings at the \og geometric\fg points of $X$ which, from the viewpoint of the \'etale topology, are meant to reflect the local properties of the prescheme $X$. They also play, in many respects, the role formerly assigned to the \emph{completions} of the local rings of $X$ (say, at points with algebraically closed residue field), while remaining \og closer\fg to $X$ and allowing an easier passage to \og neighboring points\fg.

One should therefore reconsider for noetherian henselian local rings (\resp noetherian strictly local rings) many questions generally posed for complete local rings (possibly restricted\pageoriginale to those having an algebraically closed residue field). Thus the topological problems raised in n$^{\textup{os}}$~\Ref{XIII.2} and \Ref{XIII.3} arise more generally for strictly local rings. One can moreover state conjecturally, for \og good\fg strictly local rings, certain simple-connectedness and acyclicity properties for the geometric fibers of the canonical morphism $\Spec(\widehat{A})\to\Spec(A)$; these would show that, for many properties of a \og topological\fg nature, it amounts to the same thing to prove them for the ring $A$ or for its completion $\widehat{A}$. Certain results already obtained in this direction\sfootnote{\Cf M\ptbl\sisi{ARTIN}{Artin} in \SGA 4~XIX} give reason to hope that complete results in this direction will soon be available.

\refstepcounter{toto}\label{commpara}%
The notion of \'etale localization provides a seemingly reasonable definition of a \og \emph{geometrically parafactorial}\fg or \og \emph{geometrically factorial}\fg local ring\sisi{,}{} (the need for which was indicated in \numero\Ref{XIII.5}, p\ptbl\sisi{20}{\pageref{pXIII20}}): this will mean a local ring whose strict local closure is parafactorial, \resp factorial. Hypotheses of this nature arise naturally in the study of the \'etale cohomology of preschemes\nde{See, for example, (Strano~R., {\og The Brauer group of a scheme\fg}, \emph{Ann. Mat. Pura Appl. (4)} \textbf{121} (1979), p\ptbl 157--169), where the hypothesis that the local rings of a scheme $X$ are geometrically parafactorial sometimes makes it possible to prove that the Brauer groups of $X$ (computed in terms of Azumaya algebras) coincide with the cohomological Brauer group of $X$.}. Thus, if $X$ is a locally noetherian prescheme whose strictly local rings are factorial (\ie whose ordinary local rings are \og geometrically factorial\fg), one proves that the $\H^i(X_{\et},\GG_m)$ are torsion groups for $i\geq 2$ (which sometimes makes it possible to express these groups in terms of cohomology groups with coefficients in the groups $\bbmu_n$ of $n$\sisi{th}{-th} roots of unity), and if $X$ is integral with fraction field $K$, the natural homomorphism $\H^2(X_{\et},\GG_m)\to\H^2(K,\GG_m)=\Br(K)$ is injective\nde{The connection between the Brauer group and the Picard group is a close one. In this connection, let us cite the following results of Saito (Saito~S., {\og Arithmetic on two-dimensional local rings\fg}, \emph{Invent. Math.} \textbf{85} (1986), \numero 2, p\ptbl 379--414) in the case of surfaces, the first local and the other global. Let $A$ be an excellent normal henselian local ring of dimension $2$ with \textit{finite} residue field, and let $X$ be the complement of the closed point in $\Spec(A)$. Then there is a perfect duality of \emph{torsion} groups $\Pic(X)\times\Br(X)\to\QQ/\ZZ$ --- by the Brauer group of $X$ we mean the \emph{cohomological} Brauer group $\textrm{Br}(X)=H^2_\et(X,G_m)$. In the global case one has the following generalization of a result of Lichtenbaum (Lichtenbaum~S., {\og Duality theorems for curves over $p$-adic fields\fg}, \emph{Invent. Math.} \textbf{7} (1969), p\ptbl 120--136): let $k$ be the fraction field of a complete discrete valuation ring $\Oo$ with finite residue field, and let $X$ be a projective, smooth, and geometrically complete curve over $k$. The group $\Pic^0(X)$ is endowed with the topology induced by the adic topology of $k$, and $\Pic(X)$ is the topological group for which $\Pic^0(X)$ is an open subgroup. Then there is a perfect duality of topological groups $\Pic(X)\times\Br(X)\to\QQ/\ZZ$. Note that this statement, which concerns curves, is of course proved by considering a regular (proper and flat) model of $X$ over $\Oo$: it is a result about surfaces.}; examples show that these conclusions may fail, even for local $X$, if one assumes only that $X$ is factorial rather than geometrically factorial\sfootnote{\Cf A\ptbl \sisi{GROTHENDIECK}{Grothendieck}, le groupe de Brauer~II (S\'eminaire Bourbaki \numero 297, Nov.~1965), notably 1.8 and 1.11~b.}.

As regards the problems of local and global Lefschetz type raised in \Ref{XIII.3.4}, and their analogues in scheme theory, the homological version of these questions has been considerably clarified: everything follows formally from three general theorems. The first concerns the cohomological dimension of\pageoriginale certain affine schemes (\resp Stein spaces), such as affine schemes $X$ of finite type over an algebraically closed field: their cohomological dimension is $\leq\dim X$ (\og affine Lefschetz theorem\fg)\sfootnote{\Cf \SGA 4~XIV.}; the second is a duality theorem for the cohomology (with discrete coefficients) of a projective morphism\sfootnote{\Cf \SGA 4~XVIII.}; and the last is a \emph{local duality} theorem of an analogous nature\sfootnote{\Cf \SGA 5~I.}. In Algebraic Geometry, only the last of these has not been proved at the time of writing (it is, however, proved in \sisi{char.~$0$}{characteristic $0$}, using resolution of singularities by \sisi{HIRONAKA}{Hironaka}). Moreover, in the transcendental setting, global and local duality are already available, having recently been proved by \sisi{VERDIER}{Verdier}\nde{See Verdier~J.-L., {\og Dualit\'e dans la cohomologie des espaces localement compacts\fg}, in \emph{S\'eminaire Bourbaki}, vol.~9, Soci\'et\'e math\'ematique de France, Paris, 1995, \Exp 300, p\ptbl 337--349.}. Let us merely indicate that, in the statements of the homological versions of Problems \Ref{XIII.4.2} and \Ref{XIII.4.3} (which henceforth deserve the name of conjectures), the \og at infinity\fg conditions a) and c) are certainly superfluous; only the \emph{local cohomological structure} of $X-Y$ is important, and one may assume, for example, that it is locally a complete intersection of dimension $\geq n$. Moreover, in \Ref{XIII.4.3}, say, the fact that $Y$ is a hyperplane section should play no role and should be replaceable by the sole hypothesis that $X$ is compact and $X-Y$ is Stein (\ie in Algebraic Geometry, $X$ is proper over $k$ and $X-Y$ is affine; as we said, the homological version of this conjecture is proved for algebraic spaces over the field $\CC$)\sfootnote{\Cf \Exp \Ref{XIV}.}.

\refstepcounter{toto}\label{p25y}%
In the definition (p\ptbl\sisi{15}{\pageref{pip15}}) of the $\pi_i^x(X)$, one must assume $i\geq 2$. For $i=0,1$, there is no reasonable definition of the $\pi_i^x(X)$\sisi{,}{;} they should be replaced by $\H_0^x(X)$ and $\H_1^x(X)$, defined respectively as the co\sisi{-}{}kernel and kernel in the natural homomorphism
$$
\varprojlim \H_0(U-(x),\ZZ)\to\varprojlim\H_0(U,\ZZ).
$$

\sisi{At}{At} a pinch, and for convenience of formulation, one may set $\pi_i^x(X)=\H_i^x(X)$ for $i\leq 1$; otherwise, subsequent assertions concerning the $\pi_i^x$ must be supplemented by the corresponding assertions for $\H_0^x,\H_1^x$. If $x$ is an isolated point of $X$, one should set $\pi_i^x(X)=0$ for $i\neq 0$ and $\pi_0^x(X)=\H_0^x(X)=\ZZ$.

The assertion\pageoriginale that the $\pi_i^x(U,f)$ are mutually isomorphic is true only when $X$ is not disconnected by $x$ in a neighborhood of $x$, \ie if $\pi_i^x(X)=0$ for $i=0,1$. In the general case $\pi_i^x(X)$ can designate only a \emph{family} of groups, not necessarily mutually isomorphic; nevertheless the notation $\pi_i^x(X)=0$ retains an evident meaning.

\refstepcounter{toto}\label{p26y}\sisi{Page~15, line~-~11}{Page~\pageref{p15l}}, where I predict that the vanishing of the local homotopy invariants $\pi_i^x(X)$ for $x\in Y$, $i\leq n$, should imply that $\pi_i(X-Y)\to\pi_i(X)$ is bijective for $i<n-d$ and surjective for $i=n-d$, one must be cautious, since in the present context (as in Algebraic Geometry) one does not have at one's disposal \og general\fg points at which the local conditions would also have to apply. For this reason it will doubtless be necessary to use relative local homotopy invariants
$$
\pi_i^Y(X,f)=\pi_i^Y(X,x)=\varprojlim_U\pi_{i-1}(U-U\cap Y,f(t))\qquad\text{for }i\geq 2\quoi,
$$
(with an {\sisi{ad-hoc}{\textit{ad hoc}}} definition as above for $i=0,1$), where $Y$ is a closed subset of~$X$; or else to compensate for the absence of general points by expressing the hypotheses on $X$ in terms of properties, topological in nature (for the \'etale topology), of the spectra of the local rings of $X$, which makes it possible to recover general points. The same reservation applies to the generalization of Conjectures \Ref{XIII.4.2} and \Ref{XIII.4.3} to the case in which $X-Y$ is not assumed to be locally a complete intersection, a generalization suggested in the statements of conditions~b) of those conjectures.

To formulate the expurgated versions of Conjectures \Ref{XIII.4.2} and \Ref{XIII.4.3} suggested by the results mentioned above, we make the following

\begin{supdefinition} \label{XIII.6.1}
Let $X$ be a topological space, let $Y$ be a locally closed subset of $X$, and let $n$ be an integer. One says that $X$ has \emph{homotopical depth} $\geq n$ along $Y$, and writes $\mathrm{prof\:hpt}_Y(X)\geq n$, if for every $x\in Y$ one has $\pi_i^Y(X,x)=0$ for $i<n$.
\end{supdefinition}

It should\pageoriginale be equivalent to say that, for every open subset $X'$ of $X$ and every $x\in X'\cap U=U'$ (where $U=X-Y$), the canonical homomorphism
$$
\pi_i(U',x)\to\pi_i(X',x)
$$
is an isomorphism for $i<n-1$ and a monomorphism for $i=n-1$\nde{When the pair $(X,Y)$ is moreover polyhedral, this equivalence holds; \cf (Eyral~C., {\og Profondeur homotopique et conjecture de Grothendieck\fg}, \emph{Ann. Sci. \'Ec. Norm. Sup. (4)} \textbf{33} (2000), \numero 6, p\ptbl 823--836).}.

\begin{supdefinition} \label{XIII.6.2}
Let $X$ be a complex analytic space and let $n$ be an integer. One says that the \emph{rectified homotopical depth} of $X$ is $n$\sfootnote{In the first edition of these notes, we used the term \og true homotopical depth\fg. In the present version we follow \EGA IV~10.8.1.} if, for every locally closed analytic subset $Y$ of $X$, one has
\begin{equation*} \label{eq:XIII.6.x} \tag{$x$} {\mathrm{prof\:hpt}_Y(X)\geq n-\dim Y}
\end{equation*}
(where, of course, $\dim Y$ denotes the \emph{complex} dimension of $Y$).
\end{supdefinition}

It should be equivalent to say that, for every irreducible locally closed analytic subset $Y$ of $X$, there exists a closed analytic subset $Z$ of $Y$, of dimension $<\dim Y$, such that relation~(\Ref{eq:XIII.6.x}) holds with $Y-Z$ in place of $Y$. This would make it possible, for example, to restrict in Definition~\Ref{XIII.6.2} to the case where $Y$ is nonsingular.\refstepcounter{toto}\nde{\label{HL}All the conjectures that follow, suitably rectified, if I may say so, became theorems through the work of Hamm and L\^e D\~ung Tr\'ang (Hamm~H.A. \& L\^e D\~ung Tr\'ang, {\og Rectified homotopical depth and Grothendieck conjectures\fg}, in \emph{The Grothendieck Festschrift, Vol. II}, Progr. Math., vol.~87, Birkh\"auser, Boston, 1990, p\ptbl 311--351), cited as \cite{HL} below. As for the two conjecturally equivalent definitions of rectified depth, they are in fact equivalent to a third, expressed in terms of Whitney stratification (\cf \loccit, Theorem 1.4).}

\enlargethispage{\baselineskip}%
The following conjecture, purely topological in nature, is in the spirit of a \og \emph{local Hurewicz theorem}\fg.

\refstepcounter{aconjecture} \label{XIII.6.A}
\begin{enonce*}{Conjecture A}[\og local Hurewic\sisi{s}{z} \sisi{th.}{theorem}\fg\ndemark]\ndetext{As observed in \cite{HL}, Example 3.1.3, this conjecture is already false for $X=\{\boldsymbol{z}\in\nobreak\CC^n\mid\nobreak z_1^2+z_2^3+\cdots+z_n^3=0\}, n\geq 4$, and $Y$ reduced to the origin. Suitably modified, however, it is true (Theorem 3.1.4 of {\loccit}).}
Let $X$ be a topological space and $Y$ a locally closed subset, subject if necessary to \og smoothness\fg conditions such as local triangulability of the pair $(X,Y)$, and let $n$ be an integer $\geq 3$. In order that $\mathrm{prof\:hpt}_Y(X)\geq n$, it is necessary and sufficient that\pageoriginale
$$
\H^i_Y(\ZZ_X)=0\quad\text{for }i<n
$$
(one then says that $X$ has \emph{cohomological depth} $\geq n$ along $Y$), and that the local fundamental groups
$$
\sisi{\pi^2_Y}{\pi_2^Y}(X,x)=\varprojlim_{U\ni x}\pi_1(U-U\cap Y)
$$
vanish (one then says that $X$ is \og pure\fg along $Y$).
\end{enonce*}

Note that if $X$ is an analytic space, $Y$ an analytic subspace, and $X$ is pure along $Y$, then for every $x\in Y$ the local ring $\OXx$, as well as its localizations with respect to prime ideals containing the ideal defining the germ $Y$ at $x$ (\ie in the inverse image $Y_x$ of $Y$ under $\Spec(\OXx)=X_x\to X$), are pure in the sense of \Exp \Ref{X}; it seems plausible that the converse is also true. Analogous remarks apply to cohomological depth, it being understood that one works with the \'etale topology on the $\Spec(\OXx)$.

Conjecture \Ref{XIII.4.1} then generalizes to

\refstepcounter{aconjecture} \label{XIII.6.B}
\begin{enonce*}{Conjecture B}[\og Purity\fg\ndemark]\ndetext{This conjecture is proved, even when $E$ is singular, in \cite{HL}: it is Theorem 3.2.1.}
Let $E$ be an analytic space and $X$ an analytic subset of $E$. Suppose that $E$ is nonsingular of dimension $N$ at $x\sisi{}{\in X}$, and that $X$ can be described by $p$ analytic equations in a neighborhood of every point. Then the rectified homotopical depth of $X$ is $\geq N-p$.
\end{enonce*}

In particular, a local complete intersection of dimension $n$ at every point would have rectified homotopical depth $\geq n$, which is precisely Conjecture \Ref{XIII.4.1}.

Conjectures \Ref{XIII.4.2} and \Ref{XIII.4.3} generalize respectively to:

\refstepcounter{aconjecture} \label{XIII.6.C}
\begin{enonce*}{Conjecture C}[\og Local Lefschetz\fg\ndemark]\ndetext{This conjecture is proved in \cite{HL}, indeed even in the strong form of the remark that follows; \cf Theorem 3.3.1 of~\loccit.}
Let $X$ be an analytic space, let $Y$ be a closed analytic subset, and let $x$ be a point of $Y$. Suppose that $X-Y$ is Stein in a neighborhood of $x$ (for example, that $Y$ is defined by one equation at $x$), and that $X-Y$ has rectified homotopical depth $\geq n$ in a neighborhood of $x$ (for example, that at every point of $X-Y$ near $x$ it is a complete intersection of dimension $\geq n$; \cf Conjecture \Ref{XIII.6.B}). Then the canonical homomorphism
$$
\pi_i^x(Y)\to\pi_i^x(X)
$$
is an isomorphism for $i<n-1$ and an epimorphism for $i=n-1$.
\end{enonce*}

\refstepcounter{aconjecture} \label{XIII.6.D}
\begin{enonce*}{Conjecture D}[\og Global Lefschetz\fg\ndemark]
\ndetext{This conjecture is again proved in \cite{HL}, indeed even in the strong form of the remark that follows; \cf Theorem 3.4.1 of~\loccit.}
Let\pageoriginaled $X$ be a \emph{compact} analytic space and let $Y$ be an analytic subspace of $X$ such that $U=X-Y$ is Stein and has rectified homotopical depth $\geq n$ (for example, is a complete intersection of dimension $\geq n$ at every point). Then the canonical homomorphism
$$
\pi_i(Y)\to\pi_i(X)
$$
is an isomorphism for $i<n-1$ and an epimorphism for $i=n-1$.
\end{enonce*}

\begin{remarquestar}
When, in the statements \Ref{XIII.6.C} and \Ref{XIII.6.D}, the hypothesis that $X-Y$ is Stein is replaced by the hypothesis that $X-Y$ is a union of $c+1$ Stein open subsets (which plays the role of a topological \og concavity\fg hypothesis), the conclusions should simply be modified by replacing $n$ there by $n-c$.\nde{Finally, let us mention the following result of Fulton, to be compared with the Fulton--Hansen result cited in the editor's note~\eqref{FH} on page~\pageref{FH}: let $X$ and $H$ be closed subschemes of $\PP^m_\CC$, let $n$ be the dimension of $X$, and let $d$ be the codimension of $H$. Then the map
$$
\pi_i(X,X\cap H)\to\pi_i(\PP^m_\CC,H)
$$
is an isomorphism if $i\leq n-d$ and is surjective if $i=n-d-1$; see (Fulton~W. \& Lazarsfeld.~R., {\og Connectivity and its applications in algebraic geometry\fg}, in \emph{Algebraic geometry (Chicago, Ill., 1980)}, Lect. Notes in Math., vol.~862, Springer, Berlin--New York, 1981, p\ptbl 26--92).}
\end{remarquestar}

Finally, let us make explicit, in the \og global case\fg \Ref{XIII.6.D}, the conjecture concerning the fundamental group (obtained by taking $n=3$):

\refstepcounter{aconjecture}\label{XIII.6.D'}
\begin{enonce*}{Conjecture D$\boldsymbol'$}[Global Lefschetz for the fundamental group\ndemark]
\ndetext{This conjecture is proved in \cite{HL}; \cf Theorem 3.5.1 of~\loccit.}
Let $X$ be a compact analytic space over the field of complex numbers and let $Y$ be a closed analytic subset such that $U=X-Y$ is Stein. Suppose moreover that the following conditions hold:
\begin{enumeratei}
\item
For every $x\in U$, the local fundamental group $\pi_2^x(X,x)$ vanishes (\ie $X$ \emph{is \og pure at }$x$\fg), or merely the local ring $\OXx$ is pure.
\item
The local rings at the points of $U$ are \og connected in dimension $\geq 2$\fg.
\item
The local rings at the points of $U$ have dimension $\geq 3$.
\end{enumeratei}

Under these conditions, for every $x\in Y$, the homomorphism
$$
\pi_1(Y,x)\to\pi_1(X,x)
$$
is\pageoriginale an isomorphism (and $\pi_2(Y,x)\to\pi_2(X,x)$ is an epimorphism).
\end{enonce*}

Note that the local conditions (i), (ii), and (iii) on $U$ hold if $U$ is locally a complete intersection of dimension $\geq 3$. From the viewpoint of Algebraic Geometry (when $U$ arises from a scheme, again denoted by $U$), conditions (i) through (iii) correspond to hypotheses on the local invariants $\pi_i^x(U)$, \sisi{\ignorespaces}{namely} \hbox{$\pi_i^x(U)=0$} for $i<3-\degtr k(x)/k$, at points $x$ for which one has respectively $\degtr k(x)/k=0,1,2$. The global condition on $U$ ($U$ Stein) is satisfied if $X$ is projective and $Y$ is a hyperplane section.

\begin{thebibliography}{HL}\refstepcounter{toto}\pageoriginale\label{XIII.7}

\bibitem[1]{XIII.1} \sisi{%
S. Abhyankar, Local uniformisation on algebraic surfaces, Annals of Math. vol.~63, \No~3, 1956}
{%
\textsc{S. Abhyankar} -- {\og Local uniformisation on algebraic surfaces over ground fields of characteristic $p\neq0$\fg}, \emph{Annals of Math.} \textbf{63} (1956), p\ptbl 491--526.
}

\bibitem[2]{XIII.2} \sisi{%
W.L. Chow, On the theorem of Bertini for local domains, Proc. Nat. Acad. Sc. Vol.~44, \No~6, pp\ptbl 580--584, 1958}
{%
\textsc{W.L. Chow} -- {\og On the theorem of Bertini for local domains\fg}, \emph{Proc. Nat. Acad. Sci. U.S.A.} \textbf{44} (1958), p\ptbl 580--584.
}

\bibitem[3]{XIII.3} \sisi{%
M. Greenberg, Schemata over local rings, Annals of Math., vol~73, 1961}
{%
\textsc{M. Greenberg} -- {\og Schemata over local rings\fg}, \emph{Annals of Math.} \textbf{73} (1961), p\ptbl 624--648.
}

\bibitem[4]{XIII.4} \sisi{%
M. Kneser, Uber die Darstellung algebraischer Raumkurven als Durchschnitte von Fl\"achen, Archiv der Math. Vol.~XI, fasc.~3, 1960}
{%
\textsc{M. Kneser} -- {\og \"Uber die Darstellung algebraischer Raumkurven als Durchschnitte von Fl\"achen\fg}, \emph{Archiv der Math.} \textbf{XI} (1960), p\ptbl 157--158.
}

\bibitem[5]{XIII.5} \sisi{%
D. Mumford, The topology of normal singularities of an algebraic surface, and a criterion for simplicity, Publications Math\'ematiques \No~9, 1961}
{%
\textsc{D. Mumford} -- {\og The topology of normal singularities of an algebraic surface, and a criterion for simplicity\fg}, \emph{Publ. Math. Inst. Hautes \'Etudes Sci.} \textbf{9} (1961), p\ptbl 5--22.
}

\bibitem[6]{XIII.6} \sisi{%
A. N\'eron, Mod\`eles minimaux des vari\'et\'es ab\'eliennes sur les corps locaux et globaux, to appear in Publications Math\'ematiques (available as mimeographed notes in limited circulation...)}
{%
\textsc{A. N\'eron} -- {\og Mod\`eles minimaux des vari\'et\'es ab\'eliennes sur les corps locaux et globaux\fg}, \emph{Publ. Math. Inst. Hautes \'Etudes Sci.} \textbf{21} (1964), p\ptbl 5-128.
}

\bibitem[7]{XIII.7} \sisi{%
A.H. Wallace, Homology Theory of algebraic Varieties, Pergamon Press, 1958.}
{%
\textsc{A.H. Wallace} -- \emph{Homology Theory of algebraic Varieties}, Pergamon Press, 1958.
}

\bibitem[8]{XIII.8} \sisi{%
S. Abhyankar, Resolution of singularities of arithmetical surfaces, p\ptbl 111--152, Arithmetical Algebraic Geometry, Harper and Row, New-York 1963.}
{%
\textsc{S. Abhyankar} -- {\og Resolution of singularities of arithmetical surfaces\fg}, in \emph{Arithmetical Algebraic Geometry}, Harper and Row, New-York, 1965, p\ptbl 111--152.
}

\sisi{}{%
\bibitem[HL]{HL}
\textsc{H.A. Hamm \& L\^e D\~ung Tr\'ang} -- {\og Rectified homotopical depth and Grothendieck conjectures\fg}, in \emph{The Grothendieck Festschrift, Vol.~II}, Progr. Math., vol.~87, Birkhäuser, Boston, 1990, p\ptbl 311--351.
}
\end{thebibliography}

\chapterspace{-4}

% END INCLUDED active_source/parts/13_expose_XIII.tex

% BEGIN INCLUDED active_source/parts/14_expose_XIV.tex
\chapter[Depth and Lefschetz theorems]
{Depth and Lefschetz theorems in~\'{e}tale cohomology} \label{XIV}
\begin{center}
by Mrs M. Raynaud\sfootnote{Based on unpublished notes by A.~Grothendieck.}
\end{center}

\vspace*{1cm}
In\pageoriginale \numero \Ref{XIV.1}, we define a notion of \og \textit{\'{e}tale depth}\fg that is the analogue in \'{e}tale cohomology of the notion of depth studied in \Ref{III}, in the cohomology of coherent sheaves. After a technical part, we prove \og Lefschetz theorems\fg in \numero \Ref{XIV.4}, the central theorem being \Ref{XIV.4.2}. Let $X$ be a scheme, $Y$ a closed subset of $X$, $U$ the complementary open subset $X-Y$, and $F$ an abelian sheaf on $X$ for the \'{e}tale topology; in general, the purpose of Lefschetz theorems is to show that, if $F$ satisfies certain local conditions on $U$, expressible in terms of \'{e}tale depth at the points of $U$, then, under certain additional conditions of a global nature on $U$ (for example, $U$ \textit{affine}), the natural map of \'{e}tale cohomology groups
$$
\H^i(X, F)\to \H^i(Y, \sisi{F{|Y}}{F_{|Y}})$$
is an isomorphism for values $i<n$, where $n$ is a certain specified integer. By taking $F$ to be a constant sheaf, one thereby obtains conditions for $\pi_0(X)$ to equal $\pi_0(Y)$ and conditions for the abelianized fundamental groups of $X$ and $Y$ to be the same. In \numero \Ref{XIV.5}, the introduction of a notion of \og geometric depth\fg makes it possible to give useful special cases of the Lefschetz theorems (\Ref{XIV.5.7}). Finally, in \numero \Ref{XIV.6}, we mention some conjectures, concerning in particular \og noncommutative\fg variants of the theorems obtained.

\section{Cohomological and homotopical depth}\label{XIV.1}

\setcounter{subsection}{-1}

\subsection{\texorpdfstring{}{XIV.1.0}}\label{XIV.1.0}
Fix\pageoriginale the following notation. Let $X$ be a scheme\sfootnote{In accordance with the new terminology (\cf the reissue of \EGA I), here we shall call \og scheme\fg what was formerly called \og prescheme\fg, and \og separated scheme\fg what was called \og scheme\fg.}, $Y$ a closed subset of~$X$, $U$ the complementary open subset, and $i:Y=X-U\to X$ the canonical immersion. Let~$\underline{\Gamma}_Y$ be the functor that associates to an abelian sheaf on $X$ the \og sheaf of sections with support in $Y$\fg, that is, $\underline{\Gamma}_Y=i_*i^!$ (\cf \SGA 4 IV 3.8 and VIII 6.6), and let $\Gamma_Y$ be the functor $\Gamma\cdot \underline{\Gamma}_Y$ (where $\Gamma$ is the \og global sections\fg functor). Consider the derived category $D^+(X)$ and the derived functor $\R\underline{\Gamma}_Y$ (\resp $\R\Gamma_Y$) of $\underline{\Gamma}_Y$ (\resp of $\Gamma_Y$) (\cf \cite{XIV.3}). Given a complex of abelian sheaves $F$ on $X$, bounded below, it may be regarded as an element of $D^+(X)$; we then denote by ${\h}_Y^p(F)$ the $p$-th cohomology sheaf of $\R\underline{\Gamma}_Y(F)$ and by $\H_Y^p(X, F)$ the $p$-th cohomology group of $\R\Gamma_Y(F)$. The results of (\SGA 4 V 4.3 and 4.4) extend trivially to ${\h}_Y^p(F)$ and $\H_Y^p(X, F)$.

\enlargethispage{\baselineskip}%
\begin{proposition} \label{XIV.1.1}
Let $X$ be a scheme, $Y$ a closed subset of $X$, $U$ the complementary open subset, and $i:U\to X$ the canonical immersion. Let $F$ denote either a sheaf of sets on $X$, a sheaf of groups on $X$, or a complex of abelian sheaves on $X$, bounded below. Fix the following notation: if $X'\to X$ is a morphism, $U'$ and $F'$ denote the inverse images of $U$ and $F$ on $X'$; moreover, if $\bar{y}$ is a geometric point of $X$, $\bar{X}$ denotes the strict localization of $X$ at $\bar{y}$, and $\bar{U}$ and $\bar{F}$ the inverse images of $U$ and $F$ in $\bar{X}$.

$1^\circ)$ Let\pageoriginale $F$ be a sheaf of sets on $X$ and $n$ an integer $\leq 2$; then the following conditions are equivalent:

\textup{(i)} The canonical morphism
$$F\to i_*i^*F$$
is injective if $n\geq 1$, and bijective if $n\geq 2$.

\textup{(ii)} For every scheme $X'$ \'{e}tale over $X$, the canonical morphism
$$
\H^0(X', F')\to \H^0(U', F')
$$
is injective if $n\geq 1$, and bijective if $n\geq 2$.

Suppose in addition that the open subset $U$ is retrocompact in $X$; then the preceding conditions are equivalent to the following one:

\textup{(iii)} For every geometric point $\bar{y}$ of $Y$, the canonical morphism
$$
\H^0(\bar{X}, \bar{F})\to \H^0(\bar{U}, \bar{F})$$
is injective if $n\geq 1$, and bijective if $n\geq 2$.

$2^\circ)$ Let $F$ be a sheaf of groups on $X$ and $n$ an integer $\leq 3$; then the following conditions are equivalent:

\textup{(i)} The canonical morphism
$$F\to i_*i^*F$$
is injective if $n\geq 1$, bijective if $n\geq 2$, and if $n\geq 3$, in addition to the preceding conditions, the sheaf of pointed sets $\R^1i_*(i^*F)$ is zero.

\textup{(ii)} For every scheme $X'$ \'{e}tale over $X$, the canonical morphism
$$
\H^0(X', F')\to \H^0(U', F')
$$
is\pageoriginale injective if $n\geq 1$, bijective if $n\geq 2$, and moreover the canonical morphism
$$
\H^1(X', F')\to \H^1(U', F')
$$
is injective if $n\geq 2$, and bijective if $n\geq 3$.

\textup{(ii bis)} The same as \textup{(ii)}, except in the case $n=2$, where it is only assumed that $\H^0(X', F')\to \H^0(U', F')$ is bijective.

Suppose in addition that $U$ is retrocompact in $X$; then the preceding conditions are also equivalent to the following one:

\textup{(iii)} For every geometric point $\bar{y}$ of $Y$, the canonical morphism
$$
\H^0(\bar{X}, \bar{F})\to \H^0(\bar{U}, \bar{F})
$$
is injective if $n\geq 1$, and bijective if $n\geq 2$; finally, if $n\geq 3$, in addition to the preceding conditions, $\H^1(\bar{U}, \bar{F})$ is zero.

$3^\circ)$ Let $F$ be a complex of abelian sheaves, bounded below, and let $n$~be an integer; then the following conditions are equivalent:

\textup{(i)} We have ${\h}_Y^p(F)=0$ for $p<n$ (\cf \hyperref[XIV.1.0]{1.0}).

\textup{(ii)} For every scheme $X'$ \'{e}tale over $X$, the canonical morphism
$$
\H^p(X', F')\to \H^p(U', F')
$$
is bijective for $p<n-1$, and injective for $p=n-1$.

Suppose that $U$ is retrocompact in $X$; then the preceding conditions are also equivalent to the following one:

\textup{(iii)} For every geometric point $\bar{y}$ of $Y$, the canonical morphism
$$
\H^p(\bar{X}, \bar{F})\to \H^p(\bar{U}, \bar{F})$$
is bijective for $p<n-1$, and injective for $p=n-1$.

In\pageoriginale the case where $F$ is an abelian sheaf and $n\geq 2$, conditions \textup{(i)} and \textup{(ii)} are also equivalent to the following one:

\textup{(ii bis)} For every scheme $X'$ \'{e}tale over $X$, the canonical morphism
$$
\H^p(X', F')\to \H^p(U', F')$$
is bijective for $p<n-1$.
\end{proposition}

\skpt
\begin{proof}
$1^\circ)$ It is clear that \textup{(i)}\SSI \textup{(ii)}.
Let us show that, if $U$ is quasi-compact in $X$, \textup{(i)}\SSI
\textup{(iii)}. Indeed, \textup{(i)} is equivalent to saying that, for
every geometric point $\bar{y}$ of $X$, the morphism
$F_{\bar{y}}\to (i_*i^*F)_{\bar{y}}$ is injective if $n\geq 1$ and
bijective if $n\geq 2$ (\SGA 4 VIII 3.6). Since this morphism is in
any case bijective when $\bar{y}$ is a geometric point of $U$, it is
enough to consider the geometric points $\bar{y}$ of $Y$. Now it
follows from the fact that $i$ is quasi-compact and from
(\SGA 4 VIII 5.3) that the morphism
$$F_{\bar{y}}\to (i_*i^*F)_{\bar{y}}$$
is canonically identified with the morphism
$$
\H^0(\bar{X}, \bar{F})\to \H^0(\bar{U}, \bar{F}),
$$
which proves the equivalence of \textup{(i)} and \textup{(iii)}.

$2^\circ)$ \textup{(i)}\ALORS \textup{(ii)}. The assertions concerning $\H^0$ follow from $1^\circ)$. Let $i'$ be the canonical immersion of $U'$ into $X'$; the assertions concerning $\H^1$ follow from the exact sequence (\SGA 4 XII 3.2)
$$0\to \H^1(X', i'_*i'^*F')\to \H^1(U', F')\to \H^0(X', \R^1i'_*(i'^*F')).
$$

(ii bis)\ALORS \textup{(i)}\pageoriginale. By $1^\circ)$, it is enough to show that, for $n\geq 3$, one has $\R^1i_*(i^*F)=0$. Now $\R^1i_*(i^*F)$ is the sheaf associated with the presheaf $X'\mto \H^1(U', F')$, that is, by hypothesis, the sheaf associated with the presheaf $X'\mto \H^1(X', F')$, and the latter is zero.

\textup{(i)}\SSI \textup{(iii)}. Taking $1^\circ)$ into account, the only thing left to verify is that the relation $\R^1i_*(i^*F)=0$ is equivalent to the condition that $\H^1(\bar U, \bar F)=0$ for every geometric point $\bar y$ of $Y$. Since $\R^1i_*(i^*F)$ is zero outside $Y$, saying that $\R^1i_*(i^*F)=0$ is equivalent to saying that $(\R^1i_*(i^*F))_{\bar y}=0$ for every geometric point $\bar y$ of $Y$. It is then enough to note that, since $i$ is quasi-compact, one has $(\R^1i_*(i^*F))_{\bar y}=\H^1(\bar U, \bar F)$ (\SGA 4 VIII 5.3).

$3^\circ)$ \textup{(i)}\ALORS \textup{(ii)}. Let $X'$ be a scheme \`etale over $X$; one has the exact sequence (\SGA 4 V 4.5)
\begin{equation*} \label{eq:XIV.1.1.*} \tag{$*$} {\to \H^p_{Y'}(X', F')\to H^p(X', F')\to H^p(U', F')\to};
\end{equation*}
hence \textup{(ii)} is equivalent to $H^p_{Y'}(X', F')=0$ for $p<n$ and for every scheme $X'$ \`etale over $X$. Consider now the spectral sequence
$$
\E_2^{pq}=\H^p(X', {\h}_Y^q(F))\To \H_{Y'}^*(X', F');$$
by hypothesis, ${\h}_Y^q(F)=0$ for $q<n$, whence $\E_2^{pq}=0$ for $p+q<n$ and consequently $\H_{Y'}^p(X', F')=0$ for $p<n$.

\textup{(ii)}\ALORS \textup{(i)}. The sheaf ${\h}_Y^p(F)$ is associated with the presheaf $X'\mto \H_{Y'}^p(X', F')$; since we have already observed that \textup{(ii)} is equivalent to the relation $\H_{Y'}^p(X', F')=0$ for $p<n$ and for every scheme $X'$ \`etale over $X$, it follows that ${\h}_Y^p(F)=0$ for $p<n$.

\textup{(i)}\SSI \textup{(iii)}. The sheaves ${\h}_Y^p(F)$ are concentrated on $Y$; consequently, saying that ${\h}_Y^p(F)=0$ is equivalent to saying that, for every geometric point\pageoriginale $\bar y$ of $Y$, the fiber $({\h}_Y^p(F))_{\bar y}$ is zero. Since $i$ is quasi-compact, (\SGA 4 VIII 5.2) gives $({\h}_Y^p(F))_{\bar y}={\h}_{\bar Y}^p(\bar X, \bar F)$. The equivalence of \textup{(i)} and \textup{(iii)} follows, in view of the analogue on $\bar X$ of the exact sequence (\Ref{eq:XIV.1.1.*}).

(ii bis)\ALORS \textup{(ii)} when $F$ is an abelian sheaf. The only thing left to show is that ${\h}_Y^{n-1}(F)=0$. Now, for $n>2$, the sheaf ${\h}_Y^{n-1}(F)$ is associated with the cokernel of the morphism of presheaves $X'\mto \H^{n-2}(X', F')\to\H^{n-2}(U', F')$, and is therefore zero. The case $n=2$ follows from the fact that ${\h}_Y^1(F)$ is the cokernel of the morphism $F\to i_*i^*F$.
\skipqed
\end{proof}

\begin{definition} \label{XIV.1.2}
The notation is that of \Ref{XIV.1.1}. One says that $F$ has \`etale $Y$-depth $\geq n$ and writes
$$
\prof_Y(F)\geq n$$
if $F$ satisfies the equivalent conditions \textup{(i)} and \textup{(ii)} of \Ref{XIV.1.1}. If $F$ is a complex of abelian sheaves, the \`etale $Y$-depth of $F$ is the supremum of the~$n$ for which $\prof_Y(F)\geq n$; the same notation will be used if $F$ is a sheaf of sets \resp of not necessarily commutative groups (so that one then has $0\leq \prof_Y(F)\leq 2$\sisi{}{,} \resp $0\leq \prof_Y(F)\leq 3$, when the context does not make clear which of the three variants considered here is being used).

If $\LL$ is a set of prime numbers, one says that the \`etale $Y$-depth of $X$ for $\LL$ is $\geq n$ and writes
$$
\prof_Y^{\LL}(X)\geq n
$$
if, for every constant sheaf of the form $\ZZ/\ell\ZZ$ with $\ell\in \LL$, one has $\prof_Y(\ZZ/\ell\ZZ)\geq n$\pageoriginale. The \`etale $Y$-depth of $X$ for $\LL$ is defined in the evident way. If $\LL=\PP$, the set of all prime numbers, and there is no risk of confusion with the notation of \textup{(\EGA IV 5.7.1)} (which concerns the case $F=\OX$), $\LL$ is omitted from the notation; otherwise one writes $\prof \et(X)$.

Finally, one says that $X$ has homotopical $Y$-depth \sisi{\ignorespaces}{$\geq 3$} for $\LL \sisi{\geq 3}{}$ and writes
$$
\prof \hop_Y^{\LL}(X)\geq 3$$
if, for every constant finite sheaf $F$ of $\LL$-groups on $X$, one has $\prof_Y(F)\geq 3$. If $\LL=\PP$, $\LL$ is omitted from the notation.
\end{definition}

\begin{corollaire} \label{XIV.1.3}
Under the conditions of \Ref{XIV.1.1}, if $\prof_Y(F)\geq n$, then, for every closed subset~$Z$ of $Y$, one has
$$
\prof_Z(F)\geq n.
$$
\end{corollaire}

For example, let us give the argument when $F$ is a complex of abelian sheaves bounded below. We use \Ref{XIV.1.1} $3^\circ)$ \textup{(ii)}. Let $V=X-Z$ and, for every integer $p$, consider the sequence of morphisms
$$
\H^p(X, F)\lto{f} \H^p(V, F)\lto{g} \H^p(U, F).
$$

By hypothesis, $g$ and $g\circ f$ are bijective for $p<n-1$ and injective for $p=n-1$; hence the same is true of $f$. Since the argument remains valid when $X$ is replaced by a scheme $X'$ \`etale over $X$, this proves \Ref{XIV.1.3}.

\begin{corollaire} \label{XIV.1.4} The notation is that of \Ref{XIV.1.1} $2^\circ)$. If $X'$ is a scheme over\pageoriginale $X$, let $\Phi'$ denote the functor that assigns to an \`etale covering of $X'$ its restriction to~$U'$, and let $\Phi'_{F'}$ denote the functor that assigns to a torsor\sfootnote{That is, a \og homogeneous principal bundle\fg in older terminology.} under $F'$ its restriction to~$U'$. Then the following conditions are equivalent:

\textup{(i)} One has $\prof_Y(F)\geq 1$ (\resp $\prof_Y(F)\geq 2$, \resp $\prof_Y(F)\geq 3$).

\textup{(ii)} For every scheme $X'$ \`etale over $X$, the functor $\Phi'_{F'}$ is faithful (\resp fully faithful, \resp an equivalence of categories).

In particular, one has $\prof_Y({X})\geq 1$ (resp.\ $\prof_Y({X})\geq 2$, resp.\ $\prof \hop_Y(X)\geq 3$) if and only if the functor $\Phi'$ is faithful (resp.\ fully faithful, \resp an equivalence of categories).
\end{corollaire}

This follows from \Ref{XIV.1.1} $2^\circ)$ \textup{(ii)}, in view of the interpretation of $\H^1(X', F')$ as the set of classes (up to isomorphism) of torsors under $F'$ (\SGA 4 VII~2), and of \`etale coverings $Z$ of degree $n$ of a scheme as being associated with principal Galois coverings whose group is the symmetric group $\mathfrak{S}_n$: to $Z$ one associates the covering $\underline{\Isom}_X(Z_0, Z)$, where $Z_0$ is the trivial covering of $X$ of degree $n$.

\begin{corollaire} \label{XIV.1.5} Under the conditions of \Ref{XIV.1.1} $3^\circ)$, suppose that $\prof_Y(F)\geq n$; then
$$
\H_Y^n(X, F)\simeq \H^0(X, {\h}_Y^n(F)).
$$
\end{corollaire}

The corollary follows from the spectral sequence
$$
\E_2^{pq}=H^p(X, {\h}_Y^q(F))\To H_Y^*(X, F).
$$

Indeed, by hypothesis $\E_2^{pq}=0$ for $q<n$, and hence
$$
\H_Y^n(X, F)=\E_2^{0n}=\H^0(X, {\h}_Y^n(F)).
$$

\skpt
\begin{remarques} \label{XIV.1.6}
a) The notion\pageoriginale of $Y$-depth, in the form of the equivalent conditions \textup{(i)} and \textup{(ii)} of \Ref{XIV.1.1}, makes sense on any site. In the particular case in which $X$ is a locally Noetherian scheme equipped with the Zariski topology and $F$ is a sheaf of coherent $\OX$-modules, one recovers the usual notion of $Y$-depth as the infimum of the depths at the points of $Y$ (\Ref{III}).

b) For $n\leq 2$, the notion of \`etale $Y$-depth of $X$ is independent of $\LL$. For $n=1$, it simply means that $U$ is dense in $X$. Indeed, this condition is necessary for $\prof_Y(F)\geq 1$, and it is also sufficient because one may assume that $X$ is reduced, in which case the condition that $U$ be dense in $X$ is preserved by \`etale base change \textup{(\EGA IV 11.10.5) \textup{(ii)} b))}. If $U$ is quasi-compact in $X$, the relation $\prof_Y(X)\geq 1$ is also equivalent to saying that $Y$ contains no generic point of~$X$ (\EGA I 6.6.5). For $n=2$ and $U$ quasi-compact in $X$, the condition $\prof_Y(X)\geq 2$ is equivalent to the condition that, for every geometric point $\bar y$ of $Y$, $\bar U$ be nonempty and connected, that is, that \og $Y$ does not disconnect $X$ locally for the \`etale topology\fg.

c) If $X$ has $Y$-depth $\geq n$ for $\LL$ and $U$ is quasi-compact in $X$, then for every locally constant abelian sheaf of $\LL$-torsion $F$ on $X$, one has $\prof_Y(F)\geq n$. Indeed, since the property $\prof_Y(F)\geq n$ is local for the \`etale topology, one may assume that $F$ is constant; then $F$ is a filtered direct limit of sheaves that are finite direct sums of sheaves of the form $\ZZ/p^m\ZZ$, where $m$ is an integer $>0$ and $p\in \LL$. Using \Ref{XIV.1.1} \textup{(iii)} and \textup{(\SGA 4 VII 3.3)}, one sees that one may reduce to the case $F=\ZZ/p^m\ZZ$, and then\pageoriginale, by induction on $m$, to the case $F=\ZZ/p\ZZ$, for which the assertion follows from the definition.

d) By \Ref{XIV.1.4}, if $\prof_Y(X)\geq 3$, the pair $(X, Y)$ is \textit{pure} in the sense of \Ref{X}~\Ref{X.3.1}. In fact, the pure pairs encountered in practice (\cf \Ref{X}~\Ref{X.3.4}) satisfy the stronger condition of homotopical depth $\geq 3$, and this notion can therefore advantageously replace that of a pure pair.

e) Let $F$ be a complex of abelian sheaves and $T(F)$ the complex obtained by applying the translation functor to $F$ (\cite{XIV.3}); then one plainly has:
$$
\prof_Y(T(F))=\prof_Y(F)-1.
$$

f) Let us point out that the recent work of Artin--Mazur (\cite{XIV.1}) makes it possible to define the notion of homotopical depth $\geq n$ for every integer $n$ (not only for $n\geq 3$).

g) Under the conditions of \Ref{XIV.1.1} 3), one has $\prof_Y(F)=\infty$ if and only if $F\isomto Ri_*(i^*F)$ in $D^+(X)$. Indeed, the ${\h}_Y^p(F)$ are the cohomology sheaves of the cone (= mapping cylinder) of the canonical morphism $F\to Ri_*i^*(F)$.
\end{remarques}

\begin{definition} \label{XIV.1.7}
Let $X$ be a scheme, $x$ a point of $X$, $\bar x$ a geometric point above $x$, and $\bar X$ the strict localization of $X$ at $\bar x$. As before, $F$ denotes either a sheaf of sets on $X$, a sheaf of groups on $X$, or a bounded-below complex of abelian sheaves on $X$; $\bar F$ denotes its inverse image on~$\bar X$, and $\LL$ a set of prime numbers. We say that $F$ has \emph{étale depth $\geq n$ at the point $x$}\pageoriginale (respectively that the \emph{étale depth for $\LL$ of $X$ at $x$ is $\geq n$}, respectively that the \emph{homotopical depth for $\LL$ of $X$ at $x$ is $\geq 3$}) and write $\prof_x(F)\geq n$ (respectively $\prof_x^{\LL}(X)\geq n$, respectively $\prof \hop_x^{\LL}(X)\geq \sisi{3}{3}$) if $\prof_{\bar x}(\bar F)\geq n$ (respectively $\prof_{\bar x}^{\LL}(\bar X)\geq\nobreak n$, respectively $\prof \hop_{\bar x}^{\LL}(\bar X)\geq \sisi{3}{3}$). The integer $\prof_x^{\LL}(X)$ and, if $F$ is a complex of abelian sheaves, the integer $\prof_x(F)$ are defined in the evident way. If $\LL$ is the set of all prime numbers, we omit $\LL$ from the notation $\prof_x^{\LL}(X)$, unless there is a risk of confusion with the notation of \textup{(\EGA IV 5.7.1)}, in which case we write $\prof \et_x(X)$.
\end{definition}

We then have the following \textit{pointwise characterization} of depth:

\begin{theoreme} \label{XIV.1.8} Let $X$ be a scheme and $Y$ a closed subset of $X$ such that the open subset $U=X-Y$ is retrocompact in $X$. If $F$ is either a sheaf of sets on $X$, a sheaf of groups on $X$, or a bounded-below complex of abelian sheaves on $X$, then
$$
\prof_Y(F)=\inf_{y\in Y} \prof_y(F).
$$
\end{theoreme}

\subsubsection{} \label{XIV.1.8.1} Let us first show that, for every point $y$ of $Y$, one has the inequality $\prof_y(F)\geq \prof_Y(F)$. Indeed, let $\bar y$ be a geometric point above $y$, $\bar X$ the strict localization of $X$ at $\bar y$, and $\bar F$ and $\bar Y$ the inverse images of $F$ and $Y$ on $\bar X$. By \Ref{XIV.1.7} and \Ref{XIV.1.3}, one has
$$
\prof_y(F)=\prof_{\bar y}(\bar F)\geq \prof_{\bar Y}(\bar F)\geq \prof_Y(F),
$$
the last inequality using the hypothesis \og $U$ retrocompact\fg, through conditions \textup{(iii)} in \Ref{XIV.1.1} and transitivity in the formation of strict localizations.

\subsubsection{} \label{XIV.1.8.2}
Conversely, suppose that, for every point $y$ of $Y$, one has $\prof_y(\sisi{Y}{F})\geq n$ ($n$~an integer), and let us show that $\prof_Y(F)\geq n$.

First recall\pageoriginale the following well-known results \textup{(\SGA 4 VIII)}:

\begin{subsublemme} \label{XIV.1.8.2.1}
Let $X$ be a scheme and $F$ a sheaf of sets on $X$ (respectively let $G\to F$ be a monomorphism of sheaves of sets on $X$). Then two sections $s$ and $s'$ of $F$ coincide (respectively a section $s$ of~$F$ comes from a section of~$G$) if and only if this is so locally. In particular, if~$s$ and~$s'$ are two sections of $F$, there is a largest open subset $V$ of $X$ on which they coincide (respectively, if $s$ is a section of $F$ on $X$, there is a largest open subset $V$ of $X$ such that $s|V$ comes from a section of $G$ on $V$). This open subset is also the set of points~$x$ of $X$ such that, denoting by $\bar x$ a geometric point above~$x$, the sections~$s$ and~$s'$ have the same image in the fiber $F_{\bar x}$ (respectively the image of $s$ in $F_{\bar x}$ comes from an element of $G_{\bar x}$).
\end{subsublemme}

Let us return to the proof of \Ref{XIV.1.8}.

$1^\circ)$ Case where $F$ is a sheaf of sets. If $n=1$, it is enough to show that the canonical morphism
$$
\H^0(X, F)\to \H^0(U, F)$$
is injective, since the result also applies when $X$ is replaced by a scheme étale over $X$. Let $s$ and $s'$ be two sections of $F$ over $X$ which have the same image in $\H^0(U, F)$, and let $V$ be the largest open subset over which they are equal; clearly $V\supset U$. Suppose $V\neq X$, and let $y$ be a maximal point of $X-V$, $\bar y$ a geometric point above\pageoriginale~$y$, $\bar X$ the strict localization of $X$ at $\bar y$, and $\bar V$ and $\bar F$ the inverse images of $V$ and $F$ on $\bar X$. By the choice of $y$, one has $\bar X-\bar y=\bar V$, and by hypothesis the morphism
$$
\H^0(\bar X, \bar F)\to \H^0(\bar X-\bar y, \bar F)=\H^0(\bar V, \bar F)$$
is injective. It follows that $s$ and $s'$ coincide at the point $y$, which is absurd. If $n=2$, it is enough, in view of the preceding, to show that the morphism
$$
\H^0(X, F)\to \H^0(U, F)=\H^0(X, i_*i^*F)$$
is surjective (where $i$ is the canonical immersion of $U$ into $X$). Let $s$ be a section of $i_*i^*F$ over $X$, and let $V$ be the largest open subset over which it comes from a section of $F$. Suppose $V\neq X$, and let $y$ be a maximal point of $X-V$; with the preceding notation, it follows from the hypothesis that the canonical morphism
$$
\H^0(\bar X, \bar F)\to \H^0(\bar X-\bar y, \bar F)=\H^0(\bar V, \bar F)
$$
is bijective; consequently $s|\bar V$ extends to the point $\bar y$, which is absurd and completes the proof in case $1^\circ)$.

$2^\circ)$ Case where $F$ is a sheaf of groups. In view of $1^\circ)$, the only thing left to show is that, in the case $n=3$, the morphism
$$
\H^1(X, F)=\H^1(X, i_*i^*F)\to \H^1(U, F)$$
is bijective. We already know that it is injective, by $1^\circ)$ and \Ref{XIV.1.1} $2^\circ)$ (ii bis). For surjectivity, we use the exact sequence (\SGA 4 XII 3.2)
$$
0\to \H^1(X, i_*i^*F)\to \H^1(U, F)\lto{d}\H^0(X, \R^1i_*(i^*F)).
$$
Let $s\in \H^1(U, F)$, and let $V\supset U$ be the largest open subset over which $d(s)=0$; this is\pageoriginale also the largest open subset such that $s$ comes from an element of $\H^1(V, F)$. Suppose $V\neq X$, and let $y$ be a maximal point of $X-V$; if $\bar X$ is the strict localization of $X$ at a geometric point $\bar y$ above~$y$, then, with the evident notation, we have the exact sequence
$$
0\to \H^1(\bar X, \bar i_*(\bar i^*\bar F))\to \H^1(\bar U, \bar F)\lto{\bar d}\H^0(\bar X, \R^1\bar i_*(\bar i^*\bar F)).
$$
Since $i:U\to X$ is quasi-compact, $\R^1\bar i_*(\bar i^*\bar F)$ is the inverse image of $\R^1i_*(i^*F)$ under the morphism $\bar X\to X$, whence $\H^0(\bar X, \R^1\bar i_*(\bar i^*\bar F))=(\R^1i_*(i^*F))_{\bar y}$. By hypothesis, and since $y\in Y$, the morphism
$$
\H^1(\bar X, \bar F)\to \H^1(\bar V, \bar F)
$$
is bijective. The image $\bar s$ of $s$ in $\H^1(\bar U, \bar F)$, which extends to $\bar V$ by the definition of $V$, therefore also extends to $\bar X$; it follows that $\bar d(\bar s)=0$, and hence the image of $d(s)$ in the geometric fiber $(\R^1i_*(i^*F))_{\bar y}$ is zero; but this contradicts the definition of $V$, proving case~$2^\circ)$.

$3^\circ)$ Case where $F$ is a bounded-below complex of abelian sheaves. We argue by induction on $n$. The conclusion holds for sufficiently small $n$, since $F$ is bounded below. Suppose, therefore, that $\prof_Y(F)\geq n-1$, and let us show that $\prof_Y(F)\geq n$, knowing that for every point $y$ of $Y$ one has $\prof_y(F)\geq n$. It is enough to see that the canonical morphism
\begin{equation*} \label{eq:XIV.1.8.*} \tag{$*$} {\H^{n-2}(X, F)\to \H^{n-2}(U, F)}
\end{equation*}
is surjective and that
\begin{equation*} \label{eq:XIV.1.8.**} \tag{$**$} {\H^{n-1}(X, F)\to \H^{n-1}(U, F)}
\end{equation*}
is injective (the result also applying when $X$ is replaced by a scheme étale over\pageoriginale $X$).

a) Surjectivity of (\Ref{eq:XIV.1.8.*}). The proof is analogous to that of $2^\circ)$. In view of \Ref{XIV.1.5} and (\SGA 4 V 4.5), one has the exact sequence
$$
\H^{n-2}(X, F)\to \H^{n-2}(U, F)\lto{d}\H^{n-1}_Y(X, F)=\H^0(X, {\h}_Y^{n-1}(F)).
$$
Let $s\in \H^{n-2}(U, F)$, and let $V\supset U$ be the largest open subset over which $d(s)=0$, which is also the largest open subset such that $s$ extends to $\H^{n-2}(V, F)$. Suppose $V\neq X$, and let $y$ be a maximal point of $X-V$ and $\bar X$ the strict localization of $X$ at a geometric point $\bar y$ above~$y$. Since $i:U\to X$ is quasi-compact, the formation of ${\h}_Y^{n-1}(F)$ commutes with the base change $\bar X\to X$, and hence one has (with the evident notation) the exact sequence
$$
\H^{n-2}(\bar X, \bar F)\to \H^{n-2}(\bar U, \bar F)\lto{\bar d}\H^{n-1}_{\bar Y}(\bar X, \bar F)=({\h}_Y^{n-1}(F))_{\bar y},
$$
the last equality following from the retrocompactness hypothesis on $U$.

But by hypothesis one has the isomorphism
$$
\H^{n-2}(\bar X, \bar F)\isomto \H^{n-2}(\bar X-\bar y, \bar F)=\H^{n-2}(\bar V, \bar F);
$$
consequently the image $\bar s$ of $s$ in $\H^{n-2}(\bar U, \bar F)$, which extends (by definition of $V$) to $\H^{n-2}(\bar V, \bar F)$, also extends to $\H^{n-2}(\bar X, \bar F)$; but this shows that $\bar d(\bar s)=0$, that is, that $d(s)$ is zero at $\bar y$, which is absurd.

b) Injectivity of (\Ref{eq:XIV.1.8.**}). Using the surjectivity of (\Ref{eq:XIV.1.8.*}), one obtains the exact sequence
$$0\to \H^0(X, {\h}_Y^{n-1}(F))\to \H^{n-1}(X, F)\to \H^{n-1}(U, F)$$
and we must show that every element $s\in \H^0(X, {\h}_Y^{n-1}(F))$ is zero. Let $V$\pageoriginale be the largest open subset over which $s=0$. Suppose that $V\neq X$, and let $y$ be a maximal point of $X-V$ and $\bar X$ a strict localization of $X$ at a geometric point $\bar y$ above~$y$. By the induction hypothesis and \Ref{XIV.1.8.1}, one has $\prof_{\bar Y}(\bar F)\geq \prof_Y(F)\geq n-1$, and hence the map $e$ in the following diagram is injective:
$$
\H^0(\bar X, {\h}_Y^{n-1}(\bar F))= ({\h}_Y^{n-1}(F))_{\bar y}\lto{e} \H^{n-1}(\bar X, \bar F) \lto{f} \H^{n-1}(\bar V, \bar F).
$$
The same is true of $f$ by the hypothesis; the equality on the left follows from the retrocompactness hypothesis on $U$. Let $\bar s$ be the image of $s$ in $({\h}_Y^{n-1}(F))_{\bar y}$; since $s$ vanishes over $V$, one has $f\cdot e(\bar s)=0$, whence $\bar s=0$, which contradicts the choice of $y$ and completes the proof.

\begin{remarque} \label{XIV.1.9}
A result analogous to \Ref{XIV.1.8} is probably valid when the étale topos of a scheme $X$ is replaced by a \og locally finite-type topos\fg, that is, one definable by a locally finite-type site (\SGA 4 VI 1.1). To see this, one must use a result of P. Deligne (\SGA \sisi{4 XVII}{4 VI.9}) asserting that there are \og enough fiber functors\fg in such a topos.
\end{remarque}

We shall deduce from \Ref{XIV.1.8} important cases in which the étale depth can be determined.

\refstepcounter{subsection}\label{XIV.1.10}
\begin{enonce*}{Theorem 1.10}[Cohomological semipurity theorem\ndemark]\ndetext{\label{gabpur} Gabber subsequently proved --- in 1994 --- Grothendieck's absolute cohomological purity conjecture: if $Y$ is a closed subscheme of absolute Noetherian schemes of pure codimension $c$ and $n$ is an integer invertible on $X$, then $\SheafH_Y^q(\Lambda)$ is zero if $q\neq 2c$ and is $\Lambda_Y(-c)$ (Tate twist) otherwise, where $\Lambda=\ZZ/n\ZZ$. See (Fujiwara~K., {\og A Proof of the Absolute Purity Conjecture (after Gabber)\fg}, in \emph{Algebraic geometry 2000, Azumino (Hotaka)}, Adv. Stud. in Pure Math., vol.~36, 2002, p\ptbl 153-183). For applications to the existence of the dualizing complex, see (\SGA 5, Lect. Notes in Math., vol.~589, Springer-Verlag, 1977, p\ptbl 16--72), Exposé 1 and \loccit, \S8. This conjecture had been proved in the case $n=\ell^\nu$, with $\ell$ a sufficiently large prime invertible on $X$, by making crucial use of $K$-theory (Thomason~R.W., {\og Absolute cohomological purity\fg}, \emph{Bull. Soc. Math. France} \textbf{112} (1984), \numero 3, p\ptbl 397--406). $K$-theory reappears in Gabber's proof through the Atiyah-Hirzebruch-Thomason spectral sequence relating étale cohomology and $K$-theory, a method already used in Thomason's approach. Besides this result, the other fundamental argument is the generalization of the Lefschetz theorem cited in note~\eqref{gabaf}, page~\pageref{gabaf}.}
Let~$X$ denote either a scheme smooth over a field $k$, or an excellent regular scheme \textup{(\EGA IV 7.8.2)} of characteristic zero (N.B. if resolution of singularities in the sense of \textup{(\SGA 4 XIX)} is assumed, it is enough, more generally, to suppose that $X$ is an excellent regular scheme of equal characteristic).\pageoriginaled Let $Y$ be a closed subset of $X$ and $\LL$ the set of prime numbers distinct from the characteristic of $X$. Then
$$
\prof_Y^{\LL}(X)=2\codim(Y, X).
$$
\end{enonce*}

\begin{proof}
It follows from \Ref{XIV.1.8} that
$$
\prof_Y^{\LL}(X)=\inf_{y\in Y} \prof_y(X).
$$

On the other hand, since $\codim(Y, X)=\inf_{y\in Y} \dim \Oo_{X, y}$, we are reduced to showing that
$$
\prof_y^{\LL}(X)=2 \dim \Oo_{X, y},
$$
which follows from (\SGA 4 XVI 3.7 and XIX 3.2).
\skipqed
\end{proof}

\begin{theoreme}[Homotopical purity theorem\ndemark] \label{XIV.1.11}
If $X$ is a locally noetherian scheme that is regular (\resp whose local rings are complete intersections), and $Y$ is a closed subset of $X$ such that $\codim(Y, X)\geq 2$ (resp.\ \hbox{$\codim(Y, X)\geq3$}), then one has\ndetext{\label{jongoort} recently, de Jong and Oort obtained the following purity statement: let $\tilde S\to S$ be a resolution of the singularities of the spectrum $S$ of a normal noetherian local ring of dimension $2$, and let~$U$ be the complement of the closed point $s$ in $S$. Suppose in addition that $k(s)$ is algebraically closed. Then, for every prime number $p$, in particular if $S$ has characteristic $p$, the restriction morphism $H^1_\et(\tilde S,\QQ_p)\to H^1_\et(U,\QQ_p)$ is bijective (de Jong A.J. \& Oort~F, {\og Purity of the stratification by Newton polygons\fg}, \emph{J.~Amer. Math. Soc.} \textbf{13} (2000), \numero 1, p\ptbl 209--241, Theorem 3.2). If $k=\CC$ and $A$ is the completion of a surface singularity, this result is due to Mumford (see page~\pageref{XIII.7}, \cite{XIII.5}).}
$$
\prof \hop_Y(X)\geq 3.
$$
\end{theoreme}

Indeed, it follows from \Ref{XIV.1.8} that $\prof \hop_Y(X)=\inf_{y\in Y} \prof \hop_y(X)$. But the strictly local rings of $X$ at the various points of $Y$ are regular rings of dimension $\geq 2$ (\resp complete intersections of dimension $\geq 3$). It then follows from the purity theorem \Ref{X}~\Ref{X.3.4} that $\prof \hop_y(X)\geq 3$, which proves the theorem.

\begin{exemple} \label{XIV.1.12}
Let $X$ be a locally noetherian scheme, $Y$ a closed subset of~$X$, and $n=1$ or $2$. Then, if $\prof_Y(\OX)\geq n$ ($\prof_Y(\OX)$ denoting the $Y$-depth in the sense of coherent sheaves (\cf \hyperref[XIV.1.6]{1.6} a)), one also has $\prof_Y(X)\geq n$; this is obvious for $n=1$, and for $n=2$ it is\pageoriginale precisely Hartshorne's theorem (\Ref{III}~\Ref{III.1}). On the other hand, the analogous assertion is false for $n\geq 3$. For example, take an affine space of dimension $\geq 3$ over a field of characteristic $\neq 2$ and let the group $\ZZ/2\ZZ$ act by reflection through the origin. Let $X$ be the quotient and $Y=\{x\}$ the image of the origin in $X$. Then $\Oo_{X, x}$ is a Cohen--Macaulay ring, so $\prof_x(\OX)\geq 3$; but the affine space with the origin removed is an \'{e}tale covering of $X-\{x\}$ that does not extend to an \'{e}tale covering of $X$; hence, by \Ref{XIV.1.4}, $\prof_Y(X)=2$.
\end{exemple}

The following theorem is the analogue of (\EGA IV 6.3.1):

\begin{theoreme} \label{XIV.1.13}
Let $f:X\to S$ be a morphism of schemes, $Y$ a closed subset of $X$, and $Z$ a closed subset of $S$, such that $f(Y)\subset Z$. Suppose that the local rings of $X$ at the various points of $Y$ are noetherian and that the open subsets $X-Y$ and $S-Z$ are retrocompact in $X$ and $S$, respectively. Let $p$, $q$, $r$ be integers such that $p\geq -r$, $q\geq 0$, let $\LL$ be a set of prime numbers, and let $F$ be a complex of abelian sheaves of $\LL$-torsion on $S$, such that the cohomology sheaves $\H^i(F)$ vanish for $i<-r$. Suppose that
\begin{enumeratea}
\item
The morphism $f$ is locally ($p+q+r-2$)-acyclic for $\LL$ \textup{(\SGA 4 XV 1.11)}.
\item
One has
$$
\prof_Z(F)\geq p.
$$
\item
For every point $s$ of $Z$, one has
$$
\prof_{Y_s}^{\LL}(X_s)\geq q.
$$
\end{enumeratea}
\noindent
Then one has\pageoriginale
$$
\prof_Y(f^*F)\geq p+q.
$$
\end{theoreme}

We shall need the following lemma:

\begin{sublemme} \label{XIV.1.13.1}
Let $\LL$ be a set of prime numbers, let $n$ and $r$ be integers, and let $f:X\to S$ be a morphism locally $n$-acyclic for $\LL$. Let $F$ be a complex of abelian sheaves whose cohomology sheaves are of $\LL$-torsion, such that $\H^i(F)=0$ for $i<-r$, let $Z$ be a closed subset of $S$ such that $S-Z$ is retrocompact in $S$, and let $T=f^{-1}(Z)$. Then the canonical morphism
$$f^*({\h}_Z^i(F))\to {\h}_T^i(f^*F)$$
is bijective for $i<n-r+2$ and injective for $i=n-r+2$.
\end{sublemme}

Set $U=S-Z$ and $V=X-T$, so that one has the cartesian square
$$
\xymatrix{
V \ar[r]^-g\ar[d]_k& U\ar[d]^j\\
X \ar[r]^-f&S
}$$
Consider the following commutative diagram whose rows are exact
$$
\xymatrix{ \ar[r]&f^*({\h}_Z^i(F))\ar[r]\ar[d]& f^*({\h}^i(F))\ar[r]\ar[d]^{\wr} &f^*({\h}^i(\R j_*(j^*F)))\ar[r]\ar[d]& \\
\ar[r] &{\h}_T^i(f^*F)\ar[r]& {\h}^i(f^*F) \ar[r]&{\h}^i(\R k_*(k^*f^*F)) \ar[r]&\text{;}}$$
it follows that one is reduced to showing that the morphism
$$f^*({\h}^i(\R j_*(j^*F)))\to {\h}^i(\R k_*(k^*f^*F))$$
is\pageoriginale bijective for $i<n-r+1$ and injective for $i=n-r+1$. But such a morphism arises from the following morphism between hypercohomology spectral sequences
$$
\xymatrix{f^*E_2^{p, q}=f^*(\R^pj_*({\h}^q(j^*F)))\ar@{=>} [r]\ar[d]&f^*({\h}^*(\R j_*(j^*F)))\ar[d]\\
E'^{p, q}_2=\R^pk_*({\h}^q(k^*f^*F))\ar@{=>}[r]&{\h}^*(\R k_*(k^*f^*F)).}
$$
Since $j$ is quasi-compact, it follows from (\SGA 4 XV 1.10) that the morphism $f^*(E^{p, q}_2)\to E'^{p, q}_2$ is bijective for $p\leq n$ and injective for $p=n+1$; in particular, it is bijective for $p+q\leq n-r$ and injective for $p+q=n-r+1$. The conclusion follows at once.

Let us return to the proof of \Ref{XIV.1.13}. Let $T=f^{-1}(Z)$. By \Ref{XIV.1.13.1} and condition a), the canonical morphism $f^*(\h_Z^i(F))\to\h_T^i(f^*F)$ is an isomorphism for $i\leq p+q$. It therefore follows from b) that $\h^i_T(f^*F)=0$ for $i<p$ and, for $i<p+q$, the restriction of $\h^i_T(f^*F)$ to $T$ is the inverse image of a sheaf $G^i$ on $Z$. Let
$$f_T:T\to Z$$
be the restriction of $f$ to $T$. It then follows from c) and the corollary below that $\h^j_Y(f^*_T(G^i))=0$ for $j<q$. We conclude that
$$
\h^j_Y(\h^i_T(f^*F))=0\text{ for } i+j<p+q,
$$
because the inequality $i+j<p+q$ implies either $i<p$, in which case $\h^i_T(f^*F)=0$, or $j<q$, in which case $\h_Y^j(\h_T^i(f^*F))=0$. Given that, with the notation of \Ref{XIV.1.0}, $\underline{\Gamma}_Y=\underline{\Gamma}_Y.\underline{\Gamma}_T$, one has the spectral sequence
\begin{equation*} \label{eq:XIV.1.13.2} \tag{1.13.2}
E_2^{i, j}=\h_Y^j(\h_T^i(f^*F))\To\h^*_Y(f^*F);
\end{equation*}
since\pageoriginale $E_2^{i, j}=0$ for $i+j<p+q$, one sees that $\h_Y^k(f^*F)=0$ for $k<p+q$. The theorem will therefore be proved if we prove the following corollary (which is the special case of \hyperref[XIV.1.13]{1.13} obtained by taking $Z=S, \ r=p=0$ and $F$ concentrated in degree $0$).

\begin{corollaire}\label{XIV.1.14}
Let $f:X\to S$ be a morphism, let $Y$ be a closed subset such that the complementary open subset $X-Y$ is quasi-compact in $X$ and the local rings of~$X$ at the various points of~$Y$ are Noetherian. Let $\LL$ be a set of \hbox{prime numbers}, let $q$ be an integer, and let $F$ be an abelian sheaf of $\LL$-torsion on $S$. Suppose that~$f$ is locally $(q-2)$-acyclic for $\LL$ and that, for every point $s$ of $S$, one has \hbox{$\prof_{Y_s}^\LL(X_s)\!\geq\! q$}. Then $\prof_Y(f^*F)\geq q$.
\end{corollaire}

$1^\circ)$ Reduction to the case in which $X$ and $S$ are strictly local schemes, $f$ is a local morphism, and $Y$ consists of a closed point of $X$.

By \Ref{XIV.1.8}, to prove \Ref{XIV.1.14}, one must show that, for every point $y$ of $Y$, one has:$$
\prof_y(f^*F)\geq q.
$$

Let $s=f(y)$, let $\bar s$ be a geometric point above $s$, let $\bar y$ be a geometric point above~$y$ and $\bar s$, let $\bar X$ and $\bar S$ be the strict localizations of $X$ and $S$ at $\bar y$ and $\bar s$, respectively, let $\bar f:\bar X\to\bar S$ be the canonical morphism, and let $\bar F$ be the inverse image of $F$ on $\bar S$. Since $\prof_y(f^*F)=\prof_{\bar y}(\bar f^*\bar F)$, it is enough to show that the hypotheses of \Ref{XIV.1.14} are preserved when $f$ (\resp $Y$, \resp $F$) is replaced by $\bar f$ (\resp $\{\bar y\}$, \resp $\bar F$). The quasi-compactness condition follows from the Noetherian hypothesis on $\Oo_{X,y}$, which implies that $\bar X$ is Noetherian\pageoriginale. By (\SGA 4 XV~1.10 \textup{(i)}), $\bar f$ is still locally $(q-2)$-acyclic for $\LL$. Moreover, the fiber $(\bar X)_{\bar s}$ of $\bar X$ above $\bar s$ identifies with the strict localization of $X_s$ at $\bar y$, and therefore satisfies $\prof_{\bar y}((\bar X)_{\bar s})\geq q$. Since an analogous relation is trivially satisfied for the fibers of the $\bar S$-scheme $\bar X$ other than the closed fiber, this completes the reduction.

$2^\circ)$ The case in which $X$ and $S$ are strictly local, $f$ is a local morphism, and $Y$ consists of the closed point of $X$. Let
$$g:U=X-\{y\}\to S$$
be the structural morphism of $U$. We must show that the canonical morphism
$$u_i:H^i(X, f^*F)\to H^i(U, f^*F)$$
is bijective for $i\leq q-2$ and injective for $i=q-1$. Consider the commutative diagram
$$
\xymatrix{H^i(X, f^*F)\ar[rr]^{u_i}&& H^i(U, f^*F)\\&H^i(S, F)\ar[lu]^{v_i}\ar[ru]_{w_i}}$$
The morphism $v_i$ is plainly bijective for every $i$. Moreover, $g$ is locally $(q-2)$-acyclic for $\LL$; its fibers are also $(q-2)$-acyclic for $\LL$, as follows from the fact that $\prof_y(X_s)\geq q$ and that the fibers of $f$ are $(q-2)$-acyclic for~$\LL$. Since $g$ is quasi-compact because $X$ is Noetherian, (\SGA 4 XV 1.16) shows that $g$ is $(q-2)$-acyclic for $\LL$. Consequently $w_i$, and hence also $u_i$, is bijective for $i\leq q-2$ and injective for $i=q-1$, which completes the proof\pageoriginale of \Ref{XIV.1.14}.

\begin{corollaire} \label{XIV.1.15}
Let $f:X\to S$ be a morphism of schemes, let $\LL$ be a set of prime numbers, let $m$ and $r$ be integers, and let $F$ be a complex of abelian sheaves of $\LL$-torsion on $S$ such that $\h^i(F)=0$ for $i<-r$. Let $x$ be a point of $X$, let $s=f(x)$, and suppose that the local ring $\Oo_{X,x}$ is Noetherian. Then, if $f$ is locally $m$-acyclic for $\LL$, one has
\begin{equation*} \label{eq:XIV.115..*} \tag{$*$} {\prof_x(f^*F)\geq\inf(\prof_s(F)+\prof^\LL_x(X_s), n)\text{ where }n=m-r+2.}
\end{equation*}
In particular, if $n\geq\prof_s(F)+\prof_x^\LL(X_s)$, for example if $f$ is locally acyclic for $\LL$, then
\begin{equation*} \label{eq:XIV.15.**} \tag{$**$} {\prof_x(f^*F)\geq\prof_s(F)+\prof_x^\LL(X_s).}
\end{equation*}
If $\LL$ consists of a single element $\ell$ and $n\geq \prof_s(F)+\prof_x^\LL(X_s)$, the preceding inequality is an equality.
\end{corollaire}

One reduces to the case in which $s$ and $x$ are closed points by taking the strict localizations of $S$ and $X$ at geometric points $\bar s$ above $s$ and $\bar x$ above $x$ and $\bar s$. If $n\geq\prof_s(F)+\prof_x^\LL(X_s)$, then ($*$) follows from \Ref{XIV.1.13} by taking $p=\prof_s(F)$ and $q=\prof_x^\LL(X_s)$ there (the hypothesis that $S-\{s\}$ is quasi-compact in $S$ follows from the fact that $X-X_s$ is quasi-compact in $X$ and that $f$ is surjective, since it is $(-1)$-acyclic, except possibly when the conclusion of \hyperref[XIV.1.15]{1.15} is empty). If $n<\prof_s(F)+\prof_x^\LL(X_s)$, inequality \eqref{eq:XIV.115..*} again follows from \Ref{XIV.1.13}, for example by taking $p=\prof_s(F)$ and $q=n-p$. It remains to prove\pageoriginale the last assertion. Let $p=\prof_s(F)$ and $q=\prof_x(X_s)$; it follows from \eqref{eq:XIV.1.13.2} that
$$
\h_x^{p+q}(f^*(F))\simeq \h_x^q(f^*(\h^p_s(F))).
$$
Since $\prof_s(F)=p$, the sheaf $\h^p_s(F)$ is a nonzero sheaf of $\ell$-torsion, constant on~$s$. Consequently the sheaf $G=f^*(\h_s^p(F))$ is a nonzero sheaf of $\ell$-torsion, constant on $X_s$, and therefore contains a subsheaf isomorphic to $\ZZ/\ell\ZZ$; since $\h_x^q(\ZZ/\ell\ZZ)$ is nonzero, one indeed has $\h_x^q(G)\neq 0$.

\begin{corollaire} \label{XIV.1.16}
Let $f:X\to S$ be a regular morphism of excellent schemes (\textup{\EGA IV 7.8.2}) of characteristic zero, let $\ell$ be a prime number, and let $F$ be a complex of sheaves of $\ell$-torsion on $S$. Let $x\in X$ and $s=f(x)$; then
$$
\prof_x(f^*F)=\prof_s(F)+2\dim(\Oo_{X_s,x}).
$$
\end{corollaire}

Indeed, $f$ \sisi{is regular by the definition of \og excellent\fg and hence is}{is} locally acyclic (\SGA 4 XIX 4.1). It then follows from \Ref{XIV.1.15} that
$$
\prof_x(f^*F)=\prof_s(F)+\prof_x(X_s).
$$
But by \Ref{XIV.1.10},
$$
\prof_x(X_s)=2\dim\Oo_{X_s,x},
$$
which gives the result.

\begin{remarque} \label{XIV.1.17}
It follows from \Ref{XIV.1.15} that \Ref{XIV.1.13} remains valid when b) and c) are replaced\pageoriginale by the conditions:

b$'$) For every point $s\in f(Y)$, one has $\prof_s(F)\geq p$.

c$'$) For every point $x\in Y$, if $s=f(x)$, one has $\prof^\LL_x(X_s)\geq q$.
\end{remarque}

For a sheaf of sets or groups, one has the following theorem, analogous to \Ref{XIV.1.13}.

\begin{theoreme} \label{XIV.1.18}
Let $f:X\to S$ be a morphism of schemes and let $Y$ be a closed subset of $X$ such that $X-Y$ is quasi-compact in $X$ and, for every point $x$ of $Y$, the local ring $\Oo_{X,x}$ is Noetherian.

$1^\circ)$ Let $F$ be a sheaf of sets on $S$ and let $n$ be an integer equal to $1$ or $2$. Suppose that $f$ is locally $(n-2)$-acyclic and that, for every point $s$ of $f(Y)$,
$$
\prof_{Y_s}(X_s)+\prof_s(F)\geq n.
$$
Then:$$
\prof_Y(f^*F)\geq n.
$$

$2^\circ)$ Let $\LL$ be a set of prime numbers and let $F$ be a sheaf of ind-$\LL$-groups. Suppose that $f$ is locally $1$-aspherical for $\LL$ (\textup{\SGA 4 XV 1.11}) and that, for every point $s$ of $f(Y)$,
$$
\prof\hop^\LL_{Y_s}(X_s)+\prof_s(F)\geq 3.
$$
Then:$$
\prof_Y(f^*F)\geq 3.
$$
\end{theoreme}

\begin{proof}
We\pageoriginale reduce, as in \Ref{XIV.1.14} and \Ref{XIV.1.15}, to the case where $X$ and $S$ are strictly local schemes, $f$ is a local morphism, and $Y$ is the closed point $x$ of $X$. Let $s=f(x)$ be the closed point of $S$; we have the commutative diagram:
$$
\xymatrix{ X-X_s\ar[r]^i\ar[d]_h&X-\{x\}\ar[rd]_g\ar[rr]^j&&X\ar[ld]^f\\
S-\{s\}\ar[rr]_k&&S. }$$

$1^\circ)$ a) \textit{Case} $n=1$.

If $\prof_s(F)\geq 1$, then the morphism $F\to k_*k^*F$ is injective, and hence the morphism $f^*F\to f^*(k_*k^*F)$ is also injective. On the other hand, it follows from the fact that $f$ is locally $(-1)$-acyclic that the morphism $f^*(k_*k^*F)\to (j.i)_*(f^*\sisi{F{|X-X_s}}{F_{|X-X_s}}))$ is injective. Finally, the composite morphism $f^*F\to (j.i)_*(f^*\sisi{F{|X-X_s}}{F_{|X-X_s}}))$ is injective, which shows that $\prof_{X_s}(f^*F)\geq 1$, and consequently also $\prof_x(f^*F)\geq 1$.

If $\prof_x(X_s)\geq 1$, consider the commutative diagram
\begin{equation*} \label{eq:XIV.18.*} \tag{$*$}
\begin{array}{c}
\xymatrix{
H^0(X, f^*F)\ar[r]^-{v}\ar[d]_\wr&H^0(X-\{x\}, f^*F)\ar[d]\\
H^0(X_s, f^*F)\ar[r]^-{v'}&H^0(X-\{x\}, f^*F);
}
\end{array}
\end{equation*}
By hypothesis, $v'$ is injective, and therefore so is $v$.

b) \textit{Case} $n=2$. Consider the commutative diagram
\begin{equation*} \label{eq:XIV.1.18.**} \tag{$**$}
\begin{array}{c}
\xymatrix{ H^0(S, F)\ar[d]_{m}^{\wr}\ar[rr]^u&&H^0(S-\{s\}, F)\ar[d]^{\wr}_{ n}\\
H^0(X, f^*F)\ar[r]^-v&H^0(X-\{x\}, f^*F)\ar[r]^w&H^0(X-X_s, f^*F);
}
\end{array}
\end{equation*}
we must\pageoriginaled show that $v$ is bijective. The morphism $m$ is obviously bijective, and, since $f$ is $0$-acyclic, $n$ is also bijective.

If $\prof_s(F)\geq 2$, $u$ is bijective. As we saw in a), the sole hypothesis $\prof_s(F)\geq 1$ implies the relation $\prof_{X_s}(f^*F)\geq 1$; consequently $v$ and $w$ are injective, and it then follows from \eqref{eq:XIV.1.18.**} that $v$ is bijective.

If $\prof_x(X_s)\geq 2$, then $g$ is $0$-acyclic (because it is locally $0$-acyclic and its fibers are $0$-acyclic). It follows that $v\cdot m$ is bijective, and therefore $v$ is bijective.

If $\prof_s(F)\geq 1$ and $\prof_x(X_s)\geq 1$, then we already know that $v$ and $w$ are injective. Let $z$ be a maximal point of $X_s-\{x\}$ (such a point exists by the hypothesis $\prof_x(X_s)\geq 1$), let $\bar Z$ be the strict localization of $X$ at a geometric point above $z$, and let $\overline{f^*F}$ be the inverse image of $f^*F$ on $\bar Z$. Consider the commutative diagram
\begin{equation*}
\xymatrix@C=0mm{
&H^0(S, F)\ar[rr]\ar[ld]_{m}^{\sim}&&H^0(S-\{s\}, F)\ar[rd]_{\sim}^{n}\\
H^0(X, f^*F)\ar[rr]^-{v}\ar[rd]_{m'}^{\sim}&&H^0(X-\{x\}, f^*F)\ar[rr]^-{w}\ar[ld]^r&& H^0(X-X_s, f^*F)\ar[ld]^{n'}_{\sim}\\
&H^0(\bar Z, f^*F)\ar[rr]&&H^0(\bar Z-\{\bar z\}, \overline{f^*F})}
\end{equation*}
the morphism $m'\cdot m$ is obviously bijective, and it follows from the fact that $f$ is locally $0$-acyclic that $n'\cdot n$ is bijective; consequently $m'$ and $n'$ are also bijective. Since $w$ is injective, $r$ is also injective, and therefore $v$ is bijective.

$2^\circ)$ In view of b), we already know that $\prof_x(f^*F)\geq 2$.

If\pageoriginale $\prof_s(F)\geq 3$, then $\R^1k_*(k^*F)=1$\nde{the trivial torsor is subsequently denoted successively by $0$ or $1$; we have retained this double notation, which, on reflection, causes no ambiguity.}. Since $f$ is locally $1$-aspherical, one has $\R^1(j\cdot i)_*(f^*\sisi{F{|X-X_s}}{F_{|X-X_s}})\!=\!f^*(\R^1k_*(k^*F))\!=\!1$. Thus \hbox{$\prof_{X_s}(f^*F)\!\geq\! 3$} and consequently $\prof_x(f^*F)\geq 3$.

If $\prof_x(X_s)\geq 3$, then $g$ is $1$-aspherical (because $g$ is locally $1$-aspherical and its fibers are $1$-aspherical). Thus $H^1(X-\{x\}, f^*F)=H^1(S, F)=1$, and consequently $\prof_x(f^*F)\geq 3$.

If $\prof_s(F)\geq 2$ and $\prof_x(X_s)\geq 1$, we use the exact sequence (\SGA 4 XII 3.2):
$$1\to\R^1j_*(i_*(f^*\sisi{F{|X-X_s}}{F_{|X-X_s}}))\to\R^1(j.i)_*(f^*\sisi{F{|X-X_s}}{F_{|X-X_s}}) \to j_*(\R^1i_*(f^*\sisi{F{|X-X_s}}{F_{|X-X_s}})).
$$
Since $f$ and $g$ are locally $1$-aspherical, one has
\begin{align*}
\R^1(j\cdot i)_*(f^*\sisi{F{|X-X_s}}{F_{|X-X_s}})&\simeq f^*(\R^1k_*(k^*F)) \\
\R^1i_*(f^*\sisi{F{|X-X_s}}{F_{|X-X_s}})&\simeq g^*(\R^1k_*(k^*F));
\end{align*}
the preceding exact sequence can then be written in the form
\enlargethispage{\baselineskip}%
\begin{equation} \label{eq:XIV.18.***}\tag{${*}{*}{*}$}
1\to\R^1j_*(i_*(f^*\sisi{F{|X-X_s}}{F_{|X-X_s}}))\to f^*(\R^1k_*(k^*F))\lto{a} j_*(j^*(f^*(\R^1k_*(k^*F))).
\end{equation}
The hypothesis $\prof_s(F)\geq 2$ shows that the morphism $F\to k_*k^*F$ is bijective; applying $g^*$, and taking into account the fact that $g$ is locally $0$-acyclic, we find $f^*\sisi{F{|X-\{x\}}}{F_{|X-\{x\}}}=i_*(f^*\sisi{F{|X-X_s}}{F_{|X-X_s}})$. The hypothesis $\prof_x(X_s)\geq 1$ shows that the morphism~$a$ is injective (note that $f^*(\R^1k_*(k^*F))$ is a sheaf equal to $1$ outside $X_s$ and constant on $X_s$). It then follows from \eqref{eq:XIV.18.***} that $\R^1j_*(f^*\sisi{F{|X-\{x\}}}{F_{|X-\{x\}}})=1$, and hence $\prof_x(f^*F)\geq 3$\pageoriginale.

\enlargethispage{.8\baselineskip}%
If $\prof_s(F)\geq 1$ and $\prof_x(X_s)\geq 2$, consider the sheaf of homogeneous spaces $G$ defined by the exact sequence
$$
1\to F\to k_*k^*F\to G\to 1.
$$
Applying the exact functor $g^*$ to this exact sequence and using (\SGA 4 XII 3.1), we obtain the following commutative diagram with exact rows:
$$
\xymatrix@R=5mm{f^*(k_*k^*F)\ar[r]\ar[d]&f^*G\ar[r]\ar[d]^b&1\\
j_*(g^*(k_*k^*F))\sisi{\ar[r]}{\ar[r]^-{u}}&j_*(g^*G)\ar[r]&\R^1j_*(g^*F) \ar[r]&\R^1j_*(g^*(k_*k^*F)).}
$$
Since $\prof_x(X_s)\geq 2$, the morphism $b$ is bijective, so $u$ is surjective, and we thus have a map whose kernel is reduced to the neutral element:
$$
1\to\R^1j_*(g^*F)\to\R^1j_*(g^*(k_*k^*F))=R.
$$
Since $g^*(k_*k^*F)\simeq i_*(f^*\sisi{F{|X-X_s}}{F_{|X-X_s}})$ (because $g$ is locally $0$-acyclic), $R$ is identified with the first term of the exact sequence \eqref{eq:XIV.18.***}; but in the preceding case we saw that $R=1$ as soon as $\prof_x(X_s)\geq 1$, which proves that $\prof_x(f^*F)\geq 3$ and completes the proof of \Ref{XIV.1.18}.
\skipqed
\end{proof}

The following corollaries are generalizations of (\SGA 4 XVI 3.2 and 3.3).

\begin{corollaire} \label{XIV.1.19}
Let $f:X\to S$ be a flat morphism with separable fibers between locally Noetherian schemes, and let $Y$ be a closed subset of $X$. Suppose\pageoriginale that, for every point $s\in f(Y)$, the fiber $Y_s$ is rare\nde{\og rare\fg = \og having empty interior\fg, \cf Bourbaki TG~IX.52.} in $X_s$ and that one of the following two conditions holds:
\begin{enumeratea}
\item
the closure of $f(Y)$ is rare in $S$.
\item
$X_s$ is geometrically unibranch at the points of $Y_s$.
\par\noindent
Then
$$
\prof_Y(X)\geq 2.
$$
\end{enumeratea}
\end{corollaire}

Indeed, it follows from the hypothesis on $f$ that $f$ is locally $0$-acyclic (\SGA 4 XV 4.1). We then apply \Ref{XIV.1.13}. The hypothesis that $Y_s$ is rare in $X_s$ (respectively that $\overline{f(Y)}$ is rare in~$S$) is equivalent, by \Ref{XIV.1.6} b), to the relation $\prof_{Y_s}(X_s)\geq 1$ (respectively $\prof_{\overline{f(Y)}}(S)\geq 1$). The hypothesis that $X_s$ is geometrically unibranch at every point of $Y_s$ is equivalent to saying that the strict localization of $X_s$ at a geometric point of $Y_s$ is irreducible; knowing that $Y_s$ is rare in $X_s$, this clearly implies $\prof_{Y_s}(X_s)\geq 2$, by \Ref{XIV.1.8}. In either case \Ref{XIV.1.13} indeed gives $\prof_Y(X)\geq 2$.

\begin{corollaire} \label{XIV.1.20}
Let $f:X\to S$ be a regular morphism (\textup{\EGA IV 6.8.1}) between locally Noetherian schemes, and let $Y$ be a closed subset of $X$. Suppose that, for every point $s\in f(Y)$, one of the following conditions holds:
\begin{enumeratea}
\item
One has $\codim(Y_s, X_s)\geq 2$.
\item
One has $\codim(Y_s, X_s)\geq 1$ and $\prof_s(S)\geq 1$.
\item
One has $\prof\hop_s(S)\geq 3$.
\end{enumeratea}
\noindent
Then\pageoriginale
$$
\prof\hop_Y(X)\geq 3.
$$
\end{corollaire}

Indeed, this follows from \Ref{XIV.1.18}, since hypothesis a) implies $\prof\hop_{Y_s}(X_s)\geq 3$ (\cf \Ref{XIV.1.11}), and the condition $\codim(Y_s, X_s)\geq 1$ clearly implies $\prof_Y(X)\geq 2$.

\section{Technical lemmas} \label{XIV.2}

\subsection{\texorpdfstring{}{XIV.2.1}} \label{XIV.2.1}
Let $S$ be a locally noetherian scheme, $f:X\to S$ a morphism locally of finite type, and $t$ a point of $S$. If $x\in X$ is such that $s=f(x)\in\Spec\Oo_{S, t}$, set
$$
\delta_t(x)=\degtr k(x)/k(s)+\dim(\overline{\{s\}}),
$$
where $\overline{\{s\}}$ denotes the closure of $s$ in $\Spec\Oo_{S, t}$, and $k(x)$ and $k(s)$ the residue fields of $x$ and $s$, respectively. If $S$ is a local ring with closed point $t$, we also write $\delta(x)$ instead of $\delta_t(x)$ (\cf \SGA 4 XIV 2.2).
\begin{sublemme} \label{XIV.2.1.1} Let there be a cartesian square
$$
\xymatrix{X'\ar[r]^{h}\ar[d]_{f'}&X\ar[d]^f\\S'\ar[r]^g&S},
$$
where $S$ and $S'$ are noetherian local rings, with closed points $t$ and $t'$, respectively, $g$ is a faithfully flat morphism such that $g^{-1}(t)=t'$\pageoriginale, and $f$ is a morphism locally of finite type. Let $x'\in X',\ x=h(x'),\ s=f(x),\ s'=f'(x')$; then one has
$$
\delta(x')\leq\delta(x).
$$
Moreover, the preceding inequality is an equality if and only if one has:
$$
\degtr k(x)/k(s)=\degtr k(x')/k(s')\text{ and }\dim(\overline{\{s\}})=\dim(\overline{\{s'\}}).
$$
In particular, given $x\in X$, one can find $x'$ such that $\delta(x)=\delta(x')$.
\end{sublemme}

Indeed, one has (\EGA IV 6.11)
$$
\dim(\overline{\{s\}})=\dim g^{-1}(\overline{\{s\}}).
$$
It follows that, for every point $s'$ of $g^{-1}(s)$, one has the relation $\dim(\overline{\{s'\}})\leq \dim(\overline{\{s\}})$, and that, given $s$, one can find $s'\in g^{-1}(s)$ for which equality holds. Let $Z$ then denote the schematic closure of $x$ in the fiber $X_s$ of $X$ at~$s$, and let $Z'=Z\times_{\Spec k(s)}\Spec k(s')$. Then $Z'$ is equidimensional of dimension $\degtr k(x)/k(s)$; therefore, for every point $x'\in Z'_{x}$, one has
$$
\degtr k(x')/k(s')\leq\degtr k(x)/k(s)\text{, with equality}$$
when $x'$ is a maximal point of $Z'_x$. The asserted conclusion follows at once.

\subsection{\texorpdfstring{}{XIV.2.2}} \label{XIV.2.2} Let $f:X\to S$ be a morphism locally of finite type and $T$ a closed subset of $S$. Let $x\in X, s=f(x)$; we set
$$
\delta_T(x)=\degtr k(x)/k(s)+\codim(\overline{\{s\}}\cap T, \overline{\{s\}})=\inf_{t\in T\cap\overline{\{s\}}}\delta_t(x).
$$

\begin{sublemme} \label{XIV.2.2.1}
Let\pageoriginale there be a cartesian square
$$
\xymatrix{X'\ar[r]^{h}\ar[d]_{f'}&X\ar[d]^f\\S'\ar[r]^g&S},
$$
where the schemes $S$ and $S'$ are locally noetherian and catenary, the morphism $f$ is locally of finite type, and $g$ is faithfully flat. Let $T$ be a closed subset of $S$ and $T'$ a closed subset of $S'$, such that $g(T')\subset T$, let $x'$ be an element of $X'$, let $x=h(x')$, and let
$$h_{x'}:\Spec\Oo_{X', x'}\to\Spec\Oo_{X, x}$$
be the morphism induced by $h$. Then one has:
$$
\delta_T(x)-\delta_{T'}(x')\leq\dim h_{x'}^{-1}(x).
$$
\end{sublemme}

Let $s'=f'(x'), s=f(x)$. By definition, one has:
\begin{multline*}
\delta_T(x)-\delta_{T'}(x')=\degtr k(x)/k(s)-\degtr k(x')/k(s')\\
+\codim(\overline{\{s\}}\cap T, \overline{\{s\}})-\codim(\overline{\{s'\}}\cap T', \overline{\{s'\}}).
\end{multline*}
Since $g$ is faithfully flat, it follows from (\EGA IV 6.1.4) that one has
\begin{align*} \label{eq:XIV.2.2.*} \tag{$*$}
\codim(\overline{\{s\}}\cap T, \overline{\{s\}})&=\codim(g^{-1}(\overline{\{s\}})\cap g^{-1}(T), g^{-1}(\overline{\{s\}}))\\&\leq\codim(g^{-1}(\overline{\{s\}})\cap T', g^{-1}(\overline{\{s\}}));
\end{align*}
since $S'$ is catenary, one has, by ($\EGA 0_{\textup{IV}}$ 14.3.2 b)):
\begin{multline*}
\codim(\overline{\{s'\}}\cap T', g^{-1}(\overline{\{s\}}))= \codim(\overline{\{s'\}}\cap T', \overline{\{s'\}})+ \codim(\overline{\{s'\}}, g^{-1}(\overline{\{s\}}))\\
= \codim(\overline{\{s'\}}\cap T', g^{-1}(\overline{\{s\}})\cap T')+ \codim(g^{-1}(\overline{\{s\}})\cap T', g^{-1}(\overline{\{s\}})).
\end{multline*}
From\pageoriginale this relation and \eqref{eq:XIV.2.2.*}, one deduces
$$
\delta_T(x)-\delta_{T'}(x')\leq\degtr k(x)/k(s)-\degtr k(x')/k(s')+\codim(\overline{\{s'\}}, g^{-1}(\overline{\{s\}})).
$$
Let us calculate $\codim(\overline{\{s'\}}, g^{-1}(\overline{\{s\}}))=\dim\Oo_{S'_s, s'}$ (where $S'_s$ is the fiber of $S'$ at $s$). Let $Z$ be the closed image of $x$ in $X_s$ and let $Z'\subset X'_s$ be the scheme defined by the cartesian square
$$
\xymatrix{Z'\ar[r]\ar[d]&Z\ar[d]\\S'_s\ar[r]&\Spec k(s)}.
$$
The morphism $Z\to \Spec k(s)$ is flat, locally of finite type, and $\dim Z=\degtr k(x)/k(s)$. It then follows from (\EGA IV 6.1.2) that

$$
\dim(\Oo_{Z', x'})=\dim(\Oo_{S'_s, s'})+\degtr k(x)/k(s)-\degtr k(x')/k(s');$$
taking into account the fact that $Z'_{s'}\simeq Z\otimes_{k(s)}k(s')$, one then obtains:
$$
\delta_T(x)-\delta_{T'}(x')\leq\dim(\Oo_{Z', x'}).
$$
But $\Spec(\Oo_{Z', x'})$ identifies with the fiber at $x$ of the morphism
$$(\Spec(\Oo_{X', x'}))_s\to(\Spec(\Oo_{X, x}))_s,
$$
and hence also with the fiber at $x$ of $h_{x'}$, which proves the sublemma.

\subsection{\texorpdfstring{}{XIV.2.3}} \label{XIV.2.3}
The proofs of the theorems in \numero \Ref{XIV.4} are based on duality theory; they use the following lemmas. Let $m$ be an\pageoriginale integer that is a power of a prime number $\ell$; if $X$ is a scheme, all the sheaves considered on $X$ are sheaves of $\ZZ/m\ZZ$-modules; there is then a notion of a dualizing complex on $X$ (\SGA 5 I 1.7). Suppose that such a complex $K$ exists on $X$; then, for every geometric point~$\bar x$ \sisi{over}{over} a point $x$ of $X$, one deduces from $K$ (\cf \SGA 5 I 4.5) a dualizing complex~$K_{\bar x}$ on $\Spec k(\bar x)$, so that $K_{\bar x}\simeq\ZZ/m\ZZ[n]$ (the brackets denoting the translation functor) for a certain integer $n$ depending only on $x$. We set
$$
\delta^K(x)=n.
$$
If $K$ is normalized at the point $x$ (\SGA 5 I 4.5), then $n=0$.

\begin{sublemme} \label{XIV.2.3.1} Let $X$ be a locally noetherian scheme endowed with a dualizing complex $K$. If $x$ and $x'$ are two points of $X$, such that $x$ is a specialization of $x'$ and $\codim(\overline{\{x\}}, \overline{\{x'\}})=1$, then one has
$$
\delta^K(x)=\delta^{K}(x')-2.
$$
\end{sublemme}

One may first reduce to the case where $X$ is a strictly local scheme. Indeed, let $\bar X$ be the strict localization of $X$ at a geometric point $\bar x$ over $x$, let $i:\bar X\to X$ be the canonical morphism, and let $\bar x'$ be a geometric point of $\bar X$ over $x'$. Then $i^*K$ is a dualizing complex on $\bar X$, and one has (\SGA 5 I 4.5)
$$
(i^*K)_{\bar x}\simeq K_{\bar x}\text{ and } (i^*K)_{\bar x'}\simeq K_{\bar x'},
$$
which\pageoriginale completes the reduction to the strictly local case.

If $j:\overline{\{x'\}}\to X$ denotes the immersion of the reduced closed subscheme of $X$ whose underlying space is $\overline{\{x'\}}$, then $\R^!j(K)$ is a dualizing complex on $\overline{\{x'\}}$, and one sees at once, using (\SGA 5 I 4.5), that it suffices to prove the lemma for $\overline{\{x'\}}$. One is thus reduced to the case where $X$ is an integral strictly local scheme of dimension $1$.

Let $X'$ then be the normalization of $X$ and let $f:X'\to X$ be the canonical morphism; $f$ is integral, surjective, and radicial, and it follows that $f^*K$ is a dualizing complex on $X'$ and that it suffices to prove the lemma for $X'$ and for the points \sisi{over}{over}~$x$ and~$x'$. One is thus reduced to the case where $X$ is a local, integral, regular scheme of dimension $1$, but it is known (\cf \SGA 5 I 4.6.2 and 5.1) that then $\mu_m[2]$ and $\ZZ/m\ZZ$ are dualizing complexes, normalized at the points $x$ and $x'$, respectively; the lemma follows at once.

\begin{sublemme} \label{XIV.2.3.2} Let $S$ be a noetherian local scheme and $f:X\to S$ a morphism of finite type. If $K$ is a dualizing complex on $S$, normalized at the closed point $t$ of $S$, and if $\R^!f(K)=K'$ is a dualizing complex on $X$ (\cf \SGA 5 I 3.4.3), one has, for every point $x$ of $X$:$$
\delta^{\sisi{K}{{K'}}}(x)=2\delta(x).
$$
\end{sublemme}

Indeed, let $s=f(x)$ and let $x'$ be a closed point of the fiber $X_s$; then $\delta^{K'}(x')=\delta^K(s)$ and, by \Ref{XIV.2.3.1},
$$
\delta^K(s)=2\codim(\overline{\{t\}}, \overline{\{s\}})=2\dim(\overline{\{s\}}).
$$
Since\pageoriginaled $x'$ may be chosen to be a specialization of $x$, one has, by \Ref{XIV.2.3.1},
$$
\delta^{K'}(x)=\delta^{K'}(x')+2\codim(\overline{\{x'\}}, \overline{\{x\}})=\delta^{K'}(x')+2\degtr k(x)/k(s);$$
the lemma follows at once.

The following lemma will be used only for the converse of the Lefschetz theorem, in \numero\Ref{XIV.4}:

\begin{sublemme} \label{XIV.2.3.3}
Let there be a cartesian square
$$
\xymatrix{X'\ar[r]^{h}\ar[d]_{f'}&X\ar[d]^f\\S'\ar[r]^g&S},
$$
where $S$ is an excellent strictly local scheme of characteristic zero, $S'$ is the completion of $S$, and $f$ is a morphism of finite type. Let $\ell$ be a prime number, $x\in X$, let $Z$ be the schematic closure of $X'_x$ in $X'$, and let $i:X'_x\to Z$, $j:Z\to X'$ be the canonical morphisms. Then, if $k:X'\to R$ is a closed immersion of $X'$ into an excellent regular scheme $R$ of characteristic zero, the complex
$$K'=i^*(R^!(k.j)(\ZZ/\ell\ZZ))$$
is a constant dualizing complex on $X'_x$ (that is, it has only one nonzero cohomology sheaf, isomorphic to $\ZZ/\ell\ZZ$).
\end{sublemme}

\enlargethispage{\baselineskip}%
In view of (\SGA 5 I 3.4.3), the only thing to prove is that $K'$ is constant. But, since $Z$ is excellent, the set of points of $Z$ whose\pageoriginale local rings are regular is an open subset $U$ (\EGA IV 7.8.3 \textup{(iv)}), and $U$ plainly contains $X'_x$, which is regular. Let
$$u:U\to R$$
then be the canonical immersion of $U$ into $R$; it follows from the purity theorem (\SGA 4 XIX 3.2 and 3.4) and the isomorphism
$$(\mu_\ell)_S\simeq(\ZZ/\ell\ZZ)_S$$
($S$ strictly local) that one has
$$
\R^!u(\ZZ/\ell\ZZ)\simeq\ZZ/\ell\ZZ[2c],
$$
where $c$ is a locally constant function on $U$, necessarily constant in a neighborhood of $X'_x$, because the fibers of $g$ are geometrically integral by (\EGA IV 18.9.1), and hence $X'_x$ is integral. The lemma follows at once.

\section{Converse to the affine Lefschetz theorem} \label{XIV.3}

The present section will be used in \numero \Ref{XIV.4} to prove a converse to the \og Lefschetz theorem\fg; a reader interested only in the direct part of that theorem may therefore omit the present section.

\subsection{\texorpdfstring{}{XIV.3.1}} \label{XIV.3.1}

Recall the statement of the affine Lefschetz theorem\nde{\label{gabaf}Gabber proved the following generalization. Let $Y$ be a strictly local scheme of arithmetic type over a regular Noetherian scheme $S$ of dimension $\leq 1$. Let $f:X\to Y$ be an affine morphism of finite type, $\Lambda=\ZZ/n\ZZ$ with $n$ invertible on $X$, and $F$ a $\Lambda$-sheaf. Then $H^q(X,F)=0$ if $q>\delta(F)$. The following local Lefschetz theorem follows. Let $\Oo$ be strictly local of arithmetic type over $S$. For every non-zero-divisor $f\in \Oo$ and every $\Lambda$-sheaf $F$ on $\Spec(\Oo[f^{-1}])$, one has $H^q(\Spec(\Oo[f^{-1}]),F)=0$ for $q>\dim(\Oo)$. \Cf (Fujiwara~K., {\og A Proof of the Absolute Purity Conjecture (after Gabber)\fg}, in \emph{Algebraic geometry 2000, Azumino (Hotaka)}, Adv. Stud. in Pure Math., vol.~36, 2002, p\ptbl 153-183, $\S$5), and especially Illusie's article (Illusie~L., {\og Perversité et variation\fg}, \emph{Manuscripta Math.} \textbf{112} (2003), p\ptbl 271-295). This result is one of the crucial points used by Gabber in his proof of Grothendieck's purity theorem (\cf note~\eqref{gabpur}, page~\pageref{gabpur}).} (\SGA 4~XIX~6.1~bis):

Let $S$ be a strictly local excellent scheme of characteristic zero, $f:X\to S$ an affine morphism of finite type, and $F$ a torsion sheaf\pageoriginale on $X$. Then, setting
$$
\delta(F)=\sup\{\delta(x)|x\in X \text{ and } F_{\bar{x}} \neq 0\},
$$
one has
$$
H^q(X, F)=0 \text{ for } q>\delta(F).
$$

Before stating the converse, let us prove a few lemmas.

\begin{lemme} \label{XIV.3.2} Let $K$ be a field, $\ell$ a prime number distinct from the characteristic of $K$, and $F$ a nonzero constructible sheaf of $\ell$-torsion on $K$. Suppose that the $\ell$-cohomological dimension of $K$ (\textup{\SGA 4 X~1}) is equal to $n$ (this holds, for example, if $K$ is the field of fractions of an integral strictly local excellent ring of characteristic zero and dimension $n$ (\textup{\SGA 4~XIX~6.3}), or if $K$ is an extension of finite type and transcendence degree $n$ over a separably closed field (\textup{\SGA 4~X~2.1})). Then one can find a finite separable extension $L$ of $K$ such that:
$$H^n(L, \sisi{F{|L}}{F_{|L}}) \neq 0.
$$
\end{lemme}

One can find a finite extension $K'$ of $K$ such that the restrictions of $F$ and $\mu_\ell$ to $\Spec K'$ are constant sheaves. Then $\cd_{\ell}(K') = \cd_{\ell}(K) = n$ (\SGA 4~X~2.1), and it follows from (\cite{XIV.2}~II~\S 3~Prop\ptbl4~\textup{(iii)}) that one can find a finite extension $L$ of $K'$ such that\pageoriginale
$$
H^n(L, \mu_{\ell})\neq 0, \quad\text {\ie } H^n(L, \ZZ/\ell\ZZ) \neq 0.
$$
Now the functor $H^n(L,\cdot)$ is right exact on the category of sheaves of $\ell$-torsion, since $\cd_{\ell}(L)=n$; as $F$ has a quotient isomorphic to $\ZZ/\ell\ZZ$, one also has $H^n(L, \sisi{F{|L}}{F_{|L}}) \neq 0$.

\begin{corollaire} \label{XIV.3.3}
Let $k$ be a field, $K$ an extension of finite type and transcendence degree $n$ over $k$, and $F$ a nonzero constructible sheaf of $\ell$-torsion on $K$, with $\ell$ prime to the characteristic of $k$. Then one can find a finite separable extension~$L$ of~$K$ such that, if $u: \Spec L \to \Spec k$ denotes the canonical morphism, one has
$$R^n u_*(\sisi{F{|\Spec L}}{F_{|\Spec L}}) \neq 0.
$$
\end{corollaire}

When the field $k$ is separab\sisi{le}{ly} closed, the corollary is a special case of \Ref{XIV.3.2}. In the general case, one can find a finite separable extension $k_1$ of $k\sisi{'}{}$ such that the irreducible components of $K \otimes_k k_1$ are geometrically irreducible (\EGA IV~4.5.11); let $K_1$ be one of them. If $k'$ is a separable closure of $k_1$, then $K' = K_1 \otimes_{k_1} k'$ is a field, and by (\EGA IV~4.2) one has
$$
\deg \tr K'/k' = \deg \tr K/k = n.
$$
It then follows from \Ref{XIV.3.2} that one can find a finite separable extension $L'$ of $K'$ such that $H^n(L', \sisi{F{|L'}}{F_{|L'}}) \neq 0$. But $k' = \varinjlim_i k_i$, where $k_i$ runs through the finite extensions of $k_1$ contained in $k'$, and consequently\pageoriginale $K' = \varinjlim_i (k_i \otimes_{k_1} K_1)$. It follows that one can find an index $i$ and a finite separable extension $L$ of $k_i \otimes_{k_1} K_1 =K_i$ such that $L' \simeq L \otimes_{K_i} K'$. The extension $L$ of $K$ has the required property; indeed, it follows from the commutative diagram
$$
\xymatrix{
&{\Spec L}\ar[dl]_v\ar[dr]^u&\\
{\Spec k_i} \ar[rr]^w &{}&{\Spec k}\,,
}
$$
with $w$ finite and hence $R^q w_* = 0$ if $q>0$, that
$$R^n u_*(\sisi{F{|\Spec L}}{F_{|\Spec L}}) \simeq w_*(R^n v_*(\sisi{F{|\Spec L}}{F_{|\Spec L}})).
$$
Now $R^nv_*(\sisi{F{|\Spec L}}{F_{|\Spec L}}) \neq 0$, since $H^n(L', \sisi{F{|L'}}{F_{|L'}}) \neq 0$; therefore one also has $R^nu_*(\sisi{F{|\Spec L}}{F_{|\Spec L}}) \neq 0$.

Recall the following familiar lemma (\cf $\EGA 0_{\textup{III}}$~10.3.1.2 and \EGA IV~18.2.3):

\begin{lemme} \label{XIV.3.4}
Let $X$ be a scheme, $x$ a point of $X$, and $K$ a finite separable extension of $k(x)$. Then there exist a scheme $X_1$ étale over $X$ and affine, and a point $x_1 \in X_1$ above $x$, such that $k(x_1)$ is $k(x)$-isomorphic to $K$.
\end{lemme}

In \numero \Ref{XIV.4} we shall use the following technical form of the converse to \Ref{XIV.3.1}.

\begin{proposition} \label{XIV.3.5}
Let\pageoriginale there be a Cartesian square
$$
\xymatrix{
{X'}\ar[r]^h\ar[d]_{f'} &{X}\ar[d]^f\\
{S'}\ar[r]^g &{S}\,,
}
$$
where the schemes $S$ and $S'$ are strictly local excellent schemes of characteristic zero, the morphism $f$ is locally of finite type, $g$ is regular (\textup{\EGA IV~6.8.1}) and surjective, and the closed fiber of $g$ is reduced to the closed point of $S'$. Given an $S$-scheme $X_1$ (respectively an $S$-morphism $f_1$, etc.), we denote by $X_1'$ (respectively $f_1'$, etc.) the scheme $X_1 \times_S S'$ (respectively the morphism $(f_1)_{(S')}$, etc.). Let $F$ be a sheaf of $\ZZ/m\ZZ$-modules on $X'$ ($m$ a power of a prime number $\ell$), constructible and satisfying the following conditions:
\begin{enumeratei}
\item
For every point $x \in X$, one can find a finite separable extension~$K$ of $k(x)$ such that the restriction of $F$ to the fiber $(X')_{(\Spec K)}$ comes by inverse image from a constructible sheaf on $\Spec K$.
\item
For every morphism $f_1: X_1 \to S$, with $X_1$ affine and étale over $X$, for every point $s \in S$, and for every integer $q>0$, one can find a finite separable extension~$K$ of $k(s)$ such that the restriction of $R^q f'_{1*}(F|X'_1)$ to the fiber $S'_{(\Spec K)}$ comes by inverse image from a constructible sheaf on $\Spec K$.
\end{enumeratei}

Let $n$ be an integer, and suppose that for every affine scheme $X_1$ étale over $X$ one has\pageoriginale
$$
H^i(X_1', F)=0 \text{ for } i>n.
$$
Then, if $\bar{x}'$ is a geometric point above the point $x' \in X'$ such that $F_{\bar{x}'} \neq 0$, one has
$$
\delta(x') \leq n.
$$
\end{proposition}

Let $Z'$ be the set of points $x'$ of $X'$ such that $F_{\bar{x}'} \neq 0$. Then, if $Z=h(Z')$, one has, by \textup{(i)}, $Z' = h^{-1}(Z)$; let $x' \in X'$, $x=h(x')$, $s'=f'(x')$, $s=f(x)$. It follows from \Ref{XIV.2.1.1} and the fact that the function $\delta$ decreases under specialization that it suffices to prove the inequality $\delta(x') \leq n$ when $x$ is a maximal point of $Z$ and $x'$ is such that
$$r=\deg \tr k(x)/k(s) = \deg \tr k(x')/k(s') \text{ and }d=\dim \overline{\{s\}} = \dim\overline{\{s'\}}$$
Let $x'$ be such a point; it suffices to show that one can find an affine scheme $X_1$ \'{e}tale over~$X$ such that
$$
H^{d+r}(X_1', F) \neq 0.
$$

The set $Z'$ is constructible (\SGA 4~IX~2.4), and hence so is $Z$ (\EGA IV~1.9.12); one may then suppose, after restricting $X$ to a neighborhood of~$x$, that $Z$ is an irreducible closed subset with generic point $x$. Let $T=f(Z)$; $T$ is a constructible set contained in $\overline{\{s\}}$; one can therefore find an affine open subset $U$ of~$S$ such that $s \in U$ and $T \cap U = T_U$ is an irreducible closed subset of $U$ with generic point~$s$.

Let\pageoriginale $V$ then be an affine scheme \'{e}tale over $X$, whose image in $X$ contains $x$ and whose image in $S$ is contained in $U$; let $Z_V$ be the inverse image of $Z$ in $V$ and let $u: Z_V \to T_U$ be the canonical morphism. Let\sisi{}{} $W$ be an affine scheme \'{e}tale over $U$; denote by $T_W$ the inverse image of $T_U$ in $W$ and let $X_1 = W\times_U V$. Since $F$ is zero outside $Z'$, one has the spectral sequence
$$
E^{pq}_2 = H^p((T_W)', R^qu_*'(F\sisi{}{_{|(Z_V)'}})) \To H^*(X'_1, F).
$$
We shall show that $V$ and $W$ may be chosen so that
\begin{enumeratea}
\item
$E^{pq}_2=0$ for $p>d$ and for $q>r$.
\item
$E^{dr}_2 \neq 0$.
\end{enumeratea}
\noindent
The relation $H^{d+r}(X_1', F) \neq 0$ will then follow from the spectral sequence.
\begin{enumerate}
\item[$1^\circ)$]
Set $G_q = R^q u'_*(F\sisi{}{_{|(Z_V)'}})$; then one has:
$$(G_q)_{\bar{s}'}=H^q((Z_V)'_{\bar{s}'}, F\sisi{}{{_{|(Z_V)'}}}_{\bar{s}'}),
$$
because $s'$ is a maximal point of $(T_U)'$. Since the fiber $(Z_V)'_{\bar{s}'}$ is an affine scheme of finite type of dimension $r$ over a separably closed field, it follows from \Ref{XIV.3.1} that
$$(G_q)_{\bar{s}'}=0 \text{ for } q>r.
$$

For $q>r$, let $Y'_q$ be the set of points of $(T_U)'$ where the geometric fiber\pageoriginale of $G_q$ is $\neq 0$, and let $Y_q=g(Y'_q)$; then $Y'_q = g^{-1}(Y_q)$ by \textup{(ii)}, so $Y_q$ is a constructible subset of $T_U$ (\SGA 4~XIX~5.1 and \EGA 1.9.12) that does not contain $s$; after restricting $U$ to an open neighborhood of $s$, one may suppose that $G_q=0$ for $q>r$, and hence that $E_2^{pq}=0$ for $q>r$.

Moreover, since $(T_W)'$ is an affine scheme of finite type over $g^{-1}(\overline{\{s\}})$, one has, for every $q$ (\cf \Ref{XIV.3.1}):
$$H^p((T_W)', G_q)=0 \text{ for } p>\dim g^{-1}(\overline{\{s\}})=d,
$$
whence condition a).

\item[$2^\circ)$]
Let us show that $V$ may be chosen so that $(G_r)_{\bar{s}'} \neq 0$. By~\textup{(i)}, there exists a constructible sheaf $I$, defined over a finite separable extension $K$ of $k(x)$, whose inverse image on $(X')_{(\Spec K)}$ is isomorphic to $\sisi{F{|(X')_{(\Spec K)}}}{F_{|(X')_{(\Spec K)}}}$. By \Ref{XIV.3.3}, one finds a finite separable extension $L$ of $K$ such that, if $v:\Spec L \to \Spec k(s)$ denotes the canonical morphism, then $R^r v_*(I) \neq 0$. Since the morphism $S'_s \to \Spec k(s)$ is regular, one has, by (\SGA 4~XIX~4.2):
$$R^rv_*'(\sisi{F{|(\Spec L)'}}{F_{|(\Spec L)'}}) \simeq (R^rv_*(I))' \neq 0.
$$
By Lemma \Ref{XIV.3.4}, one can find an affine scheme $X_2$ \'{e}tale over $X$, and a point~$x_2$ of $X_2$ over $x$, such that $L$ is $k(x)$-isomorphic to $k(x_2)$, and one may suppose that $X_2$ lies over $U$. Since $x$ is a maximal point of $Z$, one has\pageoriginale
$$
\Spec L \simeq \varprojlim_V Z_V,
$$
where $V$ ranges over the affine open neighborhoods of $x_2$. By passage to the limit (\SGA 4~VII~5.8), after restriction to the geometric fiber at $\bar{s}'$, one deduces:
$$
(R^rv_*'(\sisi{F{|\Spec L}}{F_{|\Spec L}})')_{\bar{s}'} = \varinjlim_V (R^ru_*'(F\sisi{}{_{|(Z_V)'}})\sisi{}{)}_{\bar{s}'},
$$
which indeed shows that one can find $V$ such that $(G_r)_{\bar{s}'} \neq 0$.

\item[$3^\circ)$]
With the scheme $V$ having been chosen in $2^\circ)$, let us show that the scheme~$W$ may be chosen so that
$$
E^{dr}_2 = H^d((T_W)', G_r) \neq 0.
$$
By \textup{(ii)}, there exists a constructible sheaf $J$, defined over a finite separable extension $K$ of $k(s)$, whose inverse image on $(S')_{(\Spec K)}$ is isomorphic to $\sisi{{G_r}{|(S')_{(\Spec K)}}}{{G_r}_{|(S')_{(\Spec K)}}}$. By Lemma \Ref{XIV.3.2}, one can find a finite separable extension $L$ of $K$ such that $H^d(\Spec L, J) \neq 0$. Since the morphism $(S')_{(\Spec L)} \to \Spec L$ is acyclic (\SGA 4~XIX~4.1 and XV~1.10 and 1.16), one has
$$
H^d((\Spec L)', \sisi{{G_r}{|(S')_{(\Spec L)'}}}{{G_r}_{|(S')_{(\Spec L)'}}}) = H^d(\Spec L, J) \neq 0.
$$
By \Ref{XIV.3.4}, one can find an affine scheme $U_1$ \'{e}tale over $U$, and a point~$s_1$ over $s$, such that $k(s_1)$ is $k(s)$-isomorphic to $L$. Since $s$ is a maximal point of~$T_U$, one has
$$
\Spec L \simeq \varprojlim_W T_W,
$$
where\pageoriginaled $W$ ranges over the affine open neighborhoods of $s_1$. One deduces that $(\Spec L)' \simeq \varprojlim_W(T_W)'$, and by passage to the limit (\SGA 4~VII~5.8):
$$
H^d((\Spec L)', \sisi{{G_r}{|(\Spec L)'}}{{G_r}_{|(\Spec L)'}}) \simeq \varinjlim_W H^d((T_W)', \sisi{{G_r}{|(T_W)'}}{{G_r}_{|(T_W)'}});
$$
consequently, one can find $W$ such that
$$
H^d((T_W)', \sisi{{G_r}{|(T_W)'}}{{G_r}_{|(T_W)'}})\neq 0,
$$
which completes the proof of the theorem.
\end{enumerate}

\begin{corollaire} \label{XIV.3.6}
The hypotheses concerning $S, S', f, f', m$ are those of \Ref{XIV.3.5}. Let $F$ now denote a complex of sheaves of $\ZZ/m\ZZ$-modules on $X'$, bounded below and with constructible cohomology, and whose cohomology sheaves satisfy conditions \textup{(i)} and \textup{(ii)} of \Ref{XIV.3.5}. Let $n$ be an integer, and suppose that, for every affine scheme $X_1$ \'{e}tale over $X$, one has
$$
H^i(X'_1, F)=0\text{ for } i>n.
$$
Then, if $\bar{x}'$ is a geometric point lying over a point $x'$ of $X'$ such that, for an integer $j$, $(\underline{H}^j(F))_{\bar{x}'} \neq 0$, one has
$$
\delta(x')\leq n-j.
$$
\end{corollaire}

Let $T'$ be the set of points of $X'$ where the conclusion of \Ref{XIV.3.7} fails, and suppose $T' \neq \emptyset$; let $T=f(T')$, let $x$ be a maximal point of $T$, and let $x'$ be a point of $X'$ lying over~$x$. Let $j$ be the largest\pageoriginale integer such that $(\underline{H}^j(F))_{\bar{x}'} \neq 0$; thus $r=\delta(x)>n-j$. Let $Z'_q$ be the set of points where the geometric fiber of $\underline{H}^q(F)$ is $\neq 0$, and let $Z_q =h(Z'_q)$; as in the proof of \Ref{XIV.3.5}, one sees that $Z_q$ is constructible. Clearly $Z'_q=\emptyset$ for $q>n$ and for $q$ sufficiently small. The other values of $q$ fall into three subsets. Let
$$
Q_1=\{q\mid x\in Z_q\text{ and a generization of }x, \text{ distinct from }x, \not\in Z_q\}.
$$
One has $j \in Q_1$, and one can find an affine open neighborhood $U_1$ of $x$ such that, for every $q \in Q_1$, $U_1 \cap Z_q$ is an irreducible closed subset with generic point $x$. If $q \in Q_1$, one has
\begin{equation*} \label{eq:XIV.3.6.*} \tag{$*$} \delta(\sisi{\underline{H}^q(F){|U_1'}}{\underline{H}^q(F)_{|U_1'}})=\delta(x)\quad\text{(for the definition of $\delta(\underline{H}^q(F))$, \cf \Ref{XIV.3.1})}.
\end{equation*}
Let
$$
Q_2=\{q\mid\text{no generization of $x$ belongs to }Z_q\}.
$$
Then, if $j<q\leq n$, one has $q\in Q_2$, and one can find an affine open neighborhood $U_2$ of $x$ such that, for every $q \in Q_2$, $Z_q \cap U_2 = \emptyset$; thus one has
\begin{equation*} \label{eq:XIV.3.6.**} \tag{$**$} \sisi{\underline{H}^q(F){|U'_2}}{\underline{H}^q(F)_{|U'_2}}=0\text{ for }q\in Q_2.
\end{equation*}
Finally, let
$$
Q_3= \{q\mid Z_q \text{ contains strict generizations of }x\}.
$$
Then one can find an affine open neighborhood of $x$, $U_3$, such that, for every $q \in Q_3$, all the maximal points of $Z_q \cap U_3$ are generizations of $x$. If $q \in Q_3$, one has
\begin{equation*} \label{eq:XIV.3.6.***} \tag{${*}{*}{*}$} \delta(\sisi{\underline{H}^q(F){|U'_3}}{\underline{H}^q(F)_{|U'_3}})\leq n-q.
\end{equation*}

For\pageoriginale every affine scheme $X_1$ \'{e}tale over $U_1 \cap U_2 \cap \sisi{{}}{U}_3$, consider the \sisi{}{h}ypercohomology spectral sequence
$$
E^{pq}_2= H^p(X_1', \underline{H}^q(F)) \To H^*(X'_1, F).
$$
One has $E^{pq}_2=0$ for $q\sisi{>}{\in Q_2}$ by \eqref{eq:XIV.3.6.**}. One has $E^{pq}_2=0$ for $p+q \geq r+j$ except possibly for $p=r, q=j$. Indeed, this is clear if $q\in Q_2$; if $q\in Q_1$, then $p>r$ except when $p=r, q=j$, and this follows from \Ref{XIV.3.1} in view of (\Ref{eq:XIV.3.6.*}); finally, if $q\in Q_3$, since $r>n-j$, one has $p>n-q$, and the assertion follows from \Ref{XIV.3.1} in view of (\Ref{eq:XIV.3.6.***}). Since $H^{r+j}(X_1', F)=0$, it follows from the spectral sequence that
$$H^r(X_1', \underline{H}^j(F))=0;$$
but, by \Ref{XIV.3.5}, this implies $\delta(x)<r$, which is absurd.

\begin{corollaire} \label{XIV.3.7}
Let $S$ be an excellent strictly local scheme of characteristic zero, $f:X \to S$ a morphism locally of finite type, $m$ a power of a prime number, $F$ a complex of sheaves of $\ZZ/m\ZZ$-modules on $X$, bounded\sisi{}{} below, with constructible cohomology, and $n$ an integer. Then the following conditions are equivalent:
\begin{enumeratei}
\item
For every affine scheme $X_1$ \'{e}tale over $X$, one has
$$
H^i(X_1, F)=0\text{ for }i>n.
$$

\item
For\pageoriginale every geometric point $\bar{x}$ lying over the point $x$ of $X$, and for every integer~$j$ such that $(\underline{H}^j(F))_{\bar{x}} \neq 0$, one has
$$
\delta(x) \leq n-j.
$$
\end{enumeratei}
\end{corollaire}

\textup{(i)} \ALORS \textup{(ii)} is the special case of \Ref{XIV.3.6} obtained by taking $S=S'$.

\textup{(ii)} \ALORS \textup{(i)} follows immediately from \Ref{XIV.3.1}, using the hypercohomology spectral sequence
$$
H^p(X_1, \underline{H}^q(F)) \To H^*(X_1, F).
$$

\section{Main theorem and variants} \label{XIV.4}

\setcounter{subsection}{-1}

\subsection{\texorpdfstring{}{XIV.4.0}} \label{XIV.4.0}
Let $g\colon X\to S$ be a separated morphism of finite type, let $T$ be a closed subset of $S$, let \hbox{$Z=g^{-1}(T)$}, and let $F$ be a complex of abelian sheaves on $X$, bounded below. We call the group
$$
\H_{Z!}^{i}(X/S, F)=\H_{T}^{i}(S, \R_{!}g(F))
$$
the \emph{$i$-th cohomology group of $F$, with proper support and support in~$Z$}, where $\R_{!}g$ denotes the \og direct image with proper support\fg (\SGA 4~XVII). In the particular case in which $g$ is proper, one simply has
$$
\H_{Z!}^{i}(X/S, F)=\H_{Z}^{i}(X, F).
$$

\begin{proposition} \label{XIV.4.1} Let $f\colon U\to S$ be a morphism of finite type and let $F$ be a complex of abelian sheaves on $U$, bounded below. Suppose\pageoriginale that a factorization of $f$ is given:
$$
\xymatrix@C=15pt{ U\ar[rr]^-{i}\ar[rd]_-{f} & & X\ar[ld]^-{g}\\
& S & }$$
where $i$ is an open immersion and $g$ is a separated morphism of finite type, and let $G$ be a complex of abelian sheaves on $X$, bounded below, that extends $F$. Let $Y$ be a closed subscheme of $X$ whose underlying space is $X-U$, so that one has a commutative diagram:
$$
\xymatrix@C=15pt{ Y\ar[rr]^-{j}\ar[rd]_-{h} & & X\ar[ld]^-{g}\\
& S & }$$
Finally, let $n$ be an integer and let $T$ be a closed subset of $S$. Then the following conditions are equivalent:
\begin{enumeratei}
\item
One has $\prof_{T}(\R_{!}f(F))\geq n$.
\item
The canonical morphism
$$
\underline{H}_{T}^{i}(\R_{!}g(G))\to\underline{H}_{T}^{i}(\R_{!}h(j^{\ast}G))
$$
is bijective for $i<n-1$ and injective for $i=n-1$.
\item
For every scheme $S^{\prime}$ \`etale over $S$, if $X^{\prime}$ (\resp $f^{\prime}$, \resp etc.) denotes the scheme $X\times_{S}S^{\prime}$ (\resp the morphism $f_{(S^{\prime})}$, \resp etc.), the canonical morphism\pageoriginale
$$
\H_{g^{\prime -1}(T^{\prime})!}^{i}(X^{\prime}/S^{\prime}, G^{\prime})\to\H_{h^{\prime -1}(T^{\prime})!}^{i}(Y^{\prime}/S^{\prime}, j^{\prime\ast}G^{\prime})
$$
is bijective for $i<n-1$ and injective for $i=n-1$.
\end{enumeratei}
\end{proposition}

In the derived category $\D^{+}(X)$ (\cf \cite{XIV.3}), consider the distinguished triangle
$$
\xymatrix@C=15pt{ & j_{\ast}j^{\ast}G\ar[ld] & \\
i_{!}F\ar[rr] & & G.\ar[lu]}$$
Applying the functor $\R_{!}g$ to this triangle gives the triangle
\begin{equation*} \label{eq:XIV.4.1.*} \tag{$*$}
\begin{array}{c}
\xymatrix@C=0pt{
& \R_{!}h(j^{\ast}G)\ar[ld] & \\
\R_{!}f(F)\ar[rr] & & \R_{!}g(G).\ar[lu]
}
\end{array}
\end{equation*}

Let us show that \textup{(i)}\SSI\textup{(ii)}. Indeed, by Definition \Ref{XIV.1.2}, \textup{(i)} is equivalent to the relation
$$
\underline{H}_{T}^{i}(\R_{!}f(F))=0\quad\textrm{ for }i<n\textrm{;}$$
while \eqref{eq:XIV.4.1.*} gives the exact sequence of sheaves
$$
\to\underline{H}_{T}^{i}(\R_{!}f(F))\to\underline{H}_{T}^{i}(\R_{!}g(G))\to\underline{H}_{T}^{i}(\R_{!}h(j^{\ast}G))\to,
$$
which proves\pageoriginale the equivalence of \textup{(i)} and \textup{(ii)}.

\textup{(i)}\SSI\textup{(iii)}. Indeed, \textup{(i)} is equivalent to saying that, for every scheme $S^{\prime}$ \`etale over $S$, one has
\begin{equation*} \label{eq:XIV.4.1.**} \tag{$**$} {\H_{T^{\prime}}^{i}(S^{\prime}, \R_{!}f^{\prime}(F^{\prime}))=0\quad\textrm{ for }i<n.}
\end{equation*}
But \eqref{eq:XIV.4.1.*} gives the exact sequence of abelian groups
$$
\to\H_{T^{\prime}}^{i}(S^{\prime}, \R_{!}f^{\prime}(F^{\prime}))\to\H_{T^{\prime}}^{i}(S^{\prime}, \R_{!}g^{\prime}(G^{\prime}))\to \H_{T^{\prime}}^{i}(S^{\prime}, \R_{!}h^{\prime}(j^{\prime\ast}G^{\prime}))\to\;;$$
in view of \Ref{XIV.4.0}, this exact sequence takes the form
$$
\to\H_{T^{\prime}}^{i}(S^{\prime}, \R_{!}f^{\prime}(F^{\prime}))\to\H_{g^{\prime -1}(T^{\prime})!}^{i}(X^{\prime}/S^{\prime}, G^{\prime})\to\H_{h^{\prime -1}(T^{\prime})!}^{i}(Y^{\prime}/S^{\prime}, j^{\prime\ast}G^{\prime})\sisi{}{\to.}$$
The equivalence of \textup{(i)} and \textup{(iii)} follows, taking into account the form \eqref{eq:XIV.4.1.**} of \textup{(i)}.

\setcounter{subsubsection}{-1} \setcounter{subsection}{2}

\subsubsection{} \label{XIV.4.2.0}
When $f\colon U\to S$ is affine, we shall give \emph{local} conditions on $F$ under which conditions \textup{(i)} through \textup{(iii)} of \Ref{XIV.4.1} hold. In what follows, all schemes considered are excellent schemes of characteristic zero, and all sheaves are sheaves of $\ZZ/m\ZZ$-modules, where $m$ is a power of a prime number. If resolution of singularities in the sense of (\SGA 4~XIX) were available, the stated results and their proofs would remain valid for excellent schemes of equal characteristic, with $m$ prime to the characteristic.

\setcounter{subsection}{1}

\begin{theoreme} \label{XIV.4.2}
Let\pageoriginale $S$ be an excellent scheme of characteristic zero and let \hbox{$f\colon U\!\to\! S$} be a separated morphism of finite type. Let $F$ be a bounded-below complex of sheaves of $\ZZ/m\ZZ$-modules on $U$ with constructible cohomology, let $n$ be an integer, and let~$T$ be a closed subset of $S$. Then the following conditions are equivalent:
\begin{enumeratei}
\item
For every scheme $U_{1}$ \`etale over $U$ and affine over $S$, denoting by~$f_{1}$ the structural morphism of $U_{1}$ and by $F_{1}$ the restriction of $F\sisi{_{1}}{}$ to $U_{1}$, one has
$$
\prof_{T}(\R_{!}f_{1}(F_{1}))\geq n$$
(\cf Proposition\ptbl \Ref{XIV.4.1} for the meaning of this relation).
\item
For every point $u$ of $U$, one has
$$
\prof_{u}(F)\geq n-\delta_{T}(u),
$$
where, writing $s=f(u)$, one sets (\cf \Ref{XIV.2.2}): $\delta_{T}(u)=\degtr(k(u)/k(s))+\codim\big(\overline{\{ s\} }\cap T, \overline{\{ s\} } \big)$.
\end{enumeratei}
\end{theoreme}

\begin{proof}
$1^\circ)$ Let $t$ be a point of $T$, let $\overline{S}$ be the strict localization of $S$ at a geometric point above $t$, and let $S^{\prime}$ be the completion of $\overline{S}$, with closed point $t^{\prime}$; then $\overline{S}$ is excellent by (\EGA IV~7.9.5), so $S^{\prime}$ is an excellent complete strictly local scheme. Given a scheme $U$ over $S$ (respectively an $S$-morphism $f$, respectively etc.), we shall denote by $U^{\prime}$ (respectively $f^{\prime}$, respectively etc.) the scheme $U\times_{S}S^{\prime}$ (respectively the morphism $f_{(S^{\prime})}$, respectively etc.). We have the Cartesian square
$$
\xymatrix{ U^{\prime}\ar[r]^-{h}\ar[d]_-{f^{\prime}} & U\ar[d]^-{f}\\
S^{\prime}\ar[r]^-{g} & S, }
$$
in\pageoriginale which the morphism $g$ is regular (\EGA IV~7.8.2). Let us show that it is enough to prove that (for every point $t\in T$) the following two properties are equivalent:
\begin{itemize}
\item[$\textup{(i)}_{t}$] For every scheme $U_{1}$ étale over $U$ and affine over $S$, setting $f_{1}\colon U_{1}\to S$, one has
$$
\prof_{t^{\prime}}(\R_{!}f_{1}^{\prime}(F_{1}^{\prime}))\geq n.
$$

\item[$\textup{(ii)}_{t}$] For every point $u^{\prime}$ of $U^{\prime}$, one has
$$
\prof_{u^{\prime}}(F^{\prime})\geq n-\delta_{t^{\prime}}(u^{\prime}).
$$
\end{itemize}
It is enough to prove the following lemma:

\begin{sublemme} \label{XIV.4.2.1}
One has \textup{(i)}\SSI$\textup{(i)}_{t}$ for every $t\in T$ and \textup{(ii)}\SSI$\textup{(ii)}_{t}$ for every $t\in T$.
\end{sublemme}

\textup{(i)}\SSI$\textup{(i)}_{t}$ for every $t\in T$. Indeed, \textup{(i)} is equivalent to saying that, for every scheme $U_{1}$ étale over $U$ and affine over $S$, one has
$$
\prof_{T}(\R_{!}f_{1}(F_{1}))\geq n\;;$$
now, by \Ref{XIV.1.8},
$$
\prof_{T}(\R_{!}f_{1}(F_{1}))=\inf_{t\in T}\prof_{t}(\R_{!}f_{1}(F_{1})).
$$
Since $g^{\ast}(\R_{!}f_{1}(F_{1}))\simeq\R_{!}f_{1}^{\prime}(F_{1}^{\prime})$ (\SGA 4~XVII), one has, by \Ref{XIV.1.16},
$$
\prof_{t}(\R_{!}f_{1}(F_{1}))=\prof_{t^{\prime}}(\R_{!}f_{1}^{\prime}(F_{1}^{\prime})),
$$
and hence\pageoriginale \textup{(i)} is equivalent to saying that, for every $t\in T$, one has $\prof_{t^{\prime}}(\R_{!}f_{1}^{\prime}(F_{1}^{\prime}))\geq n$, which is precisely $\textup{(i)}_{t}$.

\textup{(ii)}$_{t}$ for every $t\in T$\ALORS\textup{(ii)}. Indeed, let $u\in U$; we must show the relation
$$
\prof_{u}(F)\geq n-\delta_{T}(u),
$$
where $\delta_{T}(u)=\inf_{t\in T\cap\overline{\{ s\} }}\delta_{t}(u)$ (\cf \Ref{XIV.2.2}); we are therefore reduced to showing that, for every $t\in T\cap\overline{\{ s\} }$, one has
$$
\prof_{u}(F)\geq n-\delta_{t}(u).
$$
Let $u^{\prime}$ be a point of $U^{\prime}$ such that $h(u^{\prime})=u$ and $\delta_{t^{\prime}}(u^{\prime})=\delta_{t}(u)$ (\cf \Ref{XIV.2.1.1}). Since $h$ is locally acyclic (\SGA 4~XIX~4.1), it follows from \Ref{XIV.1.16} and the fact that~$u^{\prime}$ is a generic point of $U_{u}^{\prime}$ that
$$
\prof_{u^{\prime}}(F^{\prime})=\prof_{u}(F).
$$
But by $\textup{(ii)}_{t}$ one has $\prof_{u}(F)=\prof_{u^{\prime}}(F^{\prime})\geq n-\delta_{t}(u)$, which proves \textup{(ii)}.

\textup{(ii)}\ALORS$\textup{(ii)}_{t}$ for every $t$. With the notation of \Ref{XIV.2.2.1}, for every point $u^{\prime}$ of $U^{\prime}$ one has, by \Ref{XIV.1.16},
$$
\prof_{u^{\prime}}(F^{\prime})\geq\prof_{u}(F)+2\dim h_{u^{\prime}}^{-1}(u)\geq\prof_{u}(F)+\dim h_{u^{\prime}}^{-1}(u).
$$
Taking into account \Ref{XIV.2.2.1} and \textup{(ii)}, we obtain
$$
\prof_{u^{\prime}}(F^{\prime})\geq n-\delta_{T}(u)+\dim h_{u^{\prime}}^{-1}(u)\geq n-\delta_{t^{\prime}}(u^{\prime}),
$$
which\pageoriginaled is precisely $\textup{(ii)}_{t}$.

$2^\circ)$ $\textup{(ii)}_{t}$\SSI$\textup{(i)}_{t}$. We reduce immediately to the case where $F$ is bounded, by truncating $F$ at a sufficiently high degree. We can realize $S^{\prime}$ as a closed subscheme of a complete regular local scheme, hence an excellent one; it then follows from (\SGA 5~I~3.4.3) that there exists a dualizing complex $K$ on $S^{\prime}$ and that $\R^{!}f^{\prime}(K)=K^{\prime}$ is a dualizing complex on $U^{\prime}$. We choose $K$ so that $\delta^{K}(t^{\prime})=0$ (for the definition of $\delta^{K}(t^{\prime})$, \cf \Ref{XIV.2.3}), and denote by $\D F^{\prime}$ the dual of $F^{\prime}$ with respect to $K^{\prime}$. We can reformulate the hypothesis $\textup{(ii)}_{t}$ as follows:

\begin{sublemme} \label{XIV.4.2.2}
Let $u^{\prime}$ be a point of $U^{\prime}$; then the following conditions are equivalent:
\begin{enumeratei}
\item
One has $\prof_{u^{\prime}}(F^{\prime})\geq n-\delta_{t^{\prime}}(u^{\prime})$.
\item
One has $\big(\underline{H}^{q}(\D F^{\prime}) \big)_{\overline{u}^{\prime}}=0$ for $q>-n-\delta_{t^{\prime}}(u^{\prime})$ ($\overline{u}^{\prime}$ a geometric point above $u^{\prime}$).
\end{enumeratei}
\end{sublemme}

Let $\overline{U}^{\prime}$ be the strict localization of $U^{\prime}$ at $\overline{u}^{\prime}$, and let $\overline{F}^{\prime}$ be the inverse image of $F^{\prime}$ under the morphism $\overline{U}^{\prime}\to U^{\prime}$. The relation $\prof_{u^{\prime}}(F^{\prime})\geq n-\delta_{t^{\prime}}(u^{\prime})$ is, by definition, equivalent to the following:
\begin{equation*} \label{eq:XIV.4.2.*} \tag{$*$} {\H_{\overline{u}^{\prime}}^{i}(\overline{F}^{\prime})=0\quad\textrm{ for }i<n-\delta_{t^{\prime}}(u^{\prime}).}
\end{equation*}
Let $\D\big(\H_{\overline{u}^{\prime}}^{i}(\overline{F}^{\prime})\big)$ be the dual of the abelian group $\H_{\overline{u}^{\prime}}^{i}(\overline{F}^{\prime})$ with respect to $\ZZ/m\ZZ$. By \Ref{XIV.2.3.2}, $K^{\prime}[-2\delta_{t^{\prime}}(u^{\prime})]=K^{\prime\prime}$ satisfies $\delta^{K^{\prime\prime}}(u^{\prime})=0$; since $F^{\prime}$ has constructible cohomology, one has $\overline{\D F^{\prime}}=\D(\overline{F}^{\prime})$, and the local\pageoriginale duality theorem (\SGA 5~I~4.5.3) then shows that
$$
\D\big(\H_{\overline{u}^{\prime}}^{i}(\overline{F}^{\prime})\big) \simeq\big(\H^{-i-2\delta_{t^{\prime}}(u^{\prime})}(\D F^{\prime})\big)_{\overline{u}^{\prime}}.
$$
Thus \eqref{eq:XIV.4.2.*} is equivalent to the relation
\begin{equation*} \label{eq:XIV.4.2.**} \tag{$**$}
\big(\underline{H}^{q}(\D F^{\prime})\big)_{\overline{u}^{\prime}}=0\quad\text{for }q>-n-\delta_{t^{\prime}}(u^{\prime}).
\end{equation*}

We are now in a position to prove the theorem. The relation \textup{(ii)}$_{t}$ is equivalent to relation \eqref{eq:XIV.4.2.**}. Let $G^{q}=\underline{H}^{q}(\D F^{\prime})$; the affine Lefschetz theorem (\Ref{XIV.3.1}) implies in particular that, for every scheme $U_{1}$ étale over $U$ and affine over $S$, one has
$$
\H^{p}(U_{1}^{\prime}, G^{q})=0\quad\textrm{ for }p>\delta(G^{q}),
$$
where $\delta(G^{q})$ is the upper bound of the $\delta_{t^{\prime}}(u^{\prime})$ for those $u^{\prime}$ such that $G_{\overline{u}^{\prime}}^{q}\neq 0$; by \eqref{eq:XIV.4.2.**}, one has $\delta(G^{q})\leq -n-q$, so $\textup{(ii)}_{t}$ implies the relation
$$
\H^{p}(U_{1}^{\prime}, \underline{H}^{q}(\D F^{\prime}))=0\quad\textrm{ for }p>-q-n.
$$
Taking into account the hypercohomology spectral sequence for the functor \og sections on $U_{1}^{\prime}$\fg with respect to the complex $\D F^{\prime}$:
$$
\E_{2}^{pq}=\H^{p}(U_{1}^{\prime}, \underline{H}^{q}(\D F^{\prime}))\To\H^{\ast}(U_{1}^{\prime}, \D F^{\prime}),
$$
we obtain the relation
\begin{equation*} \label{eq:XIV.4.2.***} \tag{${*}{*}{*}$}
\H^{i}(U_{1}^{\prime}, \D F^{\prime})=0\quad\textrm{ for }i>-n.
\end{equation*}

Conversely,\pageoriginale suppose that the preceding relation holds for every $U_{1}$ étale over $U$ and affine over $S$. Apply Proposition \Ref{XIV.3.6}, replacing $S$ there by $\overline{S}^{\prime}$; the hypotheses of \Ref{XIV.3.6} concerning $\sisi{S}{\overline{S}}$ are satisfied, since, for every affine scheme $\overline{U}_{1}$ étale over $\overline{U}$, one can find a scheme above $\overline{U}_{1}$ which comes by inverse image from a scheme étale over $U$ and affine over $S$; as for the hypotheses concerning $F$, they are satisfied by \Ref{XIV.2.3.3}. Thus, for every point $u^{\prime}$ of $U^{\prime}$ such that $\big(\underline{H}^{q}(\D F^{\prime})\big)_{\overline{u}^{\prime}}\neq 0$, one has:
$$
\delta_{t^{\prime}}(u^{\prime})\leq -n-q,
$$
which is precisely relation \eqref{eq:XIV.4.2.**}; we have therefore proved the equivalence
$$
\text{\textup{(ii)}$_{t}$\SSI\eqref{eq:XIV.4.2.***}}.
$$
We shall transform relation \eqref{eq:XIV.4.2.***}; first, one has
$$
\H^{i}(U_{1}^{\prime}, \D F^{\prime})=\big(\underline{H}^{i}(\R f_{1\ast}^{\prime}(\D F_{1}^{\prime}))\big)_{t^{\prime}}\;;$$
but by (\SGA 5~I~1.12), there is a canonical isomorphism
$$
\R f_{1\ast}^{\prime}(\D F_{1}^{\prime})\simeq\D\big(\R_{!}f_{1}^{\prime}(F_{1}^{\prime})\big),
$$
where $\D\big(\R_{!}f_{1}^{\prime}(F_{1}^{\prime})\big)$ denotes the dual of $\R_{!}f_{1}^{\prime}(F_{1}^{\prime})$ with respect to $K$. We thus see that $\textup{(ii)}_{t}$ is equivalent to
$$
\big(\underline{H}^{i}(\D(\R_{!}f_{1}^{\prime}(F_{1}^{\prime})))\big)_{t^{\prime}}=0\quad\textrm{ for }i>-n.
$$
Applying\pageoriginale the local duality theorem (\SGA 5~I~4.5.3) once again, but this time at the point $t^{\prime}$, we find that
$$
\big(\underline{H}^{i}(\D(\R_{!}f_{1}^{\prime}(F_{1}^{\prime})))\big)_{t^{\prime}}\simeq\D\big(\H_{t^{\prime}}^{-i}(\R_{!}f_{1}^{\prime}(F_{1}^{\prime}))\big),
$$
and finally $\textup{(ii)}_{t}$ is equivalent to the relation
$$
\H_{t^{\prime}}^{i}(\R_{!}f_{1}^{\prime}(F_{1}^{\prime}))=0\quad\textrm{ for }i<n,
$$
that is, $\prof_{t^{\prime}}(\R_{!}f_{1}^{\prime}(F_{1}^{\prime}))\geq n$, which completes the proof of the theorem.
\skipqed
\end{proof}

\begin{subremarque} \label{XIV.4.2.3}
The argument simplifies quite considerably when one assumes that $S$ admits (at least locally) a dualizing complex (for example, that it is locally immersible in a regular scheme). This avoids recourse to a completion (the passage to the case where $S$ is strictly local being immediate), to \Ref{XIV.2.3.3}, and to the rather unappealing technical statement \Ref{XIV.3.6}, which may then be replaced by the more congenial reference \Ref{XIV.3.7}.
\end{subremarque}

\begin{corollaire} \label{XIV.4.3}
Let $S$ be an excellent scheme of characteristic zero and $f\colon U\to S$ a separated morphism of finite type such that $U$ is the \emph{union of $c+1$ open subsets affine over~$S$}. Let $F$ be a complex of sheaves of $\ZZ/m\ZZ$-modules, bounded below and with constructible cohomology, let $n$ be an integer, and let $T$ be a closed subset of $S$. Suppose that, for every point $u\in U$, one has
$$
\prof_{u}(F)\geq n-\delta_{T}(u).
$$
Then one has\pageoriginale
$$
\prof_{T}(\R_{!}f(F))\geq n-c.
$$
\end{corollaire}

Indeed, let $U_{\sisi{i}{j}}$, $0\leq \sisi{i}{j}\leq c$, be a covering of $U$ by open subsets $U_{\sisi{i}{j}}$ affine over~$S$. Adopting again the notation of the proof of \Ref{XIV.4.2}, one has, for every $\sisi{i}{j}$,
$$
\H^{i}(U_{\sisi{i}{j}}^{\prime}, \underline{H}^{q}(\D F^{\prime}))=0\quad\textrm{ for }i>-n-q.
$$
Using the spectral sequence that relates the cohomology of $U$ to that of the covering formed by the $U_{\sisi{i}{j}}$ (\SGA 4~V~2.4), the preceding relation shows that
$$
\H^{i}(U^{\prime}, \underline{H}^{q}(\D F^{\prime}))=0\quad\textrm{ for }i>-n-q+c.
$$
The corollary then follows from the end of the proof of \Ref{XIV.4.2}.

\begin{corollaire} \label{XIV.4.4}
Let $S$ be an excellent scheme of characteristic zero, $g\colon X\to S$ a morphism, $U$ an open subset of $X$ that is the union of $c+1$ open subsets affine over $S$, $Y$ a closed subscheme with underlying space $X-U$, and $j\colon Y\to X$ the natural morphism. Let~$F$ be a complex of sheaves of $\ZZ/m\ZZ$-modules on $X$, bounded below and with constructible cohomology, let $T$ be a closed subset of $S$, and let $n$ be an integer. Suppose that, for every point $u$ of $U$, one has
$$
\prof_{u}(F)\geq n-\delta_{T}(u).
$$
Then\pageoriginale the canonical morphism
$$
\H_{g^{-1}(T)!}^{i}(X/S, F)\to\H_{(g^{-1}(T)\cap Y)!}^{i}(Y/S, j^{\ast}F)$$
is bijective for $i<n-c-1$, and injective for $i=n-c-1$.
\end{corollaire}

This follows immediately from \Ref{XIV.4.1} and \Ref{XIV.4.3}.

\begin{corollaire}[Local Lefschetz theorem] \label{XIV.4.5}
Let $S$ be an excellent henselian local scheme of characteristic zero, $t$ the closed point of $S$, $X$ a scheme proper over $S^{\prime}=S-\{ t\}$, and $U$ an open subset of $X$ that is the union of $c+1$ affine open subsets. Let $Y$ be a closed subscheme of $X$ with underlying space $X-U$, let $j\colon Y\to X$ be the canonical morphism, let $F$ be a complex of sheaves of $\ZZ/m\ZZ$-modules on $X$, bounded below and with constructible cohomology, and let $n$ be an integer. Suppose that, for every point $u$ of~$U$, one has
$$
\prof_{u}(F)\geq n-\delta_{t}^{\prime}(u), \quad\text{where } \delta_{t}^{\prime}(u)=\delta_{t}(u)-1.
$$
Then the canonical morphism
$$
\H^{i}(X, F)\to\H^{i}(Y, j^{\ast}F)$$
is bijective for $i<n-c-1$, and injective for $i=n-c-1$.
\end{corollaire}

Let $f\colon U\to S$ be the canonical morphism; it follows from
\Ref{XIV.4.2}, applied with $n$ replaced by $n+1$, that
$$
\prof_{t}(\R_{!}f(F\sisi{|U}{_{|U}}))\geq n+1-c.
$$
The preceding\pageoriginale relation shows that the canonical morphism
$$
\H^{i}(S, \R_{!}f(F\sisi{|U}{_{|U}}))\to\H^{i}(S^{\prime}, \R_{!}f(F\sisi{|U}{_{|U}}))$$
is bijective for $i<n-c$, and injective for $i=n-c$. Since $\R_{!}f(F\sisi{|U}{_{|U}})$ is zero outside $S^{\prime}$, one has $\H^{i}(S, \R_{!}f(F\sisi{|U}{_{|U}}))\simeq\big(\underline{H}^{i}(\R_{!}f(F\sisi{|U}{_{|U}}))\big)_{t}=0$, and consequently
\begin{equation*} \label{eq:XIV.4.5.*} \tag{$*$}
\H^{i}(S^{\prime}, \R_{!}f(F\sisi{|U}{_{|U}}))=0\quad\text{for }i<n-c.
\end{equation*}
Let $g\colon X\to S^{\prime}$, $h\colon Y\to S^{\prime}$, $f^{\prime}\colon U\to S^{\prime}$ be the canonical morphisms. It follows from the distinguished triangle
$$
\xymatrix@C=0pt{ & \R h_{\ast}(j^{\ast}F)\ar[ld] & \\
\R_{!}f^{\prime}(F\sisi{|U}{_{|U}})\ar[rr] & & \R g_{\ast}(F)\ar[lu] }$$
that condition \eqref{eq:XIV.4.5.*} is equivalent to the assertion that the morphism
$$
\H^{i}(S^{\prime}, \R g_{\ast}(F))\to\H^{i}(S^{\prime}, \R h_{\ast}(j^{\ast}F))$$
is bijective for $i<n-c-1$, and injective for $i=n-c-1$. Since this morphism identifies canonically with the morphism
$$
\H^{i}(X, F)\to\H^{i}(Y, j^{\ast}F),
$$
the conclusion follows at once.

\begin{corollaire}[Global Lefschetz theorem] \label{XIV.4.6}
Let\pageoriginale $S$ be the spectrum of a field, $X$ a scheme proper over $S$, and $U$ an open subset of $X$ that is the union of $c+1$ affine open subsets. Let~$Y$ be a closed subscheme of $X$ with underlying space $X-U$, let $j\colon Y\to X$ be the canonical morphism, let $F$ be a complex of sheaves of $\ZZ/m\ZZ$-modules on $X$, bounded below and with constructible cohomology, and let $n$ be an integer. Suppose that, for every point $u$ of $U$, one has
$$
\prof_{u}(F)\geq n-\dim\big(\overline{\{ u\} }\big).
$$
Then the canonical morphism
$$
\H^{i}(X, F)\to\H^{i}(Y, j^{\ast}F)$$
is bijective for $i<n-c-1$, and injective for $i=n-c-1$.

More generally, if $g\colon X\to S$ is a separated morphism of finite type, with the hypotheses on $S$, $U$, $Y$, and $F$ being the same as above, then the canonical morphism
$$
\H_{!}^{i}(X/S, F)\to\H_{!}^{i}(Y/S, j^{\ast}F)$$
(where $\H_{!}^{i}$ denotes cohomology with proper support, that is, $\H_{!}^{i}(X\sisi{}{/S}, F)=\H^{i}(S, \R_{!}g(F))$) is bijective for $i<n-c-1$, and injective for $i=n-c-1$.
\end{corollaire}

The corollary is a special case of \Ref{XIV.4.4}, with $T=S$.

Here is a partial converse to \Ref{XIV.4.3}:

\begin{proposition} \label{XIV.4.7}
Let\pageoriginale $S$ be a Noetherian scheme and let $f\colon U\to S$ be a morphism of finite type. Suppose that there exists a dualizing complex $K$ on $S$ and that $\R^{!}f(K)$ is a dualizing complex on $U$. Let $T$ be a closed subset of $S$ and let $c$ be an integer. Then the following conditions are equivalent:
\begin{enumeratei}
\item
For every bounded-below complex $F$ of sheaves of $\ZZ/m\ZZ$-modules on $U$ with constructible cohomology, and every integer $n$ such that, for every point $u$ of $U$,
$$
\prof_{u}(F)\geq n-\delta_{T}(u),
$$
one has
$$
\prof_{T}(\R_{!}f(F))\geq n-c.
$$

\item
For every constructible sheaf $G$ of $\ZZ/m\ZZ$-modules on $U$, and every point $t\in T$, one has
$$
\big(\R^{p}f_{\ast}(G)\big)_{\overline{t}}=0\quad\textrm{ for }p>\delta(G, f, t)+c$$
(recall from (\textup{\SGA \sisi{X~}{}4~XIX~6.0}) that $\delta(G, f, t)=\sup\{ \delta_{t}(u)|t\in\overline{\{ u\} }\textrm{ and }G_{\overline{u}}\neq 0\}$).
\end{enumeratei}
\end{proposition}

N.B. Condition \textup{(ii)} is satisfied by \Ref{XIV.3.1} if $f$ is separated and, locally on $S$ for the \`etale topology, $U$ is a union of $c+1$ open subsets affine over $S$; thus \Ref{XIV.4.7} contains \Ref{XIV.4.3}\sfootnote{At least in the \sisi{case where}{case where} $S$ locally admits a dualizing complex, \sisi{e.g.}{for example} when $S$ is locally embeddable in a regular scheme.}.

We may plainly assume that $S$ is local and that $T$ is the closed point $t$ of~$S$. The proof of \textup{(ii)}\ALORS\textup{(i)} is essentially identical\pageoriginale to part $2^\circ)$ of the proof of \Ref{XIV.4.2}. Let us briefly prove that \textup{(i)}\ALORS \textup{(ii)}. Set $n=\delta(G,f,t)$. The local duality theorem (\SGA 5~I~4.3.2), applied to $\D G$, shows that
$$
\D\big(\H_{\overline{u}}^{i}(\D G)\big) \simeq\big(\underline{H}^{-i-2\delta_{t}(u)}(G)\big)_{\overline{u}}.
$$
Since $G$ is concentrated in degree $0$, one therefore has $\H_{\overline{u}}^{i}(\D G)=0$ except possibly for $i=-2\delta_{t}(u)$; more precisely,
\sisi{\begin{align*}
\prof_{u}(\D G)=-2\delta_{t}(u) &\;\textrm{ if }\; G_{\overline{u}}\neq 0, \\
\prof_{u}(\D G)=\infty &\;\textrm{ if }\; G_{\overline{u}}=0.
\end{align*}}
{\[
\prof_{u}(\D G)=
\begin{cases}
-2\delta_{t}(u) &\text{if } G_{\overline{u}}\neq 0, \\
\infty &\text{if } G_{\overline{u}}=0.
\end{cases}
\]}
It follows that, for every $u\in U$,
$$
\prof_{u}(\D G)\geq -n-\delta_{t}(u).
$$
Hypothesis \textup{(i)} then gives $\prof_{t}(\R_{!}f(\D G))\geq -n-c$. We transform this relation using the isomorphism $\R_{!}f(\D G)\simeq\D(\R f_{\ast}(G))$ (\SGA 5 I~1.12) and applying the local duality theorem at $t$; this gives
$$
\big(\underline{H}^{i}(\R f_{\ast}(G)\big)_{\overline{t}}=0\quad\textrm{ for }i>n+c,
$$
which is precisely \textup{(ii)}.

\subsection{\texorpdfstring{}{XIV.4.8}} \label{XIV.4.8} Under the hypotheses of \Ref{XIV.4.4} with $g$ proper (\resp \Ref{XIV.4.5}, \resp \Ref{XIV.4.6} with $g$ proper), if $V$ is an open neighborhood of $Y$ in $X$, the\pageoriginaled morphism
$$
\H^{i}(V, F)\to\H^{i}(Y, j^{\ast}F)$$
is bijective for $i<n-c-1$ and injective for $i=n-c-1$. Indeed, if $\sisi{i}{\iota}\colon V\to X$ is the canonical morphism, it is enough to apply \Ref{XIV.4.4} (\resp \Ref{XIV.4.5}, \resp \Ref{XIV.4.6}) to the complex $\R \sisi{i}{\iota}_{\ast}(\sisi{F{|V}}{F_{|V}})$. One may ask whether the preceding morphism is bijective for $i=n-c-1$ and injective for $i=n-c$. It plainly suffices that the hypotheses hold with $n$ replaced by $n+1$; the following proposition shows that slightly less is enough.

\begin{proposition} \label{XIV.4.9}
Let $S$ be an excellent local scheme of characteristic zero with closed point $t$ (\resp assume in addition to the preceding conditions that $S$ is Henselian), let $f\colon X\to S$ be a proper morphism (\resp proper over $S-\{ \sisi{s}{t}\}$), and let $U$ be an open subset of $X$ that is the union of $c+1$ affine open subsets. Let $Y$ be a closed subscheme of $X$ whose underlying space is $X-U$, let $j\colon Y\to X$ be the canonical morphism, let $F$ be a bounded-below complex of sheaves of $\ZZ/m\ZZ$-modules on $X$ with constructible cohomology, and let $n$ be an integer. Suppose that, for every point $u$ of $U$,
$$
\prof_{u}(F)\geq\inf(n-1, n-\delta_{t}(u))\quad(\textrm{\resp }\;\prof_{u}(F)\geq\inf(n-1, n+1-\delta_{t}(u))).
$$
Then, for every open neighborhood $V$ of $Y$ in $X$, the canonical morphism\pageoriginale
$$
\H_{f^{-1}(t)}^{i}(V, F)\to\H_{f^{-1}(t)\cap Y}^{i}(Y, j^{\ast}F)\quad(\textrm{\resp }\;\H^{i}(V, F)\to\H^{i}(Y, j^{\ast}F))
$$
is bijective for $i<n-c-2$ and injective for $i=n-c-2$. Moreover, there exists an open neighborhood $V_{0}$ of $Y$ in $X$ such that, for every such $V$ with $V\subset V_{0}$, the canonical morphism
$$
\H_{f^{-1}(t)\cap V}^{i}(V, F)\to\H_{f^{-1}(t)\cap Y}^{i}(Y, j^{\ast}F)\quad(\textrm{\resp }\;\H^{i}(V, F)\to\H^{i}(Y, j^{\ast}F))
$$
is bijective for $i<n-c-1$ and injective for $i=n-c-1$.
\end{proposition}

\begin{proof}
For simplicity, set $\delta_{t}^{\prime}(u)=\delta_{t}(u)$ (\resp $\delta_{t}^{\prime}(u)=\delta_{t}(u)-1$). The first assertion of \Ref{XIV.4.9} follows from \Ref{XIV.4.8}, since the hypotheses of \Ref{XIV.4.4} (\resp \Ref{XIV.4.5}) hold there with $n$ replaced by $n-1$. They also hold for $n$ itself, except at points $u$ such that $\delta_{t}^{\prime}(u)=0$. Now, for $u\in U$, the condition $\delta_{t}^{\prime}(u)=0$ is equivalent to saying that $u$ is a \emph{closed point of $U_t$} (\resp a \emph{closed point of $X$}). Let $E$ be the set of points $u$ of $U$ such that $\delta_{t}^{\prime}(u)=0$; let us show that \emph{for all but finitely many points $u\in E$}, one has $\prof_{u}(F)\geq n$. Let $\overline{S}$ be the strict localization of $S$ at $t$, let $S^{\prime}$ be the completion of~$\overline{S}$, with closed point $t^{\prime}$, and consider the Cartesian square
$$
\xymatrix{ U^{\prime}\ar[r]^-{h}\ar[d]_-{f^{\prime}} & U\ar[d]^-{f}\\
S^{\prime}\ar[r]^-{g} & S.}$$
The depth hypotheses at the points of $U$ are preserved when\pageoriginale $U$ is replaced by $U^{\prime}$ and $F$ by its inverse image $F^{\prime}$ on $U^{\prime}$. Indeed, let $u^{\prime}\in U^{\prime}$ and $u=h(u^{\prime})$. If $u\not\in E$, one has $\prof_{u}(F)\geq n-\delta_{t}^{\prime}(u)$, and \Ref{XIV.4.2.1} shows that this implies $\prof_{u^{\prime}}(F^{\prime})\geq n-\delta_{t^{\prime}}^{\prime}(u^{\prime})$. If $u\in E$, then $u^{\prime}$ is a closed point of $U_{t}^{\prime}$ (\resp a point of $X^{\prime}=X\times_{S}S^{\prime}$ closed in $X^{\prime}$), and, since the fiber $U_{u}^{\prime}$ of $h$ at $u$ has dimension zero and $h$ is regular, \Ref{XIV.1.16} gives $\prof_{u^{\prime}}(F^{\prime})=\prof_{u}(F)\geq n-1$.

Let $K$ be a dualizing complex on $S^{\prime}$ normalized at the closed point $t^{\prime}$, and let $\D F^{\prime}$ be the dual of $F^{\prime}$ with respect to $\R^{!}f(K)$. By \Ref{XIV.4.2.2}, the \`etale-depth hypotheses at the points of $U^{\prime}$ translate into the relations
$$
\big(\underline{H}^{q}(\D F^{\prime})\big)_{\overline{u}^{\prime}}=0\quad\textrm{ for }q>-n-\delta_{t^{\prime}}^{\prime}(u^{\prime})\quad(\textrm{\resp }q>-n-2-\delta_{t^{\prime}}^{\prime}(u^{\prime})),
$$
if $u^{\prime}$ is not a point of $E^{\prime}=h^{-1}(E)$, and
$$
\big(\underline{H}^{q}(\D F^{\prime})\big)_{\overline{u}^{\prime}}=0\quad\textrm{ for }q>-(n-1)\quad(\textrm{\resp }q>-n-1), \textrm{ if }u^{\prime}\in E^{\prime}.
$$

Let $G=\underline{H}^{-(n-1)}(\D F^{\prime})$ (\resp $G=\underline{H}^{-n-1}(\D F^{\prime})$). Since $G$ is a constructible sheaf, the set of points at which its geometric fiber is nonzero is constructible (\SGA 4~IX~2.4~\textup{(iv)}). By hypothesis this set is contained in the set $E^{\prime}$ of closed points of $U_{t^{\prime}}^{\prime}$ (\resp of points of $U^{\prime}$ closed in $X^{\prime}$); it follows from \Ref{XIV.4.9.1} below that this set consists of finitely many points. Applying \Ref{XIV.4.2.2}, one sees that $\prof_{u^{\prime}}(F^{\prime})\geq n$ for all but finitely many points of $E^{\prime}$. By \Ref{XIV.1.16}, it follows that, for all but finitely many\pageoriginale points of $E^{\sisi{\prime}{}}$,
$$
\prof_{u}(F)\geq n.
$$

Let $V$ be an open neighborhood of $Y$ in $X$, contained in the complement in $X$ of the finite set of point\sisi{}{s} $u$ of $E$ for which $\prof_{u}(F)=n-1$. If $\sisi{i}{\iota}\colon V\to X$ is the canonical immersion, set
$$F_{1}=\R \sisi{i}{\iota}_{\ast}(\sisi{F{|V}}{F_{|V}})\;;$$
then $F_{1}$ is a bounded-below complex of sheaves on $X$ with constructible cohomology (\SGA 4~XIX~5.1). We shall see that, for every point $u$ of $U$, the complex $F_{1}$ satisfies
\begin{equation*} \label{eq:XIV.4.9.*} \tag{$*$}
\prof_{u}(F_{1})\geq n-\delta_{t}^{\prime}(u).
\end{equation*}
If $u\in U\cap V$, then $\prof_{u}(F_{1})=\prof_{u}(F)$, and relation \eqref{eq:XIV.4.9.*} holds by hypothesis at the points of $U$ that do not belong to $E$; at the latter points, it also holds by the choice of $V$. Finally, if $u\in U$ and $u\not\in V$, then by \Ref{XIV.1.6} g), $\prof_{u}(F_{1})=\infty$. Apply \Ref{XIV.4.4} (\resp \Ref{XIV.4.5}) with $F$ replaced by $F_{1}$; this gives the asserted result, since, for every $i$,
$$
\H_{f^{-1}(t)}^{i}(X, \R \sisi{i}{\iota}_{\ast}(\sisi{F{|V}}{F_{|V}}))\simeq\H_{f^{-1}(t)\cap V}^{i}(V, F)\quad(\textrm{\resp }\H^{i}(X, \R \sisi{i}{\iota}_{\ast}(\sisi{F{|V}}{F_{|V}}))\simeq\H^{i}(V, F).
$$
\skipqed
\end{proof}

\begin{sublemme} \label{XIV.4.9.1} A constructible set $E$ contained in the set of closed points of a Noetherian scheme $X$ consists of finitely many points\pageoriginale.
\end{sublemme}

Indeed, $E$ is a \sisi{finite union}{finite union} \sisi{of a set}{of sets} of the form $U\cap\complement V$, where $U$ and $V$ are open subsets of $X$; by hypothesis, all the points of $U\cap\complement V$ are generic points of this set, so there are only finitely many of them.

\section{Geometric depth} \label{XIV.5}

To apply \Ref{XIV.4.2} and its corollaries in practice, one needs a convenient criterion for verifying the étale-depth hypotheses at the points of $U$. We shall give such a criterion, using the local Lefschetz theorem~\sisi{(\Ref{XIV.4.5})}{\Ref{XIV.4.5}}.

\subsection{\texorpdfstring{}{XIV.5.1}} \label{XIV.5.1} \sisi{Let}{Let} $A$ be a Noetherian local ring; when we speak of the étale depth of $A$, we mean the depth at the closed point. We shall introduce a notion of \og geometric depth of $A$\fg and use \Ref{XIV.4.5} to compare it with the étale depth $\profet(A)$.

\begin{proposition} \label{XIV.5.2} Let $A$ be a Noetherian local ring; suppose that $A$ is isomorphic to a quotient of a regular local ring $B$ by an ideal $I$ (this is true, for example, when $A$ is complete, by Cohen's theorem (\textup{$\EGA 0_{\textup{IV}}$~19.8.8})). Let $q$ be the minimal number of generators of $I$; then the number $\dim(B)-q$ is independent of the choice of $B$.
\end{proposition}

The minimal number of generators of $I$ is also equal to the dimension of the $k$-vector space $I\otimes_{B}k$, where $k$ denotes the residue field of $A$. We reduce at once to the case where $A$ is complete, since $\widehat{A}\simeq\widehat{B}/\widehat{I}$ with $\dim \widehat{B}=\dim B$ and $\rg_{k}(I\otimes_{B}k)=\rg_{k}(\widehat{I}\otimes_{\widehat{B}}k)$; for the same reason we\pageoriginale may suppose that th\sisi{ese}{e} ring\sisi{s}{} $B$ \sisi{are}{is} complete\sisi{}{}. Let $B$ and $B^{\prime}$ be two complete regular local rings, $f\colon B\to A$ and $f^{\prime}\colon B^{\prime}\to A$ two surjective homomorphisms, and $I=\Ker(f)$, $I^{\prime}=\Ker(f^{\prime})$. We must show that
$$
\dim B-\rg_{k}(I\otimes_{B}k)=\dim B^{\prime}-\rg_{k}(I^{\prime}\otimes_{B^{\prime}}k).
$$

First consider the case in which there is a factorization
of the form
$$
\xymatrix@C=15pt{ B\ar[rr]^-{f}\ar[rd]_-{g} & & A\\
& B^{\prime}\ar[ru]_-{f^{\prime}} & }$$
with $g$ surjective. Let $J=\Ker(g)$; then $J\subset I$ and $I/J=I^{\prime}$. Since $B^{\prime}$ is regular, $\dim(B^{\prime})=\dim(B)-\rg_{k}(J\otimes_{B}k)$, and $J$ is generated by elements belonging to a regular system of parameters of $B$. It follows that there is an exact sequence
$$0\to J\otimes_{B}k\to I\otimes_{B}k\to I/J\otimes_{B^{\prime}}k\to 0,
$$
and consequently
$$
\dim B-\rg_{k}(I\otimes_{B}k)=\dim B-\rg_{k}(J\otimes_{B}k)-\rg_{k}(I/J\otimes_{B^{\prime}}k)=\dim B^{\prime}-\rg_{k}(I^{\prime}\otimes_{B^{\prime}}k).
$$

The general case reduces to the preceding one; to see this, it is enough to show that one can find a complete regular local ring $B^{\prime\prime}$ and surjective homomorphisms $g\colon B^{\prime\prime}\to B$ and $g^{\prime}\colon B^{\prime\prime}\to B^{\prime}$ making the diagram\pageoriginale commutative
\begin{equation*} \label{eq:XIV.5.2.*} \tag{$*$}
\begin{array}{c}
\xymatrix@C=10pt@R=15pt{ & B\ar[rd]^-{f} & \\
B^{\prime\prime}\ar[ru]^-{g}\ar[rd]_-{g^{\prime}} & & A\\
& B^{\prime}.\ar[ru]_-{f^{\prime}} &
}
\end{array}
\end{equation*}
Now, if $W$ is a Cohen ring with residue field $k$, there is a local morphism $W\to A$ which lifts to $B$ and $B^{\prime}$ (\EGA IV~19.8.6), so that one has the commutative diagram
$$
\xymatrix@R=15pt@C=10pt{ & B\ar[rd] & \\
W\ar[ru]\ar[rd] & & A\\
& B^{\prime}.\ar[ru] & }$$
One can find integers $n$ and $n^{\prime}$ and surjective morphisms \hbox{$h\colon W\llbracket T_{1},\dots, T_{n}\rrbracket\!\to\! B$} and $h^{\prime}\colon W\llbracket T_{1}^{\prime},\dots, T_{n^{\prime}}^{\prime}\rrbracket\to B^{\prime}$ which are morphisms of $W$-algebras ($\EGA 0_{\textup{IV}}$ 19.8.8); if we then set $B^{\prime\prime}=W\llbracket T_{1},\dots, T_{n}, T_{1}^{\prime},\dots, T_{n^{\prime}}^{\prime}\rrbracket$ and define $g$ and~$g^{\prime}$ as morphisms of $W$-algebras such that
$$g(T_{i})=h(T_{i}), \quad g(T_{i}^{\prime})=b_{i}, \quad g^{\prime}(T_{i})=b_{i}^{\prime}, \quad g^{\prime}(T_{i}^{\prime})=h^{\prime}(T_{i}^{\prime}),
$$
where $b_{i}$ (respectively $b_{i}^{\prime}$) is an element of $B$ (respectively of $B^{\prime}$) lifting \sisi{$h^{\prime}.f^{\prime}(T_{i}^{\prime})$ (respectively $h.f(T_{i})$)}{$(f^{\prime}\circ h^{\prime})(T_{i}^{\prime})$ (respectively $(f\circ\nobreak h)(T_{i})$)}, the diagram \eqref{eq:XIV.5.2.*} is indeed commutative.

Proposition \Ref{XIV.5.2} justifies the following definition:

\begin{definition} \label{XIV.5.3}
Let\pageoriginale $A$ be a Noetherian local ring and $\widehat{A}$ its completion, which is thus isomorphic to the quotient of a complete regular local ring $B$ by an ideal $I$; if $q$ is the minimal number of generators of $I$, the \emph{geometric depth} of $A$ is the number
$$
\profgeom(A)=\dim B-q
$$
\end{definition}

\begin{proposition} \label{XIV.5.4} Let $A$ be a Noetherian local ring. Then
$$
\profgeom(A)\leq\dim A,
$$
with equality if and only if $A$ is \sisi{a}{a} complete intersection.
\end{proposition}

We may suppose that $A$ is complete. Let $A=B/I$, where $B$ is a complete regular local ring and $I$ an ideal of $B$. If $(x_{1},\dots, x_{q})$ is a minimal system of generators of $I$, then $\dim A\geq\dim B-q$, and saying that $\dim A=\dim B-q$ is equivalent to saying that $(x_{1},\dots, x_{q})$ belongs to a system of parameters of $B$ ($\EGA 0_{\textup{IV}}$~16.3.7); the proposition follows immediately.

\begin{proposition} \label{XIV.5.5} Let $A$ and $A^{\prime}$ be Noetherian local rings, and $f\colon A\to A^{\prime}$ a local homomorphism. Suppose that $f$ is flat and that, denoting by $k$ the residue field of $A$, $A^{\prime}\otimes_{A}k$ is a field, separable over $\sisi{K}{k}$. Then
$$
\profgeom(A)=\profgeom(A^{\prime}).
$$
\end{proposition}

After\pageoriginale replacing $A$ and $A^{\prime}$ by their completions, we may suppose that $A$ and $A^{\prime}$ are complete (it follows from ($\EGA 0_{\textup{III}}$~10.2.1) that the flatness hypothesis is preserved, and this is evident for the other hypotheses). Let $A=B/I$, where $B$ is a regular local ring and $I$ an ideal of $B$. Since $A^{\prime}$ is formally smooth over $A$ ($\EGA 0_{\textup{IV}}$~19.8.2), it follows from ($\EGA 0_{\textup{IV}}$~19.7.2) that one can find a complete Noetherian local ring $B^{\prime}$ and a local homomorphism $B\to B^{\prime}$ such that $B^{\prime}$ is a flat $B$-module and $B^{\prime}\otimes_{B}A\simeq A^{\prime}$. Thus $A^{\prime}\simeq B^{\prime}/IB^{\prime}$; moreover the ring $B^{\prime}$ is regular: indeed, if $\mm$ is the maximal ideal of $B$, $\mm B^{\prime}$ is the maximal ideal of $B^{\prime}$, and, since $\mm$ is generated by a regular sequence by the definition of \og regular\fg, $\sisi{B^{\prime}}{\mm B^{\prime}}$ is generated by a $B^{\prime}$-regular sequence ($\EGA 0_{\textup{IV}}$~15.1.14). Since one clearly has $\dim B=\dim B^{\prime}$, and since $I$ and $IB^{\prime}$ have the same minimal number of generators, the assertion follows.

\refstepcounter{subsection}\label{XIV.5.6}
\begin{enonce*}{Theorem 5.6\ndemark}\ndetext{Illusie subsequently proved the inequality $\prof_x(\ZZ/\ell^\nu\ZZ)\geq \profgeom_x(X/S)-\delta(x)+1$ for $x$ a point of $X$, where $X$ is a scheme of finite type over a trait $S$ whose residual characteristic is prime to $\ell$, and $\nu\geq 1$. If $S$ has characteristic zero, this is a consequence of Theorem~\Ref{XIV.5.6}; see (Illusie~L., {\og Perversité et variation\fg}, \emph{Manuscripta Math.} \textbf{112} (2003), p\ptbl 271-295).}
Let $A$ be an excellent local ring of characteristic zero. Then
$$
\profet(A)\geq\profgeom(A).
$$
\end{enonce*}

One may suppose that $A$ is strictly local and complete, since geometric depth and \'{e}tale depth are preserved under passage to the strict henselization and to the completion by \Ref{XIV.5.5} and \Ref{XIV.1.16}. Let $A\simeq B/I$, where $B$ is a complete regular local ring, and let $(f_{1},\dots, f_{q})$ be a minimal system of generators of the ideal $I$. Thus one has\pageoriginale
$$
\pi=\profgeom(A)=\dim B-q.
$$
Consider the closed immersion
$$Y=\Spec A\to X=\Spec B,
$$
and let\sisi{}{} $U=X-Y=\bigcup_{1\leq i\leq q}X_{f_{i}}$. If $a$ denotes the closed point of $X$, one must show that, for every prime number $p$, one has
$$
\H_{a}^{i}(Y, \ZZ/p\ZZ)=0\quad\textrm{ for }i<\pi.
$$
Since $B$ is regular and excellent, one has $\profet(B)=2\dim X$ (\cf \Ref{XIV.1.10}) and consequently $\H_{a}^{i}(X, \ZZ/p\ZZ)=0$ for $i<2\dim X$. Thus, to prove the theorem, it suffices to prove that the morphism
\begin{equation*} \label{eq:XIV.5.6.*} \tag{$*$}
\H_{a}^{i}(X, \ZZ/p\ZZ)\to\H_{a}^{i}(Y, \ZZ/p\ZZ)
\end{equation*}
is bijective for $i<\pi$. For this, apply the local Lefschetz theorem \sisi{(\Ref{XIV.4.5})}{\Ref{XIV.4.5}} with $n=\pi+q$, $c=q-1$, and hence $n-c-1=\pi$. Note that $U=X-Y$ is the union of the $q$ affine open subsets $X_{f_{i}}$. Let us show that, for every point $u$ of $U$, one has:
$$
\profet_{u}(X)\geq\pi+q-\dim\big(\overline{\{ u\} }\big)=\dim\Oo_{X, u}$$
(where $\overline{\{ u\} }$ denotes the closure of $u$ \emph{in} $X-\overline{\{ a\} }$). Indeed, it follows from \Ref{XIV.1.10} that
$$
\profet_{u}(X)=2\dim\Oo_{X, u}\geq\dim\Oo_{X, u}.
$$
Using\pageoriginaled \Ref{XIV.4.5}, one sees that \eqref{eq:XIV.5.6.*} is bijective for $i<\pi$, which completes the proof of the theorem.

\begin{corollaire} \label{XIV.5.7} Let $S$ be the spectrum of a field of characteristic zero (\resp an excellent henselian local scheme of characteristic zero), and let $f\colon X\to S$ exhibit $X$ as a scheme proper over $S$ (\resp over $S-\{ s\}$). Let $U$ be a union of $c+1$ open subsets of $X$ that are \emph{affine}, let $Y$ be a closed subscheme with underlying space $X-U$, and let $n$ and $m$ be integers $>0$. Suppose that, for every point $u$ of $U$, one has
$$
\profgeom(\Oo_{X, u})\geq n-\dim\big(\overline{\{ u\} }\big)$$
($\overline{\{ u\} }$ being the closure of $u$ in $X$). Then the canonical morphism
$$
\H^{i}(X, \ZZ/m\ZZ)\to\H^{i}(Y, \ZZ/m\ZZ)$$
is bijective for $i<n-c-1$, and injective for $i=n-c-1$.
\end{corollaire}

Apply \Ref{XIV.4.5} and \Ref{XIV.4.6}. The hypotheses on \'{e}tale depth at the points of $U$ are satisfied because, by \Ref{XIV.5.6}, one has \sisi{$$
\profet_{u}(X)\geq\profgeom\Oo_{X, u}\geq n-\dim\big(\overline{\{ u\} }\big).
$$
}{$$
\profet_{u}(X)\geq\profgeom(\Oo_{X, u})\geq n-\dim\big(\overline{\{ u\} }\big).
$$
}
\enlargethispage{\baselineskip}%

\section{Open questions} \label{XIV.6}

\subsection{\texorpdfstring{}{XIV.6.1}} \label{XIV.6.1} One may ask whether the implication \textup{(ii)} \ALORS \textup{(i)} of \Ref{XIV.4.2} holds more generally for torsion sheaves $F$, not necessarily annihilated by a given integer $m$ and not necessarily constructible\refstepcounter{pageoriginal}\unskip\marginpar{\raggedright\footnotesize\bfseries\thepageoriginal}\nobreak. In the case where $S$ is not of characteristic zero, it seems possible that this implication remains valid, even for sheaves of $p$-torsion ($p$ being the residue characteristic). Finally, it is also unclear whether the hypothesis that $S$ is excellent can be removed.

\subsection{\texorpdfstring{}{XIV.6.2}} \label{XIV.6.2}
Let $X$ be a scheme proper over a field $k$, or else the complement of the closed point of a henselian local scheme, and let $j:Y \to X$ be a closed subscheme of $X$ whose complement $U$ is affine. Then, if $F$ is a sheaf of sets on $X$ or a sheaf of not necessarily commutative groups, the statements \Ref{XIV.4.5} or \Ref{XIV.4.6} and \Ref{XIV.4.9} still make sense for such an $F$, provided that one restricts to small values of $n$. If $u$ is a point of $U$, denote by $\bar{u}$ a geometric point lying over~$u$, by $X(\bar{u})$ the strict localization of $X$ at $\bar{u}$, and by $F_{\bar{u}}$ the fiber of $F$ at $\bar{u}$. Then, possibly making certain hypotheses on $X$ and $F$, for example by assuming $X$ excellent (possibly of characteristic zero, or of equal characteristic using resolution of singularities) and $F$ ind-finite (or, if necessary, even $\LL$-ind-finite with $\LL$ prime to the characteristic of $X$), one would like to prove the following statements:
\begin{enumeratea}
\item
Let $F$ be a sheaf of sets (\resp a sheaf of groups), and suppose that, for every point $u$ of $U$, one has
$$F_{\bar{u}} \to H^0(X(\bar{u})-\bar{u}, F)\text{ injective if } \dim(\overline{\{u\}}) \leq 1$$
(that is, for such a $u$, one has $\prof_u(F) \geq 1$). Then, as $V$ ranges over the set of open neighborhoods of $Y$, the canonical morphism\pageoriginale
$$
\varinjlim_V H^0(V, F) \to H^0(Y, j^*F)
$$
is bijective (\resp one has the preceding conclusion and, moreover, the morphism $\varinjlim_V H^1(V, F) \to H^1(Y, j^* F)$ is injective). If $F$ is constructible, one may replace the $\varinjlim$ by the cohomology of $V$ for $V$ \og sufficiently small\fg.
\item
Let $F$ be a sheaf of sets (\resp a sheaf of groups), and suppose that, for every point $u$ of $U$, one has $\prof_u(F) \geq 2 -\dim(\overline{\{u\}})$, which is also expressed by the relations
\begin{gather*}
F_{\bar{u}}\to H^0(X(\bar{u})-\bar{u}, F)\text{ is bijective if } \dim(\overline{\{u\}})=0\\
F_{\bar{u}}\to H^0(X(\bar{u})-\bar{u}, F)\text{ is injective if } \dim(\overline{\{u\}})=1.
\end{gather*}
Then the canonical morphism
$$H^0(X, F) \to H^0(Y, j^*F)$$
is bijective (\resp one has the preceding conclusion and, moreover, the morphism $H^1(X, F) \to H^1(Y, j^*F)$ is injective).
\item
Let $F$ be an ind-finite sheaf of groups. Suppose that, for every point $u$ of $U$, one has
\begin{gather*}
F_{\bar{u}}\to H^0(X(\bar{u})-\bar{u}, F)\text{ bijective if } \dim(\overline{\{u\}})=0\text{ or }1,\\
F_{\bar{u}}\to H^0(X(\bar{u})-\bar{u}, F)\text{ injective if } \dim(\overline{\{u\}})=2.
\end{gather*}
Then\pageoriginale, as $V$ ranges over the set of open neighborhoods of $Y$, the canonical morphisms
$$
\varinjlim_V H^0(V, F) \to H^0(Y, j^*F) \text{ and } \varinjlim_V H^1(V, F) \to H^1(Y, j^*F)
$$
are bijective. If $F$ is constructible, one may replace the $\varinjlim$ by the cohomology of~$V$ for $V$ sufficiently small.
\item
Let $F$ be a sheaf of groups. Suppose that, for every point $u$ of $U$, one has $\prof_u(F) \geq 3-\dim(\overline{\{u\}})$, which is also expressed by the conditions
\begin{gather*}
F_{\bar{u}}\to H^0(X(\bar{u})-\bar{u}, F)\text{ bijective, and } H^1(X(\bar{u})-\bar{u}, F)=0 \text{ if } \dim(\overline{\{u\}})=0,\\
F_{\bar{u}}\to H^0(X(\bar{u})-\bar{u}, F)\text{ bijective if } \dim(\overline{\{u\}})=1,\\
F_{\bar{u}}\to H^0(X(\bar{u})-\bar{u}, F)\text{ injective if } \dim(\overline{\{u\}})=2.
\end{gather*}
Then the canonical morphisms
$$H^0(X, F) \to H^0(Y, j^*F)\text{ and } H^1(X, F) \to H^1(Y, j^*F)$$
are bijective.
\end{enumeratea}

As evidence in favor of these statements\refstepcounter{toto}\nde{\label{MiR}all the statements of~\Ref{XIV.6.2}, apart from the constructible variants, were proved by Mrs Raynaud; see (Raynaud~M., {\og Th\'eor\`emes de Lefschetz en cohomologie des faisceaux coh\'erents et en cohomologie \'etale. Application au groupe fondamental\fg}, \emph{Ann. Sci. \'Ec. Norm. Sup. (4)} \textbf{7} (1974), p\ptbl 29--52, Corollary III.1.3).}, let us mention \Ref{XIII}~\Ref{XIII.2.1}, \Ref{X}~\Ref{X.3.4}, and \Ref{XII}~\Ref{XII.3.5}. Note that, by the argument of \Ref{XIV.4.8} and \Ref{XIV.4.9}, statement a) (\resp c)) would follow from b) (\resp d)).

\subsection{\texorpdfstring{}{XIV.6.3}} \label{XIV.6.3}
The following statement, analogous to \Ref{XIV.5.6}, would follow from d): if $A$ is\pageoriginale a noetherian local ring (possibly excellent) and $\prof \geom (A) \geq 3$, then $\prof \hop (A) \geq 3$. To see this, realize $Y'=\Spec A$ as a closed subscheme of a regular local scheme $X'=\Spec B$, whose complement is the union of $q$ affine open subsets, with the relation $\dim B -q=\prof \geom(A)$. For every point $x$ of $X'$, if $n=\dim B$, one has $\prof \hop_x(X') \geq \inf(3, n-\dim(\overline{\{x\}}))$ (\cf \Ref{XIV.1.11}), and one deduces from d) that this implies $\prof \hop_y(Y') \geq \inf(3, n-q-\dim(\overline{\{y\}}))$ for every point $y$ of $Y'$. The result is then obtained by taking $y$ to be the closed point of $Y'$.

\subsection{\texorpdfstring{}{XIV.6.4}} \label{XIV.6.4}
A variant of \Ref{XIV.4.2}, at least of the implication \textup{(ii)} \ALORS \textup{(i)}, should also hold in the complex analytic case, provided one works with \og analytically constructible\fg sheaves (\cf \Ref{XIII}); the proof would be analogous to that of \Ref{XIV.4.2}, using the duality theory of J.-L. Verdier. Note, however, that for the complex analytic analogue of the noncommutative variants mentioned in \Ref{XIV.6.2}, not even a method of approach is available for the statements concerning the fundamental group suggested by the results of Expos\'{e}s \Ref{X}, \Ref{XII}, and \Ref{XIII}, recalled at the end of \Ref{XIV.6.2}. Indeed, the methods of the Seminar seem inextricably tied to the case of \emph{finite} coverings (which may be studied in terms of coherent sheaves of \sisi{Algebras}{algebras}).

\begin{thebibliography}{9}
\bibitem[1]{XIV.1}\sisi{%
M\ptbl Artin et B\ptbl Mazur, Homotopy of Varieties in Etale Topology, in Proceedings of a Conference on Local Fields, Springer (1967).}
{%
\textsc{M\ptbl Artin \& B\ptbl Mazur} -- {\og Homotopy of Varieties in Etale Topology\fg}, in \emph{Proceedings of a Conference on Local Fields}, Springer, 1967.
}
\bibitem[2]{XIV.2}\sisi{%
J.-P\ptbl Serre, Cohomologie Galoisienne, Springer-Verlag, 1964.}
{%
\textsc{J.-P\ptbl Serre} -- \emph{Cohomologie Galoisienne}, Springer-Verlag, 1964.
}
\bibitem[3]{XIV.3}\sisi{%
J.-L\ptbl Verdier, Alg\`ebre homologique et Cat\'egories d\'eriv\'ees, North.\ Holland Pub.\ Cie.\, to appear}
{%
\textsc{J.-L\ptbl Verdier} -- \emph{Des cat\'egories d\'eriv\'ees des cat\'egories ab\'eliennes}, with a preface by Luc Illusie, edited by Georges Maltsiniotis, Ast\'erisque, vol.~239, Soci\'et\'e math\'ematique de France, Paris, 1996.
}
\end{thebibliography}

\chapter*{Index of notation}
\markboth{INDEX OF NOTATION}{INDEX OF NOTATION}

\begin{tabular}{p{9cm}p{3cm}}
$\Gamma_{Z}, \ \underline{\Gamma}_{Z}$\ \dotfill & \Ref{I}~\Ref{I.1}\\
$i_{!}, \ i^{!}$\ \dotfill & \Ref{I}~\Ref{I.1}\\
$\ZZ_{Z, X}$\ \dotfill & \Ref{I}~\Ref{I.1.6}\\
$H^{i}_{Z}(X, F), \underline{H}^{i}_{Z}(F)$\ \dotfill & \Ref{I}~\Ref{I.2}\\
$H^{i}_{J}(M), H^{i}((f), M), H^{i}(M)$\ \dotfill & \Ref{IV}~\Ref{IV.5.4}, \Ref{V}~\Ref{V.2} (p\ptbl\sisi{4}{\pageref{pV2.4}})\\
$\Ass M$, $\Supp M$, $\Ann \ M\sisi{}{,} \rr(\aaa)$\ \dotfill & \Ref{III}~\Ref{III.1.1}\\
$\prof_{I}(M), \prof_{Y}(F)$\ \dotfill & \Ref{III}~\Ref{III.2.3}, \Ref{III.2.8}\\
$\Ab$\ \dotfill & \Ref{IV}~\Ref{IV.1} (p\ptbl\sisi{1}{\pageref{pIV.1}})\\
$\Hom^{.}(F^{.}, G^{.})$\ \dotfill & \Ref{V}~\Ref{V.1.1}\\
$\Ext^{i}_{Z}(X;F, G), \underline{\Ext}^{i}_{Z}(F, G)$\ \dotfill & \Ref{VI}~\Ref{VI.1.1}\\
${\Et}(X), \LL(Z)$\ \dotfill & \Ref{X} (p\ptbl\sisi{1}{\pageref{pX.1}})\\
$\Lef(X, Y), \Leff(X, Y)$\ \dotfill & \Ref{X}~\Ref{X.2} (p\ptbl\sisi{2}{\pageref{pX.2}})\\
$\underline{P}(Z), \Pic(Z)$\ \dotfill & \Ref{XI} (p\ptbl\sisi{1}{\pageref{pXI.1}})\\
$\Pi_{i}^{x}(X)$\ \dotfill & \Ref{XIII} (p\ptbl\sisi{15}{\pageref{pXIII.15}})\\
$D^{+}(X)$\ \dotfill & \Ref{XIV}~\Ref{XIV.1.0}\\
$\prof_{Y}(F) \geq n$\ \dotfill & \Ref{XIV}~\Ref{XIV.1.2}\\
$\prof_{Y}^{\LL}(X) \geq n$\ \dotfill & \Ref{XIV}~\Ref{XIV.1.2}\\
$\profet_{Y}(X)$\ \dotfill & \Ref{XIV}~\Ref{XIV.1.2}\\
$\LL$\ \dotfill & \Ref{XIV}~\Ref{XIV.1.2}\\
$\prof_{x}(F) \geq n$\ \dotfill & \Ref{XIV}~\Ref{XIV.1.7}\\
$\prof_{x}^{\LL}(X)\geq n$\ \dotfill & \Ref{XIV}~\Ref{XIV.1.7}\\
$\prof\ \hop^{\LL}_{x}(X) \geq 3$\ \dotfill & \Ref{XIV}~\Ref{XIV.1.7}\\
$\profet_{x}(X)$\ \dotfill& \Ref{XIV}~\Ref{XIV.1.7}\\
$\delta_{t}(x), \ \delta (n)$\ \dotfill & \Ref{XIV}~\Ref{XIV.2.1}\\
$H^{i}_{Z!}(X/S, F)$\ \dotfill & \Ref{XIV}~\Ref{XIV.4.0}
\end{tabular}

\chapter*{Terminological index}
\markboth{TERMINOLOGICAL INDEX}{TERMINOLOGICAL INDEX}

\begin{footnotesize}
\noindent
\begin{tabular}{lll}
Auslander-Buchsbaum (th\'eor\`eme de -) & & \Ref{XI}~\Ref{XI.3.13}\\
comparaison (th\'eor\`eme de -) & comparison theorem & \Ref{IX}~\Ref{IX.1.1}, \Ref{XII}~\Ref{XII.2.1}\\
dualit\'e locale (th\'eor\`eme de -) & local duality theorem & \Ref{V}~\Ref{V.2.1}\\
dualit\'e projective (th\'eor\`eme de -) & projective duality theorem & \Ref{XII}~\Ref{XII.1.1} \\
enveloppe injective & injective envelope & \Ref{IV} (p\ptbl\sisi{15}{\pageref{pIV.13}})\\
excision (th\'eor\`eme d' -) & excision theorem & \Ref{I}~\Ref{I.2.2}, \Ref{VI}~\Ref{VI.1.3}\\
existence (th\'eor\`eme d' -) & existence theorem & \Ref{IX}~\Ref{IX.2.1}, \Ref{XII}~\Ref{XII.3.1}\\
extension essentielle & essential extension & \Ref{IV} (p\ptbl\sisi{15}{\pageref{pIV.13}})\\
finitude (th\'eor\`eme de -) & finiteness theorem & \Ref{VIII}~\Ref{VIII.2.1}, \Ref{XII}~\Ref{XII.1.5}\\
foncteur dualisant & dualizing functor & \Ref{IV}~\Ref{IV.4.1}, \Ref{IV.4.2}\\
forme r\'esidu & residue form & \Ref{IV}~\Ref{IV.5.5}\\
g\'eom\'etriquement factoriel, \resp & geometrically factorial (\resp &\\
g\'eom. parafactoriel (anneau local -) & geom. parafactorial) local ring & \Ref{XIII} (p\ptbl20, 24)\\
groupes d'homotopie locale & local homotopy groups & \Ref{XIII} (p\ptbl\sisi{15}{\pageref{pXIII.15}})\\
Gysin (homomorphisme de -) & Gysin homomorphism & \Ref{I} (p\ptbl\sisi{13}{\pageref{hgysin}})\\
Hartshorne (th\'eor\`eme de -) & & \Ref{III}~\Ref{III.3.6}\\
Lefschetz (condition de -) & Lefschetz condition & \Ref{X}~\Ref{X.2} (p\ptbl\sisi{2}{\pageref{pX.2}})\\
Lefschetz (condition de - effective) & effective Lefschetz condition & \Ref{X}~\Ref{X.2} (p\ptbl\sisi{2}{\pageref{pX.2}})\\
Lefschetz (th\'eor\`eme de - affine) & affine Lefschetz theorem & \Ref{XIV}~\Ref{XIV.3.1}\\
module dualisant & dualizing module & \Ref{IV}~\Ref{IV.4.1}\\
\sisi{module orthogonal d'un module}{orthogonal d'un sous-module} & \sisi{module orthogonal of a module}{orthogonal of a submodule} & \Ref{IV}~\Ref{IV.5}\\
parafactoriel (couple - de pr\'e\-sch\'emas) & parafactorial pair of preschemes & \Ref{XI}~\Ref{XI.3.1}\\
parafactoriel (anneau local -) & parafactorial local ring & \Ref{XI}~\Ref{XI.3.2}\\
profondeur & depth & \Ref{III}~\Ref{III.2.3}\\
profondeur \'etale & etale depth & \Ref{XIV}~\Ref{XIV.1.2}\\
profondeur g\'eom\'etrique d'un anneau & geometrical depth of a noetherian &\\
local noeth\'erien & local ring & \Ref{XIV}~\Ref{XIV.5.3}\\
profondeur homotopique & homotopical depth & \Ref{XIII}~\Ref{XIII.6} def\ptbl1 (p\ptbl\sisi{4}{\pageref{XIII.6.1}})\\
profondeur homotopique (pour $\LL$) & homotopical depth (with respect to $\LL$) & \Ref{XIV}~\Ref{XIV.1.2}\\
profondeur homotopique rectifi\'ee & rectified homotopical depth & \Ref{XIII} def\ptbl2 (p\ptbl\sisi{5}{\pageref{XIII.6.2}})\\
\end{tabular}
\newpage
\begin{tabular}{lll}
pur (anneau local -) & pure local ring & \Ref{X}~\Ref{X.3.2} \\
pure (couple - de pr\'esch\'emas) & pure pair of preschemes & \Ref{X}~\Ref{X.3.1}\\
puret\'e (th\'eor\`eme de - de Zariski-Nagata) & purity theorem of Zariski-Nagata & \Ref{X}~\Ref{X.3.4}\\
puret\'e cohomologique (th\'eor\`eme de -) & cohomological purity theorem & \Ref{XIV}~\Ref{XIV.1.11}\\
r\'egulier ($M$-r\'egulier) & $M$-regular & \Ref{III.2.1}\\
Samuel (conjecture de -) & & \Ref{XI}~\Ref{XI.3.14}\\
semi-puret\'e cohomologique (th\'eor\`eme de -) & cohomological semi-purity theorem & \Ref{XIV}~\Ref{XIV.1.10}\\
strictement local (anneau -) & strictly local ring & \Ref{XIII}~\Ref{XIII.6} (p\ptbl\sisi{23}{\pageref{pXIII.6}})\\
Zariski-Nagata (th\'eor\`eme de puret\'e de -) & purity theorem of Zariski-Nagata & \Ref{X}~\Ref{X.3.4}
\end{tabular}
\end{footnotesize}

\backmatter

% END INCLUDED active_source/parts/14_expose_XIV.tex


\end{document}
